% 2026-10-26 DRAFT appendix subsection: market exposure of the sign strategy and its allocation with the market.
% Not \input by paper_rehaul.tex.
% Pre-registration: results/econ_beta_2026-10-26/RULE.txt (written 2026-10-26T02:30Z, before any strategy-vs-market
%   statistic was computed; its ADDENDUM lists what was added after the results were seen -- descriptive only).
% Generator: writeup/make_rehaul_econ_beta.py (run from the repository root with /c/Users/james/miniconda3/python.exe;
%   it sets REHAUL_FINAL_RUN=1, gates the 14 strategies' Sharpe ratios against \rhemSh<tag> / \rhemShX<tag> and the
%   best-day union against \rhedBestUnion / \rhedTopDate before writing).
% Files:   generated/rehaul_econ_beta_numbers.tex  (\rheb... macros; macro-only; no name collides with a macro of
%                                                    another writeup/generated/*.tex file, checked by the generator)
%          generated/rehaul_econ_beta_table.tex    (Table tab:rh_econ_beta)
%          generated/rehaul_econ_beta_fig.tex      (Figure fig:rh_econ_beta_cum, generated/rehaul_econ_beta_cum.pdf)
% To use:  % AUTO-GENERATED by writeup/make_rehaul_econ_beta.py -- do not edit; re-run the script.
% \rheb... macros: market exposure of the sign strategy and its allocation with the market (appendix,
% tab:rh_econ_beta). Pre-registration: results/econ_beta_2026-10-26/RULE.txt. Row stems <Est><Set> (Est = Ridge,
% Lasso, Enet, Lgbm, Xgb, Rf; Set = Full, Har), CiHar, Short; X = quoted bid/ask prices, else mid-quotes; Ex =
% without the union of the full-set forecasts' ten best days. Market: S&P 500 close-to-close price return minus
% the DTB3 T-bill return (no dividends); Thirty = the index's 15:30-16:00 return. STANDARD Sharpe ratio =
% sqrt(q) mean / sd of the pooled daily returns, q = T / sum_y c_y of the sample's days (2026-10-26, D19: the
% paper's Sharpe ratio; was sqrt(252)); D6 = the equal-year mean of SR_y (hac_inference.d6_sharpe; robustness).
\providecommand{\rhebNDays}{}\renewcommand{\rhebNDays}{614}  % post-warmup option days of the exposure and spanning statistics
\providecommand{\rhebNBurn}{}\renewcommand{\rhebNBurn}{126}  % burn-in days of the real-time rules (expanding window)
\providecommand{\rhebNTrail}{}\renewcommand{\rhebNTrail}{63}  % trailing days of the inverse-volatility rule
\providecommand{\rhebNRt}{}\renewcommand{\rhebNRt}{488}  % evaluation days of the real-time rules (post-warmup days after the burn-in)
\providecommand{\rhebRtFirst}{}\renewcommand{\rhebRtFirst}{17~May~2022}  % first evaluation day of the real-time rules
\providecommand{\rhebCap}{}\renewcommand{\rhebCap}{90}  % cap: largest share (%) of the portfolio volatility either leg may carry
\providecommand{\rhebNHarDistinct}{}\renewcommand{\rhebNHarDistinct}{4}  % distinct HAR-set forecasts (the three linear ones coincide)
\providecommand{\rhebNHarDistinctWord}{}\renewcommand{\rhebNHarDistinctWord}{four}  % ... in words
\providecommand{\rhebNBest}{}\renewcommand{\rhebNBest}{11}  % days removed in the robustness check (= \rhedBestUnion)
\providecommand{\rhebNEx}{}\renewcommand{\rhebNEx}{603}  % post-warmup days without them
\providecommand{\rhebNExRt}{}\renewcommand{\rhebNExRt}{477}  % real-time evaluation days without them
\providecommand{\rhebTopDate}{}\renewcommand{\rhebTopDate}{22~November~2021}  % the largest of them (= \rhedTopDate)
\providecommand{\rhebBootBlock}{}\renewcommand{\rhebBootBlock}{10}  % block length of the paired circular block bootstrap
\providecommand{\rhebBootN}{}\renewcommand{\rhebBootN}{4{,}999}  % bootstrap resamples
\providecommand{\rhebMktMean}{}\renewcommand{\rhebMktMean}{0.011}  % market excess return, mean, % per day, option days
\providecommand{\rhebMktSd}{}\renewcommand{\rhebMktSd}{1.12}  % ... standard deviation, % per day
\providecommand{\rhebRfMean}{}\renewcommand{\rhebRfMean}{0.0143}  % risk-free return (DTB3, actual/360), mean, % per day
\providecommand{\rhebRateMean}{}\renewcommand{\rhebRateMean}{3.25}  % DTB3, mean over the trading days of the span, % a year
\providecommand{\rhebSRm}{}\renewcommand{\rhebSRm}{0.15}  % market: Sharpe ratio, STANDARD (D19: sqrt(q), q = 225.5), 614 option days
\providecommand{\rhebSRmDSix}{}\renewcommand{\rhebSRmDSix}{0.57}  % market: Sharpe ratio, D6, 614 option days
\providecommand{\rhebSRmAll}{}\renewcommand{\rhebSRmAll}{0.17}  % market: Sharpe ratio, STANDARD (D19: sqrt(q) of these days), every trading day of the span
\providecommand{\rhebQAll}{}\renewcommand{\rhebQAll}{251.3}  % D19: q of every trading day of the span
\providecommand{\rhebQRt}{}\renewcommand{\rhebQRt}{249.2}  % D19: q of the real-time evaluation days
\providecommand{\rhebNAll}{}\renewcommand{\rhebNAll}{684}  % ... trading days of the span 2021-08-11 .. 2024-04-30
\providecommand{\rhebSRmRt}{}\renewcommand{\rhebSRmRt}{0.50}  % market: Sharpe ratio, STANDARD, real-time evaluation days
\providecommand{\rhebSRmRtDSix}{}\renewcommand{\rhebSRmRtDSix}{0.68}  % market: Sharpe ratio, D6, real-time evaluation days
\providecommand{\rhebSRmEx}{}\renewcommand{\rhebSRmEx}{0.23}  % market: Sharpe ratio, STANDARD, without the best-day union
\providecommand{\rhebSRmExRt}{}\renewcommand{\rhebSRmExRt}{0.54}  % market: Sharpe ratio, STANDARD, real-time days without the union
\providecommand{\rhebWorstCut}{}\renewcommand{\rhebWorstCut}{\ensuremath{-}1.29}  % worst-decile threshold: 10% quantile of the market excess return, %
\providecommand{\rhebNWorst}{}\renewcommand{\rhebNWorst}{62}  % worst-decile market days
\providecommand{\rhebNDown}{}\renewcommand{\rhebNDown}{307}  % market-down days (excess return < 0)
\providecommand{\rhebWorstCutThirty}{}\renewcommand{\rhebWorstCutThirty}{\ensuremath{-}0.33}  % ... 10% quantile of the 15:30-16:00 return, %
\providecommand{\rhebCorrRidgeFull}{}\renewcommand{\rhebCorrRidgeFull}{0.01}  % poolpt13_ridge_all_features, mid: correlation with the market excess return
\providecommand{\rhebBetaRidgeFull}{}\renewcommand{\rhebBetaRidgeFull}{1.27}  % ... CAPM beta (% of premium per % of market)
\providecommand{\rhebBetaTRidgeFull}{}\renewcommand{\rhebBetaTRidgeFull}{0.34}  % ... its HAC t (S = 33)
\providecommand{\rhebBetaPRidgeFull}{}\renewcommand{\rhebBetaPRidgeFull}{0.746}  % ... its fixed-b p
\providecommand{\rhebRsqRidgeFull}{}\renewcommand{\rhebRsqRidgeFull}{0.00}  % ... CAPM R^2
\providecommand{\rhebAlphaRidgeFull}{}\renewcommand{\rhebAlphaRidgeFull}{12.3}  % ... CAPM alpha, % of premium per day
\providecommand{\rhebAlphaTRidgeFull}{}\renewcommand{\rhebAlphaTRidgeFull}{2.93}  % ... its HAC t
\providecommand{\rhebAlphaPRidgeFull}{}\renewcommand{\rhebAlphaPRidgeFull}{0.008}  % ... its fixed-b p
\providecommand{\rhebAlphaAnnRidgeFull}{}\renewcommand{\rhebAlphaAnnRidgeFull}{2776}  % ... alpha x q (D19; was x 252), % of premium a year
\providecommand{\rhebARRidgeFull}{}\renewcommand{\rhebARRidgeFull}{1.64}  % ... appraisal ratio alpha / sigma_eps, annualized sqrt(q) (D19)
\providecommand{\rhebSRsRidgeFull}{}\renewcommand{\rhebSRsRidgeFull}{1.64}  % ... strategy Sharpe ratio, STANDARD = the paper's (= \rhemSh / X)
\providecommand{\rhebSRsDSixRidgeFull}{}\renewcommand{\rhebSRsDSixRidgeFull}{1.75}  % ... strategy Sharpe ratio, D6 (= \rhemShDSix)
\providecommand{\rhebSRmaxRidgeFull}{}\renewcommand{\rhebSRmaxRidgeFull}{1.65}  % ... max Sharpe ratio of market + strategy, STANDARD
\providecommand{\rhebDSRmaxRidgeFull}{}\renewcommand{\rhebDSRmaxRidgeFull}{1.50}  % ... its gain over the market alone, STANDARD
\providecommand{\rhebDSRmaxLoRidgeFull}{}\renewcommand{\rhebDSRmaxLoRidgeFull}{0.24}  % ... bootstrap 95% interval, lower end
\providecommand{\rhebDSRmaxHiRidgeFull}{}\renewcommand{\rhebDSRmaxHiRidgeFull}{2.46}  % ... upper end
\providecommand{\rhebDSRmaxPRidgeFull}{}\renewcommand{\rhebDSRmaxPRidgeFull}{0.002}  % ... bootstrap p (descriptive; Delta >= 0)
\providecommand{\rhebGrsRidgeFull}{}\renewcommand{\rhebGrsRidgeFull}{7.31}  % ... one-asset GRS statistic, F(1, T - 2)
\providecommand{\rhebGrsPRidgeFull}{}\renewcommand{\rhebGrsPRidgeFull}{0.007}  % ... its p (i.i.d. errors)
\providecommand{\rhebTanShareRidgeFull}{}\renewcommand{\rhebTanShareRidgeFull}{99}  % ... tangency: strategy's share of volatility, %
\providecommand{\rhebTanDSixRidgeFull}{}\renewcommand{\rhebTanDSixRidgeFull}{1.76}  % ... tangency (fixed in-sample weights): Sharpe ratio, D6
\providecommand{\rhebRtMvRidgeFull}{}\renewcommand{\rhebRtMvRidgeFull}{1.70}  % ... real-time mean-variance: Sharpe ratio, STANDARD
\providecommand{\rhebRtMvDRidgeFull}{}\renewcommand{\rhebRtMvDRidgeFull}{1.20}  % ... minus the market's, STANDARD
\providecommand{\rhebRtMvPRidgeFull}{}\renewcommand{\rhebRtMvPRidgeFull}{0.146}  % ... Ledoit-Wolf p of the difference (D19)
\providecommand{\rhebRtMvDSixRidgeFull}{}\renewcommand{\rhebRtMvDSixRidgeFull}{1.76}  % ... real-time mean-variance: Sharpe ratio, D6
\providecommand{\rhebRtMvDSixDRidgeFull}{}\renewcommand{\rhebRtMvDSixDRidgeFull}{1.07}  % ... minus the market's, D6
\providecommand{\rhebRtMvDSixPRidgeFull}{}\renewcommand{\rhebRtMvDSixPRidgeFull}{0.243}  % ... bootstrap p of the D6 difference
\providecommand{\rhebExRtMvRidgeFull}{}\renewcommand{\rhebExRtMvRidgeFull}{1.11}  % ... w/o best days, real-time mean-variance: SR, STANDARD
\providecommand{\rhebExRtMvDRidgeFull}{}\renewcommand{\rhebExRtMvDRidgeFull}{0.57}  % ... minus the market's, STANDARD
\providecommand{\rhebExRtMvPRidgeFull}{}\renewcommand{\rhebExRtMvPRidgeFull}{0.506}  % ... its Ledoit-Wolf p (D19)
\providecommand{\rhebRtIvRidgeFull}{}\renewcommand{\rhebRtIvRidgeFull}{1.46}  % ... real-time inverse-volatility: Sharpe ratio, STANDARD
\providecommand{\rhebRtIvDRidgeFull}{}\renewcommand{\rhebRtIvDRidgeFull}{0.96}  % ... minus the market's, STANDARD
\providecommand{\rhebRtIvPRidgeFull}{}\renewcommand{\rhebRtIvPRidgeFull}{0.103}  % ... Ledoit-Wolf p of the difference (D19)
\providecommand{\rhebRtIvDSixRidgeFull}{}\renewcommand{\rhebRtIvDSixRidgeFull}{1.76}  % ... real-time inverse-volatility: Sharpe ratio, D6
\providecommand{\rhebRtIvDSixDRidgeFull}{}\renewcommand{\rhebRtIvDSixDRidgeFull}{1.07}  % ... minus the market's, D6
\providecommand{\rhebRtIvDSixPRidgeFull}{}\renewcommand{\rhebRtIvDSixPRidgeFull}{0.092}  % ... bootstrap p of the D6 difference
\providecommand{\rhebExRtIvRidgeFull}{}\renewcommand{\rhebExRtIvRidgeFull}{1.18}  % ... w/o best days, real-time inverse-volatility: SR, STANDARD
\providecommand{\rhebExRtIvDRidgeFull}{}\renewcommand{\rhebExRtIvDRidgeFull}{0.64}  % ... minus the market's, STANDARD
\providecommand{\rhebExRtIvPRidgeFull}{}\renewcommand{\rhebExRtIvPRidgeFull}{0.286}  % ... its Ledoit-Wolf p (D19)
\providecommand{\rhebRtMvShareRidgeFull}{}\renewcommand{\rhebRtMvShareRidgeFull}{90}  % ... real-time MV: median strategy vol share, %
\providecommand{\rhebRtMvCapDaysRidgeFull}{}\renewcommand{\rhebRtMvCapDaysRidgeFull}{100}  % ... % of days the 90% cap binds
\providecommand{\rhebRtMvWsRidgeFull}{}\renewcommand{\rhebRtMvWsRidgeFull}{3.4}  % ... real-time MV: median strategy dollar weight, %
\providecommand{\rhebRtIvWsRidgeFull}{}\renewcommand{\rhebRtIvWsRidgeFull}{0.9}  % ... real-time IV: median strategy dollar weight, %
\providecommand{\rhebExBetaRidgeFull}{}\renewcommand{\rhebExBetaRidgeFull}{1.59}  % ... without the best-day union: CAPM beta
\providecommand{\rhebExBetaTRidgeFull}{}\renewcommand{\rhebExBetaTRidgeFull}{0.47}  % ... its HAC t
\providecommand{\rhebExAlphaRidgeFull}{}\renewcommand{\rhebExAlphaRidgeFull}{6.7}  % ... without the best-day union: alpha, % per day
\providecommand{\rhebExAlphaTRidgeFull}{}\renewcommand{\rhebExAlphaTRidgeFull}{1.98}  % ... its HAC t
\providecommand{\rhebExAlphaPRidgeFull}{}\renewcommand{\rhebExAlphaPRidgeFull}{0.064}  % ... its fixed-b p
\providecommand{\rhebExARRidgeFull}{}\renewcommand{\rhebExARRidgeFull}{1.02}  % ... appraisal ratio, sqrt(q) (D19)
\providecommand{\rhebExSRmaxRidgeFull}{}\renewcommand{\rhebExSRmaxRidgeFull}{1.05}  % ... max Sharpe ratio, STANDARD
\providecommand{\rhebExDSRmaxPRidgeFull}{}\renewcommand{\rhebExDSRmaxPRidgeFull}{0.065}  % ... bootstrap p of its gain
\providecommand{\rhebExGrsPRidgeFull}{}\renewcommand{\rhebExGrsPRidgeFull}{0.093}  % ... GRS p
\providecommand{\rhebTmGammaRidgeFull}{}\renewcommand{\rhebTmGammaRidgeFull}{\ensuremath{-}1.76}  % poolpt13_ridge_all_features, mid: Treynor-Mazuy gamma on (r^m)^2 (% per %^2)
\providecommand{\rhebTmGammaTRidgeFull}{}\renewcommand{\rhebTmGammaTRidgeFull}{\ensuremath{-}1.31}  % ... its HAC t
\providecommand{\rhebTmGammaPRidgeFull}{}\renewcommand{\rhebTmGammaPRidgeFull}{0.213}  % ... its fixed-b p
\providecommand{\rhebTmBetaRidgeFull}{}\renewcommand{\rhebTmBetaRidgeFull}{1.20}  % ... Treynor-Mazuy beta on r^m
\providecommand{\rhebTmBetaTRidgeFull}{}\renewcommand{\rhebTmBetaTRidgeFull}{0.31}  % ... its HAC t
\providecommand{\rhebNLongRidgeFull}{}\renewcommand{\rhebNLongRidgeFull}{123}  % ... long days
\providecommand{\rhebTmGammaLongRidgeFull}{}\renewcommand{\rhebTmGammaLongRidgeFull}{\ensuremath{-}1.25}  % ... TM gamma on its long days only
\providecommand{\rhebTmGammaLongTRidgeFull}{}\renewcommand{\rhebTmGammaLongTRidgeFull}{\ensuremath{-}0.20}  % ... its HAC t (no test: RULE section 4)
\providecommand{\rhebNShortRidgeFull}{}\renewcommand{\rhebNShortRidgeFull}{491}  % ... short days
\providecommand{\rhebTmGammaShortRidgeFull}{}\renewcommand{\rhebTmGammaShortRidgeFull}{\ensuremath{-}1.75}  % ... TM gamma on its short days only
\providecommand{\rhebTmGammaShortTRidgeFull}{}\renewcommand{\rhebTmGammaShortTRidgeFull}{\ensuremath{-}1.31}  % ... its HAC t (no test: RULE section 4)
\providecommand{\rhebDownMeanRidgeFull}{}\renewcommand{\rhebDownMeanRidgeFull}{16.2}  % ... mean return on market-down days, % of premium
\providecommand{\rhebDownRestRidgeFull}{}\renewcommand{\rhebDownRestRidgeFull}{8.5}  % ... on the other days
\providecommand{\rhebDownDiffTRidgeFull}{}\renewcommand{\rhebDownDiffTRidgeFull}{0.71}  % ... HAC t of the difference
\providecommand{\rhebDownDiffPRidgeFull}{}\renewcommand{\rhebDownDiffPRidgeFull}{0.497}  % ... its fixed-b p
\providecommand{\rhebWorstMeanRidgeFull}{}\renewcommand{\rhebWorstMeanRidgeFull}{\ensuremath{-}10.2}  % ... mean return on worst-decile market days, % of premium
\providecommand{\rhebWorstRestRidgeFull}{}\renewcommand{\rhebWorstRestRidgeFull}{14.9}  % ... on the other days
\providecommand{\rhebWorstDiffTRidgeFull}{}\renewcommand{\rhebWorstDiffTRidgeFull}{\ensuremath{-}1.70}  % ... HAC t of the difference
\providecommand{\rhebWorstDiffPRidgeFull}{}\renewcommand{\rhebWorstDiffPRidgeFull}{0.109}  % ... its fixed-b p
\providecommand{\rhebThirtyCorrRidgeFull}{}\renewcommand{\rhebThirtyCorrRidgeFull}{0.03}  % ... correlation with the 15:30-16:00 index return
\providecommand{\rhebThirtyBetaRidgeFull}{}\renewcommand{\rhebThirtyBetaRidgeFull}{11.08}  % ... beta on the 15:30-16:00 index return
\providecommand{\rhebThirtyBetaTRidgeFull}{}\renewcommand{\rhebThirtyBetaTRidgeFull}{0.37}  % ... its HAC t
\providecommand{\rhebThirtyBetaPRidgeFull}{}\renewcommand{\rhebThirtyBetaPRidgeFull}{0.725}  % ... its fixed-b p
\providecommand{\rhebThirtyAlphaRidgeFull}{}\renewcommand{\rhebThirtyAlphaRidgeFull}{12.5}  % ... alpha against it, % per day
\providecommand{\rhebThirtyAlphaTRidgeFull}{}\renewcommand{\rhebThirtyAlphaTRidgeFull}{3.02}  % ... its HAC t
\providecommand{\rhebThirtyTmGammaRidgeFull}{}\renewcommand{\rhebThirtyTmGammaRidgeFull}{\ensuremath{-}111.64}  % ... TM gamma on the squared 15:30-16:00 return
\providecommand{\rhebThirtyTmGammaTRidgeFull}{}\renewcommand{\rhebThirtyTmGammaTRidgeFull}{\ensuremath{-}2.26}  % ... its HAC t
\providecommand{\rhebThirtyTmGammaPRidgeFull}{}\renewcommand{\rhebThirtyTmGammaPRidgeFull}{0.036}  % ... its fixed-b p
\providecommand{\rhebThirtyTmGammaLongRidgeFull}{}\renewcommand{\rhebThirtyTmGammaLongRidgeFull}{374.16}  % ... ... on its long days only
\providecommand{\rhebThirtyTmGammaLongTRidgeFull}{}\renewcommand{\rhebThirtyTmGammaLongTRidgeFull}{2.59}  % ... its HAC t (no test)
\providecommand{\rhebThirtyTmGammaShortRidgeFull}{}\renewcommand{\rhebThirtyTmGammaShortRidgeFull}{\ensuremath{-}272.01}  % ... ... on its short days only
\providecommand{\rhebThirtyTmGammaShortTRidgeFull}{}\renewcommand{\rhebThirtyTmGammaShortTRidgeFull}{\ensuremath{-}4.35}  % ... its HAC t (no test)
\providecommand{\rhebThirtyDownMeanRidgeFull}{}\renewcommand{\rhebThirtyDownMeanRidgeFull}{11.9}  % ... mean on 15:30-16:00 down days, %
\providecommand{\rhebThirtyDownDiffTRidgeFull}{}\renewcommand{\rhebThirtyDownDiffTRidgeFull}{\ensuremath{-}0.11}  % ... HAC t of the difference with the other days
\providecommand{\rhebThirtyWorstMeanRidgeFull}{}\renewcommand{\rhebThirtyWorstMeanRidgeFull}{\ensuremath{-}51.7}  % ... mean on 15:30-16:00 worst days, %
\providecommand{\rhebThirtyWorstDiffTRidgeFull}{}\renewcommand{\rhebThirtyWorstDiffTRidgeFull}{\ensuremath{-}2.60}  % ... HAC t of the difference with the other days
\providecommand{\rhebCorrXRidgeFull}{}\renewcommand{\rhebCorrXRidgeFull}{0.01}  % poolpt13_ridge_all_features, bid/ask: correlation with the market excess return
\providecommand{\rhebBetaXRidgeFull}{}\renewcommand{\rhebBetaXRidgeFull}{1.34}  % ... CAPM beta (% of premium per % of market)
\providecommand{\rhebBetaTXRidgeFull}{}\renewcommand{\rhebBetaTXRidgeFull}{0.35}  % ... its HAC t (S = 33)
\providecommand{\rhebBetaPXRidgeFull}{}\renewcommand{\rhebBetaPXRidgeFull}{0.737}  % ... its fixed-b p
\providecommand{\rhebRsqXRidgeFull}{}\renewcommand{\rhebRsqXRidgeFull}{0.00}  % ... CAPM R^2
\providecommand{\rhebAlphaXRidgeFull}{}\renewcommand{\rhebAlphaXRidgeFull}{9.7}  % ... CAPM alpha, % of premium per day
\providecommand{\rhebAlphaTXRidgeFull}{}\renewcommand{\rhebAlphaTXRidgeFull}{2.32}  % ... its HAC t
\providecommand{\rhebAlphaPXRidgeFull}{}\renewcommand{\rhebAlphaPXRidgeFull}{0.032}  % ... its fixed-b p
\providecommand{\rhebAlphaAnnXRidgeFull}{}\renewcommand{\rhebAlphaAnnXRidgeFull}{2183}  % ... alpha x q (D19; was x 252), % of premium a year
\providecommand{\rhebARXRidgeFull}{}\renewcommand{\rhebARXRidgeFull}{1.28}  % ... appraisal ratio alpha / sigma_eps, annualized sqrt(q) (D19)
\providecommand{\rhebSRsXRidgeFull}{}\renewcommand{\rhebSRsXRidgeFull}{1.29}  % ... strategy Sharpe ratio, STANDARD = the paper's (= \rhemSh / X)
\providecommand{\rhebSRsDSixXRidgeFull}{}\renewcommand{\rhebSRsDSixXRidgeFull}{1.40}  % ... strategy Sharpe ratio, D6 (= \rhemShDSix)
\providecommand{\rhebSRmaxXRidgeFull}{}\renewcommand{\rhebSRmaxXRidgeFull}{1.29}  % ... max Sharpe ratio of market + strategy, STANDARD
\providecommand{\rhebDSRmaxXRidgeFull}{}\renewcommand{\rhebDSRmaxXRidgeFull}{1.14}  % ... its gain over the market alone, STANDARD
\providecommand{\rhebDSRmaxLoXRidgeFull}{}\renewcommand{\rhebDSRmaxLoXRidgeFull}{0.05}  % ... bootstrap 95% interval, lower end
\providecommand{\rhebDSRmaxHiXRidgeFull}{}\renewcommand{\rhebDSRmaxHiXRidgeFull}{2.11}  % ... upper end
\providecommand{\rhebDSRmaxPXRidgeFull}{}\renewcommand{\rhebDSRmaxPXRidgeFull}{0.011}  % ... bootstrap p (descriptive; Delta >= 0)
\providecommand{\rhebGrsXRidgeFull}{}\renewcommand{\rhebGrsXRidgeFull}{4.48}  % ... one-asset GRS statistic, F(1, T - 2)
\providecommand{\rhebGrsPXRidgeFull}{}\renewcommand{\rhebGrsPXRidgeFull}{0.035}  % ... its p (i.i.d. errors)
\providecommand{\rhebTanShareXRidgeFull}{}\renewcommand{\rhebTanShareXRidgeFull}{99}  % ... tangency: strategy's share of volatility, %
\providecommand{\rhebTanDSixXRidgeFull}{}\renewcommand{\rhebTanDSixXRidgeFull}{1.41}  % ... tangency (fixed in-sample weights): Sharpe ratio, D6
\providecommand{\rhebRtMvXRidgeFull}{}\renewcommand{\rhebRtMvXRidgeFull}{1.35}  % ... real-time mean-variance: Sharpe ratio, STANDARD
\providecommand{\rhebRtMvDXRidgeFull}{}\renewcommand{\rhebRtMvDXRidgeFull}{0.85}  % ... minus the market's, STANDARD
\providecommand{\rhebRtMvPXRidgeFull}{}\renewcommand{\rhebRtMvPXRidgeFull}{0.293}  % ... Ledoit-Wolf p of the difference (D19)
\providecommand{\rhebRtMvDSixXRidgeFull}{}\renewcommand{\rhebRtMvDSixXRidgeFull}{1.41}  % ... real-time mean-variance: Sharpe ratio, D6
\providecommand{\rhebRtMvDSixDXRidgeFull}{}\renewcommand{\rhebRtMvDSixDXRidgeFull}{0.73}  % ... minus the market's, D6
\providecommand{\rhebRtMvDSixPXRidgeFull}{}\renewcommand{\rhebRtMvDSixPXRidgeFull}{0.420}  % ... bootstrap p of the D6 difference
\providecommand{\rhebExRtMvXRidgeFull}{}\renewcommand{\rhebExRtMvXRidgeFull}{0.75}  % ... w/o best days, real-time mean-variance: SR, STANDARD
\providecommand{\rhebExRtMvDXRidgeFull}{}\renewcommand{\rhebExRtMvDXRidgeFull}{0.20}  % ... minus the market's, STANDARD
\providecommand{\rhebExRtMvPXRidgeFull}{}\renewcommand{\rhebExRtMvPXRidgeFull}{0.807}  % ... its Ledoit-Wolf p (D19)
\providecommand{\rhebRtIvXRidgeFull}{}\renewcommand{\rhebRtIvXRidgeFull}{1.17}  % ... real-time inverse-volatility: Sharpe ratio, STANDARD
\providecommand{\rhebRtIvDXRidgeFull}{}\renewcommand{\rhebRtIvDXRidgeFull}{0.68}  % ... minus the market's, STANDARD
\providecommand{\rhebRtIvPXRidgeFull}{}\renewcommand{\rhebRtIvPXRidgeFull}{0.241}  % ... Ledoit-Wolf p of the difference (D19)
\providecommand{\rhebRtIvDSixXRidgeFull}{}\renewcommand{\rhebRtIvDSixXRidgeFull}{1.48}  % ... real-time inverse-volatility: Sharpe ratio, D6
\providecommand{\rhebRtIvDSixDXRidgeFull}{}\renewcommand{\rhebRtIvDSixDXRidgeFull}{0.79}  % ... minus the market's, D6
\providecommand{\rhebRtIvDSixPXRidgeFull}{}\renewcommand{\rhebRtIvDSixPXRidgeFull}{0.209}  % ... bootstrap p of the D6 difference
\providecommand{\rhebExRtIvXRidgeFull}{}\renewcommand{\rhebExRtIvXRidgeFull}{0.90}  % ... w/o best days, real-time inverse-volatility: SR, STANDARD
\providecommand{\rhebExRtIvDXRidgeFull}{}\renewcommand{\rhebExRtIvDXRidgeFull}{0.36}  % ... minus the market's, STANDARD
\providecommand{\rhebExRtIvPXRidgeFull}{}\renewcommand{\rhebExRtIvPXRidgeFull}{0.547}  % ... its Ledoit-Wolf p (D19)
\providecommand{\rhebRtMvShareXRidgeFull}{}\renewcommand{\rhebRtMvShareXRidgeFull}{90}  % ... real-time MV: median strategy vol share, %
\providecommand{\rhebRtMvCapDaysXRidgeFull}{}\renewcommand{\rhebRtMvCapDaysXRidgeFull}{100}  % ... % of days the 90% cap binds
\providecommand{\rhebRtMvWsXRidgeFull}{}\renewcommand{\rhebRtMvWsXRidgeFull}{3.4}  % ... real-time MV: median strategy dollar weight, %
\providecommand{\rhebRtIvWsXRidgeFull}{}\renewcommand{\rhebRtIvWsXRidgeFull}{0.8}  % ... real-time IV: median strategy dollar weight, %
\providecommand{\rhebExBetaXRidgeFull}{}\renewcommand{\rhebExBetaXRidgeFull}{1.60}  % ... without the best-day union: CAPM beta
\providecommand{\rhebExBetaTXRidgeFull}{}\renewcommand{\rhebExBetaTXRidgeFull}{0.46}  % ... its HAC t
\providecommand{\rhebExAlphaXRidgeFull}{}\renewcommand{\rhebExAlphaXRidgeFull}{4.3}  % ... without the best-day union: alpha, % per day
\providecommand{\rhebExAlphaTXRidgeFull}{}\renewcommand{\rhebExAlphaTXRidgeFull}{1.25}  % ... its HAC t
\providecommand{\rhebExAlphaPXRidgeFull}{}\renewcommand{\rhebExAlphaPXRidgeFull}{0.235}  % ... its fixed-b p
\providecommand{\rhebExARXRidgeFull}{}\renewcommand{\rhebExARXRidgeFull}{0.64}  % ... appraisal ratio, sqrt(q) (D19)
\providecommand{\rhebExSRmaxXRidgeFull}{}\renewcommand{\rhebExSRmaxXRidgeFull}{0.68}  % ... max Sharpe ratio, STANDARD
\providecommand{\rhebExDSRmaxPXRidgeFull}{}\renewcommand{\rhebExDSRmaxPXRidgeFull}{0.165}  % ... bootstrap p of its gain
\providecommand{\rhebExGrsPXRidgeFull}{}\renewcommand{\rhebExGrsPXRidgeFull}{0.293}  % ... GRS p
\providecommand{\rhebCorrLassoFull}{}\renewcommand{\rhebCorrLassoFull}{\ensuremath{-}0.01}  % poolpt13_lasso_all_features, mid: correlation with the market excess return
\providecommand{\rhebBetaLassoFull}{}\renewcommand{\rhebBetaLassoFull}{\ensuremath{-}0.52}  % ... CAPM beta (% of premium per % of market)
\providecommand{\rhebBetaTLassoFull}{}\renewcommand{\rhebBetaTLassoFull}{\ensuremath{-}0.13}  % ... its HAC t (S = 33)
\providecommand{\rhebBetaPLassoFull}{}\renewcommand{\rhebBetaPLassoFull}{0.899}  % ... its fixed-b p
\providecommand{\rhebRsqLassoFull}{}\renewcommand{\rhebRsqLassoFull}{0.00}  % ... CAPM R^2
\providecommand{\rhebAlphaLassoFull}{}\renewcommand{\rhebAlphaLassoFull}{13.0}  % ... CAPM alpha, % of premium per day
\providecommand{\rhebAlphaTLassoFull}{}\renewcommand{\rhebAlphaTLassoFull}{3.34}  % ... its HAC t
\providecommand{\rhebAlphaPLassoFull}{}\renewcommand{\rhebAlphaPLassoFull}{0.003}  % ... its fixed-b p
\providecommand{\rhebAlphaAnnLassoFull}{}\renewcommand{\rhebAlphaAnnLassoFull}{2923}  % ... alpha x q (D19; was x 252), % of premium a year
\providecommand{\rhebARLassoFull}{}\renewcommand{\rhebARLassoFull}{1.73}  % ... appraisal ratio alpha / sigma_eps, annualized sqrt(q) (D19)
\providecommand{\rhebSRsLassoFull}{}\renewcommand{\rhebSRsLassoFull}{1.73}  % ... strategy Sharpe ratio, STANDARD = the paper's (= \rhemSh / X)
\providecommand{\rhebSRsDSixLassoFull}{}\renewcommand{\rhebSRsDSixLassoFull}{1.75}  % ... strategy Sharpe ratio, D6 (= \rhemShDSix)
\providecommand{\rhebSRmaxLassoFull}{}\renewcommand{\rhebSRmaxLassoFull}{1.73}  % ... max Sharpe ratio of market + strategy, STANDARD
\providecommand{\rhebDSRmaxLassoFull}{}\renewcommand{\rhebDSRmaxLassoFull}{1.58}  % ... its gain over the market alone, STANDARD
\providecommand{\rhebDSRmaxLoLassoFull}{}\renewcommand{\rhebDSRmaxLoLassoFull}{0.29}  % ... bootstrap 95% interval, lower end
\providecommand{\rhebDSRmaxHiLassoFull}{}\renewcommand{\rhebDSRmaxHiLassoFull}{2.55}  % ... upper end
\providecommand{\rhebDSRmaxPLassoFull}{}\renewcommand{\rhebDSRmaxPLassoFull}{0.002}  % ... bootstrap p (descriptive; Delta >= 0)
\providecommand{\rhebGrsLassoFull}{}\renewcommand{\rhebGrsLassoFull}{8.11}  % ... one-asset GRS statistic, F(1, T - 2)
\providecommand{\rhebGrsPLassoFull}{}\renewcommand{\rhebGrsPLassoFull}{0.005}  % ... its p (i.i.d. errors)
\providecommand{\rhebTanShareLassoFull}{}\renewcommand{\rhebTanShareLassoFull}{99}  % ... tangency: strategy's share of volatility, %
\providecommand{\rhebTanDSixLassoFull}{}\renewcommand{\rhebTanDSixLassoFull}{1.76}  % ... tangency (fixed in-sample weights): Sharpe ratio, D6
\providecommand{\rhebRtMvLassoFull}{}\renewcommand{\rhebRtMvLassoFull}{1.84}  % ... real-time mean-variance: Sharpe ratio, STANDARD
\providecommand{\rhebRtMvDLassoFull}{}\renewcommand{\rhebRtMvDLassoFull}{1.35}  % ... minus the market's, STANDARD
\providecommand{\rhebRtMvPLassoFull}{}\renewcommand{\rhebRtMvPLassoFull}{0.114}  % ... Ledoit-Wolf p of the difference (D19)
\providecommand{\rhebRtMvDSixLassoFull}{}\renewcommand{\rhebRtMvDSixLassoFull}{1.85}  % ... real-time mean-variance: Sharpe ratio, D6
\providecommand{\rhebRtMvDSixDLassoFull}{}\renewcommand{\rhebRtMvDSixDLassoFull}{1.16}  % ... minus the market's, D6
\providecommand{\rhebRtMvDSixPLassoFull}{}\renewcommand{\rhebRtMvDSixPLassoFull}{0.212}  % ... bootstrap p of the D6 difference
\providecommand{\rhebExRtMvLassoFull}{}\renewcommand{\rhebExRtMvLassoFull}{1.28}  % ... w/o best days, real-time mean-variance: SR, STANDARD
\providecommand{\rhebExRtMvDLassoFull}{}\renewcommand{\rhebExRtMvDLassoFull}{0.73}  % ... minus the market's, STANDARD
\providecommand{\rhebExRtMvPLassoFull}{}\renewcommand{\rhebExRtMvPLassoFull}{0.390}  % ... its Ledoit-Wolf p (D19)
\providecommand{\rhebRtIvLassoFull}{}\renewcommand{\rhebRtIvLassoFull}{1.60}  % ... real-time inverse-volatility: Sharpe ratio, STANDARD
\providecommand{\rhebRtIvDLassoFull}{}\renewcommand{\rhebRtIvDLassoFull}{1.10}  % ... minus the market's, STANDARD
\providecommand{\rhebRtIvPLassoFull}{}\renewcommand{\rhebRtIvPLassoFull}{0.061}  % ... Ledoit-Wolf p of the difference (D19)
\providecommand{\rhebRtIvDSixLassoFull}{}\renewcommand{\rhebRtIvDSixLassoFull}{1.82}  % ... real-time inverse-volatility: Sharpe ratio, D6
\providecommand{\rhebRtIvDSixDLassoFull}{}\renewcommand{\rhebRtIvDSixDLassoFull}{1.14}  % ... minus the market's, D6
\providecommand{\rhebRtIvDSixPLassoFull}{}\renewcommand{\rhebRtIvDSixPLassoFull}{0.072}  % ... bootstrap p of the D6 difference
\providecommand{\rhebExRtIvLassoFull}{}\renewcommand{\rhebExRtIvLassoFull}{1.32}  % ... w/o best days, real-time inverse-volatility: SR, STANDARD
\providecommand{\rhebExRtIvDLassoFull}{}\renewcommand{\rhebExRtIvDLassoFull}{0.78}  % ... minus the market's, STANDARD
\providecommand{\rhebExRtIvPLassoFull}{}\renewcommand{\rhebExRtIvPLassoFull}{0.183}  % ... its Ledoit-Wolf p (D19)
\providecommand{\rhebRtMvShareLassoFull}{}\renewcommand{\rhebRtMvShareLassoFull}{90}  % ... real-time MV: median strategy vol share, %
\providecommand{\rhebRtMvCapDaysLassoFull}{}\renewcommand{\rhebRtMvCapDaysLassoFull}{100}  % ... % of days the 90% cap binds
\providecommand{\rhebRtMvWsLassoFull}{}\renewcommand{\rhebRtMvWsLassoFull}{3.3}  % ... real-time MV: median strategy dollar weight, %
\providecommand{\rhebRtIvWsLassoFull}{}\renewcommand{\rhebRtIvWsLassoFull}{0.9}  % ... real-time IV: median strategy dollar weight, %
\providecommand{\rhebExBetaLassoFull}{}\renewcommand{\rhebExBetaLassoFull}{\ensuremath{-}0.22}  % ... without the best-day union: CAPM beta
\providecommand{\rhebExBetaTLassoFull}{}\renewcommand{\rhebExBetaTLassoFull}{\ensuremath{-}0.06}  % ... its HAC t
\providecommand{\rhebExAlphaLassoFull}{}\renewcommand{\rhebExAlphaLassoFull}{7.4}  % ... without the best-day union: alpha, % per day
\providecommand{\rhebExAlphaTLassoFull}{}\renewcommand{\rhebExAlphaTLassoFull}{2.31}  % ... its HAC t
\providecommand{\rhebExAlphaPLassoFull}{}\renewcommand{\rhebExAlphaPLassoFull}{0.032}  % ... its fixed-b p
\providecommand{\rhebExARLassoFull}{}\renewcommand{\rhebExARLassoFull}{1.12}  % ... appraisal ratio, sqrt(q) (D19)
\providecommand{\rhebExSRmaxLassoFull}{}\renewcommand{\rhebExSRmaxLassoFull}{1.15}  % ... max Sharpe ratio, STANDARD
\providecommand{\rhebExDSRmaxPLassoFull}{}\renewcommand{\rhebExDSRmaxPLassoFull}{0.049}  % ... bootstrap p of its gain
\providecommand{\rhebExGrsPLassoFull}{}\renewcommand{\rhebExGrsPLassoFull}{0.064}  % ... GRS p
\providecommand{\rhebTmGammaLassoFull}{}\renewcommand{\rhebTmGammaLassoFull}{\ensuremath{-}1.73}  % poolpt13_lasso_all_features, mid: Treynor-Mazuy gamma on (r^m)^2 (% per %^2)
\providecommand{\rhebTmGammaTLassoFull}{}\renewcommand{\rhebTmGammaTLassoFull}{\ensuremath{-}1.36}  % ... its HAC t
\providecommand{\rhebTmGammaPLassoFull}{}\renewcommand{\rhebTmGammaPLassoFull}{0.197}  % ... its fixed-b p
\providecommand{\rhebTmBetaLassoFull}{}\renewcommand{\rhebTmBetaLassoFull}{\ensuremath{-}0.59}  % ... Treynor-Mazuy beta on r^m
\providecommand{\rhebTmBetaTLassoFull}{}\renewcommand{\rhebTmBetaTLassoFull}{\ensuremath{-}0.15}  % ... its HAC t
\providecommand{\rhebNLongLassoFull}{}\renewcommand{\rhebNLongLassoFull}{121}  % ... long days
\providecommand{\rhebTmGammaLongLassoFull}{}\renewcommand{\rhebTmGammaLongLassoFull}{\ensuremath{-}3.45}  % ... TM gamma on its long days only
\providecommand{\rhebTmGammaLongTLassoFull}{}\renewcommand{\rhebTmGammaLongTLassoFull}{\ensuremath{-}0.48}  % ... its HAC t (no test: RULE section 4)
\providecommand{\rhebNShortLassoFull}{}\renewcommand{\rhebNShortLassoFull}{493}  % ... short days
\providecommand{\rhebTmGammaShortLassoFull}{}\renewcommand{\rhebTmGammaShortLassoFull}{\ensuremath{-}1.71}  % ... TM gamma on its short days only
\providecommand{\rhebTmGammaShortTLassoFull}{}\renewcommand{\rhebTmGammaShortTLassoFull}{\ensuremath{-}1.35}  % ... its HAC t (no test: RULE section 4)
\providecommand{\rhebDownMeanLassoFull}{}\renewcommand{\rhebDownMeanLassoFull}{18.7}  % ... mean return on market-down days, % of premium
\providecommand{\rhebDownRestLassoFull}{}\renewcommand{\rhebDownRestLassoFull}{7.2}  % ... on the other days
\providecommand{\rhebDownDiffTLassoFull}{}\renewcommand{\rhebDownDiffTLassoFull}{1.02}  % ... HAC t of the difference
\providecommand{\rhebDownDiffPLassoFull}{}\renewcommand{\rhebDownDiffPLassoFull}{0.330}  % ... its fixed-b p
\providecommand{\rhebWorstMeanLassoFull}{}\renewcommand{\rhebWorstMeanLassoFull}{\ensuremath{-}4.4}  % ... mean return on worst-decile market days, % of premium
\providecommand{\rhebWorstRestLassoFull}{}\renewcommand{\rhebWorstRestLassoFull}{14.9}  % ... on the other days
\providecommand{\rhebWorstDiffTLassoFull}{}\renewcommand{\rhebWorstDiffTLassoFull}{\ensuremath{-}1.31}  % ... HAC t of the difference
\providecommand{\rhebWorstDiffPLassoFull}{}\renewcommand{\rhebWorstDiffPLassoFull}{0.213}  % ... its fixed-b p
\providecommand{\rhebThirtyCorrLassoFull}{}\renewcommand{\rhebThirtyCorrLassoFull}{0.04}  % ... correlation with the 15:30-16:00 index return
\providecommand{\rhebThirtyBetaLassoFull}{}\renewcommand{\rhebThirtyBetaLassoFull}{14.76}  % ... beta on the 15:30-16:00 index return
\providecommand{\rhebThirtyBetaTLassoFull}{}\renewcommand{\rhebThirtyBetaTLassoFull}{0.49}  % ... its HAC t
\providecommand{\rhebThirtyBetaPLassoFull}{}\renewcommand{\rhebThirtyBetaPLassoFull}{0.636}  % ... its fixed-b p
\providecommand{\rhebThirtyAlphaLassoFull}{}\renewcommand{\rhebThirtyAlphaLassoFull}{13.1}  % ... alpha against it, % per day
\providecommand{\rhebThirtyAlphaTLassoFull}{}\renewcommand{\rhebThirtyAlphaTLassoFull}{3.46}  % ... its HAC t
\providecommand{\rhebThirtyTmGammaLassoFull}{}\renewcommand{\rhebThirtyTmGammaLassoFull}{\ensuremath{-}109.48}  % ... TM gamma on the squared 15:30-16:00 return
\providecommand{\rhebThirtyTmGammaTLassoFull}{}\renewcommand{\rhebThirtyTmGammaTLassoFull}{\ensuremath{-}2.26}  % ... its HAC t
\providecommand{\rhebThirtyTmGammaPLassoFull}{}\renewcommand{\rhebThirtyTmGammaPLassoFull}{0.036}  % ... its fixed-b p
\providecommand{\rhebThirtyTmGammaLongLassoFull}{}\renewcommand{\rhebThirtyTmGammaLongLassoFull}{374.13}  % ... ... on its long days only
\providecommand{\rhebThirtyTmGammaLongTLassoFull}{}\renewcommand{\rhebThirtyTmGammaLongTLassoFull}{2.38}  % ... its HAC t (no test)
\providecommand{\rhebThirtyTmGammaShortLassoFull}{}\renewcommand{\rhebThirtyTmGammaShortLassoFull}{\ensuremath{-}270.34}  % ... ... on its short days only
\providecommand{\rhebThirtyTmGammaShortTLassoFull}{}\renewcommand{\rhebThirtyTmGammaShortTLassoFull}{\ensuremath{-}4.36}  % ... its HAC t (no test)
\providecommand{\rhebThirtyDownMeanLassoFull}{}\renewcommand{\rhebThirtyDownMeanLassoFull}{12.0}  % ... mean on 15:30-16:00 down days, %
\providecommand{\rhebThirtyDownDiffTLassoFull}{}\renewcommand{\rhebThirtyDownDiffTLassoFull}{\ensuremath{-}0.25}  % ... HAC t of the difference with the other days
\providecommand{\rhebThirtyWorstMeanLassoFull}{}\renewcommand{\rhebThirtyWorstMeanLassoFull}{\ensuremath{-}51.7}  % ... mean on 15:30-16:00 worst days, %
\providecommand{\rhebThirtyWorstDiffTLassoFull}{}\renewcommand{\rhebThirtyWorstDiffTLassoFull}{\ensuremath{-}2.60}  % ... HAC t of the difference with the other days
\providecommand{\rhebCorrXLassoFull}{}\renewcommand{\rhebCorrXLassoFull}{0.00}  % poolpt13_lasso_all_features, bid/ask: correlation with the market excess return
\providecommand{\rhebBetaXLassoFull}{}\renewcommand{\rhebBetaXLassoFull}{\ensuremath{-}0.45}  % ... CAPM beta (% of premium per % of market)
\providecommand{\rhebBetaTXLassoFull}{}\renewcommand{\rhebBetaTXLassoFull}{\ensuremath{-}0.11}  % ... its HAC t (S = 33)
\providecommand{\rhebBetaPXLassoFull}{}\renewcommand{\rhebBetaPXLassoFull}{0.913}  % ... its fixed-b p
\providecommand{\rhebRsqXLassoFull}{}\renewcommand{\rhebRsqXLassoFull}{0.00}  % ... CAPM R^2
\providecommand{\rhebAlphaXLassoFull}{}\renewcommand{\rhebAlphaXLassoFull}{10.3}  % ... CAPM alpha, % of premium per day
\providecommand{\rhebAlphaTXLassoFull}{}\renewcommand{\rhebAlphaTXLassoFull}{2.68}  % ... its HAC t
\providecommand{\rhebAlphaPXLassoFull}{}\renewcommand{\rhebAlphaPXLassoFull}{0.014}  % ... its fixed-b p
\providecommand{\rhebAlphaAnnXLassoFull}{}\renewcommand{\rhebAlphaAnnXLassoFull}{2329}  % ... alpha x q (D19; was x 252), % of premium a year
\providecommand{\rhebARXLassoFull}{}\renewcommand{\rhebARXLassoFull}{1.37}  % ... appraisal ratio alpha / sigma_eps, annualized sqrt(q) (D19)
\providecommand{\rhebSRsXLassoFull}{}\renewcommand{\rhebSRsXLassoFull}{1.37}  % ... strategy Sharpe ratio, STANDARD = the paper's (= \rhemSh / X)
\providecommand{\rhebSRsDSixXLassoFull}{}\renewcommand{\rhebSRsDSixXLassoFull}{1.40}  % ... strategy Sharpe ratio, D6 (= \rhemShDSix)
\providecommand{\rhebSRmaxXLassoFull}{}\renewcommand{\rhebSRmaxXLassoFull}{1.38}  % ... max Sharpe ratio of market + strategy, STANDARD
\providecommand{\rhebDSRmaxXLassoFull}{}\renewcommand{\rhebDSRmaxXLassoFull}{1.23}  % ... its gain over the market alone, STANDARD
\providecommand{\rhebDSRmaxLoXLassoFull}{}\renewcommand{\rhebDSRmaxLoXLassoFull}{0.08}  % ... bootstrap 95% interval, lower end
\providecommand{\rhebDSRmaxHiXLassoFull}{}\renewcommand{\rhebDSRmaxHiXLassoFull}{2.21}  % ... upper end
\providecommand{\rhebDSRmaxPXLassoFull}{}\renewcommand{\rhebDSRmaxPXLassoFull}{0.009}  % ... bootstrap p (descriptive; Delta >= 0)
\providecommand{\rhebGrsXLassoFull}{}\renewcommand{\rhebGrsXLassoFull}{5.11}  % ... one-asset GRS statistic, F(1, T - 2)
\providecommand{\rhebGrsPXLassoFull}{}\renewcommand{\rhebGrsPXLassoFull}{0.024}  % ... its p (i.i.d. errors)
\providecommand{\rhebTanShareXLassoFull}{}\renewcommand{\rhebTanShareXLassoFull}{99}  % ... tangency: strategy's share of volatility, %
\providecommand{\rhebTanDSixXLassoFull}{}\renewcommand{\rhebTanDSixXLassoFull}{1.40}  % ... tangency (fixed in-sample weights): Sharpe ratio, D6
\providecommand{\rhebRtMvXLassoFull}{}\renewcommand{\rhebRtMvXLassoFull}{1.49}  % ... real-time mean-variance: Sharpe ratio, STANDARD
\providecommand{\rhebRtMvDXLassoFull}{}\renewcommand{\rhebRtMvDXLassoFull}{0.99}  % ... minus the market's, STANDARD
\providecommand{\rhebRtMvPXLassoFull}{}\renewcommand{\rhebRtMvPXLassoFull}{0.232}  % ... Ledoit-Wolf p of the difference (D19)
\providecommand{\rhebRtMvDSixXLassoFull}{}\renewcommand{\rhebRtMvDSixXLassoFull}{1.50}  % ... real-time mean-variance: Sharpe ratio, D6
\providecommand{\rhebRtMvDSixDXLassoFull}{}\renewcommand{\rhebRtMvDSixDXLassoFull}{0.82}  % ... minus the market's, D6
\providecommand{\rhebRtMvDSixPXLassoFull}{}\renewcommand{\rhebRtMvDSixPXLassoFull}{0.375}  % ... bootstrap p of the D6 difference
\providecommand{\rhebExRtMvXLassoFull}{}\renewcommand{\rhebExRtMvXLassoFull}{0.91}  % ... w/o best days, real-time mean-variance: SR, STANDARD
\providecommand{\rhebExRtMvDXLassoFull}{}\renewcommand{\rhebExRtMvDXLassoFull}{0.36}  % ... minus the market's, STANDARD
\providecommand{\rhebExRtMvPXLassoFull}{}\renewcommand{\rhebExRtMvPXLassoFull}{0.662}  % ... its Ledoit-Wolf p (D19)
\providecommand{\rhebRtIvXLassoFull}{}\renewcommand{\rhebRtIvXLassoFull}{1.31}  % ... real-time inverse-volatility: Sharpe ratio, STANDARD
\providecommand{\rhebRtIvDXLassoFull}{}\renewcommand{\rhebRtIvDXLassoFull}{0.81}  % ... minus the market's, STANDARD
\providecommand{\rhebRtIvPXLassoFull}{}\renewcommand{\rhebRtIvPXLassoFull}{0.148}  % ... Ledoit-Wolf p of the difference (D19)
\providecommand{\rhebRtIvDSixXLassoFull}{}\renewcommand{\rhebRtIvDSixXLassoFull}{1.53}  % ... real-time inverse-volatility: Sharpe ratio, D6
\providecommand{\rhebRtIvDSixDXLassoFull}{}\renewcommand{\rhebRtIvDSixDXLassoFull}{0.85}  % ... minus the market's, D6
\providecommand{\rhebRtIvDSixPXLassoFull}{}\renewcommand{\rhebRtIvDSixPXLassoFull}{0.172}  % ... bootstrap p of the D6 difference
\providecommand{\rhebExRtIvXLassoFull}{}\renewcommand{\rhebExRtIvXLassoFull}{1.04}  % ... w/o best days, real-time inverse-volatility: SR, STANDARD
\providecommand{\rhebExRtIvDXLassoFull}{}\renewcommand{\rhebExRtIvDXLassoFull}{0.49}  % ... minus the market's, STANDARD
\providecommand{\rhebExRtIvPXLassoFull}{}\renewcommand{\rhebExRtIvPXLassoFull}{0.387}  % ... its Ledoit-Wolf p (D19)
\providecommand{\rhebRtMvShareXLassoFull}{}\renewcommand{\rhebRtMvShareXLassoFull}{90}  % ... real-time MV: median strategy vol share, %
\providecommand{\rhebRtMvCapDaysXLassoFull}{}\renewcommand{\rhebRtMvCapDaysXLassoFull}{100}  % ... % of days the 90% cap binds
\providecommand{\rhebRtMvWsXLassoFull}{}\renewcommand{\rhebRtMvWsXLassoFull}{3.3}  % ... real-time MV: median strategy dollar weight, %
\providecommand{\rhebRtIvWsXLassoFull}{}\renewcommand{\rhebRtIvWsXLassoFull}{0.9}  % ... real-time IV: median strategy dollar weight, %
\providecommand{\rhebExBetaXLassoFull}{}\renewcommand{\rhebExBetaXLassoFull}{\ensuremath{-}0.21}  % ... without the best-day union: CAPM beta
\providecommand{\rhebExBetaTXLassoFull}{}\renewcommand{\rhebExBetaTXLassoFull}{\ensuremath{-}0.06}  % ... its HAC t
\providecommand{\rhebExAlphaXLassoFull}{}\renewcommand{\rhebExAlphaXLassoFull}{4.9}  % ... without the best-day union: alpha, % per day
\providecommand{\rhebExAlphaTXLassoFull}{}\renewcommand{\rhebExAlphaTXLassoFull}{1.52}  % ... its HAC t
\providecommand{\rhebExAlphaPXLassoFull}{}\renewcommand{\rhebExAlphaPXLassoFull}{0.149}  % ... its fixed-b p
\providecommand{\rhebExARXLassoFull}{}\renewcommand{\rhebExARXLassoFull}{0.74}  % ... appraisal ratio, sqrt(q) (D19)
\providecommand{\rhebExSRmaxXLassoFull}{}\renewcommand{\rhebExSRmaxXLassoFull}{0.78}  % ... max Sharpe ratio, STANDARD
\providecommand{\rhebExDSRmaxPXLassoFull}{}\renewcommand{\rhebExDSRmaxPXLassoFull}{0.131}  % ... bootstrap p of its gain
\providecommand{\rhebExGrsPXLassoFull}{}\renewcommand{\rhebExGrsPXLassoFull}{0.223}  % ... GRS p
\providecommand{\rhebCorrEnetFull}{}\renewcommand{\rhebCorrEnetFull}{0.03}  % poolpt13_enet_all_features, mid: correlation with the market excess return
\providecommand{\rhebBetaEnetFull}{}\renewcommand{\rhebBetaEnetFull}{2.52}  % ... CAPM beta (% of premium per % of market)
\providecommand{\rhebBetaTEnetFull}{}\renewcommand{\rhebBetaTEnetFull}{0.65}  % ... its HAC t (S = 33)
\providecommand{\rhebBetaPEnetFull}{}\renewcommand{\rhebBetaPEnetFull}{0.534}  % ... its fixed-b p
\providecommand{\rhebRsqEnetFull}{}\renewcommand{\rhebRsqEnetFull}{0.00}  % ... CAPM R^2
\providecommand{\rhebAlphaEnetFull}{}\renewcommand{\rhebAlphaEnetFull}{13.9}  % ... CAPM alpha, % of premium per day
\providecommand{\rhebAlphaTEnetFull}{}\renewcommand{\rhebAlphaTEnetFull}{3.27}  % ... its HAC t
\providecommand{\rhebAlphaPEnetFull}{}\renewcommand{\rhebAlphaPEnetFull}{0.003}  % ... its fixed-b p
\providecommand{\rhebAlphaAnnEnetFull}{}\renewcommand{\rhebAlphaAnnEnetFull}{3127}  % ... alpha x q (D19; was x 252), % of premium a year
\providecommand{\rhebAREnetFull}{}\renewcommand{\rhebAREnetFull}{1.85}  % ... appraisal ratio alpha / sigma_eps, annualized sqrt(q) (D19)
\providecommand{\rhebSRsEnetFull}{}\renewcommand{\rhebSRsEnetFull}{1.85}  % ... strategy Sharpe ratio, STANDARD = the paper's (= \rhemSh / X)
\providecommand{\rhebSRsDSixEnetFull}{}\renewcommand{\rhebSRsDSixEnetFull}{1.95}  % ... strategy Sharpe ratio, D6 (= \rhemShDSix)
\providecommand{\rhebSRmaxEnetFull}{}\renewcommand{\rhebSRmaxEnetFull}{1.86}  % ... max Sharpe ratio of market + strategy, STANDARD
\providecommand{\rhebDSRmaxEnetFull}{}\renewcommand{\rhebDSRmaxEnetFull}{1.71}  % ... its gain over the market alone, STANDARD
\providecommand{\rhebDSRmaxLoEnetFull}{}\renewcommand{\rhebDSRmaxLoEnetFull}{0.35}  % ... bootstrap 95% interval, lower end
\providecommand{\rhebDSRmaxHiEnetFull}{}\renewcommand{\rhebDSRmaxHiEnetFull}{2.68}  % ... upper end
\providecommand{\rhebDSRmaxPEnetFull}{}\renewcommand{\rhebDSRmaxPEnetFull}{0.001}  % ... bootstrap p (descriptive; Delta >= 0)
\providecommand{\rhebGrsEnetFull}{}\renewcommand{\rhebGrsEnetFull}{9.31}  % ... one-asset GRS statistic, F(1, T - 2)
\providecommand{\rhebGrsPEnetFull}{}\renewcommand{\rhebGrsPEnetFull}{0.002}  % ... its p (i.i.d. errors)
\providecommand{\rhebTanShareEnetFull}{}\renewcommand{\rhebTanShareEnetFull}{100}  % ... tangency: strategy's share of volatility, %
\providecommand{\rhebTanDSixEnetFull}{}\renewcommand{\rhebTanDSixEnetFull}{1.95}  % ... tangency (fixed in-sample weights): Sharpe ratio, D6
\providecommand{\rhebRtMvEnetFull}{}\renewcommand{\rhebRtMvEnetFull}{1.95}  % ... real-time mean-variance: Sharpe ratio, STANDARD
\providecommand{\rhebRtMvDEnetFull}{}\renewcommand{\rhebRtMvDEnetFull}{1.45}  % ... minus the market's, STANDARD
\providecommand{\rhebRtMvPEnetFull}{}\renewcommand{\rhebRtMvPEnetFull}{0.087}  % ... Ledoit-Wolf p of the difference (D19)
\providecommand{\rhebRtMvDSixEnetFull}{}\renewcommand{\rhebRtMvDSixEnetFull}{2.03}  % ... real-time mean-variance: Sharpe ratio, D6
\providecommand{\rhebRtMvDSixDEnetFull}{}\renewcommand{\rhebRtMvDSixDEnetFull}{1.35}  % ... minus the market's, D6
\providecommand{\rhebRtMvDSixPEnetFull}{}\renewcommand{\rhebRtMvDSixPEnetFull}{0.151}  % ... bootstrap p of the D6 difference
\providecommand{\rhebExRtMvEnetFull}{}\renewcommand{\rhebExRtMvEnetFull}{1.39}  % ... w/o best days, real-time mean-variance: SR, STANDARD
\providecommand{\rhebExRtMvDEnetFull}{}\renewcommand{\rhebExRtMvDEnetFull}{0.85}  % ... minus the market's, STANDARD
\providecommand{\rhebExRtMvPEnetFull}{}\renewcommand{\rhebExRtMvPEnetFull}{0.327}  % ... its Ledoit-Wolf p (D19)
\providecommand{\rhebRtIvEnetFull}{}\renewcommand{\rhebRtIvEnetFull}{1.61}  % ... real-time inverse-volatility: Sharpe ratio, STANDARD
\providecommand{\rhebRtIvDEnetFull}{}\renewcommand{\rhebRtIvDEnetFull}{1.11}  % ... minus the market's, STANDARD
\providecommand{\rhebRtIvPEnetFull}{}\renewcommand{\rhebRtIvPEnetFull}{0.056}  % ... Ledoit-Wolf p of the difference (D19)
\providecommand{\rhebRtIvDSixEnetFull}{}\renewcommand{\rhebRtIvDSixEnetFull}{1.95}  % ... real-time inverse-volatility: Sharpe ratio, D6
\providecommand{\rhebRtIvDSixDEnetFull}{}\renewcommand{\rhebRtIvDSixDEnetFull}{1.27}  % ... minus the market's, D6
\providecommand{\rhebRtIvDSixPEnetFull}{}\renewcommand{\rhebRtIvDSixPEnetFull}{0.047}  % ... bootstrap p of the D6 difference
\providecommand{\rhebExRtIvEnetFull}{}\renewcommand{\rhebExRtIvEnetFull}{1.34}  % ... w/o best days, real-time inverse-volatility: SR, STANDARD
\providecommand{\rhebExRtIvDEnetFull}{}\renewcommand{\rhebExRtIvDEnetFull}{0.80}  % ... minus the market's, STANDARD
\providecommand{\rhebExRtIvPEnetFull}{}\renewcommand{\rhebExRtIvPEnetFull}{0.159}  % ... its Ledoit-Wolf p (D19)
\providecommand{\rhebRtMvShareEnetFull}{}\renewcommand{\rhebRtMvShareEnetFull}{90}  % ... real-time MV: median strategy vol share, %
\providecommand{\rhebRtMvCapDaysEnetFull}{}\renewcommand{\rhebRtMvCapDaysEnetFull}{100}  % ... % of days the 90% cap binds
\providecommand{\rhebRtMvWsEnetFull}{}\renewcommand{\rhebRtMvWsEnetFull}{3.5}  % ... real-time MV: median strategy dollar weight, %
\providecommand{\rhebRtIvWsEnetFull}{}\renewcommand{\rhebRtIvWsEnetFull}{0.9}  % ... real-time IV: median strategy dollar weight, %
\providecommand{\rhebExBetaEnetFull}{}\renewcommand{\rhebExBetaEnetFull}{2.84}  % ... without the best-day union: CAPM beta
\providecommand{\rhebExBetaTEnetFull}{}\renewcommand{\rhebExBetaTEnetFull}{0.79}  % ... its HAC t
\providecommand{\rhebExAlphaEnetFull}{}\renewcommand{\rhebExAlphaEnetFull}{8.3}  % ... without the best-day union: alpha, % per day
\providecommand{\rhebExAlphaTEnetFull}{}\renewcommand{\rhebExAlphaTEnetFull}{2.47}  % ... its HAC t
\providecommand{\rhebExAlphaPEnetFull}{}\renewcommand{\rhebExAlphaPEnetFull}{0.023}  % ... its fixed-b p
\providecommand{\rhebExAREnetFull}{}\renewcommand{\rhebExAREnetFull}{1.26}  % ... appraisal ratio, sqrt(q) (D19)
\providecommand{\rhebExSRmaxEnetFull}{}\renewcommand{\rhebExSRmaxEnetFull}{1.28}  % ... max Sharpe ratio, STANDARD
\providecommand{\rhebExDSRmaxPEnetFull}{}\renewcommand{\rhebExDSRmaxPEnetFull}{0.035}  % ... bootstrap p of its gain
\providecommand{\rhebExGrsPEnetFull}{}\renewcommand{\rhebExGrsPEnetFull}{0.038}  % ... GRS p
\providecommand{\rhebTmGammaEnetFull}{}\renewcommand{\rhebTmGammaEnetFull}{\ensuremath{-}2.19}  % poolpt13_enet_all_features, mid: Treynor-Mazuy gamma on (r^m)^2 (% per %^2)
\providecommand{\rhebTmGammaTEnetFull}{}\renewcommand{\rhebTmGammaTEnetFull}{\ensuremath{-}1.53}  % ... its HAC t
\providecommand{\rhebTmGammaPEnetFull}{}\renewcommand{\rhebTmGammaPEnetFull}{0.148}  % ... its fixed-b p
\providecommand{\rhebTmBetaEnetFull}{}\renewcommand{\rhebTmBetaEnetFull}{2.43}  % ... Treynor-Mazuy beta on r^m
\providecommand{\rhebTmBetaTEnetFull}{}\renewcommand{\rhebTmBetaTEnetFull}{0.60}  % ... its HAC t
\providecommand{\rhebNLongEnetFull}{}\renewcommand{\rhebNLongEnetFull}{117}  % ... long days
\providecommand{\rhebTmGammaLongEnetFull}{}\renewcommand{\rhebTmGammaLongEnetFull}{\ensuremath{-}0.05}  % ... TM gamma on its long days only
\providecommand{\rhebTmGammaLongTEnetFull}{}\renewcommand{\rhebTmGammaLongTEnetFull}{\ensuremath{-}0.01}  % ... its HAC t (no test: RULE section 4)
\providecommand{\rhebNShortEnetFull}{}\renewcommand{\rhebNShortEnetFull}{497}  % ... short days
\providecommand{\rhebTmGammaShortEnetFull}{}\renewcommand{\rhebTmGammaShortEnetFull}{\ensuremath{-}2.03}  % ... TM gamma on its short days only
\providecommand{\rhebTmGammaShortTEnetFull}{}\renewcommand{\rhebTmGammaShortTEnetFull}{\ensuremath{-}1.47}  % ... its HAC t (no test: RULE section 4)
\providecommand{\rhebDownMeanEnetFull}{}\renewcommand{\rhebDownMeanEnetFull}{15.8}  % ... mean return on market-down days, % of premium
\providecommand{\rhebDownRestEnetFull}{}\renewcommand{\rhebDownRestEnetFull}{12.0}  % ... on the other days
\providecommand{\rhebDownDiffTEnetFull}{}\renewcommand{\rhebDownDiffTEnetFull}{0.36}  % ... HAC t of the difference
\providecommand{\rhebDownDiffPEnetFull}{}\renewcommand{\rhebDownDiffPEnetFull}{0.732}  % ... its fixed-b p
\providecommand{\rhebWorstMeanEnetFull}{}\renewcommand{\rhebWorstMeanEnetFull}{\ensuremath{-}15.7}  % ... mean return on worst-decile market days, % of premium
\providecommand{\rhebWorstRestEnetFull}{}\renewcommand{\rhebWorstRestEnetFull}{17.2}  % ... on the other days
\providecommand{\rhebWorstDiffTEnetFull}{}\renewcommand{\rhebWorstDiffTEnetFull}{\ensuremath{-}2.28}  % ... HAC t of the difference
\providecommand{\rhebWorstDiffPEnetFull}{}\renewcommand{\rhebWorstDiffPEnetFull}{0.035}  % ... its fixed-b p
\providecommand{\rhebThirtyCorrEnetFull}{}\renewcommand{\rhebThirtyCorrEnetFull}{0.02}  % ... correlation with the 15:30-16:00 index return
\providecommand{\rhebThirtyBetaEnetFull}{}\renewcommand{\rhebThirtyBetaEnetFull}{9.56}  % ... beta on the 15:30-16:00 index return
\providecommand{\rhebThirtyBetaTEnetFull}{}\renewcommand{\rhebThirtyBetaTEnetFull}{0.32}  % ... its HAC t
\providecommand{\rhebThirtyBetaPEnetFull}{}\renewcommand{\rhebThirtyBetaPEnetFull}{0.758}  % ... its fixed-b p
\providecommand{\rhebThirtyAlphaEnetFull}{}\renewcommand{\rhebThirtyAlphaEnetFull}{14.0}  % ... alpha against it, % per day
\providecommand{\rhebThirtyAlphaTEnetFull}{}\renewcommand{\rhebThirtyAlphaTEnetFull}{3.36}  % ... its HAC t
\providecommand{\rhebThirtyTmGammaEnetFull}{}\renewcommand{\rhebThirtyTmGammaEnetFull}{\ensuremath{-}122.92}  % ... TM gamma on the squared 15:30-16:00 return
\providecommand{\rhebThirtyTmGammaTEnetFull}{}\renewcommand{\rhebThirtyTmGammaTEnetFull}{\ensuremath{-}2.41}  % ... its HAC t
\providecommand{\rhebThirtyTmGammaPEnetFull}{}\renewcommand{\rhebThirtyTmGammaPEnetFull}{0.026}  % ... its fixed-b p
\providecommand{\rhebThirtyTmGammaLongEnetFull}{}\renewcommand{\rhebThirtyTmGammaLongEnetFull}{355.62}  % ... ... on its long days only
\providecommand{\rhebThirtyTmGammaLongTEnetFull}{}\renewcommand{\rhebThirtyTmGammaLongTEnetFull}{2.64}  % ... its HAC t (no test)
\providecommand{\rhebThirtyTmGammaShortEnetFull}{}\renewcommand{\rhebThirtyTmGammaShortEnetFull}{\ensuremath{-}274.55}  % ... ... on its short days only
\providecommand{\rhebThirtyTmGammaShortTEnetFull}{}\renewcommand{\rhebThirtyTmGammaShortTEnetFull}{\ensuremath{-}4.40}  % ... its HAC t (no test)
\providecommand{\rhebThirtyDownMeanEnetFull}{}\renewcommand{\rhebThirtyDownMeanEnetFull}{11.3}  % ... mean on 15:30-16:00 down days, %
\providecommand{\rhebThirtyDownDiffTEnetFull}{}\renewcommand{\rhebThirtyDownDiffTEnetFull}{\ensuremath{-}0.58}  % ... HAC t of the difference with the other days
\providecommand{\rhebThirtyWorstMeanEnetFull}{}\renewcommand{\rhebThirtyWorstMeanEnetFull}{\ensuremath{-}51.7}  % ... mean on 15:30-16:00 worst days, %
\providecommand{\rhebThirtyWorstDiffTEnetFull}{}\renewcommand{\rhebThirtyWorstDiffTEnetFull}{\ensuremath{-}2.68}  % ... HAC t of the difference with the other days
\providecommand{\rhebCorrXEnetFull}{}\renewcommand{\rhebCorrXEnetFull}{0.03}  % poolpt13_enet_all_features, bid/ask: correlation with the market excess return
\providecommand{\rhebBetaXEnetFull}{}\renewcommand{\rhebBetaXEnetFull}{2.59}  % ... CAPM beta (% of premium per % of market)
\providecommand{\rhebBetaTXEnetFull}{}\renewcommand{\rhebBetaTXEnetFull}{0.66}  % ... its HAC t (S = 33)
\providecommand{\rhebBetaPXEnetFull}{}\renewcommand{\rhebBetaPXEnetFull}{0.530}  % ... its fixed-b p
\providecommand{\rhebRsqXEnetFull}{}\renewcommand{\rhebRsqXEnetFull}{0.00}  % ... CAPM R^2
\providecommand{\rhebAlphaXEnetFull}{}\renewcommand{\rhebAlphaXEnetFull}{11.2}  % ... CAPM alpha, % of premium per day
\providecommand{\rhebAlphaTXEnetFull}{}\renewcommand{\rhebAlphaTXEnetFull}{2.66}  % ... its HAC t
\providecommand{\rhebAlphaPXEnetFull}{}\renewcommand{\rhebAlphaPXEnetFull}{0.015}  % ... its fixed-b p
\providecommand{\rhebAlphaAnnXEnetFull}{}\renewcommand{\rhebAlphaAnnXEnetFull}{2534}  % ... alpha x q (D19; was x 252), % of premium a year
\providecommand{\rhebARXEnetFull}{}\renewcommand{\rhebARXEnetFull}{1.49}  % ... appraisal ratio alpha / sigma_eps, annualized sqrt(q) (D19)
\providecommand{\rhebSRsXEnetFull}{}\renewcommand{\rhebSRsXEnetFull}{1.50}  % ... strategy Sharpe ratio, STANDARD = the paper's (= \rhemSh / X)
\providecommand{\rhebSRsDSixXEnetFull}{}\renewcommand{\rhebSRsDSixXEnetFull}{1.59}  % ... strategy Sharpe ratio, D6 (= \rhemShDSix)
\providecommand{\rhebSRmaxXEnetFull}{}\renewcommand{\rhebSRmaxXEnetFull}{1.50}  % ... max Sharpe ratio of market + strategy, STANDARD
\providecommand{\rhebDSRmaxXEnetFull}{}\renewcommand{\rhebDSRmaxXEnetFull}{1.35}  % ... its gain over the market alone, STANDARD
\providecommand{\rhebDSRmaxLoXEnetFull}{}\renewcommand{\rhebDSRmaxLoXEnetFull}{0.12}  % ... bootstrap 95% interval, lower end
\providecommand{\rhebDSRmaxHiXEnetFull}{}\renewcommand{\rhebDSRmaxHiXEnetFull}{2.34}  % ... upper end
\providecommand{\rhebDSRmaxPXEnetFull}{}\renewcommand{\rhebDSRmaxPXEnetFull}{0.007}  % ... bootstrap p (descriptive; Delta >= 0)
\providecommand{\rhebGrsXEnetFull}{}\renewcommand{\rhebGrsXEnetFull}{6.06}  % ... one-asset GRS statistic, F(1, T - 2)
\providecommand{\rhebGrsPXEnetFull}{}\renewcommand{\rhebGrsPXEnetFull}{0.014}  % ... its p (i.i.d. errors)
\providecommand{\rhebTanShareXEnetFull}{}\renewcommand{\rhebTanShareXEnetFull}{99}  % ... tangency: strategy's share of volatility, %
\providecommand{\rhebTanDSixXEnetFull}{}\renewcommand{\rhebTanDSixXEnetFull}{1.60}  % ... tangency (fixed in-sample weights): Sharpe ratio, D6
\providecommand{\rhebRtMvXEnetFull}{}\renewcommand{\rhebRtMvXEnetFull}{1.59}  % ... real-time mean-variance: Sharpe ratio, STANDARD
\providecommand{\rhebRtMvDXEnetFull}{}\renewcommand{\rhebRtMvDXEnetFull}{1.10}  % ... minus the market's, STANDARD
\providecommand{\rhebRtMvPXEnetFull}{}\renewcommand{\rhebRtMvPXEnetFull}{0.189}  % ... Ledoit-Wolf p of the difference (D19)
\providecommand{\rhebRtMvDSixXEnetFull}{}\renewcommand{\rhebRtMvDSixXEnetFull}{1.69}  % ... real-time mean-variance: Sharpe ratio, D6
\providecommand{\rhebRtMvDSixDXEnetFull}{}\renewcommand{\rhebRtMvDSixDXEnetFull}{1.00}  % ... minus the market's, D6
\providecommand{\rhebRtMvDSixPXEnetFull}{}\renewcommand{\rhebRtMvDSixPXEnetFull}{0.286}  % ... bootstrap p of the D6 difference
\providecommand{\rhebExRtMvXEnetFull}{}\renewcommand{\rhebExRtMvXEnetFull}{1.02}  % ... w/o best days, real-time mean-variance: SR, STANDARD
\providecommand{\rhebExRtMvDXEnetFull}{}\renewcommand{\rhebExRtMvDXEnetFull}{0.48}  % ... minus the market's, STANDARD
\providecommand{\rhebExRtMvPXEnetFull}{}\renewcommand{\rhebExRtMvPXEnetFull}{0.568}  % ... its Ledoit-Wolf p (D19)
\providecommand{\rhebRtIvXEnetFull}{}\renewcommand{\rhebRtIvXEnetFull}{1.33}  % ... real-time inverse-volatility: Sharpe ratio, STANDARD
\providecommand{\rhebRtIvDXEnetFull}{}\renewcommand{\rhebRtIvDXEnetFull}{0.83}  % ... minus the market's, STANDARD
\providecommand{\rhebRtIvPXEnetFull}{}\renewcommand{\rhebRtIvPXEnetFull}{0.142}  % ... Ledoit-Wolf p of the difference (D19)
\providecommand{\rhebRtIvDSixXEnetFull}{}\renewcommand{\rhebRtIvDSixXEnetFull}{1.67}  % ... real-time inverse-volatility: Sharpe ratio, D6
\providecommand{\rhebRtIvDSixDXEnetFull}{}\renewcommand{\rhebRtIvDSixDXEnetFull}{0.99}  % ... minus the market's, D6
\providecommand{\rhebRtIvDSixPXEnetFull}{}\renewcommand{\rhebRtIvDSixPXEnetFull}{0.123}  % ... bootstrap p of the D6 difference
\providecommand{\rhebExRtIvXEnetFull}{}\renewcommand{\rhebExRtIvXEnetFull}{1.06}  % ... w/o best days, real-time inverse-volatility: SR, STANDARD
\providecommand{\rhebExRtIvDXEnetFull}{}\renewcommand{\rhebExRtIvDXEnetFull}{0.52}  % ... minus the market's, STANDARD
\providecommand{\rhebExRtIvPXEnetFull}{}\renewcommand{\rhebExRtIvPXEnetFull}{0.367}  % ... its Ledoit-Wolf p (D19)
\providecommand{\rhebRtMvShareXEnetFull}{}\renewcommand{\rhebRtMvShareXEnetFull}{90}  % ... real-time MV: median strategy vol share, %
\providecommand{\rhebRtMvCapDaysXEnetFull}{}\renewcommand{\rhebRtMvCapDaysXEnetFull}{100}  % ... % of days the 90% cap binds
\providecommand{\rhebRtMvWsXEnetFull}{}\renewcommand{\rhebRtMvWsXEnetFull}{3.5}  % ... real-time MV: median strategy dollar weight, %
\providecommand{\rhebRtIvWsXEnetFull}{}\renewcommand{\rhebRtIvWsXEnetFull}{0.8}  % ... real-time IV: median strategy dollar weight, %
\providecommand{\rhebExBetaXEnetFull}{}\renewcommand{\rhebExBetaXEnetFull}{2.86}  % ... without the best-day union: CAPM beta
\providecommand{\rhebExBetaTXEnetFull}{}\renewcommand{\rhebExBetaTXEnetFull}{0.78}  % ... its HAC t
\providecommand{\rhebExAlphaXEnetFull}{}\renewcommand{\rhebExAlphaXEnetFull}{5.8}  % ... without the best-day union: alpha, % per day
\providecommand{\rhebExAlphaTXEnetFull}{}\renewcommand{\rhebExAlphaTXEnetFull}{1.72}  % ... its HAC t
\providecommand{\rhebExAlphaPXEnetFull}{}\renewcommand{\rhebExAlphaPXEnetFull}{0.106}  % ... its fixed-b p
\providecommand{\rhebExARXEnetFull}{}\renewcommand{\rhebExARXEnetFull}{0.88}  % ... appraisal ratio, sqrt(q) (D19)
\providecommand{\rhebExSRmaxXEnetFull}{}\renewcommand{\rhebExSRmaxXEnetFull}{0.91}  % ... max Sharpe ratio, STANDARD
\providecommand{\rhebExDSRmaxPXEnetFull}{}\renewcommand{\rhebExDSRmaxPXEnetFull}{0.095}  % ... bootstrap p of its gain
\providecommand{\rhebExGrsPXEnetFull}{}\renewcommand{\rhebExGrsPXEnetFull}{0.149}  % ... GRS p
\providecommand{\rhebCorrLgbmFull}{}\renewcommand{\rhebCorrLgbmFull}{0.03}  % pooltree_lgbm_all_features, mid: correlation with the market excess return
\providecommand{\rhebBetaLgbmFull}{}\renewcommand{\rhebBetaLgbmFull}{2.92}  % ... CAPM beta (% of premium per % of market)
\providecommand{\rhebBetaTLgbmFull}{}\renewcommand{\rhebBetaTLgbmFull}{0.75}  % ... its HAC t (S = 33)
\providecommand{\rhebBetaPLgbmFull}{}\renewcommand{\rhebBetaPLgbmFull}{0.472}  % ... its fixed-b p
\providecommand{\rhebRsqLgbmFull}{}\renewcommand{\rhebRsqLgbmFull}{0.00}  % ... CAPM R^2
\providecommand{\rhebAlphaLgbmFull}{}\renewcommand{\rhebAlphaLgbmFull}{11.5}  % ... CAPM alpha, % of premium per day
\providecommand{\rhebAlphaTLgbmFull}{}\renewcommand{\rhebAlphaTLgbmFull}{3.29}  % ... its HAC t
\providecommand{\rhebAlphaPLgbmFull}{}\renewcommand{\rhebAlphaPLgbmFull}{0.003}  % ... its fixed-b p
\providecommand{\rhebAlphaAnnLgbmFull}{}\renewcommand{\rhebAlphaAnnLgbmFull}{2595}  % ... alpha x q (D19; was x 252), % of premium a year
\providecommand{\rhebARLgbmFull}{}\renewcommand{\rhebARLgbmFull}{1.53}  % ... appraisal ratio alpha / sigma_eps, annualized sqrt(q) (D19)
\providecommand{\rhebSRsLgbmFull}{}\renewcommand{\rhebSRsLgbmFull}{1.54}  % ... strategy Sharpe ratio, STANDARD = the paper's (= \rhemSh / X)
\providecommand{\rhebSRsDSixLgbmFull}{}\renewcommand{\rhebSRsDSixLgbmFull}{1.67}  % ... strategy Sharpe ratio, D6 (= \rhemShDSix)
\providecommand{\rhebSRmaxLgbmFull}{}\renewcommand{\rhebSRmaxLgbmFull}{1.54}  % ... max Sharpe ratio of market + strategy, STANDARD
\providecommand{\rhebDSRmaxLgbmFull}{}\renewcommand{\rhebDSRmaxLgbmFull}{1.39}  % ... its gain over the market alone, STANDARD
\providecommand{\rhebDSRmaxLoLgbmFull}{}\renewcommand{\rhebDSRmaxLoLgbmFull}{0.18}  % ... bootstrap 95% interval, lower end
\providecommand{\rhebDSRmaxHiLgbmFull}{}\renewcommand{\rhebDSRmaxHiLgbmFull}{2.27}  % ... upper end
\providecommand{\rhebDSRmaxPLgbmFull}{}\renewcommand{\rhebDSRmaxPLgbmFull}{0.003}  % ... bootstrap p (descriptive; Delta >= 0)
\providecommand{\rhebGrsLgbmFull}{}\renewcommand{\rhebGrsLgbmFull}{6.38}  % ... one-asset GRS statistic, F(1, T - 2)
\providecommand{\rhebGrsPLgbmFull}{}\renewcommand{\rhebGrsPLgbmFull}{0.012}  % ... its p (i.i.d. errors)
\providecommand{\rhebTanShareLgbmFull}{}\renewcommand{\rhebTanShareLgbmFull}{99}  % ... tangency: strategy's share of volatility, %
\providecommand{\rhebTanDSixLgbmFull}{}\renewcommand{\rhebTanDSixLgbmFull}{1.67}  % ... tangency (fixed in-sample weights): Sharpe ratio, D6
\providecommand{\rhebRtMvLgbmFull}{}\renewcommand{\rhebRtMvLgbmFull}{1.67}  % ... real-time mean-variance: Sharpe ratio, STANDARD
\providecommand{\rhebRtMvDLgbmFull}{}\renewcommand{\rhebRtMvDLgbmFull}{1.17}  % ... minus the market's, STANDARD
\providecommand{\rhebRtMvPLgbmFull}{}\renewcommand{\rhebRtMvPLgbmFull}{0.125}  % ... Ledoit-Wolf p of the difference (D19)
\providecommand{\rhebRtMvDSixLgbmFull}{}\renewcommand{\rhebRtMvDSixLgbmFull}{1.81}  % ... real-time mean-variance: Sharpe ratio, D6
\providecommand{\rhebRtMvDSixDLgbmFull}{}\renewcommand{\rhebRtMvDSixDLgbmFull}{1.12}  % ... minus the market's, D6
\providecommand{\rhebRtMvDSixPLgbmFull}{}\renewcommand{\rhebRtMvDSixPLgbmFull}{0.197}  % ... bootstrap p of the D6 difference
\providecommand{\rhebExRtMvLgbmFull}{}\renewcommand{\rhebExRtMvLgbmFull}{1.06}  % ... w/o best days, real-time mean-variance: SR, STANDARD
\providecommand{\rhebExRtMvDLgbmFull}{}\renewcommand{\rhebExRtMvDLgbmFull}{0.52}  % ... minus the market's, STANDARD
\providecommand{\rhebExRtMvPLgbmFull}{}\renewcommand{\rhebExRtMvPLgbmFull}{0.511}  % ... its Ledoit-Wolf p (D19)
\providecommand{\rhebRtIvLgbmFull}{}\renewcommand{\rhebRtIvLgbmFull}{1.29}  % ... real-time inverse-volatility: Sharpe ratio, STANDARD
\providecommand{\rhebRtIvDLgbmFull}{}\renewcommand{\rhebRtIvDLgbmFull}{0.79}  % ... minus the market's, STANDARD
\providecommand{\rhebRtIvPLgbmFull}{}\renewcommand{\rhebRtIvPLgbmFull}{0.118}  % ... Ledoit-Wolf p of the difference (D19)
\providecommand{\rhebRtIvDSixLgbmFull}{}\renewcommand{\rhebRtIvDSixLgbmFull}{1.72}  % ... real-time inverse-volatility: Sharpe ratio, D6
\providecommand{\rhebRtIvDSixDLgbmFull}{}\renewcommand{\rhebRtIvDSixDLgbmFull}{1.04}  % ... minus the market's, D6
\providecommand{\rhebRtIvDSixPLgbmFull}{}\renewcommand{\rhebRtIvDSixPLgbmFull}{0.075}  % ... bootstrap p of the D6 difference
\providecommand{\rhebExRtIvLgbmFull}{}\renewcommand{\rhebExRtIvLgbmFull}{1.03}  % ... w/o best days, real-time inverse-volatility: SR, STANDARD
\providecommand{\rhebExRtIvDLgbmFull}{}\renewcommand{\rhebExRtIvDLgbmFull}{0.48}  % ... minus the market's, STANDARD
\providecommand{\rhebExRtIvPLgbmFull}{}\renewcommand{\rhebExRtIvPLgbmFull}{0.349}  % ... its Ledoit-Wolf p (D19)
\providecommand{\rhebRtMvShareLgbmFull}{}\renewcommand{\rhebRtMvShareLgbmFull}{90}  % ... real-time MV: median strategy vol share, %
\providecommand{\rhebRtMvCapDaysLgbmFull}{}\renewcommand{\rhebRtMvCapDaysLgbmFull}{100}  % ... % of days the 90% cap binds
\providecommand{\rhebRtMvWsLgbmFull}{}\renewcommand{\rhebRtMvWsLgbmFull}{3.5}  % ... real-time MV: median strategy dollar weight, %
\providecommand{\rhebRtIvWsLgbmFull}{}\renewcommand{\rhebRtIvWsLgbmFull}{0.9}  % ... real-time IV: median strategy dollar weight, %
\providecommand{\rhebExBetaLgbmFull}{}\renewcommand{\rhebExBetaLgbmFull}{3.26}  % ... without the best-day union: CAPM beta
\providecommand{\rhebExBetaTLgbmFull}{}\renewcommand{\rhebExBetaTLgbmFull}{0.93}  % ... its HAC t
\providecommand{\rhebExAlphaLgbmFull}{}\renewcommand{\rhebExAlphaLgbmFull}{5.9}  % ... without the best-day union: alpha, % per day
\providecommand{\rhebExAlphaTLgbmFull}{}\renewcommand{\rhebExAlphaTLgbmFull}{2.34}  % ... its HAC t
\providecommand{\rhebExAlphaPLgbmFull}{}\renewcommand{\rhebExAlphaPLgbmFull}{0.030}  % ... its fixed-b p
\providecommand{\rhebExARLgbmFull}{}\renewcommand{\rhebExARLgbmFull}{0.89}  % ... appraisal ratio, sqrt(q) (D19)
\providecommand{\rhebExSRmaxLgbmFull}{}\renewcommand{\rhebExSRmaxLgbmFull}{0.92}  % ... max Sharpe ratio, STANDARD
\providecommand{\rhebExDSRmaxPLgbmFull}{}\renewcommand{\rhebExDSRmaxPLgbmFull}{0.076}  % ... bootstrap p of its gain
\providecommand{\rhebExGrsPLgbmFull}{}\renewcommand{\rhebExGrsPLgbmFull}{0.141}  % ... GRS p
\providecommand{\rhebTmGammaLgbmFull}{}\renewcommand{\rhebTmGammaLgbmFull}{\ensuremath{-}1.67}  % pooltree_lgbm_all_features, mid: Treynor-Mazuy gamma on (r^m)^2 (% per %^2)
\providecommand{\rhebTmGammaTLgbmFull}{}\renewcommand{\rhebTmGammaTLgbmFull}{\ensuremath{-}1.28}  % ... its HAC t
\providecommand{\rhebTmGammaPLgbmFull}{}\renewcommand{\rhebTmGammaPLgbmFull}{0.225}  % ... its fixed-b p
\providecommand{\rhebTmBetaLgbmFull}{}\renewcommand{\rhebTmBetaLgbmFull}{2.85}  % ... Treynor-Mazuy beta on r^m
\providecommand{\rhebTmBetaTLgbmFull}{}\renewcommand{\rhebTmBetaTLgbmFull}{0.71}  % ... its HAC t
\providecommand{\rhebNLongLgbmFull}{}\renewcommand{\rhebNLongLgbmFull}{146}  % ... long days
\providecommand{\rhebTmGammaLongLgbmFull}{}\renewcommand{\rhebTmGammaLongLgbmFull}{4.07}  % ... TM gamma on its long days only
\providecommand{\rhebTmGammaLongTLgbmFull}{}\renewcommand{\rhebTmGammaLongTLgbmFull}{0.56}  % ... its HAC t (no test: RULE section 4)
\providecommand{\rhebNShortLgbmFull}{}\renewcommand{\rhebNShortLgbmFull}{468}  % ... short days
\providecommand{\rhebTmGammaShortLgbmFull}{}\renewcommand{\rhebTmGammaShortLgbmFull}{\ensuremath{-}1.65}  % ... TM gamma on its short days only
\providecommand{\rhebTmGammaShortTLgbmFull}{}\renewcommand{\rhebTmGammaShortTLgbmFull}{\ensuremath{-}1.28}  % ... its HAC t (no test: RULE section 4)
\providecommand{\rhebDownMeanLgbmFull}{}\renewcommand{\rhebDownMeanLgbmFull}{13.1}  % ... mean return on market-down days, % of premium
\providecommand{\rhebDownRestLgbmFull}{}\renewcommand{\rhebDownRestLgbmFull}{10.0}  % ... on the other days
\providecommand{\rhebDownDiffTLgbmFull}{}\renewcommand{\rhebDownDiffTLgbmFull}{0.27}  % ... HAC t of the difference
\providecommand{\rhebDownDiffPLgbmFull}{}\renewcommand{\rhebDownDiffPLgbmFull}{0.798}  % ... its fixed-b p
\providecommand{\rhebWorstMeanLgbmFull}{}\renewcommand{\rhebWorstMeanLgbmFull}{\ensuremath{-}14.7}  % ... mean return on worst-decile market days, % of premium
\providecommand{\rhebWorstRestLgbmFull}{}\renewcommand{\rhebWorstRestLgbmFull}{14.5}  % ... on the other days
\providecommand{\rhebWorstDiffTLgbmFull}{}\renewcommand{\rhebWorstDiffTLgbmFull}{\ensuremath{-}2.16}  % ... HAC t of the difference
\providecommand{\rhebWorstDiffPLgbmFull}{}\renewcommand{\rhebWorstDiffPLgbmFull}{0.045}  % ... its fixed-b p
\providecommand{\rhebThirtyCorrLgbmFull}{}\renewcommand{\rhebThirtyCorrLgbmFull}{0.06}  % ... correlation with the 15:30-16:00 index return
\providecommand{\rhebThirtyBetaLgbmFull}{}\renewcommand{\rhebThirtyBetaLgbmFull}{23.20}  % ... beta on the 15:30-16:00 index return
\providecommand{\rhebThirtyBetaTLgbmFull}{}\renewcommand{\rhebThirtyBetaTLgbmFull}{0.77}  % ... its HAC t
\providecommand{\rhebThirtyBetaPLgbmFull}{}\renewcommand{\rhebThirtyBetaPLgbmFull}{0.462}  % ... its fixed-b p
\providecommand{\rhebThirtyAlphaLgbmFull}{}\renewcommand{\rhebThirtyAlphaLgbmFull}{11.8}  % ... alpha against it, % per day
\providecommand{\rhebThirtyAlphaTLgbmFull}{}\renewcommand{\rhebThirtyAlphaTLgbmFull}{3.43}  % ... its HAC t
\providecommand{\rhebThirtyTmGammaLgbmFull}{}\renewcommand{\rhebThirtyTmGammaLgbmFull}{\ensuremath{-}112.12}  % ... TM gamma on the squared 15:30-16:00 return
\providecommand{\rhebThirtyTmGammaTLgbmFull}{}\renewcommand{\rhebThirtyTmGammaTLgbmFull}{\ensuremath{-}2.29}  % ... its HAC t
\providecommand{\rhebThirtyTmGammaPLgbmFull}{}\renewcommand{\rhebThirtyTmGammaPLgbmFull}{0.034}  % ... its fixed-b p
\providecommand{\rhebThirtyTmGammaLongLgbmFull}{}\renewcommand{\rhebThirtyTmGammaLongLgbmFull}{385.02}  % ... ... on its long days only
\providecommand{\rhebThirtyTmGammaLongTLgbmFull}{}\renewcommand{\rhebThirtyTmGammaLongTLgbmFull}{2.55}  % ... its HAC t (no test)
\providecommand{\rhebThirtyTmGammaShortLgbmFull}{}\renewcommand{\rhebThirtyTmGammaShortLgbmFull}{\ensuremath{-}271.45}  % ... ... on its short days only
\providecommand{\rhebThirtyTmGammaShortTLgbmFull}{}\renewcommand{\rhebThirtyTmGammaShortTLgbmFull}{\ensuremath{-}4.36}  % ... its HAC t (no test)
\providecommand{\rhebThirtyDownMeanLgbmFull}{}\renewcommand{\rhebThirtyDownMeanLgbmFull}{8.4}  % ... mean on 15:30-16:00 down days, %
\providecommand{\rhebThirtyDownDiffTLgbmFull}{}\renewcommand{\rhebThirtyDownDiffTLgbmFull}{\ensuremath{-}0.75}  % ... HAC t of the difference with the other days
\providecommand{\rhebThirtyWorstMeanLgbmFull}{}\renewcommand{\rhebThirtyWorstMeanLgbmFull}{\ensuremath{-}55.2}  % ... mean on 15:30-16:00 worst days, %
\providecommand{\rhebThirtyWorstDiffTLgbmFull}{}\renewcommand{\rhebThirtyWorstDiffTLgbmFull}{\ensuremath{-}2.67}  % ... HAC t of the difference with the other days
\providecommand{\rhebCorrXLgbmFull}{}\renewcommand{\rhebCorrXLgbmFull}{0.03}  % pooltree_lgbm_all_features, bid/ask: correlation with the market excess return
\providecommand{\rhebBetaXLgbmFull}{}\renewcommand{\rhebBetaXLgbmFull}{2.99}  % ... CAPM beta (% of premium per % of market)
\providecommand{\rhebBetaTXLgbmFull}{}\renewcommand{\rhebBetaTXLgbmFull}{0.76}  % ... its HAC t (S = 33)
\providecommand{\rhebBetaPXLgbmFull}{}\renewcommand{\rhebBetaPXLgbmFull}{0.467}  % ... its fixed-b p
\providecommand{\rhebRsqXLgbmFull}{}\renewcommand{\rhebRsqXLgbmFull}{0.00}  % ... CAPM R^2
\providecommand{\rhebAlphaXLgbmFull}{}\renewcommand{\rhebAlphaXLgbmFull}{8.9}  % ... CAPM alpha, % of premium per day
\providecommand{\rhebAlphaTXLgbmFull}{}\renewcommand{\rhebAlphaTXLgbmFull}{2.56}  % ... its HAC t
\providecommand{\rhebAlphaPXLgbmFull}{}\renewcommand{\rhebAlphaPXLgbmFull}{0.019}  % ... its fixed-b p
\providecommand{\rhebAlphaAnnXLgbmFull}{}\renewcommand{\rhebAlphaAnnXLgbmFull}{2002}  % ... alpha x q (D19; was x 252), % of premium a year
\providecommand{\rhebARXLgbmFull}{}\renewcommand{\rhebARXLgbmFull}{1.18}  % ... appraisal ratio alpha / sigma_eps, annualized sqrt(q) (D19)
\providecommand{\rhebSRsXLgbmFull}{}\renewcommand{\rhebSRsXLgbmFull}{1.18}  % ... strategy Sharpe ratio, STANDARD = the paper's (= \rhemSh / X)
\providecommand{\rhebSRsDSixXLgbmFull}{}\renewcommand{\rhebSRsDSixXLgbmFull}{1.31}  % ... strategy Sharpe ratio, D6 (= \rhemShDSix)
\providecommand{\rhebSRmaxXLgbmFull}{}\renewcommand{\rhebSRmaxXLgbmFull}{1.19}  % ... max Sharpe ratio of market + strategy, STANDARD
\providecommand{\rhebDSRmaxXLgbmFull}{}\renewcommand{\rhebDSRmaxXLgbmFull}{1.04}  % ... its gain over the market alone, STANDARD
\providecommand{\rhebDSRmaxLoXLgbmFull}{}\renewcommand{\rhebDSRmaxLoXLgbmFull}{0.03}  % ... bootstrap 95% interval, lower end
\providecommand{\rhebDSRmaxHiXLgbmFull}{}\renewcommand{\rhebDSRmaxHiXLgbmFull}{1.94}  % ... upper end
\providecommand{\rhebDSRmaxPXLgbmFull}{}\renewcommand{\rhebDSRmaxPXLgbmFull}{0.014}  % ... bootstrap p (descriptive; Delta >= 0)
\providecommand{\rhebGrsXLgbmFull}{}\renewcommand{\rhebGrsXLgbmFull}{3.77}  % ... one-asset GRS statistic, F(1, T - 2)
\providecommand{\rhebGrsPXLgbmFull}{}\renewcommand{\rhebGrsPXLgbmFull}{0.053}  % ... its p (i.i.d. errors)
\providecommand{\rhebTanShareXLgbmFull}{}\renewcommand{\rhebTanShareXLgbmFull}{99}  % ... tangency: strategy's share of volatility, %
\providecommand{\rhebTanDSixXLgbmFull}{}\renewcommand{\rhebTanDSixXLgbmFull}{1.32}  % ... tangency (fixed in-sample weights): Sharpe ratio, D6
\providecommand{\rhebRtMvXLgbmFull}{}\renewcommand{\rhebRtMvXLgbmFull}{1.32}  % ... real-time mean-variance: Sharpe ratio, STANDARD
\providecommand{\rhebRtMvDXLgbmFull}{}\renewcommand{\rhebRtMvDXLgbmFull}{0.82}  % ... minus the market's, STANDARD
\providecommand{\rhebRtMvPXLgbmFull}{}\renewcommand{\rhebRtMvPXLgbmFull}{0.279}  % ... Ledoit-Wolf p of the difference (D19)
\providecommand{\rhebRtMvDSixXLgbmFull}{}\renewcommand{\rhebRtMvDSixXLgbmFull}{1.46}  % ... real-time mean-variance: Sharpe ratio, D6
\providecommand{\rhebRtMvDSixDXLgbmFull}{}\renewcommand{\rhebRtMvDSixDXLgbmFull}{0.78}  % ... minus the market's, D6
\providecommand{\rhebRtMvDSixPXLgbmFull}{}\renewcommand{\rhebRtMvDSixPXLgbmFull}{0.363}  % ... bootstrap p of the D6 difference
\providecommand{\rhebExRtMvXLgbmFull}{}\renewcommand{\rhebExRtMvXLgbmFull}{0.71}  % ... w/o best days, real-time mean-variance: SR, STANDARD
\providecommand{\rhebExRtMvDXLgbmFull}{}\renewcommand{\rhebExRtMvDXLgbmFull}{0.16}  % ... minus the market's, STANDARD
\providecommand{\rhebExRtMvPXLgbmFull}{}\renewcommand{\rhebExRtMvPXLgbmFull}{0.826}  % ... its Ledoit-Wolf p (D19)
\providecommand{\rhebRtIvXLgbmFull}{}\renewcommand{\rhebRtIvXLgbmFull}{1.02}  % ... real-time inverse-volatility: Sharpe ratio, STANDARD
\providecommand{\rhebRtIvDXLgbmFull}{}\renewcommand{\rhebRtIvDXLgbmFull}{0.52}  % ... minus the market's, STANDARD
\providecommand{\rhebRtIvPXLgbmFull}{}\renewcommand{\rhebRtIvPXLgbmFull}{0.295}  % ... Ledoit-Wolf p of the difference (D19)
\providecommand{\rhebRtIvDSixXLgbmFull}{}\renewcommand{\rhebRtIvDSixXLgbmFull}{1.45}  % ... real-time inverse-volatility: Sharpe ratio, D6
\providecommand{\rhebRtIvDSixDXLgbmFull}{}\renewcommand{\rhebRtIvDSixDXLgbmFull}{0.77}  % ... minus the market's, D6
\providecommand{\rhebRtIvDSixPXLgbmFull}{}\renewcommand{\rhebRtIvDSixPXLgbmFull}{0.181}  % ... bootstrap p of the D6 difference
\providecommand{\rhebExRtIvXLgbmFull}{}\renewcommand{\rhebExRtIvXLgbmFull}{0.75}  % ... w/o best days, real-time inverse-volatility: SR, STANDARD
\providecommand{\rhebExRtIvDXLgbmFull}{}\renewcommand{\rhebExRtIvDXLgbmFull}{0.21}  % ... minus the market's, STANDARD
\providecommand{\rhebExRtIvPXLgbmFull}{}\renewcommand{\rhebExRtIvPXLgbmFull}{0.682}  % ... its Ledoit-Wolf p (D19)
\providecommand{\rhebRtMvShareXLgbmFull}{}\renewcommand{\rhebRtMvShareXLgbmFull}{90}  % ... real-time MV: median strategy vol share, %
\providecommand{\rhebRtMvCapDaysXLgbmFull}{}\renewcommand{\rhebRtMvCapDaysXLgbmFull}{100}  % ... % of days the 90% cap binds
\providecommand{\rhebRtMvWsXLgbmFull}{}\renewcommand{\rhebRtMvWsXLgbmFull}{3.5}  % ... real-time MV: median strategy dollar weight, %
\providecommand{\rhebRtIvWsXLgbmFull}{}\renewcommand{\rhebRtIvWsXLgbmFull}{0.9}  % ... real-time IV: median strategy dollar weight, %
\providecommand{\rhebExBetaXLgbmFull}{}\renewcommand{\rhebExBetaXLgbmFull}{3.27}  % ... without the best-day union: CAPM beta
\providecommand{\rhebExBetaTXLgbmFull}{}\renewcommand{\rhebExBetaTXLgbmFull}{0.92}  % ... its HAC t
\providecommand{\rhebExAlphaXLgbmFull}{}\renewcommand{\rhebExAlphaXLgbmFull}{3.4}  % ... without the best-day union: alpha, % per day
\providecommand{\rhebExAlphaTXLgbmFull}{}\renewcommand{\rhebExAlphaTXLgbmFull}{1.34}  % ... its HAC t
\providecommand{\rhebExAlphaPXLgbmFull}{}\renewcommand{\rhebExAlphaPXLgbmFull}{0.203}  % ... its fixed-b p
\providecommand{\rhebExARXLgbmFull}{}\renewcommand{\rhebExARXLgbmFull}{0.52}  % ... appraisal ratio, sqrt(q) (D19)
\providecommand{\rhebExSRmaxXLgbmFull}{}\renewcommand{\rhebExSRmaxXLgbmFull}{0.57}  % ... max Sharpe ratio, STANDARD
\providecommand{\rhebExDSRmaxPXLgbmFull}{}\renewcommand{\rhebExDSRmaxPXLgbmFull}{0.202}  % ... bootstrap p of its gain
\providecommand{\rhebExGrsPXLgbmFull}{}\renewcommand{\rhebExGrsPXLgbmFull}{0.396}  % ... GRS p
\providecommand{\rhebCorrXgbFull}{}\renewcommand{\rhebCorrXgbFull}{0.01}  % pooltree_xgb_all_features, mid: correlation with the market excess return
\providecommand{\rhebBetaXgbFull}{}\renewcommand{\rhebBetaXgbFull}{1.45}  % ... CAPM beta (% of premium per % of market)
\providecommand{\rhebBetaTXgbFull}{}\renewcommand{\rhebBetaTXgbFull}{0.38}  % ... its HAC t (S = 33)
\providecommand{\rhebBetaPXgbFull}{}\renewcommand{\rhebBetaPXgbFull}{0.714}  % ... its fixed-b p
\providecommand{\rhebRsqXgbFull}{}\renewcommand{\rhebRsqXgbFull}{0.00}  % ... CAPM R^2
\providecommand{\rhebAlphaXgbFull}{}\renewcommand{\rhebAlphaXgbFull}{11.6}  % ... CAPM alpha, % of premium per day
\providecommand{\rhebAlphaTXgbFull}{}\renewcommand{\rhebAlphaTXgbFull}{3.23}  % ... its HAC t
\providecommand{\rhebAlphaPXgbFull}{}\renewcommand{\rhebAlphaPXgbFull}{0.004}  % ... its fixed-b p
\providecommand{\rhebAlphaAnnXgbFull}{}\renewcommand{\rhebAlphaAnnXgbFull}{2612}  % ... alpha x q (D19; was x 252), % of premium a year
\providecommand{\rhebARXgbFull}{}\renewcommand{\rhebARXgbFull}{1.54}  % ... appraisal ratio alpha / sigma_eps, annualized sqrt(q) (D19)
\providecommand{\rhebSRsXgbFull}{}\renewcommand{\rhebSRsXgbFull}{1.54}  % ... strategy Sharpe ratio, STANDARD = the paper's (= \rhemSh / X)
\providecommand{\rhebSRsDSixXgbFull}{}\renewcommand{\rhebSRsDSixXgbFull}{1.71}  % ... strategy Sharpe ratio, D6 (= \rhemShDSix)
\providecommand{\rhebSRmaxXgbFull}{}\renewcommand{\rhebSRmaxXgbFull}{1.55}  % ... max Sharpe ratio of market + strategy, STANDARD
\providecommand{\rhebDSRmaxXgbFull}{}\renewcommand{\rhebDSRmaxXgbFull}{1.40}  % ... its gain over the market alone, STANDARD
\providecommand{\rhebDSRmaxLoXgbFull}{}\renewcommand{\rhebDSRmaxLoXgbFull}{0.19}  % ... bootstrap 95% interval, lower end
\providecommand{\rhebDSRmaxHiXgbFull}{}\renewcommand{\rhebDSRmaxHiXgbFull}{2.27}  % ... upper end
\providecommand{\rhebDSRmaxPXgbFull}{}\renewcommand{\rhebDSRmaxPXgbFull}{0.002}  % ... bootstrap p (descriptive; Delta >= 0)
\providecommand{\rhebGrsXgbFull}{}\renewcommand{\rhebGrsXgbFull}{6.46}  % ... one-asset GRS statistic, F(1, T - 2)
\providecommand{\rhebGrsPXgbFull}{}\renewcommand{\rhebGrsPXgbFull}{0.011}  % ... its p (i.i.d. errors)
\providecommand{\rhebTanShareXgbFull}{}\renewcommand{\rhebTanShareXgbFull}{99}  % ... tangency: strategy's share of volatility, %
\providecommand{\rhebTanDSixXgbFull}{}\renewcommand{\rhebTanDSixXgbFull}{1.72}  % ... tangency (fixed in-sample weights): Sharpe ratio, D6
\providecommand{\rhebRtMvXgbFull}{}\renewcommand{\rhebRtMvXgbFull}{1.66}  % ... real-time mean-variance: Sharpe ratio, STANDARD
\providecommand{\rhebRtMvDXgbFull}{}\renewcommand{\rhebRtMvDXgbFull}{1.16}  % ... minus the market's, STANDARD
\providecommand{\rhebRtMvPXgbFull}{}\renewcommand{\rhebRtMvPXgbFull}{0.134}  % ... Ledoit-Wolf p of the difference (D19)
\providecommand{\rhebRtMvDSixXgbFull}{}\renewcommand{\rhebRtMvDSixXgbFull}{1.83}  % ... real-time mean-variance: Sharpe ratio, D6
\providecommand{\rhebRtMvDSixDXgbFull}{}\renewcommand{\rhebRtMvDSixDXgbFull}{1.14}  % ... minus the market's, D6
\providecommand{\rhebRtMvDSixPXgbFull}{}\renewcommand{\rhebRtMvDSixPXgbFull}{0.216}  % ... bootstrap p of the D6 difference
\providecommand{\rhebExRtMvXgbFull}{}\renewcommand{\rhebExRtMvXgbFull}{1.02}  % ... w/o best days, real-time mean-variance: SR, STANDARD
\providecommand{\rhebExRtMvDXgbFull}{}\renewcommand{\rhebExRtMvDXgbFull}{0.48}  % ... minus the market's, STANDARD
\providecommand{\rhebExRtMvPXgbFull}{}\renewcommand{\rhebExRtMvPXgbFull}{0.553}  % ... its Ledoit-Wolf p (D19)
\providecommand{\rhebRtIvXgbFull}{}\renewcommand{\rhebRtIvXgbFull}{1.22}  % ... real-time inverse-volatility: Sharpe ratio, STANDARD
\providecommand{\rhebRtIvDXgbFull}{}\renewcommand{\rhebRtIvDXgbFull}{0.72}  % ... minus the market's, STANDARD
\providecommand{\rhebRtIvPXgbFull}{}\renewcommand{\rhebRtIvPXgbFull}{0.169}  % ... Ledoit-Wolf p of the difference (D19)
\providecommand{\rhebRtIvDSixXgbFull}{}\renewcommand{\rhebRtIvDSixXgbFull}{1.72}  % ... real-time inverse-volatility: Sharpe ratio, D6
\providecommand{\rhebRtIvDSixDXgbFull}{}\renewcommand{\rhebRtIvDSixDXgbFull}{1.04}  % ... minus the market's, D6
\providecommand{\rhebRtIvDSixPXgbFull}{}\renewcommand{\rhebRtIvDSixPXgbFull}{0.100}  % ... bootstrap p of the D6 difference
\providecommand{\rhebExRtIvXgbFull}{}\renewcommand{\rhebExRtIvXgbFull}{0.95}  % ... w/o best days, real-time inverse-volatility: SR, STANDARD
\providecommand{\rhebExRtIvDXgbFull}{}\renewcommand{\rhebExRtIvDXgbFull}{0.41}  % ... minus the market's, STANDARD
\providecommand{\rhebExRtIvPXgbFull}{}\renewcommand{\rhebExRtIvPXgbFull}{0.452}  % ... its Ledoit-Wolf p (D19)
\providecommand{\rhebRtMvShareXgbFull}{}\renewcommand{\rhebRtMvShareXgbFull}{90}  % ... real-time MV: median strategy vol share, %
\providecommand{\rhebRtMvCapDaysXgbFull}{}\renewcommand{\rhebRtMvCapDaysXgbFull}{100}  % ... % of days the 90% cap binds
\providecommand{\rhebRtMvWsXgbFull}{}\renewcommand{\rhebRtMvWsXgbFull}{3.4}  % ... real-time MV: median strategy dollar weight, %
\providecommand{\rhebRtIvWsXgbFull}{}\renewcommand{\rhebRtIvWsXgbFull}{0.9}  % ... real-time IV: median strategy dollar weight, %
\providecommand{\rhebExBetaXgbFull}{}\renewcommand{\rhebExBetaXgbFull}{1.77}  % ... without the best-day union: CAPM beta
\providecommand{\rhebExBetaTXgbFull}{}\renewcommand{\rhebExBetaTXgbFull}{0.51}  % ... its HAC t
\providecommand{\rhebExAlphaXgbFull}{}\renewcommand{\rhebExAlphaXgbFull}{5.9}  % ... without the best-day union: alpha, % per day
\providecommand{\rhebExAlphaTXgbFull}{}\renewcommand{\rhebExAlphaTXgbFull}{2.38}  % ... its HAC t
\providecommand{\rhebExAlphaPXgbFull}{}\renewcommand{\rhebExAlphaPXgbFull}{0.028}  % ... its fixed-b p
\providecommand{\rhebExARXgbFull}{}\renewcommand{\rhebExARXgbFull}{0.91}  % ... appraisal ratio, sqrt(q) (D19)
\providecommand{\rhebExSRmaxXgbFull}{}\renewcommand{\rhebExSRmaxXgbFull}{0.94}  % ... max Sharpe ratio, STANDARD
\providecommand{\rhebExDSRmaxPXgbFull}{}\renewcommand{\rhebExDSRmaxPXgbFull}{0.075}  % ... bootstrap p of its gain
\providecommand{\rhebExGrsPXgbFull}{}\renewcommand{\rhebExGrsPXgbFull}{0.135}  % ... GRS p
\providecommand{\rhebTmGammaXgbFull}{}\renewcommand{\rhebTmGammaXgbFull}{\ensuremath{-}1.81}  % pooltree_xgb_all_features, mid: Treynor-Mazuy gamma on (r^m)^2 (% per %^2)
\providecommand{\rhebTmGammaTXgbFull}{}\renewcommand{\rhebTmGammaTXgbFull}{\ensuremath{-}1.47}  % ... its HAC t
\providecommand{\rhebTmGammaPXgbFull}{}\renewcommand{\rhebTmGammaPXgbFull}{0.164}  % ... its fixed-b p
\providecommand{\rhebTmBetaXgbFull}{}\renewcommand{\rhebTmBetaXgbFull}{1.38}  % ... Treynor-Mazuy beta on r^m
\providecommand{\rhebTmBetaTXgbFull}{}\renewcommand{\rhebTmBetaTXgbFull}{0.35}  % ... its HAC t
\providecommand{\rhebNLongXgbFull}{}\renewcommand{\rhebNLongXgbFull}{143}  % ... long days
\providecommand{\rhebTmGammaLongXgbFull}{}\renewcommand{\rhebTmGammaLongXgbFull}{2.18}  % ... TM gamma on its long days only
\providecommand{\rhebTmGammaLongTXgbFull}{}\renewcommand{\rhebTmGammaLongTXgbFull}{0.32}  % ... its HAC t (no test: RULE section 4)
\providecommand{\rhebNShortXgbFull}{}\renewcommand{\rhebNShortXgbFull}{471}  % ... short days
\providecommand{\rhebTmGammaShortXgbFull}{}\renewcommand{\rhebTmGammaShortXgbFull}{\ensuremath{-}1.72}  % ... TM gamma on its short days only
\providecommand{\rhebTmGammaShortTXgbFull}{}\renewcommand{\rhebTmGammaShortTXgbFull}{\ensuremath{-}1.39}  % ... its HAC t (no test: RULE section 4)
\providecommand{\rhebDownMeanXgbFull}{}\renewcommand{\rhebDownMeanXgbFull}{15.3}  % ... mean return on market-down days, % of premium
\providecommand{\rhebDownRestXgbFull}{}\renewcommand{\rhebDownRestXgbFull}{7.9}  % ... on the other days
\providecommand{\rhebDownDiffTXgbFull}{}\renewcommand{\rhebDownDiffTXgbFull}{0.65}  % ... HAC t of the difference
\providecommand{\rhebDownDiffPXgbFull}{}\renewcommand{\rhebDownDiffPXgbFull}{0.534}  % ... its fixed-b p
\providecommand{\rhebWorstMeanXgbFull}{}\renewcommand{\rhebWorstMeanXgbFull}{\ensuremath{-}14.7}  % ... mean return on worst-decile market days, % of premium
\providecommand{\rhebWorstRestXgbFull}{}\renewcommand{\rhebWorstRestXgbFull}{14.6}  % ... on the other days
\providecommand{\rhebWorstDiffTXgbFull}{}\renewcommand{\rhebWorstDiffTXgbFull}{\ensuremath{-}2.16}  % ... HAC t of the difference
\providecommand{\rhebWorstDiffPXgbFull}{}\renewcommand{\rhebWorstDiffPXgbFull}{0.044}  % ... its fixed-b p
\providecommand{\rhebThirtyCorrXgbFull}{}\renewcommand{\rhebThirtyCorrXgbFull}{0.05}  % ... correlation with the 15:30-16:00 index return
\providecommand{\rhebThirtyBetaXgbFull}{}\renewcommand{\rhebThirtyBetaXgbFull}{20.47}  % ... beta on the 15:30-16:00 index return
\providecommand{\rhebThirtyBetaTXgbFull}{}\renewcommand{\rhebThirtyBetaTXgbFull}{0.69}  % ... its HAC t
\providecommand{\rhebThirtyBetaPXgbFull}{}\renewcommand{\rhebThirtyBetaPXgbFull}{0.510}  % ... its fixed-b p
\providecommand{\rhebThirtyAlphaXgbFull}{}\renewcommand{\rhebThirtyAlphaXgbFull}{11.8}  % ... alpha against it, % per day
\providecommand{\rhebThirtyAlphaTXgbFull}{}\renewcommand{\rhebThirtyAlphaTXgbFull}{3.38}  % ... its HAC t
\providecommand{\rhebThirtyTmGammaXgbFull}{}\renewcommand{\rhebThirtyTmGammaXgbFull}{\ensuremath{-}112.55}  % ... TM gamma on the squared 15:30-16:00 return
\providecommand{\rhebThirtyTmGammaTXgbFull}{}\renewcommand{\rhebThirtyTmGammaTXgbFull}{\ensuremath{-}2.28}  % ... its HAC t
\providecommand{\rhebThirtyTmGammaPXgbFull}{}\renewcommand{\rhebThirtyTmGammaPXgbFull}{0.035}  % ... its fixed-b p
\providecommand{\rhebThirtyTmGammaLongXgbFull}{}\renewcommand{\rhebThirtyTmGammaLongXgbFull}{389.30}  % ... ... on its long days only
\providecommand{\rhebThirtyTmGammaLongTXgbFull}{}\renewcommand{\rhebThirtyTmGammaLongTXgbFull}{2.62}  % ... its HAC t (no test)
\providecommand{\rhebThirtyTmGammaShortXgbFull}{}\renewcommand{\rhebThirtyTmGammaShortXgbFull}{\ensuremath{-}271.83}  % ... ... on its short days only
\providecommand{\rhebThirtyTmGammaShortTXgbFull}{}\renewcommand{\rhebThirtyTmGammaShortTXgbFull}{\ensuremath{-}4.38}  % ... its HAC t (no test)
\providecommand{\rhebThirtyDownMeanXgbFull}{}\renewcommand{\rhebThirtyDownMeanXgbFull}{10.5}  % ... mean on 15:30-16:00 down days, %
\providecommand{\rhebThirtyDownDiffTXgbFull}{}\renewcommand{\rhebThirtyDownDiffTXgbFull}{\ensuremath{-}0.28}  % ... HAC t of the difference with the other days
\providecommand{\rhebThirtyWorstMeanXgbFull}{}\renewcommand{\rhebThirtyWorstMeanXgbFull}{\ensuremath{-}55.2}  % ... mean on 15:30-16:00 worst days, %
\providecommand{\rhebThirtyWorstDiffTXgbFull}{}\renewcommand{\rhebThirtyWorstDiffTXgbFull}{\ensuremath{-}2.68}  % ... HAC t of the difference with the other days
\providecommand{\rhebCorrXXgbFull}{}\renewcommand{\rhebCorrXXgbFull}{0.02}  % pooltree_xgb_all_features, bid/ask: correlation with the market excess return
\providecommand{\rhebBetaXXgbFull}{}\renewcommand{\rhebBetaXXgbFull}{1.52}  % ... CAPM beta (% of premium per % of market)
\providecommand{\rhebBetaTXXgbFull}{}\renewcommand{\rhebBetaTXXgbFull}{0.39}  % ... its HAC t (S = 33)
\providecommand{\rhebBetaPXXgbFull}{}\renewcommand{\rhebBetaPXXgbFull}{0.705}  % ... its fixed-b p
\providecommand{\rhebRsqXXgbFull}{}\renewcommand{\rhebRsqXXgbFull}{0.00}  % ... CAPM R^2
\providecommand{\rhebAlphaXXgbFull}{}\renewcommand{\rhebAlphaXXgbFull}{9.0}  % ... CAPM alpha, % of premium per day
\providecommand{\rhebAlphaTXXgbFull}{}\renewcommand{\rhebAlphaTXXgbFull}{2.52}  % ... its HAC t
\providecommand{\rhebAlphaPXXgbFull}{}\renewcommand{\rhebAlphaPXXgbFull}{0.020}  % ... its fixed-b p
\providecommand{\rhebAlphaAnnXXgbFull}{}\renewcommand{\rhebAlphaAnnXXgbFull}{2020}  % ... alpha x q (D19; was x 252), % of premium a year
\providecommand{\rhebARXXgbFull}{}\renewcommand{\rhebARXXgbFull}{1.19}  % ... appraisal ratio alpha / sigma_eps, annualized sqrt(q) (D19)
\providecommand{\rhebSRsXXgbFull}{}\renewcommand{\rhebSRsXXgbFull}{1.19}  % ... strategy Sharpe ratio, STANDARD = the paper's (= \rhemSh / X)
\providecommand{\rhebSRsDSixXXgbFull}{}\renewcommand{\rhebSRsDSixXXgbFull}{1.36}  % ... strategy Sharpe ratio, D6 (= \rhemShDSix)
\providecommand{\rhebSRmaxXXgbFull}{}\renewcommand{\rhebSRmaxXXgbFull}{1.20}  % ... max Sharpe ratio of market + strategy, STANDARD
\providecommand{\rhebDSRmaxXXgbFull}{}\renewcommand{\rhebDSRmaxXXgbFull}{1.05}  % ... its gain over the market alone, STANDARD
\providecommand{\rhebDSRmaxLoXXgbFull}{}\renewcommand{\rhebDSRmaxLoXXgbFull}{0.03}  % ... bootstrap 95% interval, lower end
\providecommand{\rhebDSRmaxHiXXgbFull}{}\renewcommand{\rhebDSRmaxHiXXgbFull}{1.94}  % ... upper end
\providecommand{\rhebDSRmaxPXXgbFull}{}\renewcommand{\rhebDSRmaxPXXgbFull}{0.013}  % ... bootstrap p (descriptive; Delta >= 0)
\providecommand{\rhebGrsXXgbFull}{}\renewcommand{\rhebGrsXXgbFull}{3.84}  % ... one-asset GRS statistic, F(1, T - 2)
\providecommand{\rhebGrsPXXgbFull}{}\renewcommand{\rhebGrsPXXgbFull}{0.051}  % ... its p (i.i.d. errors)
\providecommand{\rhebTanShareXXgbFull}{}\renewcommand{\rhebTanShareXXgbFull}{99}  % ... tangency: strategy's share of volatility, %
\providecommand{\rhebTanDSixXXgbFull}{}\renewcommand{\rhebTanDSixXXgbFull}{1.37}  % ... tangency (fixed in-sample weights): Sharpe ratio, D6
\providecommand{\rhebRtMvXXgbFull}{}\renewcommand{\rhebRtMvXXgbFull}{1.31}  % ... real-time mean-variance: Sharpe ratio, STANDARD
\providecommand{\rhebRtMvDXXgbFull}{}\renewcommand{\rhebRtMvDXXgbFull}{0.81}  % ... minus the market's, STANDARD
\providecommand{\rhebRtMvPXXgbFull}{}\renewcommand{\rhebRtMvPXXgbFull}{0.289}  % ... Ledoit-Wolf p of the difference (D19)
\providecommand{\rhebRtMvDSixXXgbFull}{}\renewcommand{\rhebRtMvDSixXXgbFull}{1.48}  % ... real-time mean-variance: Sharpe ratio, D6
\providecommand{\rhebRtMvDSixDXXgbFull}{}\renewcommand{\rhebRtMvDSixDXXgbFull}{0.80}  % ... minus the market's, D6
\providecommand{\rhebRtMvDSixPXXgbFull}{}\renewcommand{\rhebRtMvDSixPXXgbFull}{0.383}  % ... bootstrap p of the D6 difference
\providecommand{\rhebExRtMvXXgbFull}{}\renewcommand{\rhebExRtMvXXgbFull}{0.64}  % ... w/o best days, real-time mean-variance: SR, STANDARD
\providecommand{\rhebExRtMvDXXgbFull}{}\renewcommand{\rhebExRtMvDXXgbFull}{0.10}  % ... minus the market's, STANDARD
\providecommand{\rhebExRtMvPXXgbFull}{}\renewcommand{\rhebExRtMvPXXgbFull}{0.899}  % ... its Ledoit-Wolf p (D19)
\providecommand{\rhebRtIvXXgbFull}{}\renewcommand{\rhebRtIvXXgbFull}{0.94}  % ... real-time inverse-volatility: Sharpe ratio, STANDARD
\providecommand{\rhebRtIvDXXgbFull}{}\renewcommand{\rhebRtIvDXXgbFull}{0.44}  % ... minus the market's, STANDARD
\providecommand{\rhebRtIvPXXgbFull}{}\renewcommand{\rhebRtIvPXXgbFull}{0.388}  % ... Ledoit-Wolf p of the difference (D19)
\providecommand{\rhebRtIvDSixXXgbFull}{}\renewcommand{\rhebRtIvDSixXXgbFull}{1.45}  % ... real-time inverse-volatility: Sharpe ratio, D6
\providecommand{\rhebRtIvDSixDXXgbFull}{}\renewcommand{\rhebRtIvDSixDXXgbFull}{0.76}  % ... minus the market's, D6
\providecommand{\rhebRtIvDSixPXXgbFull}{}\renewcommand{\rhebRtIvDSixPXXgbFull}{0.218}  % ... bootstrap p of the D6 difference
\providecommand{\rhebExRtIvXXgbFull}{}\renewcommand{\rhebExRtIvXXgbFull}{0.68}  % ... w/o best days, real-time inverse-volatility: SR, STANDARD
\providecommand{\rhebExRtIvDXXgbFull}{}\renewcommand{\rhebExRtIvDXXgbFull}{0.13}  % ... minus the market's, STANDARD
\providecommand{\rhebExRtIvPXXgbFull}{}\renewcommand{\rhebExRtIvPXXgbFull}{0.807}  % ... its Ledoit-Wolf p (D19)
\providecommand{\rhebRtMvShareXXgbFull}{}\renewcommand{\rhebRtMvShareXXgbFull}{90}  % ... real-time MV: median strategy vol share, %
\providecommand{\rhebRtMvCapDaysXXgbFull}{}\renewcommand{\rhebRtMvCapDaysXXgbFull}{100}  % ... % of days the 90% cap binds
\providecommand{\rhebRtMvWsXXgbFull}{}\renewcommand{\rhebRtMvWsXXgbFull}{3.4}  % ... real-time MV: median strategy dollar weight, %
\providecommand{\rhebRtIvWsXXgbFull}{}\renewcommand{\rhebRtIvWsXXgbFull}{0.8}  % ... real-time IV: median strategy dollar weight, %
\providecommand{\rhebExBetaXXgbFull}{}\renewcommand{\rhebExBetaXXgbFull}{1.79}  % ... without the best-day union: CAPM beta
\providecommand{\rhebExBetaTXXgbFull}{}\renewcommand{\rhebExBetaTXXgbFull}{0.51}  % ... its HAC t
\providecommand{\rhebExAlphaXXgbFull}{}\renewcommand{\rhebExAlphaXXgbFull}{3.5}  % ... without the best-day union: alpha, % per day
\providecommand{\rhebExAlphaTXXgbFull}{}\renewcommand{\rhebExAlphaTXXgbFull}{1.39}  % ... its HAC t
\providecommand{\rhebExAlphaPXXgbFull}{}\renewcommand{\rhebExAlphaPXXgbFull}{0.187}  % ... its fixed-b p
\providecommand{\rhebExARXXgbFull}{}\renewcommand{\rhebExARXXgbFull}{0.53}  % ... appraisal ratio, sqrt(q) (D19)
\providecommand{\rhebExSRmaxXXgbFull}{}\renewcommand{\rhebExSRmaxXXgbFull}{0.58}  % ... max Sharpe ratio, STANDARD
\providecommand{\rhebExDSRmaxPXXgbFull}{}\renewcommand{\rhebExDSRmaxPXXgbFull}{0.191}  % ... bootstrap p of its gain
\providecommand{\rhebExGrsPXXgbFull}{}\renewcommand{\rhebExGrsPXXgbFull}{0.384}  % ... GRS p
\providecommand{\rhebCorrRfFull}{}\renewcommand{\rhebCorrRfFull}{0.02}  % pooltree_rf_all_features, mid: correlation with the market excess return
\providecommand{\rhebBetaRfFull}{}\renewcommand{\rhebBetaRfFull}{1.65}  % ... CAPM beta (% of premium per % of market)
\providecommand{\rhebBetaTRfFull}{}\renewcommand{\rhebBetaTRfFull}{0.48}  % ... its HAC t (S = 33)
\providecommand{\rhebBetaPRfFull}{}\renewcommand{\rhebBetaPRfFull}{0.647}  % ... its fixed-b p
\providecommand{\rhebRsqRfFull}{}\renewcommand{\rhebRsqRfFull}{0.00}  % ... CAPM R^2
\providecommand{\rhebAlphaRfFull}{}\renewcommand{\rhebAlphaRfFull}{14.1}  % ... CAPM alpha, % of premium per day
\providecommand{\rhebAlphaTRfFull}{}\renewcommand{\rhebAlphaTRfFull}{3.77}  % ... its HAC t
\providecommand{\rhebAlphaPRfFull}{}\renewcommand{\rhebAlphaPRfFull}{\ensuremath{<}0.001}  % ... its fixed-b p
\providecommand{\rhebAlphaAnnRfFull}{}\renewcommand{\rhebAlphaAnnRfFull}{3170}  % ... alpha x q (D19; was x 252), % of premium a year
\providecommand{\rhebARRfFull}{}\renewcommand{\rhebARRfFull}{1.88}  % ... appraisal ratio alpha / sigma_eps, annualized sqrt(q) (D19)
\providecommand{\rhebSRsRfFull}{}\renewcommand{\rhebSRsRfFull}{1.88}  % ... strategy Sharpe ratio, STANDARD = the paper's (= \rhemSh / X)
\providecommand{\rhebSRsDSixRfFull}{}\renewcommand{\rhebSRsDSixRfFull}{1.99}  % ... strategy Sharpe ratio, D6 (= \rhemShDSix)
\providecommand{\rhebSRmaxRfFull}{}\renewcommand{\rhebSRmaxRfFull}{1.88}  % ... max Sharpe ratio of market + strategy, STANDARD
\providecommand{\rhebDSRmaxRfFull}{}\renewcommand{\rhebDSRmaxRfFull}{1.73}  % ... its gain over the market alone, STANDARD
\providecommand{\rhebDSRmaxLoRfFull}{}\renewcommand{\rhebDSRmaxLoRfFull}{0.43}  % ... bootstrap 95% interval, lower end
\providecommand{\rhebDSRmaxHiRfFull}{}\renewcommand{\rhebDSRmaxHiRfFull}{2.64}  % ... upper end
\providecommand{\rhebDSRmaxPRfFull}{}\renewcommand{\rhebDSRmaxPRfFull}{\ensuremath{<}0.001}  % ... bootstrap p (descriptive; Delta >= 0)
\providecommand{\rhebGrsRfFull}{}\renewcommand{\rhebGrsRfFull}{9.56}  % ... one-asset GRS statistic, F(1, T - 2)
\providecommand{\rhebGrsPRfFull}{}\renewcommand{\rhebGrsPRfFull}{0.002}  % ... its p (i.i.d. errors)
\providecommand{\rhebTanShareRfFull}{}\renewcommand{\rhebTanShareRfFull}{99}  % ... tangency: strategy's share of volatility, %
\providecommand{\rhebTanDSixRfFull}{}\renewcommand{\rhebTanDSixRfFull}{1.99}  % ... tangency (fixed in-sample weights): Sharpe ratio, D6
\providecommand{\rhebRtMvRfFull}{}\renewcommand{\rhebRtMvRfFull}{1.92}  % ... real-time mean-variance: Sharpe ratio, STANDARD
\providecommand{\rhebRtMvDRfFull}{}\renewcommand{\rhebRtMvDRfFull}{1.42}  % ... minus the market's, STANDARD
\providecommand{\rhebRtMvPRfFull}{}\renewcommand{\rhebRtMvPRfFull}{0.071}  % ... Ledoit-Wolf p of the difference (D19)
\providecommand{\rhebRtMvDSixRfFull}{}\renewcommand{\rhebRtMvDSixRfFull}{1.96}  % ... real-time mean-variance: Sharpe ratio, D6
\providecommand{\rhebRtMvDSixDRfFull}{}\renewcommand{\rhebRtMvDSixDRfFull}{1.28}  % ... minus the market's, D6
\providecommand{\rhebRtMvDSixPRfFull}{}\renewcommand{\rhebRtMvDSixPRfFull}{0.135}  % ... bootstrap p of the D6 difference
\providecommand{\rhebExRtMvRfFull}{}\renewcommand{\rhebExRtMvRfFull}{1.38}  % ... w/o best days, real-time mean-variance: SR, STANDARD
\providecommand{\rhebExRtMvDRfFull}{}\renewcommand{\rhebExRtMvDRfFull}{0.83}  % ... minus the market's, STANDARD
\providecommand{\rhebExRtMvPRfFull}{}\renewcommand{\rhebExRtMvPRfFull}{0.295}  % ... its Ledoit-Wolf p (D19)
\providecommand{\rhebRtIvRfFull}{}\renewcommand{\rhebRtIvRfFull}{1.52}  % ... real-time inverse-volatility: Sharpe ratio, STANDARD
\providecommand{\rhebRtIvDRfFull}{}\renewcommand{\rhebRtIvDRfFull}{1.02}  % ... minus the market's, STANDARD
\providecommand{\rhebRtIvPRfFull}{}\renewcommand{\rhebRtIvPRfFull}{0.054}  % ... Ledoit-Wolf p of the difference (D19)
\providecommand{\rhebRtIvDSixRfFull}{}\renewcommand{\rhebRtIvDSixRfFull}{1.88}  % ... real-time inverse-volatility: Sharpe ratio, D6
\providecommand{\rhebRtIvDSixDRfFull}{}\renewcommand{\rhebRtIvDSixDRfFull}{1.20}  % ... minus the market's, D6
\providecommand{\rhebRtIvDSixPRfFull}{}\renewcommand{\rhebRtIvDSixPRfFull}{0.037}  % ... bootstrap p of the D6 difference
\providecommand{\rhebExRtIvRfFull}{}\renewcommand{\rhebExRtIvRfFull}{1.23}  % ... w/o best days, real-time inverse-volatility: SR, STANDARD
\providecommand{\rhebExRtIvDRfFull}{}\renewcommand{\rhebExRtIvDRfFull}{0.69}  % ... minus the market's, STANDARD
\providecommand{\rhebExRtIvPRfFull}{}\renewcommand{\rhebExRtIvPRfFull}{0.181}  % ... its Ledoit-Wolf p (D19)
\providecommand{\rhebRtMvShareRfFull}{}\renewcommand{\rhebRtMvShareRfFull}{90}  % ... real-time MV: median strategy vol share, %
\providecommand{\rhebRtMvCapDaysRfFull}{}\renewcommand{\rhebRtMvCapDaysRfFull}{100}  % ... % of days the 90% cap binds
\providecommand{\rhebRtMvWsRfFull}{}\renewcommand{\rhebRtMvWsRfFull}{3.5}  % ... real-time MV: median strategy dollar weight, %
\providecommand{\rhebRtIvWsRfFull}{}\renewcommand{\rhebRtIvWsRfFull}{0.9}  % ... real-time IV: median strategy dollar weight, %
\providecommand{\rhebExBetaRfFull}{}\renewcommand{\rhebExBetaRfFull}{3.18}  % ... without the best-day union: CAPM beta
\providecommand{\rhebExBetaTRfFull}{}\renewcommand{\rhebExBetaTRfFull}{0.96}  % ... its HAC t
\providecommand{\rhebExAlphaRfFull}{}\renewcommand{\rhebExAlphaRfFull}{8.7}  % ... without the best-day union: alpha, % per day
\providecommand{\rhebExAlphaTRfFull}{}\renewcommand{\rhebExAlphaTRfFull}{3.22}  % ... its HAC t
\providecommand{\rhebExAlphaPRfFull}{}\renewcommand{\rhebExAlphaPRfFull}{0.004}  % ... its fixed-b p
\providecommand{\rhebExARRfFull}{}\renewcommand{\rhebExARRfFull}{1.32}  % ... appraisal ratio, sqrt(q) (D19)
\providecommand{\rhebExSRmaxRfFull}{}\renewcommand{\rhebExSRmaxRfFull}{1.34}  % ... max Sharpe ratio, STANDARD
\providecommand{\rhebExDSRmaxPRfFull}{}\renewcommand{\rhebExDSRmaxPRfFull}{0.021}  % ... bootstrap p of its gain
\providecommand{\rhebExGrsPRfFull}{}\renewcommand{\rhebExGrsPRfFull}{0.030}  % ... GRS p
\providecommand{\rhebTmGammaRfFull}{}\renewcommand{\rhebTmGammaRfFull}{\ensuremath{-}1.36}  % pooltree_rf_all_features, mid: Treynor-Mazuy gamma on (r^m)^2 (% per %^2)
\providecommand{\rhebTmGammaTRfFull}{}\renewcommand{\rhebTmGammaTRfFull}{\ensuremath{-}1.02}  % ... its HAC t
\providecommand{\rhebTmGammaPRfFull}{}\renewcommand{\rhebTmGammaPRfFull}{0.332}  % ... its fixed-b p
\providecommand{\rhebTmBetaRfFull}{}\renewcommand{\rhebTmBetaRfFull}{1.60}  % ... Treynor-Mazuy beta on r^m
\providecommand{\rhebTmBetaTRfFull}{}\renewcommand{\rhebTmBetaTRfFull}{0.45}  % ... its HAC t
\providecommand{\rhebNLongRfFull}{}\renewcommand{\rhebNLongRfFull}{132}  % ... long days
\providecommand{\rhebTmGammaLongRfFull}{}\renewcommand{\rhebTmGammaLongRfFull}{11.32}  % ... TM gamma on its long days only
\providecommand{\rhebTmGammaLongTRfFull}{}\renewcommand{\rhebTmGammaLongTRfFull}{1.11}  % ... its HAC t (no test: RULE section 4)
\providecommand{\rhebNShortRfFull}{}\renewcommand{\rhebNShortRfFull}{482}  % ... short days
\providecommand{\rhebTmGammaShortRfFull}{}\renewcommand{\rhebTmGammaShortRfFull}{\ensuremath{-}1.48}  % ... TM gamma on its short days only
\providecommand{\rhebTmGammaShortTRfFull}{}\renewcommand{\rhebTmGammaShortTRfFull}{\ensuremath{-}1.18}  % ... its HAC t (no test: RULE section 4)
\providecommand{\rhebDownMeanRfFull}{}\renewcommand{\rhebDownMeanRfFull}{15.5}  % ... mean return on market-down days, % of premium
\providecommand{\rhebDownRestRfFull}{}\renewcommand{\rhebDownRestRfFull}{12.6}  % ... on the other days
\providecommand{\rhebDownDiffTRfFull}{}\renewcommand{\rhebDownDiffTRfFull}{0.27}  % ... HAC t of the difference
\providecommand{\rhebDownDiffPRfFull}{}\renewcommand{\rhebDownDiffPRfFull}{0.796}  % ... its fixed-b p
\providecommand{\rhebWorstMeanRfFull}{}\renewcommand{\rhebWorstMeanRfFull}{\ensuremath{-}3.3}  % ... mean return on worst-decile market days, % of premium
\providecommand{\rhebWorstRestRfFull}{}\renewcommand{\rhebWorstRestRfFull}{16.0}  % ... on the other days
\providecommand{\rhebWorstDiffTRfFull}{}\renewcommand{\rhebWorstDiffTRfFull}{\ensuremath{-}1.35}  % ... HAC t of the difference
\providecommand{\rhebWorstDiffPRfFull}{}\renewcommand{\rhebWorstDiffPRfFull}{0.199}  % ... its fixed-b p
\providecommand{\rhebThirtyCorrRfFull}{}\renewcommand{\rhebThirtyCorrRfFull}{0.02}  % ... correlation with the 15:30-16:00 index return
\providecommand{\rhebThirtyBetaRfFull}{}\renewcommand{\rhebThirtyBetaRfFull}{7.03}  % ... beta on the 15:30-16:00 index return
\providecommand{\rhebThirtyBetaTRfFull}{}\renewcommand{\rhebThirtyBetaTRfFull}{0.29}  % ... its HAC t
\providecommand{\rhebThirtyBetaPRfFull}{}\renewcommand{\rhebThirtyBetaPRfFull}{0.781}  % ... its fixed-b p
\providecommand{\rhebThirtyAlphaRfFull}{}\renewcommand{\rhebThirtyAlphaRfFull}{14.2}  % ... alpha against it, % per day
\providecommand{\rhebThirtyAlphaTRfFull}{}\renewcommand{\rhebThirtyAlphaTRfFull}{3.95}  % ... its HAC t
\providecommand{\rhebThirtyTmGammaRfFull}{}\renewcommand{\rhebThirtyTmGammaRfFull}{\ensuremath{-}38.65}  % ... TM gamma on the squared 15:30-16:00 return
\providecommand{\rhebThirtyTmGammaTRfFull}{}\renewcommand{\rhebThirtyTmGammaTRfFull}{\ensuremath{-}0.55}  % ... its HAC t
\providecommand{\rhebThirtyTmGammaPRfFull}{}\renewcommand{\rhebThirtyTmGammaPRfFull}{0.595}  % ... its fixed-b p
\providecommand{\rhebThirtyTmGammaLongRfFull}{}\renewcommand{\rhebThirtyTmGammaLongRfFull}{195.08}  % ... ... on its long days only
\providecommand{\rhebThirtyTmGammaLongTRfFull}{}\renewcommand{\rhebThirtyTmGammaLongTRfFull}{3.67}  % ... its HAC t (no test)
\providecommand{\rhebThirtyTmGammaShortRfFull}{}\renewcommand{\rhebThirtyTmGammaShortRfFull}{\ensuremath{-}377.20}  % ... ... on its short days only
\providecommand{\rhebThirtyTmGammaShortTRfFull}{}\renewcommand{\rhebThirtyTmGammaShortTRfFull}{\ensuremath{-}4.46}  % ... its HAC t (no test)
\providecommand{\rhebThirtyDownMeanRfFull}{}\renewcommand{\rhebThirtyDownMeanRfFull}{12.2}  % ... mean on 15:30-16:00 down days, %
\providecommand{\rhebThirtyDownDiffTRfFull}{}\renewcommand{\rhebThirtyDownDiffTRfFull}{\ensuremath{-}0.52}  % ... HAC t of the difference with the other days
\providecommand{\rhebThirtyWorstMeanRfFull}{}\renewcommand{\rhebThirtyWorstMeanRfFull}{\ensuremath{-}59.3}  % ... mean on 15:30-16:00 worst days, %
\providecommand{\rhebThirtyWorstDiffTRfFull}{}\renewcommand{\rhebThirtyWorstDiffTRfFull}{\ensuremath{-}2.94}  % ... HAC t of the difference with the other days
\providecommand{\rhebCorrXRfFull}{}\renewcommand{\rhebCorrXRfFull}{0.02}  % pooltree_rf_all_features, bid/ask: correlation with the market excess return
\providecommand{\rhebBetaXRfFull}{}\renewcommand{\rhebBetaXRfFull}{1.72}  % ... CAPM beta (% of premium per % of market)
\providecommand{\rhebBetaTXRfFull}{}\renewcommand{\rhebBetaTXRfFull}{0.49}  % ... its HAC t (S = 33)
\providecommand{\rhebBetaPXRfFull}{}\renewcommand{\rhebBetaPXRfFull}{0.637}  % ... its fixed-b p
\providecommand{\rhebRsqXRfFull}{}\renewcommand{\rhebRsqXRfFull}{0.00}  % ... CAPM R^2
\providecommand{\rhebAlphaXRfFull}{}\renewcommand{\rhebAlphaXRfFull}{11.4}  % ... CAPM alpha, % of premium per day
\providecommand{\rhebAlphaTXRfFull}{}\renewcommand{\rhebAlphaTXRfFull}{3.09}  % ... its HAC t
\providecommand{\rhebAlphaPXRfFull}{}\renewcommand{\rhebAlphaPXRfFull}{0.005}  % ... its fixed-b p
\providecommand{\rhebAlphaAnnXRfFull}{}\renewcommand{\rhebAlphaAnnXRfFull}{2576}  % ... alpha x q (D19; was x 252), % of premium a year
\providecommand{\rhebARXRfFull}{}\renewcommand{\rhebARXRfFull}{1.52}  % ... appraisal ratio alpha / sigma_eps, annualized sqrt(q) (D19)
\providecommand{\rhebSRsXRfFull}{}\renewcommand{\rhebSRsXRfFull}{1.52}  % ... strategy Sharpe ratio, STANDARD = the paper's (= \rhemSh / X)
\providecommand{\rhebSRsDSixXRfFull}{}\renewcommand{\rhebSRsDSixXRfFull}{1.63}  % ... strategy Sharpe ratio, D6 (= \rhemShDSix)
\providecommand{\rhebSRmaxXRfFull}{}\renewcommand{\rhebSRmaxXRfFull}{1.53}  % ... max Sharpe ratio of market + strategy, STANDARD
\providecommand{\rhebDSRmaxXRfFull}{}\renewcommand{\rhebDSRmaxXRfFull}{1.37}  % ... its gain over the market alone, STANDARD
\providecommand{\rhebDSRmaxLoXRfFull}{}\renewcommand{\rhebDSRmaxLoXRfFull}{0.18}  % ... bootstrap 95% interval, lower end
\providecommand{\rhebDSRmaxHiXRfFull}{}\renewcommand{\rhebDSRmaxHiXRfFull}{2.30}  % ... upper end
\providecommand{\rhebDSRmaxPXRfFull}{}\renewcommand{\rhebDSRmaxPXRfFull}{0.004}  % ... bootstrap p (descriptive; Delta >= 0)
\providecommand{\rhebGrsXRfFull}{}\renewcommand{\rhebGrsXRfFull}{6.26}  % ... one-asset GRS statistic, F(1, T - 2)
\providecommand{\rhebGrsPXRfFull}{}\renewcommand{\rhebGrsPXRfFull}{0.013}  % ... its p (i.i.d. errors)
\providecommand{\rhebTanShareXRfFull}{}\renewcommand{\rhebTanShareXRfFull}{99}  % ... tangency: strategy's share of volatility, %
\providecommand{\rhebTanDSixXRfFull}{}\renewcommand{\rhebTanDSixXRfFull}{1.64}  % ... tangency (fixed in-sample weights): Sharpe ratio, D6
\providecommand{\rhebRtMvXRfFull}{}\renewcommand{\rhebRtMvXRfFull}{1.56}  % ... real-time mean-variance: Sharpe ratio, STANDARD
\providecommand{\rhebRtMvDXRfFull}{}\renewcommand{\rhebRtMvDXRfFull}{1.07}  % ... minus the market's, STANDARD
\providecommand{\rhebRtMvPXRfFull}{}\renewcommand{\rhebRtMvPXRfFull}{0.169}  % ... Ledoit-Wolf p of the difference (D19)
\providecommand{\rhebRtMvDSixXRfFull}{}\renewcommand{\rhebRtMvDSixXRfFull}{1.62}  % ... real-time mean-variance: Sharpe ratio, D6
\providecommand{\rhebRtMvDSixDXRfFull}{}\renewcommand{\rhebRtMvDSixDXRfFull}{0.94}  % ... minus the market's, D6
\providecommand{\rhebRtMvDSixPXRfFull}{}\renewcommand{\rhebRtMvDSixPXRfFull}{0.270}  % ... bootstrap p of the D6 difference
\providecommand{\rhebExRtMvXRfFull}{}\renewcommand{\rhebExRtMvXRfFull}{1.01}  % ... w/o best days, real-time mean-variance: SR, STANDARD
\providecommand{\rhebExRtMvDXRfFull}{}\renewcommand{\rhebExRtMvDXRfFull}{0.47}  % ... minus the market's, STANDARD
\providecommand{\rhebExRtMvPXRfFull}{}\renewcommand{\rhebExRtMvPXRfFull}{0.550}  % ... its Ledoit-Wolf p (D19)
\providecommand{\rhebRtIvXRfFull}{}\renewcommand{\rhebRtIvXRfFull}{1.24}  % ... real-time inverse-volatility: Sharpe ratio, STANDARD
\providecommand{\rhebRtIvDXRfFull}{}\renewcommand{\rhebRtIvDXRfFull}{0.74}  % ... minus the market's, STANDARD
\providecommand{\rhebRtIvPXRfFull}{}\renewcommand{\rhebRtIvPXRfFull}{0.156}  % ... Ledoit-Wolf p of the difference (D19)
\providecommand{\rhebRtIvDSixXRfFull}{}\renewcommand{\rhebRtIvDSixXRfFull}{1.60}  % ... real-time inverse-volatility: Sharpe ratio, D6
\providecommand{\rhebRtIvDSixDXRfFull}{}\renewcommand{\rhebRtIvDSixDXRfFull}{0.92}  % ... minus the market's, D6
\providecommand{\rhebRtIvDSixPXRfFull}{}\renewcommand{\rhebRtIvDSixPXRfFull}{0.105}  % ... bootstrap p of the D6 difference
\providecommand{\rhebExRtIvXRfFull}{}\renewcommand{\rhebExRtIvXRfFull}{0.95}  % ... w/o best days, real-time inverse-volatility: SR, STANDARD
\providecommand{\rhebExRtIvDXRfFull}{}\renewcommand{\rhebExRtIvDXRfFull}{0.41}  % ... minus the market's, STANDARD
\providecommand{\rhebExRtIvPXRfFull}{}\renewcommand{\rhebExRtIvPXRfFull}{0.423}  % ... its Ledoit-Wolf p (D19)
\providecommand{\rhebRtMvShareXRfFull}{}\renewcommand{\rhebRtMvShareXRfFull}{90}  % ... real-time MV: median strategy vol share, %
\providecommand{\rhebRtMvCapDaysXRfFull}{}\renewcommand{\rhebRtMvCapDaysXRfFull}{100}  % ... % of days the 90% cap binds
\providecommand{\rhebRtMvWsXRfFull}{}\renewcommand{\rhebRtMvWsXRfFull}{3.5}  % ... real-time MV: median strategy dollar weight, %
\providecommand{\rhebRtIvWsXRfFull}{}\renewcommand{\rhebRtIvWsXRfFull}{0.9}  % ... real-time IV: median strategy dollar weight, %
\providecommand{\rhebExBetaXRfFull}{}\renewcommand{\rhebExBetaXRfFull}{3.19}  % ... without the best-day union: CAPM beta
\providecommand{\rhebExBetaTXRfFull}{}\renewcommand{\rhebExBetaTXRfFull}{0.94}  % ... its HAC t
\providecommand{\rhebExAlphaXRfFull}{}\renewcommand{\rhebExAlphaXRfFull}{6.2}  % ... without the best-day union: alpha, % per day
\providecommand{\rhebExAlphaTXRfFull}{}\renewcommand{\rhebExAlphaTXRfFull}{2.28}  % ... its HAC t
\providecommand{\rhebExAlphaPXRfFull}{}\renewcommand{\rhebExAlphaPXRfFull}{0.034}  % ... its fixed-b p
\providecommand{\rhebExARXRfFull}{}\renewcommand{\rhebExARXRfFull}{0.94}  % ... appraisal ratio, sqrt(q) (D19)
\providecommand{\rhebExSRmaxXRfFull}{}\renewcommand{\rhebExSRmaxXRfFull}{0.96}  % ... max Sharpe ratio, STANDARD
\providecommand{\rhebExDSRmaxPXRfFull}{}\renewcommand{\rhebExDSRmaxPXRfFull}{0.069}  % ... bootstrap p of its gain
\providecommand{\rhebExGrsPXRfFull}{}\renewcommand{\rhebExGrsPXRfFull}{0.124}  % ... GRS p
\providecommand{\rhebCorrRidgeHar}{}\renewcommand{\rhebCorrRidgeHar}{0.02}  % poolpt13_ridge_baseline, mid: correlation with the market excess return
\providecommand{\rhebBetaRidgeHar}{}\renewcommand{\rhebBetaRidgeHar}{2.44}  % ... CAPM beta (% of premium per % of market)
\providecommand{\rhebBetaTRidgeHar}{}\renewcommand{\rhebBetaTRidgeHar}{0.71}  % ... its HAC t (S = 33)
\providecommand{\rhebBetaPRidgeHar}{}\renewcommand{\rhebBetaPRidgeHar}{0.496}  % ... its fixed-b p
\providecommand{\rhebRsqRidgeHar}{}\renewcommand{\rhebRsqRidgeHar}{0.00}  % ... CAPM R^2
\providecommand{\rhebAlphaRidgeHar}{}\renewcommand{\rhebAlphaRidgeHar}{7.2}  % ... CAPM alpha, % of premium per day
\providecommand{\rhebAlphaTRidgeHar}{}\renewcommand{\rhebAlphaTRidgeHar}{1.80}  % ... its HAC t
\providecommand{\rhebAlphaPRidgeHar}{}\renewcommand{\rhebAlphaPRidgeHar}{0.091}  % ... its fixed-b p
\providecommand{\rhebAlphaAnnRidgeHar}{}\renewcommand{\rhebAlphaAnnRidgeHar}{1615}  % ... alpha x q (D19; was x 252), % of premium a year
\providecommand{\rhebARRidgeHar}{}\renewcommand{\rhebARRidgeHar}{0.95}  % ... appraisal ratio alpha / sigma_eps, annualized sqrt(q) (D19)
\providecommand{\rhebSRsRidgeHar}{}\renewcommand{\rhebSRsRidgeHar}{0.95}  % ... strategy Sharpe ratio, STANDARD = the paper's (= \rhemSh / X)
\providecommand{\rhebSRsDSixRidgeHar}{}\renewcommand{\rhebSRsDSixRidgeHar}{0.98}  % ... strategy Sharpe ratio, D6 (= \rhemShDSix)
\providecommand{\rhebSRmaxRidgeHar}{}\renewcommand{\rhebSRmaxRidgeHar}{0.96}  % ... max Sharpe ratio of market + strategy, STANDARD
\providecommand{\rhebDSRmaxRidgeHar}{}\renewcommand{\rhebDSRmaxRidgeHar}{0.81}  % ... its gain over the market alone, STANDARD
\providecommand{\rhebDSRmaxLoRidgeHar}{}\renewcommand{\rhebDSRmaxLoRidgeHar}{0.00}  % ... bootstrap 95% interval, lower end
\providecommand{\rhebDSRmaxHiRidgeHar}{}\renewcommand{\rhebDSRmaxHiRidgeHar}{1.76}  % ... upper end
\providecommand{\rhebDSRmaxPRidgeHar}{}\renewcommand{\rhebDSRmaxPRidgeHar}{0.040}  % ... bootstrap p (descriptive; Delta >= 0)
\providecommand{\rhebGrsRidgeHar}{}\renewcommand{\rhebGrsRidgeHar}{2.46}  % ... one-asset GRS statistic, F(1, T - 2)
\providecommand{\rhebGrsPRidgeHar}{}\renewcommand{\rhebGrsPRidgeHar}{0.118}  % ... its p (i.i.d. errors)
\providecommand{\rhebTanShareRidgeHar}{}\renewcommand{\rhebTanShareRidgeHar}{98}  % ... tangency: strategy's share of volatility, %
\providecommand{\rhebTanDSixRidgeHar}{}\renewcommand{\rhebTanDSixRidgeHar}{0.98}  % ... tangency (fixed in-sample weights): Sharpe ratio, D6
\providecommand{\rhebRtMvRidgeHar}{}\renewcommand{\rhebRtMvRidgeHar}{1.02}  % ... real-time mean-variance: Sharpe ratio, STANDARD
\providecommand{\rhebRtMvDRidgeHar}{}\renewcommand{\rhebRtMvDRidgeHar}{0.52}  % ... minus the market's, STANDARD
\providecommand{\rhebRtMvPRidgeHar}{}\renewcommand{\rhebRtMvPRidgeHar}{0.524}  % ... Ledoit-Wolf p of the difference (D19)
\providecommand{\rhebRtMvDSixRidgeHar}{}\renewcommand{\rhebRtMvDSixRidgeHar}{1.10}  % ... real-time mean-variance: Sharpe ratio, D6
\providecommand{\rhebRtMvDSixDRidgeHar}{}\renewcommand{\rhebRtMvDSixDRidgeHar}{0.42}  % ... minus the market's, D6
\providecommand{\rhebRtMvDSixPRidgeHar}{}\renewcommand{\rhebRtMvDSixPRidgeHar}{0.676}  % ... bootstrap p of the D6 difference
\providecommand{\rhebExRtMvRidgeHar}{}\renewcommand{\rhebExRtMvRidgeHar}{0.95}  % ... w/o best days, real-time mean-variance: SR, STANDARD
\providecommand{\rhebExRtMvDRidgeHar}{}\renewcommand{\rhebExRtMvDRidgeHar}{0.41}  % ... minus the market's, STANDARD
\providecommand{\rhebExRtMvPRidgeHar}{}\renewcommand{\rhebExRtMvPRidgeHar}{0.616}  % ... its Ledoit-Wolf p (D19)
\providecommand{\rhebRtIvRidgeHar}{}\renewcommand{\rhebRtIvRidgeHar}{1.08}  % ... real-time inverse-volatility: Sharpe ratio, STANDARD
\providecommand{\rhebRtIvDRidgeHar}{}\renewcommand{\rhebRtIvDRidgeHar}{0.58}  % ... minus the market's, STANDARD
\providecommand{\rhebRtIvPRidgeHar}{}\renewcommand{\rhebRtIvPRidgeHar}{0.305}  % ... Ledoit-Wolf p of the difference (D19)
\providecommand{\rhebRtIvDSixRidgeHar}{}\renewcommand{\rhebRtIvDSixRidgeHar}{1.20}  % ... real-time inverse-volatility: Sharpe ratio, D6
\providecommand{\rhebRtIvDSixDRidgeHar}{}\renewcommand{\rhebRtIvDSixDRidgeHar}{0.51}  % ... minus the market's, D6
\providecommand{\rhebRtIvDSixPRidgeHar}{}\renewcommand{\rhebRtIvDSixPRidgeHar}{0.452}  % ... bootstrap p of the D6 difference
\providecommand{\rhebExRtIvRidgeHar}{}\renewcommand{\rhebExRtIvRidgeHar}{1.07}  % ... w/o best days, real-time inverse-volatility: SR, STANDARD
\providecommand{\rhebExRtIvDRidgeHar}{}\renewcommand{\rhebExRtIvDRidgeHar}{0.52}  % ... minus the market's, STANDARD
\providecommand{\rhebExRtIvPRidgeHar}{}\renewcommand{\rhebExRtIvPRidgeHar}{0.367}  % ... its Ledoit-Wolf p (D19)
\providecommand{\rhebRtMvShareRidgeHar}{}\renewcommand{\rhebRtMvShareRidgeHar}{90}  % ... real-time MV: median strategy vol share, %
\providecommand{\rhebRtMvCapDaysRidgeHar}{}\renewcommand{\rhebRtMvCapDaysRidgeHar}{100}  % ... % of days the 90% cap binds
\providecommand{\rhebRtMvWsRidgeHar}{}\renewcommand{\rhebRtMvWsRidgeHar}{3.5}  % ... real-time MV: median strategy dollar weight, %
\providecommand{\rhebRtIvWsRidgeHar}{}\renewcommand{\rhebRtIvWsRidgeHar}{0.9}  % ... real-time IV: median strategy dollar weight, %
\providecommand{\rhebExBetaRidgeHar}{}\renewcommand{\rhebExBetaRidgeHar}{1.57}  % ... without the best-day union: CAPM beta
\providecommand{\rhebExBetaTRidgeHar}{}\renewcommand{\rhebExBetaTRidgeHar}{0.49}  % ... its HAC t
\providecommand{\rhebExAlphaRidgeHar}{}\renewcommand{\rhebExAlphaRidgeHar}{4.6}  % ... without the best-day union: alpha, % per day
\providecommand{\rhebExAlphaTRidgeHar}{}\renewcommand{\rhebExAlphaTRidgeHar}{1.40}  % ... its HAC t
\providecommand{\rhebExAlphaPRidgeHar}{}\renewcommand{\rhebExAlphaPRidgeHar}{0.182}  % ... its fixed-b p
\providecommand{\rhebExARRidgeHar}{}\renewcommand{\rhebExARRidgeHar}{0.71}  % ... appraisal ratio, sqrt(q) (D19)
\providecommand{\rhebExSRmaxRidgeHar}{}\renewcommand{\rhebExSRmaxRidgeHar}{0.74}  % ... max Sharpe ratio, STANDARD
\providecommand{\rhebExDSRmaxPRidgeHar}{}\renewcommand{\rhebExDSRmaxPRidgeHar}{0.139}  % ... bootstrap p of its gain
\providecommand{\rhebExGrsPRidgeHar}{}\renewcommand{\rhebExGrsPRidgeHar}{0.244}  % ... GRS p
\providecommand{\rhebTmGammaRidgeHar}{}\renewcommand{\rhebTmGammaRidgeHar}{\ensuremath{-}0.69}  % poolpt13_ridge_baseline, mid: Treynor-Mazuy gamma on (r^m)^2 (% per %^2)
\providecommand{\rhebTmGammaTRidgeHar}{}\renewcommand{\rhebTmGammaTRidgeHar}{\ensuremath{-}0.57}  % ... its HAC t
\providecommand{\rhebTmGammaPRidgeHar}{}\renewcommand{\rhebTmGammaPRidgeHar}{0.584}  % ... its fixed-b p
\providecommand{\rhebTmBetaRidgeHar}{}\renewcommand{\rhebTmBetaRidgeHar}{2.42}  % ... Treynor-Mazuy beta on r^m
\providecommand{\rhebTmBetaTRidgeHar}{}\renewcommand{\rhebTmBetaTRidgeHar}{0.69}  % ... its HAC t
\providecommand{\rhebNLongRidgeHar}{}\renewcommand{\rhebNLongRidgeHar}{150}  % ... long days
\providecommand{\rhebTmGammaLongRidgeHar}{}\renewcommand{\rhebTmGammaLongRidgeHar}{12.79}  % ... TM gamma on its long days only
\providecommand{\rhebTmGammaLongTRidgeHar}{}\renewcommand{\rhebTmGammaLongTRidgeHar}{1.88}  % ... its HAC t (no test: RULE section 4)
\providecommand{\rhebNShortRidgeHar}{}\renewcommand{\rhebNShortRidgeHar}{464}  % ... short days
\providecommand{\rhebTmGammaShortRidgeHar}{}\renewcommand{\rhebTmGammaShortRidgeHar}{\ensuremath{-}1.06}  % ... TM gamma on its short days only
\providecommand{\rhebTmGammaShortTRidgeHar}{}\renewcommand{\rhebTmGammaShortTRidgeHar}{\ensuremath{-}0.82}  % ... its HAC t (no test: RULE section 4)
\providecommand{\rhebDownMeanRidgeHar}{}\renewcommand{\rhebDownMeanRidgeHar}{7.0}  % ... mean return on market-down days, % of premium
\providecommand{\rhebDownRestRidgeHar}{}\renewcommand{\rhebDownRestRidgeHar}{7.3}  % ... on the other days
\providecommand{\rhebDownDiffTRidgeHar}{}\renewcommand{\rhebDownDiffTRidgeHar}{\ensuremath{-}0.03}  % ... HAC t of the difference
\providecommand{\rhebDownDiffPRidgeHar}{}\renewcommand{\rhebDownDiffPRidgeHar}{0.975}  % ... its fixed-b p
\providecommand{\rhebWorstMeanRidgeHar}{}\renewcommand{\rhebWorstMeanRidgeHar}{\ensuremath{-}8.1}  % ... mean return on worst-decile market days, % of premium
\providecommand{\rhebWorstRestRidgeHar}{}\renewcommand{\rhebWorstRestRidgeHar}{8.9}  % ... on the other days
\providecommand{\rhebWorstDiffTRidgeHar}{}\renewcommand{\rhebWorstDiffTRidgeHar}{\ensuremath{-}1.21}  % ... HAC t of the difference
\providecommand{\rhebWorstDiffPRidgeHar}{}\renewcommand{\rhebWorstDiffPRidgeHar}{0.249}  % ... its fixed-b p
\providecommand{\rhebThirtyCorrRidgeHar}{}\renewcommand{\rhebThirtyCorrRidgeHar}{0.04}  % ... correlation with the 15:30-16:00 index return
\providecommand{\rhebThirtyBetaRidgeHar}{}\renewcommand{\rhebThirtyBetaRidgeHar}{14.45}  % ... beta on the 15:30-16:00 index return
\providecommand{\rhebThirtyBetaTRidgeHar}{}\renewcommand{\rhebThirtyBetaTRidgeHar}{0.52}  % ... its HAC t
\providecommand{\rhebThirtyBetaPRidgeHar}{}\renewcommand{\rhebThirtyBetaPRidgeHar}{0.616}  % ... its fixed-b p
\providecommand{\rhebThirtyAlphaRidgeHar}{}\renewcommand{\rhebThirtyAlphaRidgeHar}{7.4}  % ... alpha against it, % per day
\providecommand{\rhebThirtyAlphaTRidgeHar}{}\renewcommand{\rhebThirtyAlphaTRidgeHar}{1.87}  % ... its HAC t
\providecommand{\rhebThirtyTmGammaRidgeHar}{}\renewcommand{\rhebThirtyTmGammaRidgeHar}{\ensuremath{-}117.17}  % ... TM gamma on the squared 15:30-16:00 return
\providecommand{\rhebThirtyTmGammaTRidgeHar}{}\renewcommand{\rhebThirtyTmGammaTRidgeHar}{\ensuremath{-}2.33}  % ... its HAC t
\providecommand{\rhebThirtyTmGammaPRidgeHar}{}\renewcommand{\rhebThirtyTmGammaPRidgeHar}{0.031}  % ... its fixed-b p
\providecommand{\rhebThirtyTmGammaLongRidgeHar}{}\renewcommand{\rhebThirtyTmGammaLongRidgeHar}{372.08}  % ... ... on its long days only
\providecommand{\rhebThirtyTmGammaLongTRidgeHar}{}\renewcommand{\rhebThirtyTmGammaLongTRidgeHar}{2.92}  % ... its HAC t (no test)
\providecommand{\rhebThirtyTmGammaShortRidgeHar}{}\renewcommand{\rhebThirtyTmGammaShortRidgeHar}{\ensuremath{-}272.85}  % ... ... on its short days only
\providecommand{\rhebThirtyTmGammaShortTRidgeHar}{}\renewcommand{\rhebThirtyTmGammaShortTRidgeHar}{\ensuremath{-}4.36}  % ... its HAC t (no test)
\providecommand{\rhebThirtyDownMeanRidgeHar}{}\renewcommand{\rhebThirtyDownMeanRidgeHar}{8.8}  % ... mean on 15:30-16:00 down days, %
\providecommand{\rhebThirtyDownDiffTRidgeHar}{}\renewcommand{\rhebThirtyDownDiffTRidgeHar}{0.40}  % ... HAC t of the difference with the other days
\providecommand{\rhebThirtyWorstMeanRidgeHar}{}\renewcommand{\rhebThirtyWorstMeanRidgeHar}{\ensuremath{-}60.2}  % ... mean on 15:30-16:00 worst days, %
\providecommand{\rhebThirtyWorstDiffTRidgeHar}{}\renewcommand{\rhebThirtyWorstDiffTRidgeHar}{\ensuremath{-}2.65}  % ... HAC t of the difference with the other days
\providecommand{\rhebCorrXRidgeHar}{}\renewcommand{\rhebCorrXRidgeHar}{0.02}  % poolpt13_ridge_baseline, bid/ask: correlation with the market excess return
\providecommand{\rhebBetaXRidgeHar}{}\renewcommand{\rhebBetaXRidgeHar}{2.52}  % ... CAPM beta (% of premium per % of market)
\providecommand{\rhebBetaTXRidgeHar}{}\renewcommand{\rhebBetaTXRidgeHar}{0.72}  % ... its HAC t (S = 33)
\providecommand{\rhebBetaPXRidgeHar}{}\renewcommand{\rhebBetaPXRidgeHar}{0.491}  % ... its fixed-b p
\providecommand{\rhebRsqXRidgeHar}{}\renewcommand{\rhebRsqXRidgeHar}{0.00}  % ... CAPM R^2
\providecommand{\rhebAlphaXRidgeHar}{}\renewcommand{\rhebAlphaXRidgeHar}{4.5}  % ... CAPM alpha, % of premium per day
\providecommand{\rhebAlphaTXRidgeHar}{}\renewcommand{\rhebAlphaTXRidgeHar}{1.13}  % ... its HAC t
\providecommand{\rhebAlphaPXRidgeHar}{}\renewcommand{\rhebAlphaPXRidgeHar}{0.280}  % ... its fixed-b p
\providecommand{\rhebAlphaAnnXRidgeHar}{}\renewcommand{\rhebAlphaAnnXRidgeHar}{1022}  % ... alpha x q (D19; was x 252), % of premium a year
\providecommand{\rhebARXRidgeHar}{}\renewcommand{\rhebARXRidgeHar}{0.60}  % ... appraisal ratio alpha / sigma_eps, annualized sqrt(q) (D19)
\providecommand{\rhebSRsXRidgeHar}{}\renewcommand{\rhebSRsXRidgeHar}{0.60}  % ... strategy Sharpe ratio, STANDARD = the paper's (= \rhemSh / X)
\providecommand{\rhebSRsDSixXRidgeHar}{}\renewcommand{\rhebSRsDSixXRidgeHar}{0.62}  % ... strategy Sharpe ratio, D6 (= \rhemShDSix)
\providecommand{\rhebSRmaxXRidgeHar}{}\renewcommand{\rhebSRmaxXRidgeHar}{0.62}  % ... max Sharpe ratio of market + strategy, STANDARD
\providecommand{\rhebDSRmaxXRidgeHar}{}\renewcommand{\rhebDSRmaxXRidgeHar}{0.47}  % ... its gain over the market alone, STANDARD
\providecommand{\rhebDSRmaxLoXRidgeHar}{}\renewcommand{\rhebDSRmaxLoXRidgeHar}{0.00}  % ... bootstrap 95% interval, lower end
\providecommand{\rhebDSRmaxHiXRidgeHar}{}\renewcommand{\rhebDSRmaxHiXRidgeHar}{1.43}  % ... upper end
\providecommand{\rhebDSRmaxPXRidgeHar}{}\renewcommand{\rhebDSRmaxPXRidgeHar}{0.136}  % ... bootstrap p (descriptive; Delta >= 0)
\providecommand{\rhebGrsXRidgeHar}{}\renewcommand{\rhebGrsXRidgeHar}{0.97}  % ... one-asset GRS statistic, F(1, T - 2)
\providecommand{\rhebGrsPXRidgeHar}{}\renewcommand{\rhebGrsPXRidgeHar}{0.324}  % ... its p (i.i.d. errors)
\providecommand{\rhebTanShareXRidgeHar}{}\renewcommand{\rhebTanShareXRidgeHar}{95}  % ... tangency: strategy's share of volatility, %
\providecommand{\rhebTanDSixXRidgeHar}{}\renewcommand{\rhebTanDSixXRidgeHar}{0.62}  % ... tangency (fixed in-sample weights): Sharpe ratio, D6
\providecommand{\rhebRtMvXRidgeHar}{}\renewcommand{\rhebRtMvXRidgeHar}{0.60}  % ... real-time mean-variance: Sharpe ratio, STANDARD
\providecommand{\rhebRtMvDXRidgeHar}{}\renewcommand{\rhebRtMvDXRidgeHar}{0.10}  % ... minus the market's, STANDARD
\providecommand{\rhebRtMvPXRidgeHar}{}\renewcommand{\rhebRtMvPXRidgeHar}{0.908}  % ... Ledoit-Wolf p of the difference (D19)
\providecommand{\rhebRtMvDSixXRidgeHar}{}\renewcommand{\rhebRtMvDSixXRidgeHar}{0.61}  % ... real-time mean-variance: Sharpe ratio, D6
\providecommand{\rhebRtMvDSixDXRidgeHar}{}\renewcommand{\rhebRtMvDSixDXRidgeHar}{\ensuremath{-}0.07}  % ... minus the market's, D6
\providecommand{\rhebRtMvDSixPXRidgeHar}{}\renewcommand{\rhebRtMvDSixPXRidgeHar}{0.945}  % ... bootstrap p of the D6 difference
\providecommand{\rhebExRtMvXRidgeHar}{}\renewcommand{\rhebExRtMvXRidgeHar}{0.45}  % ... w/o best days, real-time mean-variance: SR, STANDARD
\providecommand{\rhebExRtMvDXRidgeHar}{}\renewcommand{\rhebExRtMvDXRidgeHar}{\ensuremath{-}0.09}  % ... minus the market's, STANDARD
\providecommand{\rhebExRtMvPXRidgeHar}{}\renewcommand{\rhebExRtMvPXRidgeHar}{0.912}  % ... its Ledoit-Wolf p (D19)
\providecommand{\rhebRtIvXRidgeHar}{}\renewcommand{\rhebRtIvXRidgeHar}{0.81}  % ... real-time inverse-volatility: Sharpe ratio, STANDARD
\providecommand{\rhebRtIvDXRidgeHar}{}\renewcommand{\rhebRtIvDXRidgeHar}{0.31}  % ... minus the market's, STANDARD
\providecommand{\rhebRtIvPXRidgeHar}{}\renewcommand{\rhebRtIvPXRidgeHar}{0.567}  % ... Ledoit-Wolf p of the difference (D19)
\providecommand{\rhebRtIvDSixXRidgeHar}{}\renewcommand{\rhebRtIvDSixXRidgeHar}{0.93}  % ... real-time inverse-volatility: Sharpe ratio, D6
\providecommand{\rhebRtIvDSixDXRidgeHar}{}\renewcommand{\rhebRtIvDSixDXRidgeHar}{0.25}  % ... minus the market's, D6
\providecommand{\rhebRtIvDSixPXRidgeHar}{}\renewcommand{\rhebRtIvDSixPXRidgeHar}{0.717}  % ... bootstrap p of the D6 difference
\providecommand{\rhebExRtIvXRidgeHar}{}\renewcommand{\rhebExRtIvXRidgeHar}{0.79}  % ... w/o best days, real-time inverse-volatility: SR, STANDARD
\providecommand{\rhebExRtIvDXRidgeHar}{}\renewcommand{\rhebExRtIvDXRidgeHar}{0.25}  % ... minus the market's, STANDARD
\providecommand{\rhebExRtIvPXRidgeHar}{}\renewcommand{\rhebExRtIvPXRidgeHar}{0.664}  % ... its Ledoit-Wolf p (D19)
\providecommand{\rhebRtMvShareXRidgeHar}{}\renewcommand{\rhebRtMvShareXRidgeHar}{90}  % ... real-time MV: median strategy vol share, %
\providecommand{\rhebRtMvCapDaysXRidgeHar}{}\renewcommand{\rhebRtMvCapDaysXRidgeHar}{90}  % ... % of days the 90% cap binds
\providecommand{\rhebRtMvWsXRidgeHar}{}\renewcommand{\rhebRtMvWsXRidgeHar}{3.6}  % ... real-time MV: median strategy dollar weight, %
\providecommand{\rhebRtIvWsXRidgeHar}{}\renewcommand{\rhebRtIvWsXRidgeHar}{0.8}  % ... real-time IV: median strategy dollar weight, %
\providecommand{\rhebExBetaXRidgeHar}{}\renewcommand{\rhebExBetaXRidgeHar}{1.59}  % ... without the best-day union: CAPM beta
\providecommand{\rhebExBetaTXRidgeHar}{}\renewcommand{\rhebExBetaTXRidgeHar}{0.49}  % ... its HAC t
\providecommand{\rhebExAlphaXRidgeHar}{}\renewcommand{\rhebExAlphaXRidgeHar}{2.2}  % ... without the best-day union: alpha, % per day
\providecommand{\rhebExAlphaTXRidgeHar}{}\renewcommand{\rhebExAlphaTXRidgeHar}{0.66}  % ... its HAC t
\providecommand{\rhebExAlphaPXRidgeHar}{}\renewcommand{\rhebExAlphaPXRidgeHar}{0.527}  % ... its fixed-b p
\providecommand{\rhebExARXRidgeHar}{}\renewcommand{\rhebExARXRidgeHar}{0.33}  % ... appraisal ratio, sqrt(q) (D19)
\providecommand{\rhebExSRmaxXRidgeHar}{}\renewcommand{\rhebExSRmaxXRidgeHar}{0.41}  % ... max Sharpe ratio, STANDARD
\providecommand{\rhebExDSRmaxPXRidgeHar}{}\renewcommand{\rhebExDSRmaxPXRidgeHar}{0.343}  % ... bootstrap p of its gain
\providecommand{\rhebExGrsPXRidgeHar}{}\renewcommand{\rhebExGrsPXRidgeHar}{0.583}  % ... GRS p
\providecommand{\rhebCorrLassoHar}{}\renewcommand{\rhebCorrLassoHar}{0.02}  % poolpt13_lasso_baseline, mid: correlation with the market excess return
\providecommand{\rhebBetaLassoHar}{}\renewcommand{\rhebBetaLassoHar}{2.44}  % ... CAPM beta (% of premium per % of market)
\providecommand{\rhebBetaTLassoHar}{}\renewcommand{\rhebBetaTLassoHar}{0.71}  % ... its HAC t (S = 33)
\providecommand{\rhebBetaPLassoHar}{}\renewcommand{\rhebBetaPLassoHar}{0.496}  % ... its fixed-b p
\providecommand{\rhebRsqLassoHar}{}\renewcommand{\rhebRsqLassoHar}{0.00}  % ... CAPM R^2
\providecommand{\rhebAlphaLassoHar}{}\renewcommand{\rhebAlphaLassoHar}{7.2}  % ... CAPM alpha, % of premium per day
\providecommand{\rhebAlphaTLassoHar}{}\renewcommand{\rhebAlphaTLassoHar}{1.80}  % ... its HAC t
\providecommand{\rhebAlphaPLassoHar}{}\renewcommand{\rhebAlphaPLassoHar}{0.091}  % ... its fixed-b p
\providecommand{\rhebAlphaAnnLassoHar}{}\renewcommand{\rhebAlphaAnnLassoHar}{1615}  % ... alpha x q (D19; was x 252), % of premium a year
\providecommand{\rhebARLassoHar}{}\renewcommand{\rhebARLassoHar}{0.95}  % ... appraisal ratio alpha / sigma_eps, annualized sqrt(q) (D19)
\providecommand{\rhebSRsLassoHar}{}\renewcommand{\rhebSRsLassoHar}{0.95}  % ... strategy Sharpe ratio, STANDARD = the paper's (= \rhemSh / X)
\providecommand{\rhebSRsDSixLassoHar}{}\renewcommand{\rhebSRsDSixLassoHar}{0.98}  % ... strategy Sharpe ratio, D6 (= \rhemShDSix)
\providecommand{\rhebSRmaxLassoHar}{}\renewcommand{\rhebSRmaxLassoHar}{0.96}  % ... max Sharpe ratio of market + strategy, STANDARD
\providecommand{\rhebDSRmaxLassoHar}{}\renewcommand{\rhebDSRmaxLassoHar}{0.81}  % ... its gain over the market alone, STANDARD
\providecommand{\rhebDSRmaxLoLassoHar}{}\renewcommand{\rhebDSRmaxLoLassoHar}{0.00}  % ... bootstrap 95% interval, lower end
\providecommand{\rhebDSRmaxHiLassoHar}{}\renewcommand{\rhebDSRmaxHiLassoHar}{1.76}  % ... upper end
\providecommand{\rhebDSRmaxPLassoHar}{}\renewcommand{\rhebDSRmaxPLassoHar}{0.040}  % ... bootstrap p (descriptive; Delta >= 0)
\providecommand{\rhebGrsLassoHar}{}\renewcommand{\rhebGrsLassoHar}{2.46}  % ... one-asset GRS statistic, F(1, T - 2)
\providecommand{\rhebGrsPLassoHar}{}\renewcommand{\rhebGrsPLassoHar}{0.118}  % ... its p (i.i.d. errors)
\providecommand{\rhebTanShareLassoHar}{}\renewcommand{\rhebTanShareLassoHar}{98}  % ... tangency: strategy's share of volatility, %
\providecommand{\rhebTanDSixLassoHar}{}\renewcommand{\rhebTanDSixLassoHar}{0.98}  % ... tangency (fixed in-sample weights): Sharpe ratio, D6
\providecommand{\rhebRtMvLassoHar}{}\renewcommand{\rhebRtMvLassoHar}{1.02}  % ... real-time mean-variance: Sharpe ratio, STANDARD
\providecommand{\rhebRtMvDLassoHar}{}\renewcommand{\rhebRtMvDLassoHar}{0.52}  % ... minus the market's, STANDARD
\providecommand{\rhebRtMvPLassoHar}{}\renewcommand{\rhebRtMvPLassoHar}{0.524}  % ... Ledoit-Wolf p of the difference (D19)
\providecommand{\rhebRtMvDSixLassoHar}{}\renewcommand{\rhebRtMvDSixLassoHar}{1.10}  % ... real-time mean-variance: Sharpe ratio, D6
\providecommand{\rhebRtMvDSixDLassoHar}{}\renewcommand{\rhebRtMvDSixDLassoHar}{0.42}  % ... minus the market's, D6
\providecommand{\rhebRtMvDSixPLassoHar}{}\renewcommand{\rhebRtMvDSixPLassoHar}{0.676}  % ... bootstrap p of the D6 difference
\providecommand{\rhebExRtMvLassoHar}{}\renewcommand{\rhebExRtMvLassoHar}{0.95}  % ... w/o best days, real-time mean-variance: SR, STANDARD
\providecommand{\rhebExRtMvDLassoHar}{}\renewcommand{\rhebExRtMvDLassoHar}{0.41}  % ... minus the market's, STANDARD
\providecommand{\rhebExRtMvPLassoHar}{}\renewcommand{\rhebExRtMvPLassoHar}{0.616}  % ... its Ledoit-Wolf p (D19)
\providecommand{\rhebRtIvLassoHar}{}\renewcommand{\rhebRtIvLassoHar}{1.08}  % ... real-time inverse-volatility: Sharpe ratio, STANDARD
\providecommand{\rhebRtIvDLassoHar}{}\renewcommand{\rhebRtIvDLassoHar}{0.58}  % ... minus the market's, STANDARD
\providecommand{\rhebRtIvPLassoHar}{}\renewcommand{\rhebRtIvPLassoHar}{0.305}  % ... Ledoit-Wolf p of the difference (D19)
\providecommand{\rhebRtIvDSixLassoHar}{}\renewcommand{\rhebRtIvDSixLassoHar}{1.20}  % ... real-time inverse-volatility: Sharpe ratio, D6
\providecommand{\rhebRtIvDSixDLassoHar}{}\renewcommand{\rhebRtIvDSixDLassoHar}{0.51}  % ... minus the market's, D6
\providecommand{\rhebRtIvDSixPLassoHar}{}\renewcommand{\rhebRtIvDSixPLassoHar}{0.452}  % ... bootstrap p of the D6 difference
\providecommand{\rhebExRtIvLassoHar}{}\renewcommand{\rhebExRtIvLassoHar}{1.07}  % ... w/o best days, real-time inverse-volatility: SR, STANDARD
\providecommand{\rhebExRtIvDLassoHar}{}\renewcommand{\rhebExRtIvDLassoHar}{0.52}  % ... minus the market's, STANDARD
\providecommand{\rhebExRtIvPLassoHar}{}\renewcommand{\rhebExRtIvPLassoHar}{0.367}  % ... its Ledoit-Wolf p (D19)
\providecommand{\rhebRtMvShareLassoHar}{}\renewcommand{\rhebRtMvShareLassoHar}{90}  % ... real-time MV: median strategy vol share, %
\providecommand{\rhebRtMvCapDaysLassoHar}{}\renewcommand{\rhebRtMvCapDaysLassoHar}{100}  % ... % of days the 90% cap binds
\providecommand{\rhebRtMvWsLassoHar}{}\renewcommand{\rhebRtMvWsLassoHar}{3.5}  % ... real-time MV: median strategy dollar weight, %
\providecommand{\rhebRtIvWsLassoHar}{}\renewcommand{\rhebRtIvWsLassoHar}{0.9}  % ... real-time IV: median strategy dollar weight, %
\providecommand{\rhebExBetaLassoHar}{}\renewcommand{\rhebExBetaLassoHar}{1.57}  % ... without the best-day union: CAPM beta
\providecommand{\rhebExBetaTLassoHar}{}\renewcommand{\rhebExBetaTLassoHar}{0.49}  % ... its HAC t
\providecommand{\rhebExAlphaLassoHar}{}\renewcommand{\rhebExAlphaLassoHar}{4.6}  % ... without the best-day union: alpha, % per day
\providecommand{\rhebExAlphaTLassoHar}{}\renewcommand{\rhebExAlphaTLassoHar}{1.40}  % ... its HAC t
\providecommand{\rhebExAlphaPLassoHar}{}\renewcommand{\rhebExAlphaPLassoHar}{0.182}  % ... its fixed-b p
\providecommand{\rhebExARLassoHar}{}\renewcommand{\rhebExARLassoHar}{0.71}  % ... appraisal ratio, sqrt(q) (D19)
\providecommand{\rhebExSRmaxLassoHar}{}\renewcommand{\rhebExSRmaxLassoHar}{0.74}  % ... max Sharpe ratio, STANDARD
\providecommand{\rhebExDSRmaxPLassoHar}{}\renewcommand{\rhebExDSRmaxPLassoHar}{0.139}  % ... bootstrap p of its gain
\providecommand{\rhebExGrsPLassoHar}{}\renewcommand{\rhebExGrsPLassoHar}{0.244}  % ... GRS p
\providecommand{\rhebTmGammaLassoHar}{}\renewcommand{\rhebTmGammaLassoHar}{\ensuremath{-}0.69}  % poolpt13_lasso_baseline, mid: Treynor-Mazuy gamma on (r^m)^2 (% per %^2)
\providecommand{\rhebTmGammaTLassoHar}{}\renewcommand{\rhebTmGammaTLassoHar}{\ensuremath{-}0.57}  % ... its HAC t
\providecommand{\rhebTmGammaPLassoHar}{}\renewcommand{\rhebTmGammaPLassoHar}{0.584}  % ... its fixed-b p
\providecommand{\rhebTmBetaLassoHar}{}\renewcommand{\rhebTmBetaLassoHar}{2.42}  % ... Treynor-Mazuy beta on r^m
\providecommand{\rhebTmBetaTLassoHar}{}\renewcommand{\rhebTmBetaTLassoHar}{0.69}  % ... its HAC t
\providecommand{\rhebNLongLassoHar}{}\renewcommand{\rhebNLongLassoHar}{150}  % ... long days
\providecommand{\rhebTmGammaLongLassoHar}{}\renewcommand{\rhebTmGammaLongLassoHar}{12.79}  % ... TM gamma on its long days only
\providecommand{\rhebTmGammaLongTLassoHar}{}\renewcommand{\rhebTmGammaLongTLassoHar}{1.88}  % ... its HAC t (no test: RULE section 4)
\providecommand{\rhebNShortLassoHar}{}\renewcommand{\rhebNShortLassoHar}{464}  % ... short days
\providecommand{\rhebTmGammaShortLassoHar}{}\renewcommand{\rhebTmGammaShortLassoHar}{\ensuremath{-}1.06}  % ... TM gamma on its short days only
\providecommand{\rhebTmGammaShortTLassoHar}{}\renewcommand{\rhebTmGammaShortTLassoHar}{\ensuremath{-}0.82}  % ... its HAC t (no test: RULE section 4)
\providecommand{\rhebDownMeanLassoHar}{}\renewcommand{\rhebDownMeanLassoHar}{7.0}  % ... mean return on market-down days, % of premium
\providecommand{\rhebDownRestLassoHar}{}\renewcommand{\rhebDownRestLassoHar}{7.3}  % ... on the other days
\providecommand{\rhebDownDiffTLassoHar}{}\renewcommand{\rhebDownDiffTLassoHar}{\ensuremath{-}0.03}  % ... HAC t of the difference
\providecommand{\rhebDownDiffPLassoHar}{}\renewcommand{\rhebDownDiffPLassoHar}{0.975}  % ... its fixed-b p
\providecommand{\rhebWorstMeanLassoHar}{}\renewcommand{\rhebWorstMeanLassoHar}{\ensuremath{-}8.1}  % ... mean return on worst-decile market days, % of premium
\providecommand{\rhebWorstRestLassoHar}{}\renewcommand{\rhebWorstRestLassoHar}{8.9}  % ... on the other days
\providecommand{\rhebWorstDiffTLassoHar}{}\renewcommand{\rhebWorstDiffTLassoHar}{\ensuremath{-}1.21}  % ... HAC t of the difference
\providecommand{\rhebWorstDiffPLassoHar}{}\renewcommand{\rhebWorstDiffPLassoHar}{0.249}  % ... its fixed-b p
\providecommand{\rhebThirtyCorrLassoHar}{}\renewcommand{\rhebThirtyCorrLassoHar}{0.04}  % ... correlation with the 15:30-16:00 index return
\providecommand{\rhebThirtyBetaLassoHar}{}\renewcommand{\rhebThirtyBetaLassoHar}{14.45}  % ... beta on the 15:30-16:00 index return
\providecommand{\rhebThirtyBetaTLassoHar}{}\renewcommand{\rhebThirtyBetaTLassoHar}{0.52}  % ... its HAC t
\providecommand{\rhebThirtyBetaPLassoHar}{}\renewcommand{\rhebThirtyBetaPLassoHar}{0.616}  % ... its fixed-b p
\providecommand{\rhebThirtyAlphaLassoHar}{}\renewcommand{\rhebThirtyAlphaLassoHar}{7.4}  % ... alpha against it, % per day
\providecommand{\rhebThirtyAlphaTLassoHar}{}\renewcommand{\rhebThirtyAlphaTLassoHar}{1.87}  % ... its HAC t
\providecommand{\rhebThirtyTmGammaLassoHar}{}\renewcommand{\rhebThirtyTmGammaLassoHar}{\ensuremath{-}117.17}  % ... TM gamma on the squared 15:30-16:00 return
\providecommand{\rhebThirtyTmGammaTLassoHar}{}\renewcommand{\rhebThirtyTmGammaTLassoHar}{\ensuremath{-}2.33}  % ... its HAC t
\providecommand{\rhebThirtyTmGammaPLassoHar}{}\renewcommand{\rhebThirtyTmGammaPLassoHar}{0.031}  % ... its fixed-b p
\providecommand{\rhebThirtyTmGammaLongLassoHar}{}\renewcommand{\rhebThirtyTmGammaLongLassoHar}{372.08}  % ... ... on its long days only
\providecommand{\rhebThirtyTmGammaLongTLassoHar}{}\renewcommand{\rhebThirtyTmGammaLongTLassoHar}{2.92}  % ... its HAC t (no test)
\providecommand{\rhebThirtyTmGammaShortLassoHar}{}\renewcommand{\rhebThirtyTmGammaShortLassoHar}{\ensuremath{-}272.85}  % ... ... on its short days only
\providecommand{\rhebThirtyTmGammaShortTLassoHar}{}\renewcommand{\rhebThirtyTmGammaShortTLassoHar}{\ensuremath{-}4.36}  % ... its HAC t (no test)
\providecommand{\rhebThirtyDownMeanLassoHar}{}\renewcommand{\rhebThirtyDownMeanLassoHar}{8.8}  % ... mean on 15:30-16:00 down days, %
\providecommand{\rhebThirtyDownDiffTLassoHar}{}\renewcommand{\rhebThirtyDownDiffTLassoHar}{0.40}  % ... HAC t of the difference with the other days
\providecommand{\rhebThirtyWorstMeanLassoHar}{}\renewcommand{\rhebThirtyWorstMeanLassoHar}{\ensuremath{-}60.2}  % ... mean on 15:30-16:00 worst days, %
\providecommand{\rhebThirtyWorstDiffTLassoHar}{}\renewcommand{\rhebThirtyWorstDiffTLassoHar}{\ensuremath{-}2.65}  % ... HAC t of the difference with the other days
\providecommand{\rhebCorrXLassoHar}{}\renewcommand{\rhebCorrXLassoHar}{0.02}  % poolpt13_lasso_baseline, bid/ask: correlation with the market excess return
\providecommand{\rhebBetaXLassoHar}{}\renewcommand{\rhebBetaXLassoHar}{2.52}  % ... CAPM beta (% of premium per % of market)
\providecommand{\rhebBetaTXLassoHar}{}\renewcommand{\rhebBetaTXLassoHar}{0.72}  % ... its HAC t (S = 33)
\providecommand{\rhebBetaPXLassoHar}{}\renewcommand{\rhebBetaPXLassoHar}{0.491}  % ... its fixed-b p
\providecommand{\rhebRsqXLassoHar}{}\renewcommand{\rhebRsqXLassoHar}{0.00}  % ... CAPM R^2
\providecommand{\rhebAlphaXLassoHar}{}\renewcommand{\rhebAlphaXLassoHar}{4.5}  % ... CAPM alpha, % of premium per day
\providecommand{\rhebAlphaTXLassoHar}{}\renewcommand{\rhebAlphaTXLassoHar}{1.13}  % ... its HAC t
\providecommand{\rhebAlphaPXLassoHar}{}\renewcommand{\rhebAlphaPXLassoHar}{0.280}  % ... its fixed-b p
\providecommand{\rhebAlphaAnnXLassoHar}{}\renewcommand{\rhebAlphaAnnXLassoHar}{1022}  % ... alpha x q (D19; was x 252), % of premium a year
\providecommand{\rhebARXLassoHar}{}\renewcommand{\rhebARXLassoHar}{0.60}  % ... appraisal ratio alpha / sigma_eps, annualized sqrt(q) (D19)
\providecommand{\rhebSRsXLassoHar}{}\renewcommand{\rhebSRsXLassoHar}{0.60}  % ... strategy Sharpe ratio, STANDARD = the paper's (= \rhemSh / X)
\providecommand{\rhebSRsDSixXLassoHar}{}\renewcommand{\rhebSRsDSixXLassoHar}{0.62}  % ... strategy Sharpe ratio, D6 (= \rhemShDSix)
\providecommand{\rhebSRmaxXLassoHar}{}\renewcommand{\rhebSRmaxXLassoHar}{0.62}  % ... max Sharpe ratio of market + strategy, STANDARD
\providecommand{\rhebDSRmaxXLassoHar}{}\renewcommand{\rhebDSRmaxXLassoHar}{0.47}  % ... its gain over the market alone, STANDARD
\providecommand{\rhebDSRmaxLoXLassoHar}{}\renewcommand{\rhebDSRmaxLoXLassoHar}{0.00}  % ... bootstrap 95% interval, lower end
\providecommand{\rhebDSRmaxHiXLassoHar}{}\renewcommand{\rhebDSRmaxHiXLassoHar}{1.43}  % ... upper end
\providecommand{\rhebDSRmaxPXLassoHar}{}\renewcommand{\rhebDSRmaxPXLassoHar}{0.136}  % ... bootstrap p (descriptive; Delta >= 0)
\providecommand{\rhebGrsXLassoHar}{}\renewcommand{\rhebGrsXLassoHar}{0.97}  % ... one-asset GRS statistic, F(1, T - 2)
\providecommand{\rhebGrsPXLassoHar}{}\renewcommand{\rhebGrsPXLassoHar}{0.324}  % ... its p (i.i.d. errors)
\providecommand{\rhebTanShareXLassoHar}{}\renewcommand{\rhebTanShareXLassoHar}{95}  % ... tangency: strategy's share of volatility, %
\providecommand{\rhebTanDSixXLassoHar}{}\renewcommand{\rhebTanDSixXLassoHar}{0.62}  % ... tangency (fixed in-sample weights): Sharpe ratio, D6
\providecommand{\rhebRtMvXLassoHar}{}\renewcommand{\rhebRtMvXLassoHar}{0.60}  % ... real-time mean-variance: Sharpe ratio, STANDARD
\providecommand{\rhebRtMvDXLassoHar}{}\renewcommand{\rhebRtMvDXLassoHar}{0.10}  % ... minus the market's, STANDARD
\providecommand{\rhebRtMvPXLassoHar}{}\renewcommand{\rhebRtMvPXLassoHar}{0.908}  % ... Ledoit-Wolf p of the difference (D19)
\providecommand{\rhebRtMvDSixXLassoHar}{}\renewcommand{\rhebRtMvDSixXLassoHar}{0.61}  % ... real-time mean-variance: Sharpe ratio, D6
\providecommand{\rhebRtMvDSixDXLassoHar}{}\renewcommand{\rhebRtMvDSixDXLassoHar}{\ensuremath{-}0.07}  % ... minus the market's, D6
\providecommand{\rhebRtMvDSixPXLassoHar}{}\renewcommand{\rhebRtMvDSixPXLassoHar}{0.945}  % ... bootstrap p of the D6 difference
\providecommand{\rhebExRtMvXLassoHar}{}\renewcommand{\rhebExRtMvXLassoHar}{0.45}  % ... w/o best days, real-time mean-variance: SR, STANDARD
\providecommand{\rhebExRtMvDXLassoHar}{}\renewcommand{\rhebExRtMvDXLassoHar}{\ensuremath{-}0.09}  % ... minus the market's, STANDARD
\providecommand{\rhebExRtMvPXLassoHar}{}\renewcommand{\rhebExRtMvPXLassoHar}{0.912}  % ... its Ledoit-Wolf p (D19)
\providecommand{\rhebRtIvXLassoHar}{}\renewcommand{\rhebRtIvXLassoHar}{0.81}  % ... real-time inverse-volatility: Sharpe ratio, STANDARD
\providecommand{\rhebRtIvDXLassoHar}{}\renewcommand{\rhebRtIvDXLassoHar}{0.31}  % ... minus the market's, STANDARD
\providecommand{\rhebRtIvPXLassoHar}{}\renewcommand{\rhebRtIvPXLassoHar}{0.567}  % ... Ledoit-Wolf p of the difference (D19)
\providecommand{\rhebRtIvDSixXLassoHar}{}\renewcommand{\rhebRtIvDSixXLassoHar}{0.93}  % ... real-time inverse-volatility: Sharpe ratio, D6
\providecommand{\rhebRtIvDSixDXLassoHar}{}\renewcommand{\rhebRtIvDSixDXLassoHar}{0.25}  % ... minus the market's, D6
\providecommand{\rhebRtIvDSixPXLassoHar}{}\renewcommand{\rhebRtIvDSixPXLassoHar}{0.717}  % ... bootstrap p of the D6 difference
\providecommand{\rhebExRtIvXLassoHar}{}\renewcommand{\rhebExRtIvXLassoHar}{0.79}  % ... w/o best days, real-time inverse-volatility: SR, STANDARD
\providecommand{\rhebExRtIvDXLassoHar}{}\renewcommand{\rhebExRtIvDXLassoHar}{0.25}  % ... minus the market's, STANDARD
\providecommand{\rhebExRtIvPXLassoHar}{}\renewcommand{\rhebExRtIvPXLassoHar}{0.664}  % ... its Ledoit-Wolf p (D19)
\providecommand{\rhebRtMvShareXLassoHar}{}\renewcommand{\rhebRtMvShareXLassoHar}{90}  % ... real-time MV: median strategy vol share, %
\providecommand{\rhebRtMvCapDaysXLassoHar}{}\renewcommand{\rhebRtMvCapDaysXLassoHar}{90}  % ... % of days the 90% cap binds
\providecommand{\rhebRtMvWsXLassoHar}{}\renewcommand{\rhebRtMvWsXLassoHar}{3.6}  % ... real-time MV: median strategy dollar weight, %
\providecommand{\rhebRtIvWsXLassoHar}{}\renewcommand{\rhebRtIvWsXLassoHar}{0.8}  % ... real-time IV: median strategy dollar weight, %
\providecommand{\rhebExBetaXLassoHar}{}\renewcommand{\rhebExBetaXLassoHar}{1.59}  % ... without the best-day union: CAPM beta
\providecommand{\rhebExBetaTXLassoHar}{}\renewcommand{\rhebExBetaTXLassoHar}{0.49}  % ... its HAC t
\providecommand{\rhebExAlphaXLassoHar}{}\renewcommand{\rhebExAlphaXLassoHar}{2.2}  % ... without the best-day union: alpha, % per day
\providecommand{\rhebExAlphaTXLassoHar}{}\renewcommand{\rhebExAlphaTXLassoHar}{0.66}  % ... its HAC t
\providecommand{\rhebExAlphaPXLassoHar}{}\renewcommand{\rhebExAlphaPXLassoHar}{0.527}  % ... its fixed-b p
\providecommand{\rhebExARXLassoHar}{}\renewcommand{\rhebExARXLassoHar}{0.33}  % ... appraisal ratio, sqrt(q) (D19)
\providecommand{\rhebExSRmaxXLassoHar}{}\renewcommand{\rhebExSRmaxXLassoHar}{0.41}  % ... max Sharpe ratio, STANDARD
\providecommand{\rhebExDSRmaxPXLassoHar}{}\renewcommand{\rhebExDSRmaxPXLassoHar}{0.343}  % ... bootstrap p of its gain
\providecommand{\rhebExGrsPXLassoHar}{}\renewcommand{\rhebExGrsPXLassoHar}{0.583}  % ... GRS p
\providecommand{\rhebCorrEnetHar}{}\renewcommand{\rhebCorrEnetHar}{0.02}  % poolpt13_enet_baseline, mid: correlation with the market excess return
\providecommand{\rhebBetaEnetHar}{}\renewcommand{\rhebBetaEnetHar}{2.44}  % ... CAPM beta (% of premium per % of market)
\providecommand{\rhebBetaTEnetHar}{}\renewcommand{\rhebBetaTEnetHar}{0.71}  % ... its HAC t (S = 33)
\providecommand{\rhebBetaPEnetHar}{}\renewcommand{\rhebBetaPEnetHar}{0.496}  % ... its fixed-b p
\providecommand{\rhebRsqEnetHar}{}\renewcommand{\rhebRsqEnetHar}{0.00}  % ... CAPM R^2
\providecommand{\rhebAlphaEnetHar}{}\renewcommand{\rhebAlphaEnetHar}{7.2}  % ... CAPM alpha, % of premium per day
\providecommand{\rhebAlphaTEnetHar}{}\renewcommand{\rhebAlphaTEnetHar}{1.80}  % ... its HAC t
\providecommand{\rhebAlphaPEnetHar}{}\renewcommand{\rhebAlphaPEnetHar}{0.091}  % ... its fixed-b p
\providecommand{\rhebAlphaAnnEnetHar}{}\renewcommand{\rhebAlphaAnnEnetHar}{1615}  % ... alpha x q (D19; was x 252), % of premium a year
\providecommand{\rhebAREnetHar}{}\renewcommand{\rhebAREnetHar}{0.95}  % ... appraisal ratio alpha / sigma_eps, annualized sqrt(q) (D19)
\providecommand{\rhebSRsEnetHar}{}\renewcommand{\rhebSRsEnetHar}{0.95}  % ... strategy Sharpe ratio, STANDARD = the paper's (= \rhemSh / X)
\providecommand{\rhebSRsDSixEnetHar}{}\renewcommand{\rhebSRsDSixEnetHar}{0.98}  % ... strategy Sharpe ratio, D6 (= \rhemShDSix)
\providecommand{\rhebSRmaxEnetHar}{}\renewcommand{\rhebSRmaxEnetHar}{0.96}  % ... max Sharpe ratio of market + strategy, STANDARD
\providecommand{\rhebDSRmaxEnetHar}{}\renewcommand{\rhebDSRmaxEnetHar}{0.81}  % ... its gain over the market alone, STANDARD
\providecommand{\rhebDSRmaxLoEnetHar}{}\renewcommand{\rhebDSRmaxLoEnetHar}{0.00}  % ... bootstrap 95% interval, lower end
\providecommand{\rhebDSRmaxHiEnetHar}{}\renewcommand{\rhebDSRmaxHiEnetHar}{1.76}  % ... upper end
\providecommand{\rhebDSRmaxPEnetHar}{}\renewcommand{\rhebDSRmaxPEnetHar}{0.040}  % ... bootstrap p (descriptive; Delta >= 0)
\providecommand{\rhebGrsEnetHar}{}\renewcommand{\rhebGrsEnetHar}{2.46}  % ... one-asset GRS statistic, F(1, T - 2)
\providecommand{\rhebGrsPEnetHar}{}\renewcommand{\rhebGrsPEnetHar}{0.118}  % ... its p (i.i.d. errors)
\providecommand{\rhebTanShareEnetHar}{}\renewcommand{\rhebTanShareEnetHar}{98}  % ... tangency: strategy's share of volatility, %
\providecommand{\rhebTanDSixEnetHar}{}\renewcommand{\rhebTanDSixEnetHar}{0.98}  % ... tangency (fixed in-sample weights): Sharpe ratio, D6
\providecommand{\rhebRtMvEnetHar}{}\renewcommand{\rhebRtMvEnetHar}{1.02}  % ... real-time mean-variance: Sharpe ratio, STANDARD
\providecommand{\rhebRtMvDEnetHar}{}\renewcommand{\rhebRtMvDEnetHar}{0.52}  % ... minus the market's, STANDARD
\providecommand{\rhebRtMvPEnetHar}{}\renewcommand{\rhebRtMvPEnetHar}{0.524}  % ... Ledoit-Wolf p of the difference (D19)
\providecommand{\rhebRtMvDSixEnetHar}{}\renewcommand{\rhebRtMvDSixEnetHar}{1.10}  % ... real-time mean-variance: Sharpe ratio, D6
\providecommand{\rhebRtMvDSixDEnetHar}{}\renewcommand{\rhebRtMvDSixDEnetHar}{0.42}  % ... minus the market's, D6
\providecommand{\rhebRtMvDSixPEnetHar}{}\renewcommand{\rhebRtMvDSixPEnetHar}{0.676}  % ... bootstrap p of the D6 difference
\providecommand{\rhebExRtMvEnetHar}{}\renewcommand{\rhebExRtMvEnetHar}{0.95}  % ... w/o best days, real-time mean-variance: SR, STANDARD
\providecommand{\rhebExRtMvDEnetHar}{}\renewcommand{\rhebExRtMvDEnetHar}{0.41}  % ... minus the market's, STANDARD
\providecommand{\rhebExRtMvPEnetHar}{}\renewcommand{\rhebExRtMvPEnetHar}{0.616}  % ... its Ledoit-Wolf p (D19)
\providecommand{\rhebRtIvEnetHar}{}\renewcommand{\rhebRtIvEnetHar}{1.08}  % ... real-time inverse-volatility: Sharpe ratio, STANDARD
\providecommand{\rhebRtIvDEnetHar}{}\renewcommand{\rhebRtIvDEnetHar}{0.58}  % ... minus the market's, STANDARD
\providecommand{\rhebRtIvPEnetHar}{}\renewcommand{\rhebRtIvPEnetHar}{0.305}  % ... Ledoit-Wolf p of the difference (D19)
\providecommand{\rhebRtIvDSixEnetHar}{}\renewcommand{\rhebRtIvDSixEnetHar}{1.20}  % ... real-time inverse-volatility: Sharpe ratio, D6
\providecommand{\rhebRtIvDSixDEnetHar}{}\renewcommand{\rhebRtIvDSixDEnetHar}{0.51}  % ... minus the market's, D6
\providecommand{\rhebRtIvDSixPEnetHar}{}\renewcommand{\rhebRtIvDSixPEnetHar}{0.452}  % ... bootstrap p of the D6 difference
\providecommand{\rhebExRtIvEnetHar}{}\renewcommand{\rhebExRtIvEnetHar}{1.07}  % ... w/o best days, real-time inverse-volatility: SR, STANDARD
\providecommand{\rhebExRtIvDEnetHar}{}\renewcommand{\rhebExRtIvDEnetHar}{0.52}  % ... minus the market's, STANDARD
\providecommand{\rhebExRtIvPEnetHar}{}\renewcommand{\rhebExRtIvPEnetHar}{0.367}  % ... its Ledoit-Wolf p (D19)
\providecommand{\rhebRtMvShareEnetHar}{}\renewcommand{\rhebRtMvShareEnetHar}{90}  % ... real-time MV: median strategy vol share, %
\providecommand{\rhebRtMvCapDaysEnetHar}{}\renewcommand{\rhebRtMvCapDaysEnetHar}{100}  % ... % of days the 90% cap binds
\providecommand{\rhebRtMvWsEnetHar}{}\renewcommand{\rhebRtMvWsEnetHar}{3.5}  % ... real-time MV: median strategy dollar weight, %
\providecommand{\rhebRtIvWsEnetHar}{}\renewcommand{\rhebRtIvWsEnetHar}{0.9}  % ... real-time IV: median strategy dollar weight, %
\providecommand{\rhebExBetaEnetHar}{}\renewcommand{\rhebExBetaEnetHar}{1.57}  % ... without the best-day union: CAPM beta
\providecommand{\rhebExBetaTEnetHar}{}\renewcommand{\rhebExBetaTEnetHar}{0.49}  % ... its HAC t
\providecommand{\rhebExAlphaEnetHar}{}\renewcommand{\rhebExAlphaEnetHar}{4.6}  % ... without the best-day union: alpha, % per day
\providecommand{\rhebExAlphaTEnetHar}{}\renewcommand{\rhebExAlphaTEnetHar}{1.40}  % ... its HAC t
\providecommand{\rhebExAlphaPEnetHar}{}\renewcommand{\rhebExAlphaPEnetHar}{0.182}  % ... its fixed-b p
\providecommand{\rhebExAREnetHar}{}\renewcommand{\rhebExAREnetHar}{0.71}  % ... appraisal ratio, sqrt(q) (D19)
\providecommand{\rhebExSRmaxEnetHar}{}\renewcommand{\rhebExSRmaxEnetHar}{0.74}  % ... max Sharpe ratio, STANDARD
\providecommand{\rhebExDSRmaxPEnetHar}{}\renewcommand{\rhebExDSRmaxPEnetHar}{0.139}  % ... bootstrap p of its gain
\providecommand{\rhebExGrsPEnetHar}{}\renewcommand{\rhebExGrsPEnetHar}{0.244}  % ... GRS p
\providecommand{\rhebTmGammaEnetHar}{}\renewcommand{\rhebTmGammaEnetHar}{\ensuremath{-}0.69}  % poolpt13_enet_baseline, mid: Treynor-Mazuy gamma on (r^m)^2 (% per %^2)
\providecommand{\rhebTmGammaTEnetHar}{}\renewcommand{\rhebTmGammaTEnetHar}{\ensuremath{-}0.57}  % ... its HAC t
\providecommand{\rhebTmGammaPEnetHar}{}\renewcommand{\rhebTmGammaPEnetHar}{0.584}  % ... its fixed-b p
\providecommand{\rhebTmBetaEnetHar}{}\renewcommand{\rhebTmBetaEnetHar}{2.42}  % ... Treynor-Mazuy beta on r^m
\providecommand{\rhebTmBetaTEnetHar}{}\renewcommand{\rhebTmBetaTEnetHar}{0.69}  % ... its HAC t
\providecommand{\rhebNLongEnetHar}{}\renewcommand{\rhebNLongEnetHar}{150}  % ... long days
\providecommand{\rhebTmGammaLongEnetHar}{}\renewcommand{\rhebTmGammaLongEnetHar}{12.79}  % ... TM gamma on its long days only
\providecommand{\rhebTmGammaLongTEnetHar}{}\renewcommand{\rhebTmGammaLongTEnetHar}{1.88}  % ... its HAC t (no test: RULE section 4)
\providecommand{\rhebNShortEnetHar}{}\renewcommand{\rhebNShortEnetHar}{464}  % ... short days
\providecommand{\rhebTmGammaShortEnetHar}{}\renewcommand{\rhebTmGammaShortEnetHar}{\ensuremath{-}1.06}  % ... TM gamma on its short days only
\providecommand{\rhebTmGammaShortTEnetHar}{}\renewcommand{\rhebTmGammaShortTEnetHar}{\ensuremath{-}0.82}  % ... its HAC t (no test: RULE section 4)
\providecommand{\rhebDownMeanEnetHar}{}\renewcommand{\rhebDownMeanEnetHar}{7.0}  % ... mean return on market-down days, % of premium
\providecommand{\rhebDownRestEnetHar}{}\renewcommand{\rhebDownRestEnetHar}{7.3}  % ... on the other days
\providecommand{\rhebDownDiffTEnetHar}{}\renewcommand{\rhebDownDiffTEnetHar}{\ensuremath{-}0.03}  % ... HAC t of the difference
\providecommand{\rhebDownDiffPEnetHar}{}\renewcommand{\rhebDownDiffPEnetHar}{0.975}  % ... its fixed-b p
\providecommand{\rhebWorstMeanEnetHar}{}\renewcommand{\rhebWorstMeanEnetHar}{\ensuremath{-}8.1}  % ... mean return on worst-decile market days, % of premium
\providecommand{\rhebWorstRestEnetHar}{}\renewcommand{\rhebWorstRestEnetHar}{8.9}  % ... on the other days
\providecommand{\rhebWorstDiffTEnetHar}{}\renewcommand{\rhebWorstDiffTEnetHar}{\ensuremath{-}1.21}  % ... HAC t of the difference
\providecommand{\rhebWorstDiffPEnetHar}{}\renewcommand{\rhebWorstDiffPEnetHar}{0.249}  % ... its fixed-b p
\providecommand{\rhebThirtyCorrEnetHar}{}\renewcommand{\rhebThirtyCorrEnetHar}{0.04}  % ... correlation with the 15:30-16:00 index return
\providecommand{\rhebThirtyBetaEnetHar}{}\renewcommand{\rhebThirtyBetaEnetHar}{14.45}  % ... beta on the 15:30-16:00 index return
\providecommand{\rhebThirtyBetaTEnetHar}{}\renewcommand{\rhebThirtyBetaTEnetHar}{0.52}  % ... its HAC t
\providecommand{\rhebThirtyBetaPEnetHar}{}\renewcommand{\rhebThirtyBetaPEnetHar}{0.616}  % ... its fixed-b p
\providecommand{\rhebThirtyAlphaEnetHar}{}\renewcommand{\rhebThirtyAlphaEnetHar}{7.4}  % ... alpha against it, % per day
\providecommand{\rhebThirtyAlphaTEnetHar}{}\renewcommand{\rhebThirtyAlphaTEnetHar}{1.87}  % ... its HAC t
\providecommand{\rhebThirtyTmGammaEnetHar}{}\renewcommand{\rhebThirtyTmGammaEnetHar}{\ensuremath{-}117.17}  % ... TM gamma on the squared 15:30-16:00 return
\providecommand{\rhebThirtyTmGammaTEnetHar}{}\renewcommand{\rhebThirtyTmGammaTEnetHar}{\ensuremath{-}2.33}  % ... its HAC t
\providecommand{\rhebThirtyTmGammaPEnetHar}{}\renewcommand{\rhebThirtyTmGammaPEnetHar}{0.031}  % ... its fixed-b p
\providecommand{\rhebThirtyTmGammaLongEnetHar}{}\renewcommand{\rhebThirtyTmGammaLongEnetHar}{372.08}  % ... ... on its long days only
\providecommand{\rhebThirtyTmGammaLongTEnetHar}{}\renewcommand{\rhebThirtyTmGammaLongTEnetHar}{2.92}  % ... its HAC t (no test)
\providecommand{\rhebThirtyTmGammaShortEnetHar}{}\renewcommand{\rhebThirtyTmGammaShortEnetHar}{\ensuremath{-}272.85}  % ... ... on its short days only
\providecommand{\rhebThirtyTmGammaShortTEnetHar}{}\renewcommand{\rhebThirtyTmGammaShortTEnetHar}{\ensuremath{-}4.36}  % ... its HAC t (no test)
\providecommand{\rhebThirtyDownMeanEnetHar}{}\renewcommand{\rhebThirtyDownMeanEnetHar}{8.8}  % ... mean on 15:30-16:00 down days, %
\providecommand{\rhebThirtyDownDiffTEnetHar}{}\renewcommand{\rhebThirtyDownDiffTEnetHar}{0.40}  % ... HAC t of the difference with the other days
\providecommand{\rhebThirtyWorstMeanEnetHar}{}\renewcommand{\rhebThirtyWorstMeanEnetHar}{\ensuremath{-}60.2}  % ... mean on 15:30-16:00 worst days, %
\providecommand{\rhebThirtyWorstDiffTEnetHar}{}\renewcommand{\rhebThirtyWorstDiffTEnetHar}{\ensuremath{-}2.65}  % ... HAC t of the difference with the other days
\providecommand{\rhebCorrXEnetHar}{}\renewcommand{\rhebCorrXEnetHar}{0.02}  % poolpt13_enet_baseline, bid/ask: correlation with the market excess return
\providecommand{\rhebBetaXEnetHar}{}\renewcommand{\rhebBetaXEnetHar}{2.52}  % ... CAPM beta (% of premium per % of market)
\providecommand{\rhebBetaTXEnetHar}{}\renewcommand{\rhebBetaTXEnetHar}{0.72}  % ... its HAC t (S = 33)
\providecommand{\rhebBetaPXEnetHar}{}\renewcommand{\rhebBetaPXEnetHar}{0.491}  % ... its fixed-b p
\providecommand{\rhebRsqXEnetHar}{}\renewcommand{\rhebRsqXEnetHar}{0.00}  % ... CAPM R^2
\providecommand{\rhebAlphaXEnetHar}{}\renewcommand{\rhebAlphaXEnetHar}{4.5}  % ... CAPM alpha, % of premium per day
\providecommand{\rhebAlphaTXEnetHar}{}\renewcommand{\rhebAlphaTXEnetHar}{1.13}  % ... its HAC t
\providecommand{\rhebAlphaPXEnetHar}{}\renewcommand{\rhebAlphaPXEnetHar}{0.280}  % ... its fixed-b p
\providecommand{\rhebAlphaAnnXEnetHar}{}\renewcommand{\rhebAlphaAnnXEnetHar}{1022}  % ... alpha x q (D19; was x 252), % of premium a year
\providecommand{\rhebARXEnetHar}{}\renewcommand{\rhebARXEnetHar}{0.60}  % ... appraisal ratio alpha / sigma_eps, annualized sqrt(q) (D19)
\providecommand{\rhebSRsXEnetHar}{}\renewcommand{\rhebSRsXEnetHar}{0.60}  % ... strategy Sharpe ratio, STANDARD = the paper's (= \rhemSh / X)
\providecommand{\rhebSRsDSixXEnetHar}{}\renewcommand{\rhebSRsDSixXEnetHar}{0.62}  % ... strategy Sharpe ratio, D6 (= \rhemShDSix)
\providecommand{\rhebSRmaxXEnetHar}{}\renewcommand{\rhebSRmaxXEnetHar}{0.62}  % ... max Sharpe ratio of market + strategy, STANDARD
\providecommand{\rhebDSRmaxXEnetHar}{}\renewcommand{\rhebDSRmaxXEnetHar}{0.47}  % ... its gain over the market alone, STANDARD
\providecommand{\rhebDSRmaxLoXEnetHar}{}\renewcommand{\rhebDSRmaxLoXEnetHar}{0.00}  % ... bootstrap 95% interval, lower end
\providecommand{\rhebDSRmaxHiXEnetHar}{}\renewcommand{\rhebDSRmaxHiXEnetHar}{1.43}  % ... upper end
\providecommand{\rhebDSRmaxPXEnetHar}{}\renewcommand{\rhebDSRmaxPXEnetHar}{0.136}  % ... bootstrap p (descriptive; Delta >= 0)
\providecommand{\rhebGrsXEnetHar}{}\renewcommand{\rhebGrsXEnetHar}{0.97}  % ... one-asset GRS statistic, F(1, T - 2)
\providecommand{\rhebGrsPXEnetHar}{}\renewcommand{\rhebGrsPXEnetHar}{0.324}  % ... its p (i.i.d. errors)
\providecommand{\rhebTanShareXEnetHar}{}\renewcommand{\rhebTanShareXEnetHar}{95}  % ... tangency: strategy's share of volatility, %
\providecommand{\rhebTanDSixXEnetHar}{}\renewcommand{\rhebTanDSixXEnetHar}{0.62}  % ... tangency (fixed in-sample weights): Sharpe ratio, D6
\providecommand{\rhebRtMvXEnetHar}{}\renewcommand{\rhebRtMvXEnetHar}{0.60}  % ... real-time mean-variance: Sharpe ratio, STANDARD
\providecommand{\rhebRtMvDXEnetHar}{}\renewcommand{\rhebRtMvDXEnetHar}{0.10}  % ... minus the market's, STANDARD
\providecommand{\rhebRtMvPXEnetHar}{}\renewcommand{\rhebRtMvPXEnetHar}{0.908}  % ... Ledoit-Wolf p of the difference (D19)
\providecommand{\rhebRtMvDSixXEnetHar}{}\renewcommand{\rhebRtMvDSixXEnetHar}{0.61}  % ... real-time mean-variance: Sharpe ratio, D6
\providecommand{\rhebRtMvDSixDXEnetHar}{}\renewcommand{\rhebRtMvDSixDXEnetHar}{\ensuremath{-}0.07}  % ... minus the market's, D6
\providecommand{\rhebRtMvDSixPXEnetHar}{}\renewcommand{\rhebRtMvDSixPXEnetHar}{0.945}  % ... bootstrap p of the D6 difference
\providecommand{\rhebExRtMvXEnetHar}{}\renewcommand{\rhebExRtMvXEnetHar}{0.45}  % ... w/o best days, real-time mean-variance: SR, STANDARD
\providecommand{\rhebExRtMvDXEnetHar}{}\renewcommand{\rhebExRtMvDXEnetHar}{\ensuremath{-}0.09}  % ... minus the market's, STANDARD
\providecommand{\rhebExRtMvPXEnetHar}{}\renewcommand{\rhebExRtMvPXEnetHar}{0.912}  % ... its Ledoit-Wolf p (D19)
\providecommand{\rhebRtIvXEnetHar}{}\renewcommand{\rhebRtIvXEnetHar}{0.81}  % ... real-time inverse-volatility: Sharpe ratio, STANDARD
\providecommand{\rhebRtIvDXEnetHar}{}\renewcommand{\rhebRtIvDXEnetHar}{0.31}  % ... minus the market's, STANDARD
\providecommand{\rhebRtIvPXEnetHar}{}\renewcommand{\rhebRtIvPXEnetHar}{0.567}  % ... Ledoit-Wolf p of the difference (D19)
\providecommand{\rhebRtIvDSixXEnetHar}{}\renewcommand{\rhebRtIvDSixXEnetHar}{0.93}  % ... real-time inverse-volatility: Sharpe ratio, D6
\providecommand{\rhebRtIvDSixDXEnetHar}{}\renewcommand{\rhebRtIvDSixDXEnetHar}{0.25}  % ... minus the market's, D6
\providecommand{\rhebRtIvDSixPXEnetHar}{}\renewcommand{\rhebRtIvDSixPXEnetHar}{0.717}  % ... bootstrap p of the D6 difference
\providecommand{\rhebExRtIvXEnetHar}{}\renewcommand{\rhebExRtIvXEnetHar}{0.79}  % ... w/o best days, real-time inverse-volatility: SR, STANDARD
\providecommand{\rhebExRtIvDXEnetHar}{}\renewcommand{\rhebExRtIvDXEnetHar}{0.25}  % ... minus the market's, STANDARD
\providecommand{\rhebExRtIvPXEnetHar}{}\renewcommand{\rhebExRtIvPXEnetHar}{0.664}  % ... its Ledoit-Wolf p (D19)
\providecommand{\rhebRtMvShareXEnetHar}{}\renewcommand{\rhebRtMvShareXEnetHar}{90}  % ... real-time MV: median strategy vol share, %
\providecommand{\rhebRtMvCapDaysXEnetHar}{}\renewcommand{\rhebRtMvCapDaysXEnetHar}{90}  % ... % of days the 90% cap binds
\providecommand{\rhebRtMvWsXEnetHar}{}\renewcommand{\rhebRtMvWsXEnetHar}{3.6}  % ... real-time MV: median strategy dollar weight, %
\providecommand{\rhebRtIvWsXEnetHar}{}\renewcommand{\rhebRtIvWsXEnetHar}{0.8}  % ... real-time IV: median strategy dollar weight, %
\providecommand{\rhebExBetaXEnetHar}{}\renewcommand{\rhebExBetaXEnetHar}{1.59}  % ... without the best-day union: CAPM beta
\providecommand{\rhebExBetaTXEnetHar}{}\renewcommand{\rhebExBetaTXEnetHar}{0.49}  % ... its HAC t
\providecommand{\rhebExAlphaXEnetHar}{}\renewcommand{\rhebExAlphaXEnetHar}{2.2}  % ... without the best-day union: alpha, % per day
\providecommand{\rhebExAlphaTXEnetHar}{}\renewcommand{\rhebExAlphaTXEnetHar}{0.66}  % ... its HAC t
\providecommand{\rhebExAlphaPXEnetHar}{}\renewcommand{\rhebExAlphaPXEnetHar}{0.527}  % ... its fixed-b p
\providecommand{\rhebExARXEnetHar}{}\renewcommand{\rhebExARXEnetHar}{0.33}  % ... appraisal ratio, sqrt(q) (D19)
\providecommand{\rhebExSRmaxXEnetHar}{}\renewcommand{\rhebExSRmaxXEnetHar}{0.41}  % ... max Sharpe ratio, STANDARD
\providecommand{\rhebExDSRmaxPXEnetHar}{}\renewcommand{\rhebExDSRmaxPXEnetHar}{0.343}  % ... bootstrap p of its gain
\providecommand{\rhebExGrsPXEnetHar}{}\renewcommand{\rhebExGrsPXEnetHar}{0.583}  % ... GRS p
\providecommand{\rhebCorrLgbmHar}{}\renewcommand{\rhebCorrLgbmHar}{0.01}  % pooltree_lgbm_baseline, mid: correlation with the market excess return
\providecommand{\rhebBetaLgbmHar}{}\renewcommand{\rhebBetaLgbmHar}{1.05}  % ... CAPM beta (% of premium per % of market)
\providecommand{\rhebBetaTLgbmHar}{}\renewcommand{\rhebBetaTLgbmHar}{0.31}  % ... its HAC t (S = 33)
\providecommand{\rhebBetaPLgbmHar}{}\renewcommand{\rhebBetaPLgbmHar}{0.765}  % ... its fixed-b p
\providecommand{\rhebRsqLgbmHar}{}\renewcommand{\rhebRsqLgbmHar}{0.00}  % ... CAPM R^2
\providecommand{\rhebAlphaLgbmHar}{}\renewcommand{\rhebAlphaLgbmHar}{4.9}  % ... CAPM alpha, % of premium per day
\providecommand{\rhebAlphaTLgbmHar}{}\renewcommand{\rhebAlphaTLgbmHar}{1.32}  % ... its HAC t
\providecommand{\rhebAlphaPLgbmHar}{}\renewcommand{\rhebAlphaPLgbmHar}{0.209}  % ... its fixed-b p
\providecommand{\rhebAlphaAnnLgbmHar}{}\renewcommand{\rhebAlphaAnnLgbmHar}{1110}  % ... alpha x q (D19; was x 252), % of premium a year
\providecommand{\rhebARLgbmHar}{}\renewcommand{\rhebARLgbmHar}{0.65}  % ... appraisal ratio alpha / sigma_eps, annualized sqrt(q) (D19)
\providecommand{\rhebSRsLgbmHar}{}\renewcommand{\rhebSRsLgbmHar}{0.65}  % ... strategy Sharpe ratio, STANDARD = the paper's (= \rhemSh / X)
\providecommand{\rhebSRsDSixLgbmHar}{}\renewcommand{\rhebSRsDSixLgbmHar}{0.61}  % ... strategy Sharpe ratio, D6 (= \rhemShDSix)
\providecommand{\rhebSRmaxLgbmHar}{}\renewcommand{\rhebSRmaxLgbmHar}{0.67}  % ... max Sharpe ratio of market + strategy, STANDARD
\providecommand{\rhebDSRmaxLgbmHar}{}\renewcommand{\rhebDSRmaxLgbmHar}{0.52}  % ... its gain over the market alone, STANDARD
\providecommand{\rhebDSRmaxLoLgbmHar}{}\renewcommand{\rhebDSRmaxLoLgbmHar}{0.00}  % ... bootstrap 95% interval, lower end
\providecommand{\rhebDSRmaxHiLgbmHar}{}\renewcommand{\rhebDSRmaxHiLgbmHar}{1.39}  % ... upper end
\providecommand{\rhebDSRmaxPLgbmHar}{}\renewcommand{\rhebDSRmaxPLgbmHar}{0.105}  % ... bootstrap p (descriptive; Delta >= 0)
\providecommand{\rhebGrsLgbmHar}{}\renewcommand{\rhebGrsLgbmHar}{1.16}  % ... one-asset GRS statistic, F(1, T - 2)
\providecommand{\rhebGrsPLgbmHar}{}\renewcommand{\rhebGrsPLgbmHar}{0.283}  % ... its p (i.i.d. errors)
\providecommand{\rhebTanShareLgbmHar}{}\renewcommand{\rhebTanShareLgbmHar}{95}  % ... tangency: strategy's share of volatility, %
\providecommand{\rhebTanDSixLgbmHar}{}\renewcommand{\rhebTanDSixLgbmHar}{0.62}  % ... tangency (fixed in-sample weights): Sharpe ratio, D6
\providecommand{\rhebRtMvLgbmHar}{}\renewcommand{\rhebRtMvLgbmHar}{0.79}  % ... real-time mean-variance: Sharpe ratio, STANDARD
\providecommand{\rhebRtMvDLgbmHar}{}\renewcommand{\rhebRtMvDLgbmHar}{0.29}  % ... minus the market's, STANDARD
\providecommand{\rhebRtMvPLgbmHar}{}\renewcommand{\rhebRtMvPLgbmHar}{0.717}  % ... Ledoit-Wolf p of the difference (D19)
\providecommand{\rhebRtMvDSixLgbmHar}{}\renewcommand{\rhebRtMvDSixLgbmHar}{0.72}  % ... real-time mean-variance: Sharpe ratio, D6
\providecommand{\rhebRtMvDSixDLgbmHar}{}\renewcommand{\rhebRtMvDSixDLgbmHar}{0.04}  % ... minus the market's, D6
\providecommand{\rhebRtMvDSixPLgbmHar}{}\renewcommand{\rhebRtMvDSixPLgbmHar}{0.963}  % ... bootstrap p of the D6 difference
\providecommand{\rhebExRtMvLgbmHar}{}\renewcommand{\rhebExRtMvLgbmHar}{0.69}  % ... w/o best days, real-time mean-variance: SR, STANDARD
\providecommand{\rhebExRtMvDLgbmHar}{}\renewcommand{\rhebExRtMvDLgbmHar}{0.15}  % ... minus the market's, STANDARD
\providecommand{\rhebExRtMvPLgbmHar}{}\renewcommand{\rhebExRtMvPLgbmHar}{0.868}  % ... its Ledoit-Wolf p (D19)
\providecommand{\rhebRtIvLgbmHar}{}\renewcommand{\rhebRtIvLgbmHar}{0.87}  % ... real-time inverse-volatility: Sharpe ratio, STANDARD
\providecommand{\rhebRtIvDLgbmHar}{}\renewcommand{\rhebRtIvDLgbmHar}{0.37}  % ... minus the market's, STANDARD
\providecommand{\rhebRtIvPLgbmHar}{}\renewcommand{\rhebRtIvPLgbmHar}{0.501}  % ... Ledoit-Wolf p of the difference (D19)
\providecommand{\rhebRtIvDSixLgbmHar}{}\renewcommand{\rhebRtIvDSixLgbmHar}{0.93}  % ... real-time inverse-volatility: Sharpe ratio, D6
\providecommand{\rhebRtIvDSixDLgbmHar}{}\renewcommand{\rhebRtIvDSixDLgbmHar}{0.25}  % ... minus the market's, D6
\providecommand{\rhebRtIvDSixPLgbmHar}{}\renewcommand{\rhebRtIvDSixPLgbmHar}{0.705}  % ... bootstrap p of the D6 difference
\providecommand{\rhebExRtIvLgbmHar}{}\renewcommand{\rhebExRtIvLgbmHar}{0.93}  % ... w/o best days, real-time inverse-volatility: SR, STANDARD
\providecommand{\rhebExRtIvDLgbmHar}{}\renewcommand{\rhebExRtIvDLgbmHar}{0.39}  % ... minus the market's, STANDARD
\providecommand{\rhebExRtIvPLgbmHar}{}\renewcommand{\rhebExRtIvPLgbmHar}{0.493}  % ... its Ledoit-Wolf p (D19)
\providecommand{\rhebRtMvShareLgbmHar}{}\renewcommand{\rhebRtMvShareLgbmHar}{90}  % ... real-time MV: median strategy vol share, %
\providecommand{\rhebRtMvCapDaysLgbmHar}{}\renewcommand{\rhebRtMvCapDaysLgbmHar}{95}  % ... % of days the 90% cap binds
\providecommand{\rhebRtMvWsLgbmHar}{}\renewcommand{\rhebRtMvWsLgbmHar}{3.4}  % ... real-time MV: median strategy dollar weight, %
\providecommand{\rhebRtIvWsLgbmHar}{}\renewcommand{\rhebRtIvWsLgbmHar}{0.9}  % ... real-time IV: median strategy dollar weight, %
\providecommand{\rhebExBetaLgbmHar}{}\renewcommand{\rhebExBetaLgbmHar}{0.18}  % ... without the best-day union: CAPM beta
\providecommand{\rhebExBetaTLgbmHar}{}\renewcommand{\rhebExBetaTLgbmHar}{0.06}  % ... its HAC t
\providecommand{\rhebExAlphaLgbmHar}{}\renewcommand{\rhebExAlphaLgbmHar}{2.4}  % ... without the best-day union: alpha, % per day
\providecommand{\rhebExAlphaTLgbmHar}{}\renewcommand{\rhebExAlphaTLgbmHar}{0.67}  % ... its HAC t
\providecommand{\rhebExAlphaPLgbmHar}{}\renewcommand{\rhebExAlphaPLgbmHar}{0.521}  % ... its fixed-b p
\providecommand{\rhebExARLgbmHar}{}\renewcommand{\rhebExARLgbmHar}{0.36}  % ... appraisal ratio, sqrt(q) (D19)
\providecommand{\rhebExSRmaxLgbmHar}{}\renewcommand{\rhebExSRmaxLgbmHar}{0.43}  % ... max Sharpe ratio, STANDARD
\providecommand{\rhebExDSRmaxPLgbmHar}{}\renewcommand{\rhebExDSRmaxPLgbmHar}{0.317}  % ... bootstrap p of its gain
\providecommand{\rhebExGrsPLgbmHar}{}\renewcommand{\rhebExGrsPLgbmHar}{0.552}  % ... GRS p
\providecommand{\rhebTmGammaLgbmHar}{}\renewcommand{\rhebTmGammaLgbmHar}{\ensuremath{-}0.43}  % pooltree_lgbm_baseline, mid: Treynor-Mazuy gamma on (r^m)^2 (% per %^2)
\providecommand{\rhebTmGammaTLgbmHar}{}\renewcommand{\rhebTmGammaTLgbmHar}{\ensuremath{-}0.40}  % ... its HAC t
\providecommand{\rhebTmGammaPLgbmHar}{}\renewcommand{\rhebTmGammaPLgbmHar}{0.704}  % ... its fixed-b p
\providecommand{\rhebTmBetaLgbmHar}{}\renewcommand{\rhebTmBetaLgbmHar}{1.04}  % ... Treynor-Mazuy beta on r^m
\providecommand{\rhebTmBetaTLgbmHar}{}\renewcommand{\rhebTmBetaTLgbmHar}{0.30}  % ... its HAC t
\providecommand{\rhebNLongLgbmHar}{}\renewcommand{\rhebNLongLgbmHar}{191}  % ... long days
\providecommand{\rhebTmGammaLongLgbmHar}{}\renewcommand{\rhebTmGammaLongLgbmHar}{10.71}  % ... TM gamma on its long days only
\providecommand{\rhebTmGammaLongTLgbmHar}{}\renewcommand{\rhebTmGammaLongTLgbmHar}{1.83}  % ... its HAC t (no test: RULE section 4)
\providecommand{\rhebNShortLgbmHar}{}\renewcommand{\rhebNShortLgbmHar}{423}  % ... short days
\providecommand{\rhebTmGammaShortLgbmHar}{}\renewcommand{\rhebTmGammaShortLgbmHar}{\ensuremath{-}0.93}  % ... TM gamma on its short days only
\providecommand{\rhebTmGammaShortTLgbmHar}{}\renewcommand{\rhebTmGammaShortTLgbmHar}{\ensuremath{-}0.75}  % ... its HAC t (no test: RULE section 4)
\providecommand{\rhebDownMeanLgbmHar}{}\renewcommand{\rhebDownMeanLgbmHar}{6.1}  % ... mean return on market-down days, % of premium
\providecommand{\rhebDownRestLgbmHar}{}\renewcommand{\rhebDownRestLgbmHar}{3.8}  % ... on the other days
\providecommand{\rhebDownDiffTLgbmHar}{}\renewcommand{\rhebDownDiffTLgbmHar}{0.23}  % ... HAC t of the difference
\providecommand{\rhebDownDiffPLgbmHar}{}\renewcommand{\rhebDownDiffPLgbmHar}{0.826}  % ... its fixed-b p
\providecommand{\rhebWorstMeanLgbmHar}{}\renewcommand{\rhebWorstMeanLgbmHar}{\ensuremath{-}2.3}  % ... mean return on worst-decile market days, % of premium
\providecommand{\rhebWorstRestLgbmHar}{}\renewcommand{\rhebWorstRestLgbmHar}{5.8}  % ... on the other days
\providecommand{\rhebWorstDiffTLgbmHar}{}\renewcommand{\rhebWorstDiffTLgbmHar}{\ensuremath{-}0.53}  % ... HAC t of the difference
\providecommand{\rhebWorstDiffPLgbmHar}{}\renewcommand{\rhebWorstDiffPLgbmHar}{0.610}  % ... its fixed-b p
\providecommand{\rhebThirtyCorrLgbmHar}{}\renewcommand{\rhebThirtyCorrLgbmHar}{0.04}  % ... correlation with the 15:30-16:00 index return
\providecommand{\rhebThirtyBetaLgbmHar}{}\renewcommand{\rhebThirtyBetaLgbmHar}{16.56}  % ... beta on the 15:30-16:00 index return
\providecommand{\rhebThirtyBetaTLgbmHar}{}\renewcommand{\rhebThirtyBetaTLgbmHar}{0.62}  % ... its HAC t
\providecommand{\rhebThirtyBetaPLgbmHar}{}\renewcommand{\rhebThirtyBetaPLgbmHar}{0.550}  % ... its fixed-b p
\providecommand{\rhebThirtyAlphaLgbmHar}{}\renewcommand{\rhebThirtyAlphaLgbmHar}{5.1}  % ... alpha against it, % per day
\providecommand{\rhebThirtyAlphaTLgbmHar}{}\renewcommand{\rhebThirtyAlphaTLgbmHar}{1.38}  % ... its HAC t
\providecommand{\rhebThirtyTmGammaLgbmHar}{}\renewcommand{\rhebThirtyTmGammaLgbmHar}{\ensuremath{-}97.25}  % ... TM gamma on the squared 15:30-16:00 return
\providecommand{\rhebThirtyTmGammaTLgbmHar}{}\renewcommand{\rhebThirtyTmGammaTLgbmHar}{\ensuremath{-}2.08}  % ... its HAC t
\providecommand{\rhebThirtyTmGammaPLgbmHar}{}\renewcommand{\rhebThirtyTmGammaPLgbmHar}{0.053}  % ... its fixed-b p
\providecommand{\rhebThirtyTmGammaLongLgbmHar}{}\renewcommand{\rhebThirtyTmGammaLongLgbmHar}{392.81}  % ... ... on its long days only
\providecommand{\rhebThirtyTmGammaLongTLgbmHar}{}\renewcommand{\rhebThirtyTmGammaLongTLgbmHar}{2.54}  % ... its HAC t (no test)
\providecommand{\rhebThirtyTmGammaShortLgbmHar}{}\renewcommand{\rhebThirtyTmGammaShortLgbmHar}{\ensuremath{-}266.48}  % ... ... on its short days only
\providecommand{\rhebThirtyTmGammaShortTLgbmHar}{}\renewcommand{\rhebThirtyTmGammaShortTLgbmHar}{\ensuremath{-}4.39}  % ... its HAC t (no test)
\providecommand{\rhebThirtyDownMeanLgbmHar}{}\renewcommand{\rhebThirtyDownMeanLgbmHar}{5.9}  % ... mean on 15:30-16:00 down days, %
\providecommand{\rhebThirtyDownDiffTLgbmHar}{}\renewcommand{\rhebThirtyDownDiffTLgbmHar}{0.26}  % ... HAC t of the difference with the other days
\providecommand{\rhebThirtyWorstMeanLgbmHar}{}\renewcommand{\rhebThirtyWorstMeanLgbmHar}{\ensuremath{-}55.0}  % ... mean on 15:30-16:00 worst days, %
\providecommand{\rhebThirtyWorstDiffTLgbmHar}{}\renewcommand{\rhebThirtyWorstDiffTLgbmHar}{\ensuremath{-}2.41}  % ... HAC t of the difference with the other days
\providecommand{\rhebCorrXLgbmHar}{}\renewcommand{\rhebCorrXLgbmHar}{0.01}  % pooltree_lgbm_baseline, bid/ask: correlation with the market excess return
\providecommand{\rhebBetaXLgbmHar}{}\renewcommand{\rhebBetaXLgbmHar}{1.13}  % ... CAPM beta (% of premium per % of market)
\providecommand{\rhebBetaTXLgbmHar}{}\renewcommand{\rhebBetaTXLgbmHar}{0.33}  % ... its HAC t (S = 33)
\providecommand{\rhebBetaPXLgbmHar}{}\renewcommand{\rhebBetaPXLgbmHar}{0.751}  % ... its fixed-b p
\providecommand{\rhebRsqXLgbmHar}{}\renewcommand{\rhebRsqXLgbmHar}{0.00}  % ... CAPM R^2
\providecommand{\rhebAlphaXLgbmHar}{}\renewcommand{\rhebAlphaXLgbmHar}{2.3}  % ... CAPM alpha, % of premium per day
\providecommand{\rhebAlphaTXLgbmHar}{}\renewcommand{\rhebAlphaTXLgbmHar}{0.61}  % ... its HAC t
\providecommand{\rhebAlphaPXLgbmHar}{}\renewcommand{\rhebAlphaPXLgbmHar}{0.556}  % ... its fixed-b p
\providecommand{\rhebAlphaAnnXLgbmHar}{}\renewcommand{\rhebAlphaAnnXLgbmHar}{519}  % ... alpha x q (D19; was x 252), % of premium a year
\providecommand{\rhebARXLgbmHar}{}\renewcommand{\rhebARXLgbmHar}{0.30}  % ... appraisal ratio alpha / sigma_eps, annualized sqrt(q) (D19)
\providecommand{\rhebSRsXLgbmHar}{}\renewcommand{\rhebSRsXLgbmHar}{0.31}  % ... strategy Sharpe ratio, STANDARD = the paper's (= \rhemSh / X)
\providecommand{\rhebSRsDSixXLgbmHar}{}\renewcommand{\rhebSRsDSixXLgbmHar}{0.25}  % ... strategy Sharpe ratio, D6 (= \rhemShDSix)
\providecommand{\rhebSRmaxXLgbmHar}{}\renewcommand{\rhebSRmaxXLgbmHar}{0.34}  % ... max Sharpe ratio of market + strategy, STANDARD
\providecommand{\rhebDSRmaxXLgbmHar}{}\renewcommand{\rhebDSRmaxXLgbmHar}{0.19}  % ... its gain over the market alone, STANDARD
\providecommand{\rhebDSRmaxLoXLgbmHar}{}\renewcommand{\rhebDSRmaxLoXLgbmHar}{0.00}  % ... bootstrap 95% interval, lower end
\providecommand{\rhebDSRmaxHiXLgbmHar}{}\renewcommand{\rhebDSRmaxHiXLgbmHar}{1.09}  % ... upper end
\providecommand{\rhebDSRmaxPXLgbmHar}{}\renewcommand{\rhebDSRmaxPXLgbmHar}{0.313}  % ... bootstrap p (descriptive; Delta >= 0)
\providecommand{\rhebGrsXLgbmHar}{}\renewcommand{\rhebGrsXLgbmHar}{0.25}  % ... one-asset GRS statistic, F(1, T - 2)
\providecommand{\rhebGrsPXLgbmHar}{}\renewcommand{\rhebGrsPXLgbmHar}{0.616}  % ... its p (i.i.d. errors)
\providecommand{\rhebTanShareXLgbmHar}{}\renewcommand{\rhebTanShareXLgbmHar}{81}  % ... tangency: strategy's share of volatility, %
\providecommand{\rhebTanDSixXLgbmHar}{}\renewcommand{\rhebTanDSixXLgbmHar}{0.30}  % ... tangency (fixed in-sample weights): Sharpe ratio, D6
\providecommand{\rhebRtMvXLgbmHar}{}\renewcommand{\rhebRtMvXLgbmHar}{0.31}  % ... real-time mean-variance: Sharpe ratio, STANDARD
\providecommand{\rhebRtMvDXLgbmHar}{}\renewcommand{\rhebRtMvDXLgbmHar}{\ensuremath{-}0.19}  % ... minus the market's, STANDARD
\providecommand{\rhebRtMvPXLgbmHar}{}\renewcommand{\rhebRtMvPXLgbmHar}{0.814}  % ... Ledoit-Wolf p of the difference (D19)
\providecommand{\rhebRtMvDSixXLgbmHar}{}\renewcommand{\rhebRtMvDSixXLgbmHar}{\ensuremath{-}0.02}  % ... real-time mean-variance: Sharpe ratio, D6
\providecommand{\rhebRtMvDSixDXLgbmHar}{}\renewcommand{\rhebRtMvDSixDXLgbmHar}{\ensuremath{-}0.71}  % ... minus the market's, D6
\providecommand{\rhebRtMvDSixPXLgbmHar}{}\renewcommand{\rhebRtMvDSixPXLgbmHar}{0.460}  % ... bootstrap p of the D6 difference
\providecommand{\rhebExRtMvXLgbmHar}{}\renewcommand{\rhebExRtMvXLgbmHar}{0.19}  % ... w/o best days, real-time mean-variance: SR, STANDARD
\providecommand{\rhebExRtMvDXLgbmHar}{}\renewcommand{\rhebExRtMvDXLgbmHar}{\ensuremath{-}0.35}  % ... minus the market's, STANDARD
\providecommand{\rhebExRtMvPXLgbmHar}{}\renewcommand{\rhebExRtMvPXLgbmHar}{0.665}  % ... its Ledoit-Wolf p (D19)
\providecommand{\rhebRtIvXLgbmHar}{}\renewcommand{\rhebRtIvXLgbmHar}{0.61}  % ... real-time inverse-volatility: Sharpe ratio, STANDARD
\providecommand{\rhebRtIvDXLgbmHar}{}\renewcommand{\rhebRtIvDXLgbmHar}{0.11}  % ... minus the market's, STANDARD
\providecommand{\rhebRtIvPXLgbmHar}{}\renewcommand{\rhebRtIvPXLgbmHar}{0.845}  % ... Ledoit-Wolf p of the difference (D19)
\providecommand{\rhebRtIvDSixXLgbmHar}{}\renewcommand{\rhebRtIvDSixXLgbmHar}{0.67}  % ... real-time inverse-volatility: Sharpe ratio, D6
\providecommand{\rhebRtIvDSixDXLgbmHar}{}\renewcommand{\rhebRtIvDSixDXLgbmHar}{\ensuremath{-}0.02}  % ... minus the market's, D6
\providecommand{\rhebRtIvDSixPXLgbmHar}{}\renewcommand{\rhebRtIvDSixPXLgbmHar}{0.978}  % ... bootstrap p of the D6 difference
\providecommand{\rhebExRtIvXLgbmHar}{}\renewcommand{\rhebExRtIvXLgbmHar}{0.66}  % ... w/o best days, real-time inverse-volatility: SR, STANDARD
\providecommand{\rhebExRtIvDXLgbmHar}{}\renewcommand{\rhebExRtIvDXLgbmHar}{0.12}  % ... minus the market's, STANDARD
\providecommand{\rhebExRtIvPXLgbmHar}{}\renewcommand{\rhebExRtIvPXLgbmHar}{0.835}  % ... its Ledoit-Wolf p (D19)
\providecommand{\rhebRtMvShareXLgbmHar}{}\renewcommand{\rhebRtMvShareXLgbmHar}{90}  % ... real-time MV: median strategy vol share, %
\providecommand{\rhebRtMvCapDaysXLgbmHar}{}\renewcommand{\rhebRtMvCapDaysXLgbmHar}{85}  % ... % of days the 90% cap binds
\providecommand{\rhebRtMvWsXLgbmHar}{}\renewcommand{\rhebRtMvWsXLgbmHar}{3.4}  % ... real-time MV: median strategy dollar weight, %
\providecommand{\rhebRtIvWsXLgbmHar}{}\renewcommand{\rhebRtIvWsXLgbmHar}{0.9}  % ... real-time IV: median strategy dollar weight, %
\providecommand{\rhebExBetaXLgbmHar}{}\renewcommand{\rhebExBetaXLgbmHar}{0.20}  % ... without the best-day union: CAPM beta
\providecommand{\rhebExBetaTXLgbmHar}{}\renewcommand{\rhebExBetaTXLgbmHar}{0.06}  % ... its HAC t
\providecommand{\rhebExAlphaXLgbmHar}{}\renewcommand{\rhebExAlphaXLgbmHar}{0.0}  % ... without the best-day union: alpha, % per day
\providecommand{\rhebExAlphaTXLgbmHar}{}\renewcommand{\rhebExAlphaTXLgbmHar}{\ensuremath{-}0.01}  % ... its HAC t
\providecommand{\rhebExAlphaPXLgbmHar}{}\renewcommand{\rhebExAlphaPXLgbmHar}{0.991}  % ... its fixed-b p
\providecommand{\rhebExARXLgbmHar}{}\renewcommand{\rhebExARXLgbmHar}{\ensuremath{-}0.01}  % ... appraisal ratio, sqrt(q) (D19)
\providecommand{\rhebExSRmaxXLgbmHar}{}\renewcommand{\rhebExSRmaxXLgbmHar}{0.23}  % ... max Sharpe ratio, STANDARD
\providecommand{\rhebExDSRmaxPXLgbmHar}{}\renewcommand{\rhebExDSRmaxPXLgbmHar}{0.984}  % ... bootstrap p of its gain
\providecommand{\rhebExGrsPXLgbmHar}{}\renewcommand{\rhebExGrsPXLgbmHar}{0.992}  % ... GRS p
\providecommand{\rhebCorrXgbHar}{}\renewcommand{\rhebCorrXgbHar}{0.00}  % pooltree_xgb_baseline, mid: correlation with the market excess return
\providecommand{\rhebBetaXgbHar}{}\renewcommand{\rhebBetaXgbHar}{0.14}  % ... CAPM beta (% of premium per % of market)
\providecommand{\rhebBetaTXgbHar}{}\renewcommand{\rhebBetaTXgbHar}{0.04}  % ... its HAC t (S = 33)
\providecommand{\rhebBetaPXgbHar}{}\renewcommand{\rhebBetaPXgbHar}{0.967}  % ... its fixed-b p
\providecommand{\rhebRsqXgbHar}{}\renewcommand{\rhebRsqXgbHar}{0.00}  % ... CAPM R^2
\providecommand{\rhebAlphaXgbHar}{}\renewcommand{\rhebAlphaXgbHar}{7.4}  % ... CAPM alpha, % of premium per day
\providecommand{\rhebAlphaTXgbHar}{}\renewcommand{\rhebAlphaTXgbHar}{2.05}  % ... its HAC t
\providecommand{\rhebAlphaPXgbHar}{}\renewcommand{\rhebAlphaPXgbHar}{0.056}  % ... its fixed-b p
\providecommand{\rhebAlphaAnnXgbHar}{}\renewcommand{\rhebAlphaAnnXgbHar}{1667}  % ... alpha x q (D19; was x 252), % of premium a year
\providecommand{\rhebARXgbHar}{}\renewcommand{\rhebARXgbHar}{0.98}  % ... appraisal ratio alpha / sigma_eps, annualized sqrt(q) (D19)
\providecommand{\rhebSRsXgbHar}{}\renewcommand{\rhebSRsXgbHar}{0.98}  % ... strategy Sharpe ratio, STANDARD = the paper's (= \rhemSh / X)
\providecommand{\rhebSRsDSixXgbHar}{}\renewcommand{\rhebSRsDSixXgbHar}{0.91}  % ... strategy Sharpe ratio, D6 (= \rhemShDSix)
\providecommand{\rhebSRmaxXgbHar}{}\renewcommand{\rhebSRmaxXgbHar}{0.99}  % ... max Sharpe ratio of market + strategy, STANDARD
\providecommand{\rhebDSRmaxXgbHar}{}\renewcommand{\rhebDSRmaxXgbHar}{0.84}  % ... its gain over the market alone, STANDARD
\providecommand{\rhebDSRmaxLoXgbHar}{}\renewcommand{\rhebDSRmaxLoXgbHar}{0.01}  % ... bootstrap 95% interval, lower end
\providecommand{\rhebDSRmaxHiXgbHar}{}\renewcommand{\rhebDSRmaxHiXgbHar}{1.71}  % ... upper end
\providecommand{\rhebDSRmaxPXgbHar}{}\renewcommand{\rhebDSRmaxPXgbHar}{0.028}  % ... bootstrap p (descriptive; Delta >= 0)
\providecommand{\rhebGrsXgbHar}{}\renewcommand{\rhebGrsXgbHar}{2.61}  % ... one-asset GRS statistic, F(1, T - 2)
\providecommand{\rhebGrsPXgbHar}{}\renewcommand{\rhebGrsPXgbHar}{0.106}  % ... its p (i.i.d. errors)
\providecommand{\rhebTanShareXgbHar}{}\renewcommand{\rhebTanShareXgbHar}{98}  % ... tangency: strategy's share of volatility, %
\providecommand{\rhebTanDSixXgbHar}{}\renewcommand{\rhebTanDSixXgbHar}{0.91}  % ... tangency (fixed in-sample weights): Sharpe ratio, D6
\providecommand{\rhebRtMvXgbHar}{}\renewcommand{\rhebRtMvXgbHar}{1.24}  % ... real-time mean-variance: Sharpe ratio, STANDARD
\providecommand{\rhebRtMvDXgbHar}{}\renewcommand{\rhebRtMvDXgbHar}{0.74}  % ... minus the market's, STANDARD
\providecommand{\rhebRtMvPXgbHar}{}\renewcommand{\rhebRtMvPXgbHar}{0.346}  % ... Ledoit-Wolf p of the difference (D19)
\providecommand{\rhebRtMvDSixXgbHar}{}\renewcommand{\rhebRtMvDSixXgbHar}{1.15}  % ... real-time mean-variance: Sharpe ratio, D6
\providecommand{\rhebRtMvDSixDXgbHar}{}\renewcommand{\rhebRtMvDSixDXgbHar}{0.47}  % ... minus the market's, D6
\providecommand{\rhebRtMvDSixPXgbHar}{}\renewcommand{\rhebRtMvDSixPXgbHar}{0.619}  % ... bootstrap p of the D6 difference
\providecommand{\rhebExRtMvXgbHar}{}\renewcommand{\rhebExRtMvXgbHar}{1.14}  % ... w/o best days, real-time mean-variance: SR, STANDARD
\providecommand{\rhebExRtMvDXgbHar}{}\renewcommand{\rhebExRtMvDXgbHar}{0.59}  % ... minus the market's, STANDARD
\providecommand{\rhebExRtMvPXgbHar}{}\renewcommand{\rhebExRtMvPXgbHar}{0.472}  % ... its Ledoit-Wolf p (D19)
\providecommand{\rhebRtIvXgbHar}{}\renewcommand{\rhebRtIvXgbHar}{1.11}  % ... real-time inverse-volatility: Sharpe ratio, STANDARD
\providecommand{\rhebRtIvDXgbHar}{}\renewcommand{\rhebRtIvDXgbHar}{0.61}  % ... minus the market's, STANDARD
\providecommand{\rhebRtIvPXgbHar}{}\renewcommand{\rhebRtIvPXgbHar}{0.273}  % ... Ledoit-Wolf p of the difference (D19)
\providecommand{\rhebRtIvDSixXgbHar}{}\renewcommand{\rhebRtIvDSixXgbHar}{1.26}  % ... real-time inverse-volatility: Sharpe ratio, D6
\providecommand{\rhebRtIvDSixDXgbHar}{}\renewcommand{\rhebRtIvDSixDXgbHar}{0.58}  % ... minus the market's, D6
\providecommand{\rhebRtIvDSixPXgbHar}{}\renewcommand{\rhebRtIvDSixPXgbHar}{0.368}  % ... bootstrap p of the D6 difference
\providecommand{\rhebExRtIvXgbHar}{}\renewcommand{\rhebExRtIvXgbHar}{1.13}  % ... w/o best days, real-time inverse-volatility: SR, STANDARD
\providecommand{\rhebExRtIvDXgbHar}{}\renewcommand{\rhebExRtIvDXgbHar}{0.59}  % ... minus the market's, STANDARD
\providecommand{\rhebExRtIvPXgbHar}{}\renewcommand{\rhebExRtIvPXgbHar}{0.301}  % ... its Ledoit-Wolf p (D19)
\providecommand{\rhebRtMvShareXgbHar}{}\renewcommand{\rhebRtMvShareXgbHar}{90}  % ... real-time MV: median strategy vol share, %
\providecommand{\rhebRtMvCapDaysXgbHar}{}\renewcommand{\rhebRtMvCapDaysXgbHar}{100}  % ... % of days the 90% cap binds
\providecommand{\rhebRtMvWsXgbHar}{}\renewcommand{\rhebRtMvWsXgbHar}{3.3}  % ... real-time MV: median strategy dollar weight, %
\providecommand{\rhebRtIvWsXgbHar}{}\renewcommand{\rhebRtIvWsXgbHar}{0.9}  % ... real-time IV: median strategy dollar weight, %
\providecommand{\rhebExBetaXgbHar}{}\renewcommand{\rhebExBetaXgbHar}{\ensuremath{-}0.34}  % ... without the best-day union: CAPM beta
\providecommand{\rhebExBetaTXgbHar}{}\renewcommand{\rhebExBetaTXgbHar}{\ensuremath{-}0.11}  % ... its HAC t
\providecommand{\rhebExAlphaXgbHar}{}\renewcommand{\rhebExAlphaXgbHar}{4.0}  % ... without the best-day union: alpha, % per day
\providecommand{\rhebExAlphaTXgbHar}{}\renewcommand{\rhebExAlphaTXgbHar}{1.11}  % ... its HAC t
\providecommand{\rhebExAlphaPXgbHar}{}\renewcommand{\rhebExAlphaPXgbHar}{0.290}  % ... its fixed-b p
\providecommand{\rhebExARXgbHar}{}\renewcommand{\rhebExARXgbHar}{0.61}  % ... appraisal ratio, sqrt(q) (D19)
\providecommand{\rhebExSRmaxXgbHar}{}\renewcommand{\rhebExSRmaxXgbHar}{0.65}  % ... max Sharpe ratio, STANDARD
\providecommand{\rhebExDSRmaxPXgbHar}{}\renewcommand{\rhebExDSRmaxPXgbHar}{0.171}  % ... bootstrap p of its gain
\providecommand{\rhebExGrsPXgbHar}{}\renewcommand{\rhebExGrsPXgbHar}{0.317}  % ... GRS p
\providecommand{\rhebTmGammaXgbHar}{}\renewcommand{\rhebTmGammaXgbHar}{\ensuremath{-}0.82}  % pooltree_xgb_baseline, mid: Treynor-Mazuy gamma on (r^m)^2 (% per %^2)
\providecommand{\rhebTmGammaTXgbHar}{}\renewcommand{\rhebTmGammaTXgbHar}{\ensuremath{-}0.79}  % ... its HAC t
\providecommand{\rhebTmGammaPXgbHar}{}\renewcommand{\rhebTmGammaPXgbHar}{0.451}  % ... its fixed-b p
\providecommand{\rhebTmBetaXgbHar}{}\renewcommand{\rhebTmBetaXgbHar}{0.11}  % ... Treynor-Mazuy beta on r^m
\providecommand{\rhebTmBetaTXgbHar}{}\renewcommand{\rhebTmBetaTXgbHar}{0.03}  % ... its HAC t
\providecommand{\rhebNLongXgbHar}{}\renewcommand{\rhebNLongXgbHar}{193}  % ... long days
\providecommand{\rhebTmGammaLongXgbHar}{}\renewcommand{\rhebTmGammaLongXgbHar}{10.47}  % ... TM gamma on its long days only
\providecommand{\rhebTmGammaLongTXgbHar}{}\renewcommand{\rhebTmGammaLongTXgbHar}{1.80}  % ... its HAC t (no test: RULE section 4)
\providecommand{\rhebNShortXgbHar}{}\renewcommand{\rhebNShortXgbHar}{421}  % ... short days
\providecommand{\rhebTmGammaShortXgbHar}{}\renewcommand{\rhebTmGammaShortXgbHar}{\ensuremath{-}1.21}  % ... TM gamma on its short days only
\providecommand{\rhebTmGammaShortTXgbHar}{}\renewcommand{\rhebTmGammaShortTXgbHar}{\ensuremath{-}1.01}  % ... its HAC t (no test: RULE section 4)
\providecommand{\rhebDownMeanXgbHar}{}\renewcommand{\rhebDownMeanXgbHar}{10.6}  % ... mean return on market-down days, % of premium
\providecommand{\rhebDownRestXgbHar}{}\renewcommand{\rhebDownRestXgbHar}{4.2}  % ... on the other days
\providecommand{\rhebDownDiffTXgbHar}{}\renewcommand{\rhebDownDiffTXgbHar}{0.67}  % ... HAC t of the difference
\providecommand{\rhebDownDiffPXgbHar}{}\renewcommand{\rhebDownDiffPXgbHar}{0.519}  % ... its fixed-b p
\providecommand{\rhebWorstMeanXgbHar}{}\renewcommand{\rhebWorstMeanXgbHar}{\ensuremath{-}2.3}  % ... mean return on worst-decile market days, % of premium
\providecommand{\rhebWorstRestXgbHar}{}\renewcommand{\rhebWorstRestXgbHar}{8.5}  % ... on the other days
\providecommand{\rhebWorstDiffTXgbHar}{}\renewcommand{\rhebWorstDiffTXgbHar}{\ensuremath{-}0.71}  % ... HAC t of the difference
\providecommand{\rhebWorstDiffPXgbHar}{}\renewcommand{\rhebWorstDiffPXgbHar}{0.499}  % ... its fixed-b p
\providecommand{\rhebThirtyCorrXgbHar}{}\renewcommand{\rhebThirtyCorrXgbHar}{0.04}  % ... correlation with the 15:30-16:00 index return
\providecommand{\rhebThirtyBetaXgbHar}{}\renewcommand{\rhebThirtyBetaXgbHar}{15.88}  % ... beta on the 15:30-16:00 index return
\providecommand{\rhebThirtyBetaTXgbHar}{}\renewcommand{\rhebThirtyBetaTXgbHar}{0.60}  % ... its HAC t
\providecommand{\rhebThirtyBetaPXgbHar}{}\renewcommand{\rhebThirtyBetaPXgbHar}{0.567}  % ... its fixed-b p
\providecommand{\rhebThirtyAlphaXgbHar}{}\renewcommand{\rhebThirtyAlphaXgbHar}{7.6}  % ... alpha against it, % per day
\providecommand{\rhebThirtyAlphaTXgbHar}{}\renewcommand{\rhebThirtyAlphaTXgbHar}{2.12}  % ... its HAC t
\providecommand{\rhebThirtyTmGammaXgbHar}{}\renewcommand{\rhebThirtyTmGammaXgbHar}{\ensuremath{-}96.47}  % ... TM gamma on the squared 15:30-16:00 return
\providecommand{\rhebThirtyTmGammaTXgbHar}{}\renewcommand{\rhebThirtyTmGammaTXgbHar}{\ensuremath{-}2.07}  % ... its HAC t
\providecommand{\rhebThirtyTmGammaPXgbHar}{}\renewcommand{\rhebThirtyTmGammaPXgbHar}{0.053}  % ... its fixed-b p
\providecommand{\rhebThirtyTmGammaLongXgbHar}{}\renewcommand{\rhebThirtyTmGammaLongXgbHar}{396.11}  % ... ... on its long days only
\providecommand{\rhebThirtyTmGammaLongTXgbHar}{}\renewcommand{\rhebThirtyTmGammaLongTXgbHar}{2.68}  % ... its HAC t (no test)
\providecommand{\rhebThirtyTmGammaShortXgbHar}{}\renewcommand{\rhebThirtyTmGammaShortXgbHar}{\ensuremath{-}266.36}  % ... ... on its short days only
\providecommand{\rhebThirtyTmGammaShortTXgbHar}{}\renewcommand{\rhebThirtyTmGammaShortTXgbHar}{\ensuremath{-}4.42}  % ... its HAC t (no test)
\providecommand{\rhebThirtyDownMeanXgbHar}{}\renewcommand{\rhebThirtyDownMeanXgbHar}{8.1}  % ... mean on 15:30-16:00 down days, %
\providecommand{\rhebThirtyDownDiffTXgbHar}{}\renewcommand{\rhebThirtyDownDiffTXgbHar}{0.18}  % ... HAC t of the difference with the other days
\providecommand{\rhebThirtyWorstMeanXgbHar}{}\renewcommand{\rhebThirtyWorstMeanXgbHar}{\ensuremath{-}45.2}  % ... mean on 15:30-16:00 worst days, %
\providecommand{\rhebThirtyWorstDiffTXgbHar}{}\renewcommand{\rhebThirtyWorstDiffTXgbHar}{\ensuremath{-}2.16}  % ... HAC t of the difference with the other days
\providecommand{\rhebCorrXXgbHar}{}\renewcommand{\rhebCorrXXgbHar}{0.00}  % pooltree_xgb_baseline, bid/ask: correlation with the market excess return
\providecommand{\rhebBetaXXgbHar}{}\renewcommand{\rhebBetaXXgbHar}{0.22}  % ... CAPM beta (% of premium per % of market)
\providecommand{\rhebBetaTXXgbHar}{}\renewcommand{\rhebBetaTXXgbHar}{0.07}  % ... its HAC t (S = 33)
\providecommand{\rhebBetaPXXgbHar}{}\renewcommand{\rhebBetaPXXgbHar}{0.950}  % ... its fixed-b p
\providecommand{\rhebRsqXXgbHar}{}\renewcommand{\rhebRsqXXgbHar}{0.00}  % ... CAPM R^2
\providecommand{\rhebAlphaXXgbHar}{}\renewcommand{\rhebAlphaXXgbHar}{4.8}  % ... CAPM alpha, % of premium per day
\providecommand{\rhebAlphaTXXgbHar}{}\renewcommand{\rhebAlphaTXXgbHar}{1.31}  % ... its HAC t
\providecommand{\rhebAlphaPXXgbHar}{}\renewcommand{\rhebAlphaPXXgbHar}{0.213}  % ... its fixed-b p
\providecommand{\rhebAlphaAnnXXgbHar}{}\renewcommand{\rhebAlphaAnnXXgbHar}{1076}  % ... alpha x q (D19; was x 252), % of premium a year
\providecommand{\rhebARXXgbHar}{}\renewcommand{\rhebARXXgbHar}{0.63}  % ... appraisal ratio alpha / sigma_eps, annualized sqrt(q) (D19)
\providecommand{\rhebSRsXXgbHar}{}\renewcommand{\rhebSRsXXgbHar}{0.63}  % ... strategy Sharpe ratio, STANDARD = the paper's (= \rhemSh / X)
\providecommand{\rhebSRsDSixXXgbHar}{}\renewcommand{\rhebSRsDSixXXgbHar}{0.55}  % ... strategy Sharpe ratio, D6 (= \rhemShDSix)
\providecommand{\rhebSRmaxXXgbHar}{}\renewcommand{\rhebSRmaxXXgbHar}{0.65}  % ... max Sharpe ratio of market + strategy, STANDARD
\providecommand{\rhebDSRmaxXXgbHar}{}\renewcommand{\rhebDSRmaxXXgbHar}{0.50}  % ... its gain over the market alone, STANDARD
\providecommand{\rhebDSRmaxLoXXgbHar}{}\renewcommand{\rhebDSRmaxLoXXgbHar}{0.00}  % ... bootstrap 95% interval, lower end
\providecommand{\rhebDSRmaxHiXXgbHar}{}\renewcommand{\rhebDSRmaxHiXXgbHar}{1.38}  % ... upper end
\providecommand{\rhebDSRmaxPXXgbHar}{}\renewcommand{\rhebDSRmaxPXXgbHar}{0.106}  % ... bootstrap p (descriptive; Delta >= 0)
\providecommand{\rhebGrsXXgbHar}{}\renewcommand{\rhebGrsXXgbHar}{1.09}  % ... one-asset GRS statistic, F(1, T - 2)
\providecommand{\rhebGrsPXXgbHar}{}\renewcommand{\rhebGrsPXXgbHar}{0.297}  % ... its p (i.i.d. errors)
\providecommand{\rhebTanShareXXgbHar}{}\renewcommand{\rhebTanShareXXgbHar}{95}  % ... tangency: strategy's share of volatility, %
\providecommand{\rhebTanDSixXXgbHar}{}\renewcommand{\rhebTanDSixXXgbHar}{0.57}  % ... tangency (fixed in-sample weights): Sharpe ratio, D6
\providecommand{\rhebRtMvXXgbHar}{}\renewcommand{\rhebRtMvXXgbHar}{0.89}  % ... real-time mean-variance: Sharpe ratio, STANDARD
\providecommand{\rhebRtMvDXXgbHar}{}\renewcommand{\rhebRtMvDXXgbHar}{0.40}  % ... minus the market's, STANDARD
\providecommand{\rhebRtMvPXXgbHar}{}\renewcommand{\rhebRtMvPXXgbHar}{0.605}  % ... Ledoit-Wolf p of the difference (D19)
\providecommand{\rhebRtMvDSixXXgbHar}{}\renewcommand{\rhebRtMvDSixXXgbHar}{0.80}  % ... real-time mean-variance: Sharpe ratio, D6
\providecommand{\rhebRtMvDSixDXXgbHar}{}\renewcommand{\rhebRtMvDSixDXXgbHar}{0.12}  % ... minus the market's, D6
\providecommand{\rhebRtMvDSixPXXgbHar}{}\renewcommand{\rhebRtMvDSixPXXgbHar}{0.894}  % ... bootstrap p of the D6 difference
\providecommand{\rhebExRtMvXXgbHar}{}\renewcommand{\rhebExRtMvXXgbHar}{0.54}  % ... w/o best days, real-time mean-variance: SR, STANDARD
\providecommand{\rhebExRtMvDXXgbHar}{}\renewcommand{\rhebExRtMvDXXgbHar}{\ensuremath{-}0.01}  % ... minus the market's, STANDARD
\providecommand{\rhebExRtMvPXXgbHar}{}\renewcommand{\rhebExRtMvPXXgbHar}{0.994}  % ... its Ledoit-Wolf p (D19)
\providecommand{\rhebRtIvXXgbHar}{}\renewcommand{\rhebRtIvXXgbHar}{0.84}  % ... real-time inverse-volatility: Sharpe ratio, STANDARD
\providecommand{\rhebRtIvDXXgbHar}{}\renewcommand{\rhebRtIvDXXgbHar}{0.34}  % ... minus the market's, STANDARD
\providecommand{\rhebRtIvPXXgbHar}{}\renewcommand{\rhebRtIvPXXgbHar}{0.528}  % ... Ledoit-Wolf p of the difference (D19)
\providecommand{\rhebRtIvDSixXXgbHar}{}\renewcommand{\rhebRtIvDSixXXgbHar}{0.99}  % ... real-time inverse-volatility: Sharpe ratio, D6
\providecommand{\rhebRtIvDSixDXXgbHar}{}\renewcommand{\rhebRtIvDSixDXXgbHar}{0.31}  % ... minus the market's, D6
\providecommand{\rhebRtIvDSixPXXgbHar}{}\renewcommand{\rhebRtIvDSixPXXgbHar}{0.626}  % ... bootstrap p of the D6 difference
\providecommand{\rhebExRtIvXXgbHar}{}\renewcommand{\rhebExRtIvXXgbHar}{0.86}  % ... w/o best days, real-time inverse-volatility: SR, STANDARD
\providecommand{\rhebExRtIvDXXgbHar}{}\renewcommand{\rhebExRtIvDXXgbHar}{0.32}  % ... minus the market's, STANDARD
\providecommand{\rhebExRtIvPXXgbHar}{}\renewcommand{\rhebExRtIvPXXgbHar}{0.571}  % ... its Ledoit-Wolf p (D19)
\providecommand{\rhebRtMvShareXXgbHar}{}\renewcommand{\rhebRtMvShareXXgbHar}{90}  % ... real-time MV: median strategy vol share, %
\providecommand{\rhebRtMvCapDaysXXgbHar}{}\renewcommand{\rhebRtMvCapDaysXXgbHar}{95}  % ... % of days the 90% cap binds
\providecommand{\rhebRtMvWsXXgbHar}{}\renewcommand{\rhebRtMvWsXXgbHar}{3.4}  % ... real-time MV: median strategy dollar weight, %
\providecommand{\rhebRtIvWsXXgbHar}{}\renewcommand{\rhebRtIvWsXXgbHar}{0.9}  % ... real-time IV: median strategy dollar weight, %
\providecommand{\rhebExBetaXXgbHar}{}\renewcommand{\rhebExBetaXXgbHar}{\ensuremath{-}0.31}  % ... without the best-day union: CAPM beta
\providecommand{\rhebExBetaTXXgbHar}{}\renewcommand{\rhebExBetaTXXgbHar}{\ensuremath{-}0.10}  % ... its HAC t
\providecommand{\rhebExAlphaXXgbHar}{}\renewcommand{\rhebExAlphaXXgbHar}{1.6}  % ... without the best-day union: alpha, % per day
\providecommand{\rhebExAlphaTXXgbHar}{}\renewcommand{\rhebExAlphaTXXgbHar}{0.43}  % ... its HAC t
\providecommand{\rhebExAlphaPXXgbHar}{}\renewcommand{\rhebExAlphaPXXgbHar}{0.678}  % ... its fixed-b p
\providecommand{\rhebExARXXgbHar}{}\renewcommand{\rhebExARXXgbHar}{0.24}  % ... appraisal ratio, sqrt(q) (D19)
\providecommand{\rhebExSRmaxXXgbHar}{}\renewcommand{\rhebExSRmaxXXgbHar}{0.33}  % ... max Sharpe ratio, STANDARD
\providecommand{\rhebExDSRmaxPXXgbHar}{}\renewcommand{\rhebExDSRmaxPXXgbHar}{0.457}  % ... bootstrap p of its gain
\providecommand{\rhebExGrsPXXgbHar}{}\renewcommand{\rhebExGrsPXXgbHar}{0.695}  % ... GRS p
\providecommand{\rhebCorrRfHar}{}\renewcommand{\rhebCorrRfHar}{0.02}  % pooltree_rf_baseline, mid: correlation with the market excess return
\providecommand{\rhebBetaRfHar}{}\renewcommand{\rhebBetaRfHar}{2.44}  % ... CAPM beta (% of premium per % of market)
\providecommand{\rhebBetaTRfHar}{}\renewcommand{\rhebBetaTRfHar}{0.73}  % ... its HAC t (S = 33)
\providecommand{\rhebBetaPRfHar}{}\renewcommand{\rhebBetaPRfHar}{0.483}  % ... its fixed-b p
\providecommand{\rhebRsqRfHar}{}\renewcommand{\rhebRsqRfHar}{0.00}  % ... CAPM R^2
\providecommand{\rhebAlphaRfHar}{}\renewcommand{\rhebAlphaRfHar}{2.4}  % ... CAPM alpha, % of premium per day
\providecommand{\rhebAlphaTRfHar}{}\renewcommand{\rhebAlphaTRfHar}{0.61}  % ... its HAC t
\providecommand{\rhebAlphaPRfHar}{}\renewcommand{\rhebAlphaPRfHar}{0.556}  % ... its fixed-b p
\providecommand{\rhebAlphaAnnRfHar}{}\renewcommand{\rhebAlphaAnnRfHar}{535}  % ... alpha x q (D19; was x 252), % of premium a year
\providecommand{\rhebARRfHar}{}\renewcommand{\rhebARRfHar}{0.31}  % ... appraisal ratio alpha / sigma_eps, annualized sqrt(q) (D19)
\providecommand{\rhebSRsRfHar}{}\renewcommand{\rhebSRsRfHar}{0.32}  % ... strategy Sharpe ratio, STANDARD = the paper's (= \rhemSh / X)
\providecommand{\rhebSRsDSixRfHar}{}\renewcommand{\rhebSRsDSixRfHar}{0.01}  % ... strategy Sharpe ratio, D6 (= \rhemShDSix)
\providecommand{\rhebSRmaxRfHar}{}\renewcommand{\rhebSRmaxRfHar}{0.35}  % ... max Sharpe ratio of market + strategy, STANDARD
\providecommand{\rhebDSRmaxRfHar}{}\renewcommand{\rhebDSRmaxRfHar}{0.20}  % ... its gain over the market alone, STANDARD
\providecommand{\rhebDSRmaxLoRfHar}{}\renewcommand{\rhebDSRmaxLoRfHar}{0.00}  % ... bootstrap 95% interval, lower end
\providecommand{\rhebDSRmaxHiRfHar}{}\renewcommand{\rhebDSRmaxHiRfHar}{1.22}  % ... upper end
\providecommand{\rhebDSRmaxPRfHar}{}\renewcommand{\rhebDSRmaxPRfHar}{0.302}  % ... bootstrap p (descriptive; Delta >= 0)
\providecommand{\rhebGrsRfHar}{}\renewcommand{\rhebGrsRfHar}{0.27}  % ... one-asset GRS statistic, F(1, T - 2)
\providecommand{\rhebGrsPRfHar}{}\renewcommand{\rhebGrsPRfHar}{0.605}  % ... its p (i.i.d. errors)
\providecommand{\rhebTanShareRfHar}{}\renewcommand{\rhebTanShareRfHar}{82}  % ... tangency: strategy's share of volatility, %
\providecommand{\rhebTanDSixRfHar}{}\renewcommand{\rhebTanDSixRfHar}{0.02}  % ... tangency (fixed in-sample weights): Sharpe ratio, D6
\providecommand{\rhebRtMvRfHar}{}\renewcommand{\rhebRtMvRfHar}{0.74}  % ... real-time mean-variance: Sharpe ratio, STANDARD
\providecommand{\rhebRtMvDRfHar}{}\renewcommand{\rhebRtMvDRfHar}{0.24}  % ... minus the market's, STANDARD
\providecommand{\rhebRtMvPRfHar}{}\renewcommand{\rhebRtMvPRfHar}{0.776}  % ... Ledoit-Wolf p of the difference (D19)
\providecommand{\rhebRtMvDSixRfHar}{}\renewcommand{\rhebRtMvDSixRfHar}{0.32}  % ... real-time mean-variance: Sharpe ratio, D6
\providecommand{\rhebRtMvDSixDRfHar}{}\renewcommand{\rhebRtMvDSixDRfHar}{\ensuremath{-}0.36}  % ... minus the market's, D6
\providecommand{\rhebRtMvDSixPRfHar}{}\renewcommand{\rhebRtMvDSixPRfHar}{0.696}  % ... bootstrap p of the D6 difference
\providecommand{\rhebExRtMvRfHar}{}\renewcommand{\rhebExRtMvRfHar}{0.88}  % ... w/o best days, real-time mean-variance: SR, STANDARD
\providecommand{\rhebExRtMvDRfHar}{}\renewcommand{\rhebExRtMvDRfHar}{0.33}  % ... minus the market's, STANDARD
\providecommand{\rhebExRtMvPRfHar}{}\renewcommand{\rhebExRtMvPRfHar}{0.692}  % ... its Ledoit-Wolf p (D19)
\providecommand{\rhebRtIvRfHar}{}\renewcommand{\rhebRtIvRfHar}{1.02}  % ... real-time inverse-volatility: Sharpe ratio, STANDARD
\providecommand{\rhebRtIvDRfHar}{}\renewcommand{\rhebRtIvDRfHar}{0.53}  % ... minus the market's, STANDARD
\providecommand{\rhebRtIvPRfHar}{}\renewcommand{\rhebRtIvPRfHar}{0.355}  % ... Ledoit-Wolf p of the difference (D19)
\providecommand{\rhebRtIvDSixRfHar}{}\renewcommand{\rhebRtIvDSixRfHar}{0.92}  % ... real-time inverse-volatility: Sharpe ratio, D6
\providecommand{\rhebRtIvDSixDRfHar}{}\renewcommand{\rhebRtIvDSixDRfHar}{0.24}  % ... minus the market's, D6
\providecommand{\rhebRtIvDSixPRfHar}{}\renewcommand{\rhebRtIvDSixPRfHar}{0.712}  % ... bootstrap p of the D6 difference
\providecommand{\rhebExRtIvRfHar}{}\renewcommand{\rhebExRtIvRfHar}{1.07}  % ... w/o best days, real-time inverse-volatility: SR, STANDARD
\providecommand{\rhebExRtIvDRfHar}{}\renewcommand{\rhebExRtIvDRfHar}{0.53}  % ... minus the market's, STANDARD
\providecommand{\rhebExRtIvPRfHar}{}\renewcommand{\rhebExRtIvPRfHar}{0.362}  % ... its Ledoit-Wolf p (D19)
\providecommand{\rhebRtMvShareRfHar}{}\renewcommand{\rhebRtMvShareRfHar}{90}  % ... real-time MV: median strategy vol share, %
\providecommand{\rhebRtMvCapDaysRfHar}{}\renewcommand{\rhebRtMvCapDaysRfHar}{86}  % ... % of days the 90% cap binds
\providecommand{\rhebRtMvWsRfHar}{}\renewcommand{\rhebRtMvWsRfHar}{3.5}  % ... real-time MV: median strategy dollar weight, %
\providecommand{\rhebRtIvWsRfHar}{}\renewcommand{\rhebRtIvWsRfHar}{0.9}  % ... real-time IV: median strategy dollar weight, %
\providecommand{\rhebExBetaRfHar}{}\renewcommand{\rhebExBetaRfHar}{0.68}  % ... without the best-day union: CAPM beta
\providecommand{\rhebExBetaTRfHar}{}\renewcommand{\rhebExBetaTRfHar}{0.22}  % ... its HAC t
\providecommand{\rhebExAlphaRfHar}{}\renewcommand{\rhebExAlphaRfHar}{3.2}  % ... without the best-day union: alpha, % per day
\providecommand{\rhebExAlphaTRfHar}{}\renewcommand{\rhebExAlphaTRfHar}{0.91}  % ... its HAC t
\providecommand{\rhebExAlphaPRfHar}{}\renewcommand{\rhebExAlphaPRfHar}{0.382}  % ... its fixed-b p
\providecommand{\rhebExARRfHar}{}\renewcommand{\rhebExARRfHar}{0.49}  % ... appraisal ratio, sqrt(q) (D19)
\providecommand{\rhebExSRmaxRfHar}{}\renewcommand{\rhebExSRmaxRfHar}{0.54}  % ... max Sharpe ratio, STANDARD
\providecommand{\rhebExDSRmaxPRfHar}{}\renewcommand{\rhebExDSRmaxPRfHar}{0.239}  % ... bootstrap p of its gain
\providecommand{\rhebExGrsPRfHar}{}\renewcommand{\rhebExGrsPRfHar}{0.422}  % ... GRS p
\providecommand{\rhebTmGammaRfHar}{}\renewcommand{\rhebTmGammaRfHar}{0.10}  % pooltree_rf_baseline, mid: Treynor-Mazuy gamma on (r^m)^2 (% per %^2)
\providecommand{\rhebTmGammaTRfHar}{}\renewcommand{\rhebTmGammaTRfHar}{0.09}  % ... its HAC t
\providecommand{\rhebTmGammaPRfHar}{}\renewcommand{\rhebTmGammaPRfHar}{0.933}  % ... its fixed-b p
\providecommand{\rhebTmBetaRfHar}{}\renewcommand{\rhebTmBetaRfHar}{2.45}  % ... Treynor-Mazuy beta on r^m
\providecommand{\rhebTmBetaTRfHar}{}\renewcommand{\rhebTmBetaTRfHar}{0.74}  % ... its HAC t
\providecommand{\rhebNLongRfHar}{}\renewcommand{\rhebNLongRfHar}{169}  % ... long days
\providecommand{\rhebTmGammaLongRfHar}{}\renewcommand{\rhebTmGammaLongRfHar}{10.90}  % ... TM gamma on its long days only
\providecommand{\rhebTmGammaLongTRfHar}{}\renewcommand{\rhebTmGammaLongTRfHar}{1.83}  % ... its HAC t (no test: RULE section 4)
\providecommand{\rhebNShortRfHar}{}\renewcommand{\rhebNShortRfHar}{445}  % ... short days
\providecommand{\rhebTmGammaShortRfHar}{}\renewcommand{\rhebTmGammaShortRfHar}{\ensuremath{-}0.53}  % ... TM gamma on its short days only
\providecommand{\rhebTmGammaShortTRfHar}{}\renewcommand{\rhebTmGammaShortTRfHar}{\ensuremath{-}0.40}  % ... its HAC t (no test: RULE section 4)
\providecommand{\rhebDownMeanRfHar}{}\renewcommand{\rhebDownMeanRfHar}{\ensuremath{-}0.5}  % ... mean return on market-down days, % of premium
\providecommand{\rhebDownRestRfHar}{}\renewcommand{\rhebDownRestRfHar}{5.3}  % ... on the other days
\providecommand{\rhebDownDiffTRfHar}{}\renewcommand{\rhebDownDiffTRfHar}{\ensuremath{-}0.57}  % ... HAC t of the difference
\providecommand{\rhebDownDiffPRfHar}{}\renewcommand{\rhebDownDiffPRfHar}{0.585}  % ... its fixed-b p
\providecommand{\rhebWorstMeanRfHar}{}\renewcommand{\rhebWorstMeanRfHar}{\ensuremath{-}5.6}  % ... mean return on worst-decile market days, % of premium
\providecommand{\rhebWorstRestRfHar}{}\renewcommand{\rhebWorstRestRfHar}{3.3}  % ... on the other days
\providecommand{\rhebWorstDiffTRfHar}{}\renewcommand{\rhebWorstDiffTRfHar}{\ensuremath{-}0.60}  % ... HAC t of the difference
\providecommand{\rhebWorstDiffPRfHar}{}\renewcommand{\rhebWorstDiffPRfHar}{0.568}  % ... its fixed-b p
\providecommand{\rhebThirtyCorrRfHar}{}\renewcommand{\rhebThirtyCorrRfHar}{0.14}  % ... correlation with the 15:30-16:00 index return
\providecommand{\rhebThirtyBetaRfHar}{}\renewcommand{\rhebThirtyBetaRfHar}{55.15}  % ... beta on the 15:30-16:00 index return
\providecommand{\rhebThirtyBetaTRfHar}{}\renewcommand{\rhebThirtyBetaTRfHar}{1.93}  % ... its HAC t
\providecommand{\rhebThirtyBetaPRfHar}{}\renewcommand{\rhebThirtyBetaPRfHar}{0.070}  % ... its fixed-b p
\providecommand{\rhebThirtyAlphaRfHar}{}\renewcommand{\rhebThirtyAlphaRfHar}{3.1}  % ... alpha against it, % per day
\providecommand{\rhebThirtyAlphaTRfHar}{}\renewcommand{\rhebThirtyAlphaTRfHar}{0.79}  % ... its HAC t
\providecommand{\rhebThirtyTmGammaRfHar}{}\renewcommand{\rhebThirtyTmGammaRfHar}{\ensuremath{-}149.56}  % ... TM gamma on the squared 15:30-16:00 return
\providecommand{\rhebThirtyTmGammaTRfHar}{}\renewcommand{\rhebThirtyTmGammaTRfHar}{\ensuremath{-}3.09}  % ... its HAC t
\providecommand{\rhebThirtyTmGammaPRfHar}{}\renewcommand{\rhebThirtyTmGammaPRfHar}{0.005}  % ... its fixed-b p
\providecommand{\rhebThirtyTmGammaLongRfHar}{}\renewcommand{\rhebThirtyTmGammaLongRfHar}{250.00}  % ... ... on its long days only
\providecommand{\rhebThirtyTmGammaLongTRfHar}{}\renewcommand{\rhebThirtyTmGammaLongTRfHar}{2.28}  % ... its HAC t (no test)
\providecommand{\rhebThirtyTmGammaShortRfHar}{}\renewcommand{\rhebThirtyTmGammaShortRfHar}{\ensuremath{-}291.78}  % ... ... on its short days only
\providecommand{\rhebThirtyTmGammaShortTRfHar}{}\renewcommand{\rhebThirtyTmGammaShortTRfHar}{\ensuremath{-}3.80}  % ... its HAC t (no test)
\providecommand{\rhebThirtyDownMeanRfHar}{}\renewcommand{\rhebThirtyDownMeanRfHar}{\ensuremath{-}0.5}  % ... mean on 15:30-16:00 down days, %
\providecommand{\rhebThirtyDownDiffTRfHar}{}\renewcommand{\rhebThirtyDownDiffTRfHar}{\ensuremath{-}0.67}  % ... HAC t of the difference with the other days
\providecommand{\rhebThirtyWorstMeanRfHar}{}\renewcommand{\rhebThirtyWorstMeanRfHar}{\ensuremath{-}87.3}  % ... mean on 15:30-16:00 worst days, %
\providecommand{\rhebThirtyWorstDiffTRfHar}{}\renewcommand{\rhebThirtyWorstDiffTRfHar}{\ensuremath{-}3.86}  % ... HAC t of the difference with the other days
\providecommand{\rhebCorrXRfHar}{}\renewcommand{\rhebCorrXRfHar}{0.02}  % pooltree_rf_baseline, bid/ask: correlation with the market excess return
\providecommand{\rhebBetaXRfHar}{}\renewcommand{\rhebBetaXRfHar}{2.53}  % ... CAPM beta (% of premium per % of market)
\providecommand{\rhebBetaTXRfHar}{}\renewcommand{\rhebBetaTXRfHar}{0.74}  % ... its HAC t (S = 33)
\providecommand{\rhebBetaPXRfHar}{}\renewcommand{\rhebBetaPXRfHar}{0.477}  % ... its fixed-b p
\providecommand{\rhebRsqXRfHar}{}\renewcommand{\rhebRsqXRfHar}{0.00}  % ... CAPM R^2
\providecommand{\rhebAlphaXRfHar}{}\renewcommand{\rhebAlphaXRfHar}{\ensuremath{-}0.3}  % ... CAPM alpha, % of premium per day
\providecommand{\rhebAlphaTXRfHar}{}\renewcommand{\rhebAlphaTXRfHar}{\ensuremath{-}0.07}  % ... its HAC t
\providecommand{\rhebAlphaPXRfHar}{}\renewcommand{\rhebAlphaPXRfHar}{0.949}  % ... its fixed-b p
\providecommand{\rhebAlphaAnnXRfHar}{}\renewcommand{\rhebAlphaAnnXRfHar}{\ensuremath{-}60}  % ... alpha x q (D19; was x 252), % of premium a year
\providecommand{\rhebARXRfHar}{}\renewcommand{\rhebARXRfHar}{\ensuremath{-}0.03}  % ... appraisal ratio alpha / sigma_eps, annualized sqrt(q) (D19)
\providecommand{\rhebSRsXRfHar}{}\renewcommand{\rhebSRsXRfHar}{\ensuremath{-}0.03}  % ... strategy Sharpe ratio, STANDARD = the paper's (= \rhemSh / X)
\providecommand{\rhebSRsDSixXRfHar}{}\renewcommand{\rhebSRsDSixXRfHar}{\ensuremath{-}0.32}  % ... strategy Sharpe ratio, D6 (= \rhemShDSix)
\providecommand{\rhebSRmaxXRfHar}{}\renewcommand{\rhebSRmaxXRfHar}{0.15}  % ... max Sharpe ratio of market + strategy, STANDARD
\providecommand{\rhebDSRmaxXRfHar}{}\renewcommand{\rhebDSRmaxXRfHar}{0.00}  % ... its gain over the market alone, STANDARD
\providecommand{\rhebDSRmaxLoXRfHar}{}\renewcommand{\rhebDSRmaxLoXRfHar}{0.00}  % ... bootstrap 95% interval, lower end
\providecommand{\rhebDSRmaxHiXRfHar}{}\renewcommand{\rhebDSRmaxHiXRfHar}{0.98}  % ... upper end
\providecommand{\rhebDSRmaxPXRfHar}{}\renewcommand{\rhebDSRmaxPXRfHar}{0.892}  % ... bootstrap p (descriptive; Delta >= 0)
\providecommand{\rhebGrsXRfHar}{}\renewcommand{\rhebGrsXRfHar}{0.00}  % ... one-asset GRS statistic, F(1, T - 2)
\providecommand{\rhebGrsPXRfHar}{}\renewcommand{\rhebGrsPXRfHar}{0.955}  % ... its p (i.i.d. errors)
\providecommand{\rhebTanShareXRfHar}{}\renewcommand{\rhebTanShareXRfHar}{4}  % ... tangency: strategy's share of volatility, %
\providecommand{\rhebTanDSixXRfHar}{}\renewcommand{\rhebTanDSixXRfHar}{0.74}  % ... tangency (fixed in-sample weights): Sharpe ratio, D6
\providecommand{\rhebRtMvXRfHar}{}\renewcommand{\rhebRtMvXRfHar}{0.26}  % ... real-time mean-variance: Sharpe ratio, STANDARD
\providecommand{\rhebRtMvDXRfHar}{}\renewcommand{\rhebRtMvDXRfHar}{\ensuremath{-}0.24}  % ... minus the market's, STANDARD
\providecommand{\rhebRtMvPXRfHar}{}\renewcommand{\rhebRtMvPXRfHar}{0.776}  % ... Ledoit-Wolf p of the difference (D19)
\providecommand{\rhebRtMvDSixXRfHar}{}\renewcommand{\rhebRtMvDSixXRfHar}{0.15}  % ... real-time mean-variance: Sharpe ratio, D6
\providecommand{\rhebRtMvDSixDXRfHar}{}\renewcommand{\rhebRtMvDSixDXRfHar}{\ensuremath{-}0.53}  % ... minus the market's, D6
\providecommand{\rhebRtMvDSixPXRfHar}{}\renewcommand{\rhebRtMvDSixPXRfHar}{0.408}  % ... bootstrap p of the D6 difference
\providecommand{\rhebExRtMvXRfHar}{}\renewcommand{\rhebExRtMvXRfHar}{0.64}  % ... w/o best days, real-time mean-variance: SR, STANDARD
\providecommand{\rhebExRtMvDXRfHar}{}\renewcommand{\rhebExRtMvDXRfHar}{0.10}  % ... minus the market's, STANDARD
\providecommand{\rhebExRtMvPXRfHar}{}\renewcommand{\rhebExRtMvPXRfHar}{0.905}  % ... its Ledoit-Wolf p (D19)
\providecommand{\rhebRtIvXRfHar}{}\renewcommand{\rhebRtIvXRfHar}{0.75}  % ... real-time inverse-volatility: Sharpe ratio, STANDARD
\providecommand{\rhebRtIvDXRfHar}{}\renewcommand{\rhebRtIvDXRfHar}{0.26}  % ... minus the market's, STANDARD
\providecommand{\rhebRtIvPXRfHar}{}\renewcommand{\rhebRtIvPXRfHar}{0.649}  % ... Ledoit-Wolf p of the difference (D19)
\providecommand{\rhebRtIvDSixXRfHar}{}\renewcommand{\rhebRtIvDSixXRfHar}{0.66}  % ... real-time inverse-volatility: Sharpe ratio, D6
\providecommand{\rhebRtIvDSixDXRfHar}{}\renewcommand{\rhebRtIvDSixDXRfHar}{\ensuremath{-}0.03}  % ... minus the market's, D6
\providecommand{\rhebRtIvDSixPXRfHar}{}\renewcommand{\rhebRtIvDSixPXRfHar}{0.971}  % ... bootstrap p of the D6 difference
\providecommand{\rhebExRtIvXRfHar}{}\renewcommand{\rhebExRtIvXRfHar}{0.80}  % ... w/o best days, real-time inverse-volatility: SR, STANDARD
\providecommand{\rhebExRtIvDXRfHar}{}\renewcommand{\rhebExRtIvDXRfHar}{0.26}  % ... minus the market's, STANDARD
\providecommand{\rhebExRtIvPXRfHar}{}\renewcommand{\rhebExRtIvPXRfHar}{0.641}  % ... its Ledoit-Wolf p (D19)
\providecommand{\rhebRtMvShareXRfHar}{}\renewcommand{\rhebRtMvShareXRfHar}{90}  % ... real-time MV: median strategy vol share, %
\providecommand{\rhebRtMvCapDaysXRfHar}{}\renewcommand{\rhebRtMvCapDaysXRfHar}{97}  % ... % of days the 90% cap binds
\providecommand{\rhebRtMvWsXRfHar}{}\renewcommand{\rhebRtMvWsXRfHar}{3.3}  % ... real-time MV: median strategy dollar weight, %
\providecommand{\rhebRtIvWsXRfHar}{}\renewcommand{\rhebRtIvWsXRfHar}{0.8}  % ... real-time IV: median strategy dollar weight, %
\providecommand{\rhebExBetaXRfHar}{}\renewcommand{\rhebExBetaXRfHar}{0.71}  % ... without the best-day union: CAPM beta
\providecommand{\rhebExBetaTXRfHar}{}\renewcommand{\rhebExBetaTXRfHar}{0.22}  % ... its HAC t
\providecommand{\rhebExAlphaXRfHar}{}\renewcommand{\rhebExAlphaXRfHar}{0.8}  % ... without the best-day union: alpha, % per day
\providecommand{\rhebExAlphaTXRfHar}{}\renewcommand{\rhebExAlphaTXRfHar}{0.22}  % ... its HAC t
\providecommand{\rhebExAlphaPXRfHar}{}\renewcommand{\rhebExAlphaPXRfHar}{0.832}  % ... its fixed-b p
\providecommand{\rhebExARXRfHar}{}\renewcommand{\rhebExARXRfHar}{0.12}  % ... appraisal ratio, sqrt(q) (D19)
\providecommand{\rhebExSRmaxXRfHar}{}\renewcommand{\rhebExSRmaxXRfHar}{0.26}  % ... max Sharpe ratio, STANDARD
\providecommand{\rhebExDSRmaxPXRfHar}{}\renewcommand{\rhebExDSRmaxPXRfHar}{0.706}  % ... bootstrap p of its gain
\providecommand{\rhebExGrsPXRfHar}{}\renewcommand{\rhebExGrsPXRfHar}{0.846}  % ... GRS p
\providecommand{\rhebCorrCiHar}{}\renewcommand{\rhebCorrCiHar}{0.05}  % sub_ols_baseline, mid: correlation with the market excess return
\providecommand{\rhebBetaCiHar}{}\renewcommand{\rhebBetaCiHar}{5.07}  % ... CAPM beta (% of premium per % of market)
\providecommand{\rhebBetaTCiHar}{}\renewcommand{\rhebBetaTCiHar}{1.23}  % ... its HAC t (S = 33)
\providecommand{\rhebBetaPCiHar}{}\renewcommand{\rhebBetaPCiHar}{0.243}  % ... its fixed-b p
\providecommand{\rhebRsqCiHar}{}\renewcommand{\rhebRsqCiHar}{0.00}  % ... CAPM R^2
\providecommand{\rhebAlphaCiHar}{}\renewcommand{\rhebAlphaCiHar}{8.3}  % ... CAPM alpha, % of premium per day
\providecommand{\rhebAlphaTCiHar}{}\renewcommand{\rhebAlphaTCiHar}{2.33}  % ... its HAC t
\providecommand{\rhebAlphaPCiHar}{}\renewcommand{\rhebAlphaPCiHar}{0.031}  % ... its fixed-b p
\providecommand{\rhebAlphaAnnCiHar}{}\renewcommand{\rhebAlphaAnnCiHar}{1869}  % ... alpha x q (D19; was x 252), % of premium a year
\providecommand{\rhebARCiHar}{}\renewcommand{\rhebARCiHar}{1.10}  % ... appraisal ratio alpha / sigma_eps, annualized sqrt(q) (D19)
\providecommand{\rhebSRsCiHar}{}\renewcommand{\rhebSRsCiHar}{1.11}  % ... strategy Sharpe ratio, STANDARD = the paper's (= \rhemSh / X)
\providecommand{\rhebSRsDSixCiHar}{}\renewcommand{\rhebSRsDSixCiHar}{0.99}  % ... strategy Sharpe ratio, D6 (= \rhemShDSix)
\providecommand{\rhebSRmaxCiHar}{}\renewcommand{\rhebSRmaxCiHar}{1.11}  % ... max Sharpe ratio of market + strategy, STANDARD
\providecommand{\rhebDSRmaxCiHar}{}\renewcommand{\rhebDSRmaxCiHar}{0.96}  % ... its gain over the market alone, STANDARD
\providecommand{\rhebDSRmaxLoCiHar}{}\renewcommand{\rhebDSRmaxLoCiHar}{0.01}  % ... bootstrap 95% interval, lower end
\providecommand{\rhebDSRmaxHiCiHar}{}\renewcommand{\rhebDSRmaxHiCiHar}{1.89}  % ... upper end
\providecommand{\rhebDSRmaxPCiHar}{}\renewcommand{\rhebDSRmaxPCiHar}{0.021}  % ... bootstrap p (descriptive; Delta >= 0)
\providecommand{\rhebGrsCiHar}{}\renewcommand{\rhebGrsCiHar}{3.30}  % ... one-asset GRS statistic, F(1, T - 2)
\providecommand{\rhebGrsPCiHar}{}\renewcommand{\rhebGrsPCiHar}{0.070}  % ... its p (i.i.d. errors)
\providecommand{\rhebTanShareCiHar}{}\renewcommand{\rhebTanShareCiHar}{99}  % ... tangency: strategy's share of volatility, %
\providecommand{\rhebTanDSixCiHar}{}\renewcommand{\rhebTanDSixCiHar}{0.98}  % ... tangency (fixed in-sample weights): Sharpe ratio, D6
\providecommand{\rhebRtMvCiHar}{}\renewcommand{\rhebRtMvCiHar}{1.23}  % ... real-time mean-variance: Sharpe ratio, STANDARD
\providecommand{\rhebRtMvDCiHar}{}\renewcommand{\rhebRtMvDCiHar}{0.74}  % ... minus the market's, STANDARD
\providecommand{\rhebRtMvPCiHar}{}\renewcommand{\rhebRtMvPCiHar}{0.346}  % ... Ledoit-Wolf p of the difference (D19)
\providecommand{\rhebRtMvDSixCiHar}{}\renewcommand{\rhebRtMvDSixCiHar}{1.13}  % ... real-time mean-variance: Sharpe ratio, D6
\providecommand{\rhebRtMvDSixDCiHar}{}\renewcommand{\rhebRtMvDSixDCiHar}{0.44}  % ... minus the market's, D6
\providecommand{\rhebRtMvDSixPCiHar}{}\renewcommand{\rhebRtMvDSixPCiHar}{0.623}  % ... bootstrap p of the D6 difference
\providecommand{\rhebExRtMvCiHar}{}\renewcommand{\rhebExRtMvCiHar}{1.02}  % ... w/o best days, real-time mean-variance: SR, STANDARD
\providecommand{\rhebExRtMvDCiHar}{}\renewcommand{\rhebExRtMvDCiHar}{0.48}  % ... minus the market's, STANDARD
\providecommand{\rhebExRtMvPCiHar}{}\renewcommand{\rhebExRtMvPCiHar}{0.556}  % ... its Ledoit-Wolf p (D19)
\providecommand{\rhebRtIvCiHar}{}\renewcommand{\rhebRtIvCiHar}{1.12}  % ... real-time inverse-volatility: Sharpe ratio, STANDARD
\providecommand{\rhebRtIvDCiHar}{}\renewcommand{\rhebRtIvDCiHar}{0.62}  % ... minus the market's, STANDARD
\providecommand{\rhebRtIvPCiHar}{}\renewcommand{\rhebRtIvPCiHar}{0.259}  % ... Ledoit-Wolf p of the difference (D19)
\providecommand{\rhebRtIvDSixCiHar}{}\renewcommand{\rhebRtIvDSixCiHar}{1.21}  % ... real-time inverse-volatility: Sharpe ratio, D6
\providecommand{\rhebRtIvDSixDCiHar}{}\renewcommand{\rhebRtIvDSixDCiHar}{0.53}  % ... minus the market's, D6
\providecommand{\rhebRtIvDSixPCiHar}{}\renewcommand{\rhebRtIvDSixPCiHar}{0.381}  % ... bootstrap p of the D6 difference
\providecommand{\rhebExRtIvCiHar}{}\renewcommand{\rhebExRtIvCiHar}{1.04}  % ... w/o best days, real-time inverse-volatility: SR, STANDARD
\providecommand{\rhebExRtIvDCiHar}{}\renewcommand{\rhebExRtIvDCiHar}{0.50}  % ... minus the market's, STANDARD
\providecommand{\rhebExRtIvPCiHar}{}\renewcommand{\rhebExRtIvPCiHar}{0.369}  % ... its Ledoit-Wolf p (D19)
\providecommand{\rhebRtMvShareCiHar}{}\renewcommand{\rhebRtMvShareCiHar}{90}  % ... real-time MV: median strategy vol share, %
\providecommand{\rhebRtMvCapDaysCiHar}{}\renewcommand{\rhebRtMvCapDaysCiHar}{100}  % ... % of days the 90% cap binds
\providecommand{\rhebRtMvWsCiHar}{}\renewcommand{\rhebRtMvWsCiHar}{3.6}  % ... real-time MV: median strategy dollar weight, %
\providecommand{\rhebRtIvWsCiHar}{}\renewcommand{\rhebRtIvWsCiHar}{0.9}  % ... real-time IV: median strategy dollar weight, %
\providecommand{\rhebExBetaCiHar}{}\renewcommand{\rhebExBetaCiHar}{4.67}  % ... without the best-day union: CAPM beta
\providecommand{\rhebExBetaTCiHar}{}\renewcommand{\rhebExBetaTCiHar}{1.19}  % ... its HAC t
\providecommand{\rhebExAlphaCiHar}{}\renewcommand{\rhebExAlphaCiHar}{4.9}  % ... without the best-day union: alpha, % per day
\providecommand{\rhebExAlphaTCiHar}{}\renewcommand{\rhebExAlphaTCiHar}{1.45}  % ... its HAC t
\providecommand{\rhebExAlphaPCiHar}{}\renewcommand{\rhebExAlphaPCiHar}{0.169}  % ... its fixed-b p
\providecommand{\rhebExARCiHar}{}\renewcommand{\rhebExARCiHar}{0.74}  % ... appraisal ratio, sqrt(q) (D19)
\providecommand{\rhebExSRmaxCiHar}{}\renewcommand{\rhebExSRmaxCiHar}{0.78}  % ... max Sharpe ratio, STANDARD
\providecommand{\rhebExDSRmaxPCiHar}{}\renewcommand{\rhebExDSRmaxPCiHar}{0.144}  % ... bootstrap p of its gain
\providecommand{\rhebExGrsPCiHar}{}\renewcommand{\rhebExGrsPCiHar}{0.222}  % ... GRS p
\providecommand{\rhebTmGammaCiHar}{}\renewcommand{\rhebTmGammaCiHar}{0.10}  % sub_ols_baseline, mid: Treynor-Mazuy gamma on (r^m)^2 (% per %^2)
\providecommand{\rhebTmGammaTCiHar}{}\renewcommand{\rhebTmGammaTCiHar}{0.10}  % ... its HAC t
\providecommand{\rhebTmGammaPCiHar}{}\renewcommand{\rhebTmGammaPCiHar}{0.922}  % ... its fixed-b p
\providecommand{\rhebTmBetaCiHar}{}\renewcommand{\rhebTmBetaCiHar}{5.07}  % ... Treynor-Mazuy beta on r^m
\providecommand{\rhebTmBetaTCiHar}{}\renewcommand{\rhebTmBetaTCiHar}{1.23}  % ... its HAC t
\providecommand{\rhebNLongCiHar}{}\renewcommand{\rhebNLongCiHar}{198}  % ... long days
\providecommand{\rhebTmGammaLongCiHar}{}\renewcommand{\rhebTmGammaLongCiHar}{11.65}  % ... TM gamma on its long days only
\providecommand{\rhebTmGammaLongTCiHar}{}\renewcommand{\rhebTmGammaLongTCiHar}{1.88}  % ... its HAC t (no test: RULE section 4)
\providecommand{\rhebNShortCiHar}{}\renewcommand{\rhebNShortCiHar}{416}  % ... short days
\providecommand{\rhebTmGammaShortCiHar}{}\renewcommand{\rhebTmGammaShortCiHar}{\ensuremath{-}0.66}  % ... TM gamma on its short days only
\providecommand{\rhebTmGammaShortTCiHar}{}\renewcommand{\rhebTmGammaShortTCiHar}{\ensuremath{-}0.62}  % ... its HAC t (no test: RULE section 4)
\providecommand{\rhebDownMeanCiHar}{}\renewcommand{\rhebDownMeanCiHar}{5.1}  % ... mean return on market-down days, % of premium
\providecommand{\rhebDownRestCiHar}{}\renewcommand{\rhebDownRestCiHar}{11.5}  % ... on the other days
\providecommand{\rhebDownDiffTCiHar}{}\renewcommand{\rhebDownDiffTCiHar}{\ensuremath{-}0.57}  % ... HAC t of the difference
\providecommand{\rhebDownDiffPCiHar}{}\renewcommand{\rhebDownDiffPCiHar}{0.584}  % ... its fixed-b p
\providecommand{\rhebWorstMeanCiHar}{}\renewcommand{\rhebWorstMeanCiHar}{\ensuremath{-}11.8}  % ... mean return on worst-decile market days, % of premium
\providecommand{\rhebWorstRestCiHar}{}\renewcommand{\rhebWorstRestCiHar}{10.6}  % ... on the other days
\providecommand{\rhebWorstDiffTCiHar}{}\renewcommand{\rhebWorstDiffTCiHar}{\ensuremath{-}1.68}  % ... HAC t of the difference
\providecommand{\rhebWorstDiffPCiHar}{}\renewcommand{\rhebWorstDiffPCiHar}{0.114}  % ... its fixed-b p
\providecommand{\rhebThirtyCorrCiHar}{}\renewcommand{\rhebThirtyCorrCiHar}{0.10}  % ... correlation with the 15:30-16:00 index return
\providecommand{\rhebThirtyBetaCiHar}{}\renewcommand{\rhebThirtyBetaCiHar}{42.09}  % ... beta on the 15:30-16:00 index return
\providecommand{\rhebThirtyBetaTCiHar}{}\renewcommand{\rhebThirtyBetaTCiHar}{1.42}  % ... its HAC t
\providecommand{\rhebThirtyBetaPCiHar}{}\renewcommand{\rhebThirtyBetaPCiHar}{0.177}  % ... its fixed-b p
\providecommand{\rhebThirtyAlphaCiHar}{}\renewcommand{\rhebThirtyAlphaCiHar}{8.8}  % ... alpha against it, % per day
\providecommand{\rhebThirtyAlphaTCiHar}{}\renewcommand{\rhebThirtyAlphaTCiHar}{2.57}  % ... its HAC t
\providecommand{\rhebThirtyTmGammaCiHar}{}\renewcommand{\rhebThirtyTmGammaCiHar}{\ensuremath{-}86.34}  % ... TM gamma on the squared 15:30-16:00 return
\providecommand{\rhebThirtyTmGammaTCiHar}{}\renewcommand{\rhebThirtyTmGammaTCiHar}{\ensuremath{-}2.12}  % ... its HAC t
\providecommand{\rhebThirtyTmGammaPCiHar}{}\renewcommand{\rhebThirtyTmGammaPCiHar}{0.049}  % ... its fixed-b p
\providecommand{\rhebThirtyTmGammaLongCiHar}{}\renewcommand{\rhebThirtyTmGammaLongCiHar}{374.79}  % ... ... on its long days only
\providecommand{\rhebThirtyTmGammaLongTCiHar}{}\renewcommand{\rhebThirtyTmGammaLongTCiHar}{2.59}  % ... its HAC t (no test)
\providecommand{\rhebThirtyTmGammaShortCiHar}{}\renewcommand{\rhebThirtyTmGammaShortCiHar}{\ensuremath{-}255.86}  % ... ... on its short days only
\providecommand{\rhebThirtyTmGammaShortTCiHar}{}\renewcommand{\rhebThirtyTmGammaShortTCiHar}{\ensuremath{-}4.38}  % ... its HAC t (no test)
\providecommand{\rhebThirtyDownMeanCiHar}{}\renewcommand{\rhebThirtyDownMeanCiHar}{2.6}  % ... mean on 15:30-16:00 down days, %
\providecommand{\rhebThirtyDownDiffTCiHar}{}\renewcommand{\rhebThirtyDownDiffTCiHar}{\ensuremath{-}1.34}  % ... HAC t of the difference with the other days
\providecommand{\rhebThirtyWorstMeanCiHar}{}\renewcommand{\rhebThirtyWorstMeanCiHar}{\ensuremath{-}53.1}  % ... mean on 15:30-16:00 worst days, %
\providecommand{\rhebThirtyWorstDiffTCiHar}{}\renewcommand{\rhebThirtyWorstDiffTCiHar}{\ensuremath{-}2.34}  % ... HAC t of the difference with the other days
\providecommand{\rhebCorrXCiHar}{}\renewcommand{\rhebCorrXCiHar}{0.05}  % sub_ols_baseline, bid/ask: correlation with the market excess return
\providecommand{\rhebBetaXCiHar}{}\renewcommand{\rhebBetaXCiHar}{5.15}  % ... CAPM beta (% of premium per % of market)
\providecommand{\rhebBetaTXCiHar}{}\renewcommand{\rhebBetaTXCiHar}{1.23}  % ... its HAC t (S = 33)
\providecommand{\rhebBetaPXCiHar}{}\renewcommand{\rhebBetaPXCiHar}{0.242}  % ... its fixed-b p
\providecommand{\rhebRsqXCiHar}{}\renewcommand{\rhebRsqXCiHar}{0.00}  % ... CAPM R^2
\providecommand{\rhebAlphaXCiHar}{}\renewcommand{\rhebAlphaXCiHar}{5.7}  % ... CAPM alpha, % of premium per day
\providecommand{\rhebAlphaTXCiHar}{}\renewcommand{\rhebAlphaTXCiHar}{1.60}  % ... its HAC t
\providecommand{\rhebAlphaPXCiHar}{}\renewcommand{\rhebAlphaPXCiHar}{0.132}  % ... its fixed-b p
\providecommand{\rhebAlphaAnnXCiHar}{}\renewcommand{\rhebAlphaAnnXCiHar}{1279}  % ... alpha x q (D19; was x 252), % of premium a year
\providecommand{\rhebARXCiHar}{}\renewcommand{\rhebARXCiHar}{0.75}  % ... appraisal ratio alpha / sigma_eps, annualized sqrt(q) (D19)
\providecommand{\rhebSRsXCiHar}{}\renewcommand{\rhebSRsXCiHar}{0.76}  % ... strategy Sharpe ratio, STANDARD = the paper's (= \rhemSh / X)
\providecommand{\rhebSRsDSixXCiHar}{}\renewcommand{\rhebSRsDSixXCiHar}{0.63}  % ... strategy Sharpe ratio, D6 (= \rhemShDSix)
\providecommand{\rhebSRmaxXCiHar}{}\renewcommand{\rhebSRmaxXCiHar}{0.77}  % ... max Sharpe ratio of market + strategy, STANDARD
\providecommand{\rhebDSRmaxXCiHar}{}\renewcommand{\rhebDSRmaxXCiHar}{0.62}  % ... its gain over the market alone, STANDARD
\providecommand{\rhebDSRmaxLoXCiHar}{}\renewcommand{\rhebDSRmaxLoXCiHar}{0.00}  % ... bootstrap 95% interval, lower end
\providecommand{\rhebDSRmaxHiXCiHar}{}\renewcommand{\rhebDSRmaxHiXCiHar}{1.55}  % ... upper end
\providecommand{\rhebDSRmaxPXCiHar}{}\renewcommand{\rhebDSRmaxPXCiHar}{0.077}  % ... bootstrap p (descriptive; Delta >= 0)
\providecommand{\rhebGrsXCiHar}{}\renewcommand{\rhebGrsXCiHar}{1.54}  % ... one-asset GRS statistic, F(1, T - 2)
\providecommand{\rhebGrsPXCiHar}{}\renewcommand{\rhebGrsPXCiHar}{0.215}  % ... its p (i.i.d. errors)
\providecommand{\rhebTanShareXCiHar}{}\renewcommand{\rhebTanShareXCiHar}{97}  % ... tangency: strategy's share of volatility, %
\providecommand{\rhebTanDSixXCiHar}{}\renewcommand{\rhebTanDSixXCiHar}{0.63}  % ... tangency (fixed in-sample weights): Sharpe ratio, D6
\providecommand{\rhebRtMvXCiHar}{}\renewcommand{\rhebRtMvXCiHar}{0.90}  % ... real-time mean-variance: Sharpe ratio, STANDARD
\providecommand{\rhebRtMvDXCiHar}{}\renewcommand{\rhebRtMvDXCiHar}{0.40}  % ... minus the market's, STANDARD
\providecommand{\rhebRtMvPXCiHar}{}\renewcommand{\rhebRtMvPXCiHar}{0.600}  % ... Ledoit-Wolf p of the difference (D19)
\providecommand{\rhebRtMvDSixXCiHar}{}\renewcommand{\rhebRtMvDSixXCiHar}{0.79}  % ... real-time mean-variance: Sharpe ratio, D6
\providecommand{\rhebRtMvDSixDXCiHar}{}\renewcommand{\rhebRtMvDSixDXCiHar}{0.11}  % ... minus the market's, D6
\providecommand{\rhebRtMvDSixPXCiHar}{}\renewcommand{\rhebRtMvDSixPXCiHar}{0.898}  % ... bootstrap p of the D6 difference
\providecommand{\rhebExRtMvXCiHar}{}\renewcommand{\rhebExRtMvXCiHar}{0.76}  % ... w/o best days, real-time mean-variance: SR, STANDARD
\providecommand{\rhebExRtMvDXCiHar}{}\renewcommand{\rhebExRtMvDXCiHar}{0.21}  % ... minus the market's, STANDARD
\providecommand{\rhebExRtMvPXCiHar}{}\renewcommand{\rhebExRtMvPXCiHar}{0.793}  % ... its Ledoit-Wolf p (D19)
\providecommand{\rhebRtIvXCiHar}{}\renewcommand{\rhebRtIvXCiHar}{0.86}  % ... real-time inverse-volatility: Sharpe ratio, STANDARD
\providecommand{\rhebRtIvDXCiHar}{}\renewcommand{\rhebRtIvDXCiHar}{0.36}  % ... minus the market's, STANDARD
\providecommand{\rhebRtIvPXCiHar}{}\renewcommand{\rhebRtIvPXCiHar}{0.497}  % ... Ledoit-Wolf p of the difference (D19)
\providecommand{\rhebRtIvDSixXCiHar}{}\renewcommand{\rhebRtIvDSixXCiHar}{0.95}  % ... real-time inverse-volatility: Sharpe ratio, D6
\providecommand{\rhebRtIvDSixDXCiHar}{}\renewcommand{\rhebRtIvDSixDXCiHar}{0.27}  % ... minus the market's, D6
\providecommand{\rhebRtIvDSixPXCiHar}{}\renewcommand{\rhebRtIvDSixPXCiHar}{0.660}  % ... bootstrap p of the D6 difference
\providecommand{\rhebExRtIvXCiHar}{}\renewcommand{\rhebExRtIvXCiHar}{0.78}  % ... w/o best days, real-time inverse-volatility: SR, STANDARD
\providecommand{\rhebExRtIvDXCiHar}{}\renewcommand{\rhebExRtIvDXCiHar}{0.24}  % ... minus the market's, STANDARD
\providecommand{\rhebExRtIvPXCiHar}{}\renewcommand{\rhebExRtIvPXCiHar}{0.662}  % ... its Ledoit-Wolf p (D19)
\providecommand{\rhebRtMvShareXCiHar}{}\renewcommand{\rhebRtMvShareXCiHar}{90}  % ... real-time MV: median strategy vol share, %
\providecommand{\rhebRtMvCapDaysXCiHar}{}\renewcommand{\rhebRtMvCapDaysXCiHar}{100}  % ... % of days the 90% cap binds
\providecommand{\rhebRtMvWsXCiHar}{}\renewcommand{\rhebRtMvWsXCiHar}{3.6}  % ... real-time MV: median strategy dollar weight, %
\providecommand{\rhebRtIvWsXCiHar}{}\renewcommand{\rhebRtIvWsXCiHar}{0.9}  % ... real-time IV: median strategy dollar weight, %
\providecommand{\rhebExBetaXCiHar}{}\renewcommand{\rhebExBetaXCiHar}{4.70}  % ... without the best-day union: CAPM beta
\providecommand{\rhebExBetaTXCiHar}{}\renewcommand{\rhebExBetaTXCiHar}{1.18}  % ... its HAC t
\providecommand{\rhebExAlphaXCiHar}{}\renewcommand{\rhebExAlphaXCiHar}{2.5}  % ... without the best-day union: alpha, % per day
\providecommand{\rhebExAlphaTXCiHar}{}\renewcommand{\rhebExAlphaTXCiHar}{0.72}  % ... its HAC t
\providecommand{\rhebExAlphaPXCiHar}{}\renewcommand{\rhebExAlphaPXCiHar}{0.489}  % ... its fixed-b p
\providecommand{\rhebExARXCiHar}{}\renewcommand{\rhebExARXCiHar}{0.37}  % ... appraisal ratio, sqrt(q) (D19)
\providecommand{\rhebExSRmaxXCiHar}{}\renewcommand{\rhebExSRmaxXCiHar}{0.44}  % ... max Sharpe ratio, STANDARD
\providecommand{\rhebExDSRmaxPXCiHar}{}\renewcommand{\rhebExDSRmaxPXCiHar}{0.334}  % ... bootstrap p of its gain
\providecommand{\rhebExGrsPXCiHar}{}\renewcommand{\rhebExGrsPXCiHar}{0.542}  % ... GRS p
\providecommand{\rhebCorrShort}{}\renewcommand{\rhebCorrShort}{0.04}  % always_short, mid: correlation with the market excess return
\providecommand{\rhebBetaShort}{}\renewcommand{\rhebBetaShort}{4.43}  % ... CAPM beta (% of premium per % of market)
\providecommand{\rhebBetaTShort}{}\renewcommand{\rhebBetaTShort}{1.22}  % ... its HAC t (S = 33)
\providecommand{\rhebBetaPShort}{}\renewcommand{\rhebBetaPShort}{0.244}  % ... its fixed-b p
\providecommand{\rhebRsqShort}{}\renewcommand{\rhebRsqShort}{0.00}  % ... CAPM R^2
\providecommand{\rhebAlphaShort}{}\renewcommand{\rhebAlphaShort}{2.4}  % ... CAPM alpha, % of premium per day
\providecommand{\rhebAlphaTShort}{}\renewcommand{\rhebAlphaTShort}{0.65}  % ... its HAC t
\providecommand{\rhebAlphaPShort}{}\renewcommand{\rhebAlphaPShort}{0.533}  % ... its fixed-b p
\providecommand{\rhebAlphaAnnShort}{}\renewcommand{\rhebAlphaAnnShort}{537}  % ... alpha x q (D19; was x 252), % of premium a year
\providecommand{\rhebARShort}{}\renewcommand{\rhebARShort}{0.32}  % ... appraisal ratio alpha / sigma_eps, annualized sqrt(q) (D19)
\providecommand{\rhebSRsShort}{}\renewcommand{\rhebSRsShort}{0.32}  % ... strategy Sharpe ratio, STANDARD = the paper's (= \rhemSh / X)
\providecommand{\rhebSRsDSixShort}{}\renewcommand{\rhebSRsDSixShort}{\ensuremath{-}0.12}  % ... strategy Sharpe ratio, D6 (= \rhemShDSix)
\providecommand{\rhebSRmaxShort}{}\renewcommand{\rhebSRmaxShort}{0.35}  % ... max Sharpe ratio of market + strategy, STANDARD
\providecommand{\rhebDSRmaxShort}{}\renewcommand{\rhebDSRmaxShort}{0.20}  % ... its gain over the market alone, STANDARD
\providecommand{\rhebDSRmaxLoShort}{}\renewcommand{\rhebDSRmaxLoShort}{0.00}  % ... bootstrap 95% interval, lower end
\providecommand{\rhebDSRmaxHiShort}{}\renewcommand{\rhebDSRmaxHiShort}{1.20}  % ... upper end
\providecommand{\rhebDSRmaxPShort}{}\renewcommand{\rhebDSRmaxPShort}{0.303}  % ... bootstrap p (descriptive; Delta >= 0)
\providecommand{\rhebGrsShort}{}\renewcommand{\rhebGrsShort}{0.27}  % ... one-asset GRS statistic, F(1, T - 2)
\providecommand{\rhebGrsPShort}{}\renewcommand{\rhebGrsPShort}{0.603}  % ... its p (i.i.d. errors)
\providecommand{\rhebTanShareShort}{}\renewcommand{\rhebTanShareShort}{83}  % ... tangency: strategy's share of volatility, %
\providecommand{\rhebTanDSixShort}{}\renewcommand{\rhebTanDSixShort}{\ensuremath{-}0.10}  % ... tangency (fixed in-sample weights): Sharpe ratio, D6
\providecommand{\rhebRtMvShort}{}\renewcommand{\rhebRtMvShort}{0.90}  % ... real-time mean-variance: Sharpe ratio, STANDARD
\providecommand{\rhebRtMvDShort}{}\renewcommand{\rhebRtMvDShort}{0.40}  % ... minus the market's, STANDARD
\providecommand{\rhebRtMvPShort}{}\renewcommand{\rhebRtMvPShort}{0.616}  % ... Ledoit-Wolf p of the difference (D19)
\providecommand{\rhebRtMvDSixShort}{}\renewcommand{\rhebRtMvDSixShort}{0.47}  % ... real-time mean-variance: Sharpe ratio, D6
\providecommand{\rhebRtMvDSixDShort}{}\renewcommand{\rhebRtMvDSixDShort}{\ensuremath{-}0.21}  % ... minus the market's, D6
\providecommand{\rhebRtMvDSixPShort}{}\renewcommand{\rhebRtMvDSixPShort}{0.788}  % ... bootstrap p of the D6 difference
\providecommand{\rhebExRtMvShort}{}\renewcommand{\rhebExRtMvShort}{1.98}  % ... w/o best days, real-time mean-variance: SR, STANDARD
\providecommand{\rhebExRtMvDShort}{}\renewcommand{\rhebExRtMvDShort}{1.44}  % ... minus the market's, STANDARD
\providecommand{\rhebExRtMvPShort}{}\renewcommand{\rhebExRtMvPShort}{0.108}  % ... its Ledoit-Wolf p (D19)
\providecommand{\rhebRtIvShort}{}\renewcommand{\rhebRtIvShort}{1.34}  % ... real-time inverse-volatility: Sharpe ratio, STANDARD
\providecommand{\rhebRtIvDShort}{}\renewcommand{\rhebRtIvDShort}{0.85}  % ... minus the market's, STANDARD
\providecommand{\rhebRtIvPShort}{}\renewcommand{\rhebRtIvPShort}{0.138}  % ... Ledoit-Wolf p of the difference (D19)
\providecommand{\rhebRtIvDSixShort}{}\renewcommand{\rhebRtIvDSixShort}{0.89}  % ... real-time inverse-volatility: Sharpe ratio, D6
\providecommand{\rhebRtIvDSixDShort}{}\renewcommand{\rhebRtIvDSixDShort}{0.21}  % ... minus the market's, D6
\providecommand{\rhebRtIvDSixPShort}{}\renewcommand{\rhebRtIvDSixPShort}{0.707}  % ... bootstrap p of the D6 difference
\providecommand{\rhebExRtIvShort}{}\renewcommand{\rhebExRtIvShort}{1.82}  % ... w/o best days, real-time inverse-volatility: SR, STANDARD
\providecommand{\rhebExRtIvDShort}{}\renewcommand{\rhebExRtIvDShort}{1.28}  % ... minus the market's, STANDARD
\providecommand{\rhebExRtIvPShort}{}\renewcommand{\rhebExRtIvPShort}{0.033}  % ... its Ledoit-Wolf p (D19)
\providecommand{\rhebRtMvShareShort}{}\renewcommand{\rhebRtMvShareShort}{90}  % ... real-time MV: median strategy vol share, %
\providecommand{\rhebRtMvCapDaysShort}{}\renewcommand{\rhebRtMvCapDaysShort}{91}  % ... % of days the 90% cap binds
\providecommand{\rhebRtMvWsShort}{}\renewcommand{\rhebRtMvWsShort}{3.6}  % ... real-time MV: median strategy dollar weight, %
\providecommand{\rhebRtIvWsShort}{}\renewcommand{\rhebRtIvWsShort}{0.9}  % ... real-time IV: median strategy dollar weight, %
\providecommand{\rhebExBetaShort}{}\renewcommand{\rhebExBetaShort}{2.53}  % ... without the best-day union: CAPM beta
\providecommand{\rhebExBetaTShort}{}\renewcommand{\rhebExBetaTShort}{0.78}  % ... its HAC t
\providecommand{\rhebExAlphaShort}{}\renewcommand{\rhebExAlphaShort}{9.4}  % ... without the best-day union: alpha, % per day
\providecommand{\rhebExAlphaTShort}{}\renewcommand{\rhebExAlphaTShort}{2.93}  % ... its HAC t
\providecommand{\rhebExAlphaPShort}{}\renewcommand{\rhebExAlphaPShort}{0.008}  % ... its fixed-b p
\providecommand{\rhebExARShort}{}\renewcommand{\rhebExARShort}{1.44}  % ... appraisal ratio, sqrt(q) (D19)
\providecommand{\rhebExSRmaxShort}{}\renewcommand{\rhebExSRmaxShort}{1.46}  % ... max Sharpe ratio, STANDARD
\providecommand{\rhebExDSRmaxPShort}{}\renewcommand{\rhebExDSRmaxPShort}{0.022}  % ... bootstrap p of its gain
\providecommand{\rhebExGrsPShort}{}\renewcommand{\rhebExGrsPShort}{0.018}  % ... GRS p
\providecommand{\rhebTmGammaShort}{}\renewcommand{\rhebTmGammaShort}{\ensuremath{-}0.94}  % always_short, mid: Treynor-Mazuy gamma on (r^m)^2 (% per %^2)
\providecommand{\rhebTmGammaTShort}{}\renewcommand{\rhebTmGammaTShort}{\ensuremath{-}0.70}  % ... its HAC t
\providecommand{\rhebTmGammaPShort}{}\renewcommand{\rhebTmGammaPShort}{0.505}  % ... its fixed-b p
\providecommand{\rhebTmBetaShort}{}\renewcommand{\rhebTmBetaShort}{4.39}  % ... Treynor-Mazuy beta on r^m
\providecommand{\rhebTmBetaTShort}{}\renewcommand{\rhebTmBetaTShort}{1.20}  % ... its HAC t
\providecommand{\rhebNLongShort}{}\renewcommand{\rhebNLongShort}{0}  % ... long days
\providecommand{\rhebTmGammaLongShort}{}\renewcommand{\rhebTmGammaLongShort}{--}  % ... TM gamma on its long days only
\providecommand{\rhebTmGammaLongTShort}{}\renewcommand{\rhebTmGammaLongTShort}{--}  % ... its HAC t (no test: RULE section 4)
\providecommand{\rhebNShortShort}{}\renewcommand{\rhebNShortShort}{614}  % ... short days
\providecommand{\rhebTmGammaShortShort}{}\renewcommand{\rhebTmGammaShortShort}{\ensuremath{-}0.94}  % ... TM gamma on its short days only
\providecommand{\rhebTmGammaShortTShort}{}\renewcommand{\rhebTmGammaShortTShort}{\ensuremath{-}0.70}  % ... its HAC t (no test: RULE section 4)
\providecommand{\rhebDownMeanShort}{}\renewcommand{\rhebDownMeanShort}{\ensuremath{-}1.9}  % ... mean return on market-down days, % of premium
\providecommand{\rhebDownRestShort}{}\renewcommand{\rhebDownRestShort}{6.8}  % ... on the other days
\providecommand{\rhebDownDiffTShort}{}\renewcommand{\rhebDownDiffTShort}{\ensuremath{-}0.87}  % ... HAC t of the difference
\providecommand{\rhebDownDiffPShort}{}\renewcommand{\rhebDownDiffPShort}{0.405}  % ... its fixed-b p
\providecommand{\rhebWorstMeanShort}{}\renewcommand{\rhebWorstMeanShort}{\ensuremath{-}22.4}  % ... mean return on worst-decile market days, % of premium
\providecommand{\rhebWorstRestShort}{}\renewcommand{\rhebWorstRestShort}{5.2}  % ... on the other days
\providecommand{\rhebWorstDiffTShort}{}\renewcommand{\rhebWorstDiffTShort}{\ensuremath{-}1.99}  % ... HAC t of the difference
\providecommand{\rhebWorstDiffPShort}{}\renewcommand{\rhebWorstDiffPShort}{0.063}  % ... its fixed-b p
\providecommand{\rhebThirtyCorrShort}{}\renewcommand{\rhebThirtyCorrShort}{0.07}  % ... correlation with the 15:30-16:00 index return
\providecommand{\rhebThirtyBetaShort}{}\renewcommand{\rhebThirtyBetaShort}{27.60}  % ... beta on the 15:30-16:00 index return
\providecommand{\rhebThirtyBetaTShort}{}\renewcommand{\rhebThirtyBetaTShort}{0.98}  % ... its HAC t
\providecommand{\rhebThirtyBetaPShort}{}\renewcommand{\rhebThirtyBetaPShort}{0.350}  % ... its fixed-b p
\providecommand{\rhebThirtyAlphaShort}{}\renewcommand{\rhebThirtyAlphaShort}{2.8}  % ... alpha against it, % per day
\providecommand{\rhebThirtyAlphaTShort}{}\renewcommand{\rhebThirtyAlphaTShort}{0.74}  % ... its HAC t
\providecommand{\rhebThirtyTmGammaShort}{}\renewcommand{\rhebThirtyTmGammaShort}{\ensuremath{-}274.27}  % ... TM gamma on the squared 15:30-16:00 return
\providecommand{\rhebThirtyTmGammaTShort}{}\renewcommand{\rhebThirtyTmGammaTShort}{\ensuremath{-}4.16}  % ... its HAC t
\providecommand{\rhebThirtyTmGammaPShort}{}\renewcommand{\rhebThirtyTmGammaPShort}{\ensuremath{<}0.001}  % ... its fixed-b p
\providecommand{\rhebThirtyTmGammaLongShort}{}\renewcommand{\rhebThirtyTmGammaLongShort}{--}  % ... ... on its long days only
\providecommand{\rhebThirtyTmGammaLongTShort}{}\renewcommand{\rhebThirtyTmGammaLongTShort}{--}  % ... its HAC t (no test)
\providecommand{\rhebThirtyTmGammaShortShort}{}\renewcommand{\rhebThirtyTmGammaShortShort}{\ensuremath{-}274.27}  % ... ... on its short days only
\providecommand{\rhebThirtyTmGammaShortTShort}{}\renewcommand{\rhebThirtyTmGammaShortTShort}{\ensuremath{-}4.16}  % ... its HAC t (no test)
\providecommand{\rhebThirtyDownMeanShort}{}\renewcommand{\rhebThirtyDownMeanShort}{1.7}  % ... mean on 15:30-16:00 down days, %
\providecommand{\rhebThirtyDownDiffTShort}{}\renewcommand{\rhebThirtyDownDiffTShort}{\ensuremath{-}0.19}  % ... HAC t of the difference with the other days
\providecommand{\rhebThirtyWorstMeanShort}{}\renewcommand{\rhebThirtyWorstMeanShort}{\ensuremath{-}134.4}  % ... mean on 15:30-16:00 worst days, %
\providecommand{\rhebThirtyWorstDiffTShort}{}\renewcommand{\rhebThirtyWorstDiffTShort}{\ensuremath{-}5.37}  % ... HAC t of the difference with the other days
\providecommand{\rhebCorrXShort}{}\renewcommand{\rhebCorrXShort}{0.04}  % always_short, bid/ask: correlation with the market excess return
\providecommand{\rhebBetaXShort}{}\renewcommand{\rhebBetaXShort}{4.50}  % ... CAPM beta (% of premium per % of market)
\providecommand{\rhebBetaTXShort}{}\renewcommand{\rhebBetaTXShort}{1.21}  % ... its HAC t (S = 33)
\providecommand{\rhebBetaPXShort}{}\renewcommand{\rhebBetaPXShort}{0.248}  % ... its fixed-b p
\providecommand{\rhebRsqXShort}{}\renewcommand{\rhebRsqXShort}{0.00}  % ... CAPM R^2
\providecommand{\rhebAlphaXShort}{}\renewcommand{\rhebAlphaXShort}{\ensuremath{-}0.3}  % ... CAPM alpha, % of premium per day
\providecommand{\rhebAlphaTXShort}{}\renewcommand{\rhebAlphaTXShort}{\ensuremath{-}0.08}  % ... its HAC t
\providecommand{\rhebAlphaPXShort}{}\renewcommand{\rhebAlphaPXShort}{0.939}  % ... its fixed-b p
\providecommand{\rhebAlphaAnnXShort}{}\renewcommand{\rhebAlphaAnnXShort}{\ensuremath{-}68}  % ... alpha x q (D19; was x 252), % of premium a year
\providecommand{\rhebARXShort}{}\renewcommand{\rhebARXShort}{\ensuremath{-}0.04}  % ... appraisal ratio alpha / sigma_eps, annualized sqrt(q) (D19)
\providecommand{\rhebSRsXShort}{}\renewcommand{\rhebSRsXShort}{\ensuremath{-}0.03}  % ... strategy Sharpe ratio, STANDARD = the paper's (= \rhemSh / X)
\providecommand{\rhebSRsDSixXShort}{}\renewcommand{\rhebSRsDSixXShort}{\ensuremath{-}0.46}  % ... strategy Sharpe ratio, D6 (= \rhemShDSix)
\providecommand{\rhebSRmaxXShort}{}\renewcommand{\rhebSRmaxXShort}{0.16}  % ... max Sharpe ratio of market + strategy, STANDARD
\providecommand{\rhebDSRmaxXShort}{}\renewcommand{\rhebDSRmaxXShort}{0.00}  % ... its gain over the market alone, STANDARD
\providecommand{\rhebDSRmaxLoXShort}{}\renewcommand{\rhebDSRmaxLoXShort}{0.00}  % ... bootstrap 95% interval, lower end
\providecommand{\rhebDSRmaxHiXShort}{}\renewcommand{\rhebDSRmaxHiXShort}{0.98}  % ... upper end
\providecommand{\rhebDSRmaxPXShort}{}\renewcommand{\rhebDSRmaxPXShort}{0.874}  % ... bootstrap p (descriptive; Delta >= 0)
\providecommand{\rhebGrsXShort}{}\renewcommand{\rhebGrsXShort}{0.00}  % ... one-asset GRS statistic, F(1, T - 2)
\providecommand{\rhebGrsPXShort}{}\renewcommand{\rhebGrsPXShort}{0.949}  % ... its p (i.i.d. errors)
\providecommand{\rhebTanShareXShort}{}\renewcommand{\rhebTanShareXShort}{5}  % ... tangency: strategy's share of volatility, %
\providecommand{\rhebTanDSixXShort}{}\renewcommand{\rhebTanDSixXShort}{0.89}  % ... tangency (fixed in-sample weights): Sharpe ratio, D6
\providecommand{\rhebRtMvXShort}{}\renewcommand{\rhebRtMvXShort}{0.69}  % ... real-time mean-variance: Sharpe ratio, STANDARD
\providecommand{\rhebRtMvDXShort}{}\renewcommand{\rhebRtMvDXShort}{0.20}  % ... minus the market's, STANDARD
\providecommand{\rhebRtMvPXShort}{}\renewcommand{\rhebRtMvPXShort}{0.813}  % ... Ledoit-Wolf p of the difference (D19)
\providecommand{\rhebRtMvDSixXShort}{}\renewcommand{\rhebRtMvDSixXShort}{0.13}  % ... real-time mean-variance: Sharpe ratio, D6
\providecommand{\rhebRtMvDSixDXShort}{}\renewcommand{\rhebRtMvDSixDXShort}{\ensuremath{-}0.55}  % ... minus the market's, D6
\providecommand{\rhebRtMvDSixPXShort}{}\renewcommand{\rhebRtMvDSixPXShort}{0.446}  % ... bootstrap p of the D6 difference
\providecommand{\rhebExRtMvXShort}{}\renewcommand{\rhebExRtMvXShort}{1.59}  % ... w/o best days, real-time mean-variance: SR, STANDARD
\providecommand{\rhebExRtMvDXShort}{}\renewcommand{\rhebExRtMvDXShort}{1.05}  % ... minus the market's, STANDARD
\providecommand{\rhebExRtMvPXShort}{}\renewcommand{\rhebExRtMvPXShort}{0.224}  % ... its Ledoit-Wolf p (D19)
\providecommand{\rhebRtIvXShort}{}\renewcommand{\rhebRtIvXShort}{1.06}  % ... real-time inverse-volatility: Sharpe ratio, STANDARD
\providecommand{\rhebRtIvDXShort}{}\renewcommand{\rhebRtIvDXShort}{0.56}  % ... minus the market's, STANDARD
\providecommand{\rhebRtIvPXShort}{}\renewcommand{\rhebRtIvPXShort}{0.307}  % ... Ledoit-Wolf p of the difference (D19)
\providecommand{\rhebRtIvDSixXShort}{}\renewcommand{\rhebRtIvDSixXShort}{0.62}  % ... real-time inverse-volatility: Sharpe ratio, D6
\providecommand{\rhebRtIvDSixDXShort}{}\renewcommand{\rhebRtIvDSixDXShort}{\ensuremath{-}0.06}  % ... minus the market's, D6
\providecommand{\rhebRtIvDSixPXShort}{}\renewcommand{\rhebRtIvDSixPXShort}{0.910}  % ... bootstrap p of the D6 difference
\providecommand{\rhebExRtIvXShort}{}\renewcommand{\rhebExRtIvXShort}{1.52}  % ... w/o best days, real-time inverse-volatility: SR, STANDARD
\providecommand{\rhebExRtIvDXShort}{}\renewcommand{\rhebExRtIvDXShort}{0.98}  % ... minus the market's, STANDARD
\providecommand{\rhebExRtIvPXShort}{}\renewcommand{\rhebExRtIvPXShort}{0.098}  % ... its Ledoit-Wolf p (D19)
\providecommand{\rhebRtMvShareXShort}{}\renewcommand{\rhebRtMvShareXShort}{90}  % ... real-time MV: median strategy vol share, %
\providecommand{\rhebRtMvCapDaysXShort}{}\renewcommand{\rhebRtMvCapDaysXShort}{89}  % ... % of days the 90% cap binds
\providecommand{\rhebRtMvWsXShort}{}\renewcommand{\rhebRtMvWsXShort}{3.4}  % ... real-time MV: median strategy dollar weight, %
\providecommand{\rhebRtIvWsXShort}{}\renewcommand{\rhebRtIvWsXShort}{0.8}  % ... real-time IV: median strategy dollar weight, %
\providecommand{\rhebExBetaXShort}{}\renewcommand{\rhebExBetaXShort}{2.55}  % ... without the best-day union: CAPM beta
\providecommand{\rhebExBetaTXShort}{}\renewcommand{\rhebExBetaTXShort}{0.77}  % ... its HAC t
\providecommand{\rhebExAlphaXShort}{}\renewcommand{\rhebExAlphaXShort}{7.0}  % ... without the best-day union: alpha, % per day
\providecommand{\rhebExAlphaTXShort}{}\renewcommand{\rhebExAlphaTXShort}{2.11}  % ... its HAC t
\providecommand{\rhebExAlphaPXShort}{}\renewcommand{\rhebExAlphaPXShort}{0.050}  % ... its fixed-b p
\providecommand{\rhebExARXShort}{}\renewcommand{\rhebExARXShort}{1.04}  % ... appraisal ratio, sqrt(q) (D19)
\providecommand{\rhebExSRmaxXShort}{}\renewcommand{\rhebExSRmaxXShort}{1.07}  % ... max Sharpe ratio, STANDARD
\providecommand{\rhebExDSRmaxPXShort}{}\renewcommand{\rhebExDSRmaxPXShort}{0.067}  % ... bootstrap p of its gain
\providecommand{\rhebExGrsPXShort}{}\renewcommand{\rhebExGrsPXShort}{0.087}  % ... GRS p
\providecommand{\rhebBetaFullMin}{}\renewcommand{\rhebBetaFullMin}{\ensuremath{-}0.52}  % CAPM beta, mid, six full-set: min
\providecommand{\rhebBetaFullMax}{}\renewcommand{\rhebBetaFullMax}{2.92}  % ... max
\providecommand{\rhebBetaTFullMin}{}\renewcommand{\rhebBetaTFullMin}{\ensuremath{-}0.13}  % ... its HAC t, mid, six full-set: min
\providecommand{\rhebBetaTFullMax}{}\renewcommand{\rhebBetaTFullMax}{0.75}  % ... max
\providecommand{\rhebBetaPFullMin}{}\renewcommand{\rhebBetaPFullMin}{0.472}  % ... its fixed-b p, mid, six full-set: min
\providecommand{\rhebBetaPFullMax}{}\renewcommand{\rhebBetaPFullMax}{0.899}  % ... max
\providecommand{\rhebCorrFullMin}{}\renewcommand{\rhebCorrFullMin}{\ensuremath{-}0.01}  % correlation, mid, six full-set: min
\providecommand{\rhebCorrFullMax}{}\renewcommand{\rhebCorrFullMax}{0.03}  % ... max
\providecommand{\rhebRsqFullMin}{}\renewcommand{\rhebRsqFullMin}{0.00}  % CAPM R^2, mid, six full-set: min
\providecommand{\rhebRsqFullMax}{}\renewcommand{\rhebRsqFullMax}{0.00}  % ... max
\providecommand{\rhebAlphaFullMin}{}\renewcommand{\rhebAlphaFullMin}{11.5}  % alpha % per day, mid, six full-set: min
\providecommand{\rhebAlphaFullMax}{}\renewcommand{\rhebAlphaFullMax}{14.1}  % ... max
\providecommand{\rhebAlphaTFullMin}{}\renewcommand{\rhebAlphaTFullMin}{2.93}  % alpha HAC t, mid, six full-set: min
\providecommand{\rhebAlphaTFullMax}{}\renewcommand{\rhebAlphaTFullMax}{3.77}  % ... max
\providecommand{\rhebAlphaPFullMin}{}\renewcommand{\rhebAlphaPFullMin}{\ensuremath{<}0.001}  % alpha fixed-b p, mid, six full-set: min
\providecommand{\rhebAlphaPFullMax}{}\renewcommand{\rhebAlphaPFullMax}{0.008}  % ... max
\providecommand{\rhebARFullMin}{}\renewcommand{\rhebARFullMin}{1.53}  % appraisal ratio (sqrt q, D19), mid, six full-set: min
\providecommand{\rhebARFullMax}{}\renewcommand{\rhebARFullMax}{1.88}  % ... max
\providecommand{\rhebSRmaxFullMin}{}\renewcommand{\rhebSRmaxFullMin}{1.54}  % max SR STANDARD, mid, six full-set: min
\providecommand{\rhebSRmaxFullMax}{}\renewcommand{\rhebSRmaxFullMax}{1.88}  % ... max
\providecommand{\rhebDSRmaxFullMin}{}\renewcommand{\rhebDSRmaxFullMin}{1.39}  % max SR gain STANDARD, mid, six full-set: min
\providecommand{\rhebDSRmaxFullMax}{}\renewcommand{\rhebDSRmaxFullMax}{1.73}  % ... max
\providecommand{\rhebDSRmaxPFullMin}{}\renewcommand{\rhebDSRmaxPFullMin}{\ensuremath{<}0.001}  % its bootstrap p, mid, six full-set: min
\providecommand{\rhebDSRmaxPFullMax}{}\renewcommand{\rhebDSRmaxPFullMax}{0.003}  % ... max
\providecommand{\rhebDSRmaxLoFullMin}{}\renewcommand{\rhebDSRmaxLoFullMin}{0.18}  % its 95% interval, lower end, six full-set: min
\providecommand{\rhebDSRmaxLoFullMax}{}\renewcommand{\rhebDSRmaxLoFullMax}{0.43}  % ... max
\providecommand{\rhebDSRmaxHiFullMin}{}\renewcommand{\rhebDSRmaxHiFullMin}{2.27}  % its 95% interval, upper end, six full-set: min
\providecommand{\rhebDSRmaxHiFullMax}{}\renewcommand{\rhebDSRmaxHiFullMax}{2.68}  % ... max
\providecommand{\rhebGrsPFullMin}{}\renewcommand{\rhebGrsPFullMin}{0.002}  % GRS p, mid, six full-set: min
\providecommand{\rhebGrsPFullMax}{}\renewcommand{\rhebGrsPFullMax}{0.012}  % ... max
\providecommand{\rhebTanShareFullMin}{}\renewcommand{\rhebTanShareFullMin}{99}  % tangency vol share %, mid, six full-set: min
\providecommand{\rhebTanShareFullMax}{}\renewcommand{\rhebTanShareFullMax}{100}  % ... max
\providecommand{\rhebSRsFullMin}{}\renewcommand{\rhebSRsFullMin}{1.54}  % strategy SR STANDARD, mid, six full-set: min
\providecommand{\rhebSRsFullMax}{}\renewcommand{\rhebSRsFullMax}{1.88}  % ... max
\providecommand{\rhebRtMvFullMin}{}\renewcommand{\rhebRtMvFullMin}{1.66}  % real-time Mv SR STANDARD, mid, six full-set: min
\providecommand{\rhebRtMvFullMax}{}\renewcommand{\rhebRtMvFullMax}{1.95}  % ... max
\providecommand{\rhebRtMvDFullMin}{}\renewcommand{\rhebRtMvDFullMin}{1.16}  % ... minus market, STANDARD, six full-set: min
\providecommand{\rhebRtMvDFullMax}{}\renewcommand{\rhebRtMvDFullMax}{1.45}  % ... max
\providecommand{\rhebRtMvPFullMin}{}\renewcommand{\rhebRtMvPFullMin}{0.071}  % ... Ledoit-Wolf p (D19), six full-set: min
\providecommand{\rhebRtMvPFullMax}{}\renewcommand{\rhebRtMvPFullMax}{0.146}  % ... max
\providecommand{\rhebRtMvDSixFullMin}{}\renewcommand{\rhebRtMvDSixFullMin}{1.76}  % real-time Mv SR D6, mid, six full-set: min
\providecommand{\rhebRtMvDSixFullMax}{}\renewcommand{\rhebRtMvDSixFullMax}{2.03}  % ... max
\providecommand{\rhebRtMvDSixDFullMin}{}\renewcommand{\rhebRtMvDSixDFullMin}{1.07}  % ... minus market, D6, six full-set: min
\providecommand{\rhebRtMvDSixDFullMax}{}\renewcommand{\rhebRtMvDSixDFullMax}{1.35}  % ... max
\providecommand{\rhebRtMvDSixPFullMin}{}\renewcommand{\rhebRtMvDSixPFullMin}{0.135}  % ... bootstrap p, D6, six full-set: min
\providecommand{\rhebRtMvDSixPFullMax}{}\renewcommand{\rhebRtMvDSixPFullMax}{0.243}  % ... max
\providecommand{\rhebExRtMvFullMin}{}\renewcommand{\rhebExRtMvFullMin}{1.02}  % w/o best days: RT Mv SR STANDARD: min
\providecommand{\rhebExRtMvFullMax}{}\renewcommand{\rhebExRtMvFullMax}{1.39}  % ... max
\providecommand{\rhebExRtMvPFullMin}{}\renewcommand{\rhebExRtMvPFullMin}{0.295}  % ... its Ledoit-Wolf p (D19): min
\providecommand{\rhebExRtMvPFullMax}{}\renewcommand{\rhebExRtMvPFullMax}{0.553}  % ... max
\providecommand{\rhebRtIvFullMin}{}\renewcommand{\rhebRtIvFullMin}{1.22}  % real-time Iv SR STANDARD, mid, six full-set: min
\providecommand{\rhebRtIvFullMax}{}\renewcommand{\rhebRtIvFullMax}{1.61}  % ... max
\providecommand{\rhebRtIvDFullMin}{}\renewcommand{\rhebRtIvDFullMin}{0.72}  % ... minus market, STANDARD, six full-set: min
\providecommand{\rhebRtIvDFullMax}{}\renewcommand{\rhebRtIvDFullMax}{1.11}  % ... max
\providecommand{\rhebRtIvPFullMin}{}\renewcommand{\rhebRtIvPFullMin}{0.054}  % ... Ledoit-Wolf p (D19), six full-set: min
\providecommand{\rhebRtIvPFullMax}{}\renewcommand{\rhebRtIvPFullMax}{0.169}  % ... max
\providecommand{\rhebRtIvDSixFullMin}{}\renewcommand{\rhebRtIvDSixFullMin}{1.72}  % real-time Iv SR D6, mid, six full-set: min
\providecommand{\rhebRtIvDSixFullMax}{}\renewcommand{\rhebRtIvDSixFullMax}{1.95}  % ... max
\providecommand{\rhebRtIvDSixDFullMin}{}\renewcommand{\rhebRtIvDSixDFullMin}{1.04}  % ... minus market, D6, six full-set: min
\providecommand{\rhebRtIvDSixDFullMax}{}\renewcommand{\rhebRtIvDSixDFullMax}{1.27}  % ... max
\providecommand{\rhebRtIvDSixPFullMin}{}\renewcommand{\rhebRtIvDSixPFullMin}{0.037}  % ... bootstrap p, D6, six full-set: min
\providecommand{\rhebRtIvDSixPFullMax}{}\renewcommand{\rhebRtIvDSixPFullMax}{0.100}  % ... max
\providecommand{\rhebExRtIvFullMin}{}\renewcommand{\rhebExRtIvFullMin}{0.95}  % w/o best days: RT Iv SR STANDARD: min
\providecommand{\rhebExRtIvFullMax}{}\renewcommand{\rhebExRtIvFullMax}{1.34}  % ... max
\providecommand{\rhebExRtIvPFullMin}{}\renewcommand{\rhebExRtIvPFullMin}{0.159}  % ... its Ledoit-Wolf p (D19): min
\providecommand{\rhebExRtIvPFullMax}{}\renewcommand{\rhebExRtIvPFullMax}{0.452}  % ... max
\providecommand{\rhebRtMvShareFullMin}{}\renewcommand{\rhebRtMvShareFullMin}{90}  % RT MV median vol share %, six full-set: min
\providecommand{\rhebRtMvShareFullMax}{}\renewcommand{\rhebRtMvShareFullMax}{90}  % ... max
\providecommand{\rhebRtMvCapDaysFullMin}{}\renewcommand{\rhebRtMvCapDaysFullMin}{100}  % RT MV: % of days at the cap, six full-set: min
\providecommand{\rhebRtMvCapDaysFullMax}{}\renewcommand{\rhebRtMvCapDaysFullMax}{100}  % ... max
\providecommand{\rhebRtMvWsFullMin}{}\renewcommand{\rhebRtMvWsFullMin}{3.3}  % RT MV median dollar weight %, six full-set: min
\providecommand{\rhebRtMvWsFullMax}{}\renewcommand{\rhebRtMvWsFullMax}{3.5}  % ... max
\providecommand{\rhebRtIvWsFullMin}{}\renewcommand{\rhebRtIvWsFullMin}{0.9}  % RT IV median dollar weight %, six full-set: min
\providecommand{\rhebRtIvWsFullMax}{}\renewcommand{\rhebRtIvWsFullMax}{0.9}  % ... max
\providecommand{\rhebExBetaFullMin}{}\renewcommand{\rhebExBetaFullMin}{\ensuremath{-}0.22}  % w/o best days: beta, mid, six full-set: min
\providecommand{\rhebExBetaFullMax}{}\renewcommand{\rhebExBetaFullMax}{3.26}  % ... max
\providecommand{\rhebExBetaTFullMin}{}\renewcommand{\rhebExBetaTFullMin}{\ensuremath{-}0.06}  % ... its HAC t: min
\providecommand{\rhebExBetaTFullMax}{}\renewcommand{\rhebExBetaTFullMax}{0.96}  % ... max
\providecommand{\rhebExAlphaFullMin}{}\renewcommand{\rhebExAlphaFullMin}{5.9}  % w/o best days: alpha %, mid, six full-set: min
\providecommand{\rhebExAlphaFullMax}{}\renewcommand{\rhebExAlphaFullMax}{8.7}  % ... max
\providecommand{\rhebExAlphaTFullMin}{}\renewcommand{\rhebExAlphaTFullMin}{1.98}  % ... its HAC t: min
\providecommand{\rhebExAlphaTFullMax}{}\renewcommand{\rhebExAlphaTFullMax}{3.22}  % ... max
\providecommand{\rhebExAlphaPFullMin}{}\renewcommand{\rhebExAlphaPFullMin}{0.004}  % ... its fixed-b p: min
\providecommand{\rhebExAlphaPFullMax}{}\renewcommand{\rhebExAlphaPFullMax}{0.064}  % ... max
\providecommand{\rhebExARFullMin}{}\renewcommand{\rhebExARFullMin}{0.89}  % ... appraisal ratio (sqrt q, D19): min
\providecommand{\rhebExARFullMax}{}\renewcommand{\rhebExARFullMax}{1.32}  % ... max
\providecommand{\rhebExSRmaxFullMin}{}\renewcommand{\rhebExSRmaxFullMin}{0.92}  % ... max SR STANDARD: min
\providecommand{\rhebExSRmaxFullMax}{}\renewcommand{\rhebExSRmaxFullMax}{1.34}  % ... max
\providecommand{\rhebExGrsPFullMin}{}\renewcommand{\rhebExGrsPFullMin}{0.030}  % ... GRS p: min
\providecommand{\rhebExGrsPFullMax}{}\renewcommand{\rhebExGrsPFullMax}{0.141}  % ... max
\providecommand{\rhebBetaSigFull}{}\renewcommand{\rhebBetaSigFull}{0}  % six full-set: beta sig. 5%, mid
\providecommand{\rhebBetaSigFullWord}{}\renewcommand{\rhebBetaSigFullWord}{zero}  % ... in words
\providecommand{\rhebAlphaSigFull}{}\renewcommand{\rhebAlphaSigFull}{6}  % six full-set: alpha sig. 5% (HAC), mid
\providecommand{\rhebAlphaSigFullWord}{}\renewcommand{\rhebAlphaSigFullWord}{six}  % ... in words
\providecommand{\rhebGrsSigFull}{}\renewcommand{\rhebGrsSigFull}{6}  % six full-set: GRS sig. 5%, mid
\providecommand{\rhebGrsSigFullWord}{}\renewcommand{\rhebGrsSigFullWord}{six}  % ... in words
\providecommand{\rhebBootSigFull}{}\renewcommand{\rhebBootSigFull}{6}  % six full-set: max-SR bootstrap p < 0.05
\providecommand{\rhebBootSigFullWord}{}\renewcommand{\rhebBootSigFullWord}{six}  % ... in words
\providecommand{\rhebRtMvBeatFull}{}\renewcommand{\rhebRtMvBeatFull}{6}  % six full-set: RT Mv SR > mkt (STD)
\providecommand{\rhebRtMvBeatFullWord}{}\renewcommand{\rhebRtMvBeatFullWord}{six}  % ... in words
\providecommand{\rhebRtMvSigFull}{}\renewcommand{\rhebRtMvSigFull}{0}  % ... and Ledoit-Wolf p < 0.05 (D19; was the bootstrap p)
\providecommand{\rhebRtMvSigFullWord}{}\renewcommand{\rhebRtMvSigFullWord}{zero}  % ... in words
\providecommand{\rhebRtMvDSixSigFull}{}\renewcommand{\rhebRtMvDSixSigFull}{0}  % ... D6 difference > 0 with bootstrap p < 0.05
\providecommand{\rhebRtMvDSixSigFullWord}{}\renewcommand{\rhebRtMvDSixSigFullWord}{zero}  % ... in words
\providecommand{\rhebExRtMvSigFull}{}\renewcommand{\rhebExRtMvSigFull}{0}  % ... w/o best days: STANDARD difference > 0 with p < 0.05
\providecommand{\rhebExRtMvSigFullWord}{}\renewcommand{\rhebExRtMvSigFullWord}{zero}  % ... in words
\providecommand{\rhebRtIvBeatFull}{}\renewcommand{\rhebRtIvBeatFull}{6}  % six full-set: RT Iv SR > mkt (STD)
\providecommand{\rhebRtIvBeatFullWord}{}\renewcommand{\rhebRtIvBeatFullWord}{six}  % ... in words
\providecommand{\rhebRtIvSigFull}{}\renewcommand{\rhebRtIvSigFull}{0}  % ... and Ledoit-Wolf p < 0.05 (D19; was the bootstrap p)
\providecommand{\rhebRtIvSigFullWord}{}\renewcommand{\rhebRtIvSigFullWord}{zero}  % ... in words
\providecommand{\rhebRtIvDSixSigFull}{}\renewcommand{\rhebRtIvDSixSigFull}{2}  % ... D6 difference > 0 with bootstrap p < 0.05
\providecommand{\rhebRtIvDSixSigFullWord}{}\renewcommand{\rhebRtIvDSixSigFullWord}{two}  % ... in words
\providecommand{\rhebExRtIvSigFull}{}\renewcommand{\rhebExRtIvSigFull}{0}  % ... w/o best days: STANDARD difference > 0 with p < 0.05
\providecommand{\rhebExRtIvSigFullWord}{}\renewcommand{\rhebExRtIvSigFullWord}{zero}  % ... in words
\providecommand{\rhebExAlphaSigFull}{}\renewcommand{\rhebExAlphaSigFull}{5}  % ... w/o best days: alpha sig.
\providecommand{\rhebExAlphaSigFullWord}{}\renewcommand{\rhebExAlphaSigFullWord}{five}  % ... in words
\providecommand{\rhebExGrsSigFull}{}\renewcommand{\rhebExGrsSigFull}{2}  % ... w/o best days: GRS sig.
\providecommand{\rhebExGrsSigFullWord}{}\renewcommand{\rhebExGrsSigFullWord}{two}  % ... in words
\providecommand{\rhebTmGammaFullMin}{}\renewcommand{\rhebTmGammaFullMin}{\ensuremath{-}2.19}  % TM gamma, mid, six full-set: min
\providecommand{\rhebTmGammaFullMax}{}\renewcommand{\rhebTmGammaFullMax}{\ensuremath{-}1.36}  % ... max
\providecommand{\rhebTmGammaTFullMin}{}\renewcommand{\rhebTmGammaTFullMin}{\ensuremath{-}1.53}  % ... its HAC t: min
\providecommand{\rhebTmGammaTFullMax}{}\renewcommand{\rhebTmGammaTFullMax}{\ensuremath{-}1.02}  % ... max
\providecommand{\rhebTmGammaLongFullMin}{}\renewcommand{\rhebTmGammaLongFullMin}{\ensuremath{-}3.45}  % TM gamma, long days, six full-set: min
\providecommand{\rhebTmGammaLongFullMax}{}\renewcommand{\rhebTmGammaLongFullMax}{11.32}  % ... max
\providecommand{\rhebTmGammaLongTFullMin}{}\renewcommand{\rhebTmGammaLongTFullMin}{\ensuremath{-}0.48}  % ... its HAC t: min
\providecommand{\rhebTmGammaLongTFullMax}{}\renewcommand{\rhebTmGammaLongTFullMax}{1.11}  % ... max
\providecommand{\rhebTmGammaShortFullMin}{}\renewcommand{\rhebTmGammaShortFullMin}{\ensuremath{-}2.03}  % TM gamma, short days, six full-set: min
\providecommand{\rhebTmGammaShortFullMax}{}\renewcommand{\rhebTmGammaShortFullMax}{\ensuremath{-}1.48}  % ... max
\providecommand{\rhebTmGammaShortTFullMin}{}\renewcommand{\rhebTmGammaShortTFullMin}{\ensuremath{-}1.47}  % ... its HAC t: min
\providecommand{\rhebTmGammaShortTFullMax}{}\renewcommand{\rhebTmGammaShortTFullMax}{\ensuremath{-}1.18}  % ... max
\providecommand{\rhebThirtyBetaFullMin}{}\renewcommand{\rhebThirtyBetaFullMin}{7.03}  % beta on r30, mid, six full-set: min
\providecommand{\rhebThirtyBetaFullMax}{}\renewcommand{\rhebThirtyBetaFullMax}{23.20}  % ... max
\providecommand{\rhebThirtyBetaTFullMin}{}\renewcommand{\rhebThirtyBetaTFullMin}{0.29}  % ... its HAC t: min
\providecommand{\rhebThirtyBetaTFullMax}{}\renewcommand{\rhebThirtyBetaTFullMax}{0.77}  % ... max
\providecommand{\rhebThirtyTmGammaFullMin}{}\renewcommand{\rhebThirtyTmGammaFullMin}{\ensuremath{-}122.92}  % TM gamma on r30^2, six full-set: min
\providecommand{\rhebThirtyTmGammaFullMax}{}\renewcommand{\rhebThirtyTmGammaFullMax}{\ensuremath{-}38.65}  % ... max
\providecommand{\rhebThirtyTmGammaTFullMin}{}\renewcommand{\rhebThirtyTmGammaTFullMin}{\ensuremath{-}2.41}  % ... its HAC t: min
\providecommand{\rhebThirtyTmGammaTFullMax}{}\renewcommand{\rhebThirtyTmGammaTFullMax}{\ensuremath{-}0.55}  % ... max
\providecommand{\rhebThirtyTmGammaLongFullMin}{}\renewcommand{\rhebThirtyTmGammaLongFullMin}{195.08}  % ... long days, six full-set: min
\providecommand{\rhebThirtyTmGammaLongFullMax}{}\renewcommand{\rhebThirtyTmGammaLongFullMax}{389.30}  % ... max
\providecommand{\rhebThirtyTmGammaLongTFullMin}{}\renewcommand{\rhebThirtyTmGammaLongTFullMin}{2.38}  % ... its HAC t: min
\providecommand{\rhebThirtyTmGammaLongTFullMax}{}\renewcommand{\rhebThirtyTmGammaLongTFullMax}{3.67}  % ... max
\providecommand{\rhebThirtyTmGammaShortFullMin}{}\renewcommand{\rhebThirtyTmGammaShortFullMin}{\ensuremath{-}377.20}  % ... short days, six full-set: min
\providecommand{\rhebThirtyTmGammaShortFullMax}{}\renewcommand{\rhebThirtyTmGammaShortFullMax}{\ensuremath{-}270.34}  % ... max
\providecommand{\rhebThirtyTmGammaShortTFullMin}{}\renewcommand{\rhebThirtyTmGammaShortTFullMin}{\ensuremath{-}4.46}  % ... its HAC t: min
\providecommand{\rhebThirtyTmGammaShortTFullMax}{}\renewcommand{\rhebThirtyTmGammaShortTFullMax}{\ensuremath{-}4.35}  % ... max
\providecommand{\rhebDownMeanFullMin}{}\renewcommand{\rhebDownMeanFullMin}{13.1}  % mean on down days %, six full-set: min
\providecommand{\rhebDownMeanFullMax}{}\renewcommand{\rhebDownMeanFullMax}{18.7}  % ... max
\providecommand{\rhebDownRestFullMin}{}\renewcommand{\rhebDownRestFullMin}{7.2}  % ... on the other days: min
\providecommand{\rhebDownRestFullMax}{}\renewcommand{\rhebDownRestFullMax}{12.6}  % ... max
\providecommand{\rhebDownDiffTFullMin}{}\renewcommand{\rhebDownDiffTFullMin}{0.27}  % ... HAC t of the difference: min
\providecommand{\rhebDownDiffTFullMax}{}\renewcommand{\rhebDownDiffTFullMax}{1.02}  % ... max
\providecommand{\rhebWorstMeanFullMin}{}\renewcommand{\rhebWorstMeanFullMin}{\ensuremath{-}15.7}  % mean on worst-decile days, six full-set: min
\providecommand{\rhebWorstMeanFullMax}{}\renewcommand{\rhebWorstMeanFullMax}{\ensuremath{-}3.3}  % ... max
\providecommand{\rhebWorstRestFullMin}{}\renewcommand{\rhebWorstRestFullMin}{14.5}  % ... on the other days: min
\providecommand{\rhebWorstRestFullMax}{}\renewcommand{\rhebWorstRestFullMax}{17.2}  % ... max
\providecommand{\rhebWorstDiffTFullMin}{}\renewcommand{\rhebWorstDiffTFullMin}{\ensuremath{-}2.28}  % ... HAC t of the difference: min
\providecommand{\rhebWorstDiffTFullMax}{}\renewcommand{\rhebWorstDiffTFullMax}{\ensuremath{-}1.31}  % ... max
\providecommand{\rhebTmSigFull}{}\renewcommand{\rhebTmSigFull}{0}  % six full-set: TM gamma sig. 5%
\providecommand{\rhebTmSigFullWord}{}\renewcommand{\rhebTmSigFullWord}{zero}  % ... in words
\providecommand{\rhebThirtyBetaSigFull}{}\renewcommand{\rhebThirtyBetaSigFull}{0}  % six full-set: r30 beta sig.
\providecommand{\rhebThirtyBetaSigFullWord}{}\renewcommand{\rhebThirtyBetaSigFullWord}{zero}  % ... in words
\providecommand{\rhebDownSigFull}{}\renewcommand{\rhebDownSigFull}{0}  % six full-set: down-day diff sig.
\providecommand{\rhebDownSigFullWord}{}\renewcommand{\rhebDownSigFullWord}{zero}  % ... in words
\providecommand{\rhebWorstSigFull}{}\renewcommand{\rhebWorstSigFull}{3}  % six full-set: worst-decile sig.
\providecommand{\rhebWorstSigFullWord}{}\renewcommand{\rhebWorstSigFullWord}{three}  % ... in words
\providecommand{\rhebBetaHarMin}{}\renewcommand{\rhebBetaHarMin}{0.14}  % CAPM beta, mid, four distinct HAR-set: min
\providecommand{\rhebBetaHarMax}{}\renewcommand{\rhebBetaHarMax}{2.44}  % ... max
\providecommand{\rhebBetaTHarMin}{}\renewcommand{\rhebBetaTHarMin}{0.04}  % ... its HAC t, mid, four distinct HAR-set: min
\providecommand{\rhebBetaTHarMax}{}\renewcommand{\rhebBetaTHarMax}{0.73}  % ... max
\providecommand{\rhebBetaPHarMin}{}\renewcommand{\rhebBetaPHarMin}{0.483}  % ... its fixed-b p, mid, four distinct HAR-set: min
\providecommand{\rhebBetaPHarMax}{}\renewcommand{\rhebBetaPHarMax}{0.967}  % ... max
\providecommand{\rhebCorrHarMin}{}\renewcommand{\rhebCorrHarMin}{0.00}  % correlation, mid, four distinct HAR-set: min
\providecommand{\rhebCorrHarMax}{}\renewcommand{\rhebCorrHarMax}{0.02}  % ... max
\providecommand{\rhebRsqHarMin}{}\renewcommand{\rhebRsqHarMin}{0.00}  % CAPM R^2, mid, four distinct HAR-set: min
\providecommand{\rhebRsqHarMax}{}\renewcommand{\rhebRsqHarMax}{0.00}  % ... max
\providecommand{\rhebAlphaHarMin}{}\renewcommand{\rhebAlphaHarMin}{2.4}  % alpha % per day, mid, four distinct HAR-set: min
\providecommand{\rhebAlphaHarMax}{}\renewcommand{\rhebAlphaHarMax}{7.4}  % ... max
\providecommand{\rhebAlphaTHarMin}{}\renewcommand{\rhebAlphaTHarMin}{0.61}  % alpha HAC t, mid, four distinct HAR-set: min
\providecommand{\rhebAlphaTHarMax}{}\renewcommand{\rhebAlphaTHarMax}{2.05}  % ... max
\providecommand{\rhebAlphaPHarMin}{}\renewcommand{\rhebAlphaPHarMin}{0.056}  % alpha fixed-b p, mid, four distinct HAR-set: min
\providecommand{\rhebAlphaPHarMax}{}\renewcommand{\rhebAlphaPHarMax}{0.556}  % ... max
\providecommand{\rhebARHarMin}{}\renewcommand{\rhebARHarMin}{0.31}  % appraisal ratio (sqrt q, D19), mid, four distinct HAR-set: min
\providecommand{\rhebARHarMax}{}\renewcommand{\rhebARHarMax}{0.98}  % ... max
\providecommand{\rhebSRmaxHarMin}{}\renewcommand{\rhebSRmaxHarMin}{0.35}  % max SR STANDARD, mid, four distinct HAR-set: min
\providecommand{\rhebSRmaxHarMax}{}\renewcommand{\rhebSRmaxHarMax}{0.99}  % ... max
\providecommand{\rhebDSRmaxHarMin}{}\renewcommand{\rhebDSRmaxHarMin}{0.20}  % max SR gain STANDARD, mid, four distinct HAR-set: min
\providecommand{\rhebDSRmaxHarMax}{}\renewcommand{\rhebDSRmaxHarMax}{0.84}  % ... max
\providecommand{\rhebDSRmaxPHarMin}{}\renewcommand{\rhebDSRmaxPHarMin}{0.028}  % its bootstrap p, mid, four distinct HAR-set: min
\providecommand{\rhebDSRmaxPHarMax}{}\renewcommand{\rhebDSRmaxPHarMax}{0.302}  % ... max
\providecommand{\rhebDSRmaxLoHarMin}{}\renewcommand{\rhebDSRmaxLoHarMin}{0.00}  % its 95% interval, lower end, four distinct HAR-set: min
\providecommand{\rhebDSRmaxLoHarMax}{}\renewcommand{\rhebDSRmaxLoHarMax}{0.01}  % ... max
\providecommand{\rhebDSRmaxHiHarMin}{}\renewcommand{\rhebDSRmaxHiHarMin}{1.22}  % its 95% interval, upper end, four distinct HAR-set: min
\providecommand{\rhebDSRmaxHiHarMax}{}\renewcommand{\rhebDSRmaxHiHarMax}{1.76}  % ... max
\providecommand{\rhebGrsPHarMin}{}\renewcommand{\rhebGrsPHarMin}{0.106}  % GRS p, mid, four distinct HAR-set: min
\providecommand{\rhebGrsPHarMax}{}\renewcommand{\rhebGrsPHarMax}{0.605}  % ... max
\providecommand{\rhebTanShareHarMin}{}\renewcommand{\rhebTanShareHarMin}{82}  % tangency vol share %, mid, four distinct HAR-set: min
\providecommand{\rhebTanShareHarMax}{}\renewcommand{\rhebTanShareHarMax}{98}  % ... max
\providecommand{\rhebSRsHarMin}{}\renewcommand{\rhebSRsHarMin}{0.32}  % strategy SR STANDARD, mid, four distinct HAR-set: min
\providecommand{\rhebSRsHarMax}{}\renewcommand{\rhebSRsHarMax}{0.98}  % ... max
\providecommand{\rhebRtMvHarMin}{}\renewcommand{\rhebRtMvHarMin}{0.74}  % real-time Mv SR STANDARD, mid, four distinct HAR-set: min
\providecommand{\rhebRtMvHarMax}{}\renewcommand{\rhebRtMvHarMax}{1.24}  % ... max
\providecommand{\rhebRtMvDHarMin}{}\renewcommand{\rhebRtMvDHarMin}{0.24}  % ... minus market, STANDARD, four distinct HAR-set: min
\providecommand{\rhebRtMvDHarMax}{}\renewcommand{\rhebRtMvDHarMax}{0.74}  % ... max
\providecommand{\rhebRtMvPHarMin}{}\renewcommand{\rhebRtMvPHarMin}{0.346}  % ... Ledoit-Wolf p (D19), four distinct HAR-set: min
\providecommand{\rhebRtMvPHarMax}{}\renewcommand{\rhebRtMvPHarMax}{0.776}  % ... max
\providecommand{\rhebRtMvDSixHarMin}{}\renewcommand{\rhebRtMvDSixHarMin}{0.32}  % real-time Mv SR D6, mid, four distinct HAR-set: min
\providecommand{\rhebRtMvDSixHarMax}{}\renewcommand{\rhebRtMvDSixHarMax}{1.15}  % ... max
\providecommand{\rhebRtMvDSixDHarMin}{}\renewcommand{\rhebRtMvDSixDHarMin}{\ensuremath{-}0.36}  % ... minus market, D6, four distinct HAR-set: min
\providecommand{\rhebRtMvDSixDHarMax}{}\renewcommand{\rhebRtMvDSixDHarMax}{0.47}  % ... max
\providecommand{\rhebRtMvDSixPHarMin}{}\renewcommand{\rhebRtMvDSixPHarMin}{0.619}  % ... bootstrap p, D6, four distinct HAR-set: min
\providecommand{\rhebRtMvDSixPHarMax}{}\renewcommand{\rhebRtMvDSixPHarMax}{0.963}  % ... max
\providecommand{\rhebExRtMvHarMin}{}\renewcommand{\rhebExRtMvHarMin}{0.69}  % w/o best days: RT Mv SR STANDARD: min
\providecommand{\rhebExRtMvHarMax}{}\renewcommand{\rhebExRtMvHarMax}{1.14}  % ... max
\providecommand{\rhebExRtMvPHarMin}{}\renewcommand{\rhebExRtMvPHarMin}{0.472}  % ... its Ledoit-Wolf p (D19): min
\providecommand{\rhebExRtMvPHarMax}{}\renewcommand{\rhebExRtMvPHarMax}{0.868}  % ... max
\providecommand{\rhebRtIvHarMin}{}\renewcommand{\rhebRtIvHarMin}{0.87}  % real-time Iv SR STANDARD, mid, four distinct HAR-set: min
\providecommand{\rhebRtIvHarMax}{}\renewcommand{\rhebRtIvHarMax}{1.11}  % ... max
\providecommand{\rhebRtIvDHarMin}{}\renewcommand{\rhebRtIvDHarMin}{0.37}  % ... minus market, STANDARD, four distinct HAR-set: min
\providecommand{\rhebRtIvDHarMax}{}\renewcommand{\rhebRtIvDHarMax}{0.61}  % ... max
\providecommand{\rhebRtIvPHarMin}{}\renewcommand{\rhebRtIvPHarMin}{0.273}  % ... Ledoit-Wolf p (D19), four distinct HAR-set: min
\providecommand{\rhebRtIvPHarMax}{}\renewcommand{\rhebRtIvPHarMax}{0.501}  % ... max
\providecommand{\rhebRtIvDSixHarMin}{}\renewcommand{\rhebRtIvDSixHarMin}{0.92}  % real-time Iv SR D6, mid, four distinct HAR-set: min
\providecommand{\rhebRtIvDSixHarMax}{}\renewcommand{\rhebRtIvDSixHarMax}{1.26}  % ... max
\providecommand{\rhebRtIvDSixDHarMin}{}\renewcommand{\rhebRtIvDSixDHarMin}{0.24}  % ... minus market, D6, four distinct HAR-set: min
\providecommand{\rhebRtIvDSixDHarMax}{}\renewcommand{\rhebRtIvDSixDHarMax}{0.58}  % ... max
\providecommand{\rhebRtIvDSixPHarMin}{}\renewcommand{\rhebRtIvDSixPHarMin}{0.368}  % ... bootstrap p, D6, four distinct HAR-set: min
\providecommand{\rhebRtIvDSixPHarMax}{}\renewcommand{\rhebRtIvDSixPHarMax}{0.712}  % ... max
\providecommand{\rhebExRtIvHarMin}{}\renewcommand{\rhebExRtIvHarMin}{0.93}  % w/o best days: RT Iv SR STANDARD: min
\providecommand{\rhebExRtIvHarMax}{}\renewcommand{\rhebExRtIvHarMax}{1.13}  % ... max
\providecommand{\rhebExRtIvPHarMin}{}\renewcommand{\rhebExRtIvPHarMin}{0.301}  % ... its Ledoit-Wolf p (D19): min
\providecommand{\rhebExRtIvPHarMax}{}\renewcommand{\rhebExRtIvPHarMax}{0.493}  % ... max
\providecommand{\rhebRtMvShareHarMin}{}\renewcommand{\rhebRtMvShareHarMin}{90}  % RT MV median vol share %, four distinct HAR-set: min
\providecommand{\rhebRtMvShareHarMax}{}\renewcommand{\rhebRtMvShareHarMax}{90}  % ... max
\providecommand{\rhebRtMvCapDaysHarMin}{}\renewcommand{\rhebRtMvCapDaysHarMin}{86}  % RT MV: % of days at the cap, four distinct HAR-set: min
\providecommand{\rhebRtMvCapDaysHarMax}{}\renewcommand{\rhebRtMvCapDaysHarMax}{100}  % ... max
\providecommand{\rhebRtMvWsHarMin}{}\renewcommand{\rhebRtMvWsHarMin}{3.3}  % RT MV median dollar weight %, four distinct HAR-set: min
\providecommand{\rhebRtMvWsHarMax}{}\renewcommand{\rhebRtMvWsHarMax}{3.5}  % ... max
\providecommand{\rhebRtIvWsHarMin}{}\renewcommand{\rhebRtIvWsHarMin}{0.9}  % RT IV median dollar weight %, four distinct HAR-set: min
\providecommand{\rhebRtIvWsHarMax}{}\renewcommand{\rhebRtIvWsHarMax}{0.9}  % ... max
\providecommand{\rhebExBetaHarMin}{}\renewcommand{\rhebExBetaHarMin}{\ensuremath{-}0.34}  % w/o best days: beta, mid, four distinct HAR-set: min
\providecommand{\rhebExBetaHarMax}{}\renewcommand{\rhebExBetaHarMax}{1.57}  % ... max
\providecommand{\rhebExBetaTHarMin}{}\renewcommand{\rhebExBetaTHarMin}{\ensuremath{-}0.11}  % ... its HAC t: min
\providecommand{\rhebExBetaTHarMax}{}\renewcommand{\rhebExBetaTHarMax}{0.49}  % ... max
\providecommand{\rhebExAlphaHarMin}{}\renewcommand{\rhebExAlphaHarMin}{2.4}  % w/o best days: alpha %, mid, four distinct HAR-set: min
\providecommand{\rhebExAlphaHarMax}{}\renewcommand{\rhebExAlphaHarMax}{4.6}  % ... max
\providecommand{\rhebExAlphaTHarMin}{}\renewcommand{\rhebExAlphaTHarMin}{0.67}  % ... its HAC t: min
\providecommand{\rhebExAlphaTHarMax}{}\renewcommand{\rhebExAlphaTHarMax}{1.40}  % ... max
\providecommand{\rhebExAlphaPHarMin}{}\renewcommand{\rhebExAlphaPHarMin}{0.182}  % ... its fixed-b p: min
\providecommand{\rhebExAlphaPHarMax}{}\renewcommand{\rhebExAlphaPHarMax}{0.521}  % ... max
\providecommand{\rhebExARHarMin}{}\renewcommand{\rhebExARHarMin}{0.36}  % ... appraisal ratio (sqrt q, D19): min
\providecommand{\rhebExARHarMax}{}\renewcommand{\rhebExARHarMax}{0.71}  % ... max
\providecommand{\rhebExSRmaxHarMin}{}\renewcommand{\rhebExSRmaxHarMin}{0.43}  % ... max SR STANDARD: min
\providecommand{\rhebExSRmaxHarMax}{}\renewcommand{\rhebExSRmaxHarMax}{0.74}  % ... max
\providecommand{\rhebExGrsPHarMin}{}\renewcommand{\rhebExGrsPHarMin}{0.244}  % ... GRS p: min
\providecommand{\rhebExGrsPHarMax}{}\renewcommand{\rhebExGrsPHarMax}{0.552}  % ... max
\providecommand{\rhebBetaSigHar}{}\renewcommand{\rhebBetaSigHar}{0}  % four distinct HAR-set: beta sig. 5%, mid
\providecommand{\rhebBetaSigHarWord}{}\renewcommand{\rhebBetaSigHarWord}{zero}  % ... in words
\providecommand{\rhebAlphaSigHar}{}\renewcommand{\rhebAlphaSigHar}{0}  % four distinct HAR-set: alpha sig. 5% (HAC), mid
\providecommand{\rhebAlphaSigHarWord}{}\renewcommand{\rhebAlphaSigHarWord}{zero}  % ... in words
\providecommand{\rhebGrsSigHar}{}\renewcommand{\rhebGrsSigHar}{0}  % four distinct HAR-set: GRS sig. 5%, mid
\providecommand{\rhebGrsSigHarWord}{}\renewcommand{\rhebGrsSigHarWord}{zero}  % ... in words
\providecommand{\rhebBootSigHar}{}\renewcommand{\rhebBootSigHar}{2}  % four distinct HAR-set: max-SR bootstrap p < 0.05
\providecommand{\rhebBootSigHarWord}{}\renewcommand{\rhebBootSigHarWord}{two}  % ... in words
\providecommand{\rhebRtMvBeatHar}{}\renewcommand{\rhebRtMvBeatHar}{4}  % four distinct HAR-set: RT Mv SR > mkt (STD)
\providecommand{\rhebRtMvBeatHarWord}{}\renewcommand{\rhebRtMvBeatHarWord}{four}  % ... in words
\providecommand{\rhebRtMvSigHar}{}\renewcommand{\rhebRtMvSigHar}{0}  % ... and Ledoit-Wolf p < 0.05 (D19; was the bootstrap p)
\providecommand{\rhebRtMvSigHarWord}{}\renewcommand{\rhebRtMvSigHarWord}{zero}  % ... in words
\providecommand{\rhebRtMvDSixSigHar}{}\renewcommand{\rhebRtMvDSixSigHar}{0}  % ... D6 difference > 0 with bootstrap p < 0.05
\providecommand{\rhebRtMvDSixSigHarWord}{}\renewcommand{\rhebRtMvDSixSigHarWord}{zero}  % ... in words
\providecommand{\rhebExRtMvSigHar}{}\renewcommand{\rhebExRtMvSigHar}{0}  % ... w/o best days: STANDARD difference > 0 with p < 0.05
\providecommand{\rhebExRtMvSigHarWord}{}\renewcommand{\rhebExRtMvSigHarWord}{zero}  % ... in words
\providecommand{\rhebRtIvBeatHar}{}\renewcommand{\rhebRtIvBeatHar}{4}  % four distinct HAR-set: RT Iv SR > mkt (STD)
\providecommand{\rhebRtIvBeatHarWord}{}\renewcommand{\rhebRtIvBeatHarWord}{four}  % ... in words
\providecommand{\rhebRtIvSigHar}{}\renewcommand{\rhebRtIvSigHar}{0}  % ... and Ledoit-Wolf p < 0.05 (D19; was the bootstrap p)
\providecommand{\rhebRtIvSigHarWord}{}\renewcommand{\rhebRtIvSigHarWord}{zero}  % ... in words
\providecommand{\rhebRtIvDSixSigHar}{}\renewcommand{\rhebRtIvDSixSigHar}{0}  % ... D6 difference > 0 with bootstrap p < 0.05
\providecommand{\rhebRtIvDSixSigHarWord}{}\renewcommand{\rhebRtIvDSixSigHarWord}{zero}  % ... in words
\providecommand{\rhebExRtIvSigHar}{}\renewcommand{\rhebExRtIvSigHar}{0}  % ... w/o best days: STANDARD difference > 0 with p < 0.05
\providecommand{\rhebExRtIvSigHarWord}{}\renewcommand{\rhebExRtIvSigHarWord}{zero}  % ... in words
\providecommand{\rhebExAlphaSigHar}{}\renewcommand{\rhebExAlphaSigHar}{0}  % ... w/o best days: alpha sig.
\providecommand{\rhebExAlphaSigHarWord}{}\renewcommand{\rhebExAlphaSigHarWord}{zero}  % ... in words
\providecommand{\rhebExGrsSigHar}{}\renewcommand{\rhebExGrsSigHar}{0}  % ... w/o best days: GRS sig.
\providecommand{\rhebExGrsSigHarWord}{}\renewcommand{\rhebExGrsSigHarWord}{zero}  % ... in words
\providecommand{\rhebTmGammaHarMin}{}\renewcommand{\rhebTmGammaHarMin}{\ensuremath{-}0.82}  % TM gamma, mid, four distinct HAR-set: min
\providecommand{\rhebTmGammaHarMax}{}\renewcommand{\rhebTmGammaHarMax}{0.10}  % ... max
\providecommand{\rhebTmGammaTHarMin}{}\renewcommand{\rhebTmGammaTHarMin}{\ensuremath{-}0.79}  % ... its HAC t: min
\providecommand{\rhebTmGammaTHarMax}{}\renewcommand{\rhebTmGammaTHarMax}{0.09}  % ... max
\providecommand{\rhebTmGammaLongHarMin}{}\renewcommand{\rhebTmGammaLongHarMin}{10.47}  % TM gamma, long days, four distinct HAR-set: min
\providecommand{\rhebTmGammaLongHarMax}{}\renewcommand{\rhebTmGammaLongHarMax}{12.79}  % ... max
\providecommand{\rhebTmGammaLongTHarMin}{}\renewcommand{\rhebTmGammaLongTHarMin}{1.80}  % ... its HAC t: min
\providecommand{\rhebTmGammaLongTHarMax}{}\renewcommand{\rhebTmGammaLongTHarMax}{1.88}  % ... max
\providecommand{\rhebTmGammaShortHarMin}{}\renewcommand{\rhebTmGammaShortHarMin}{\ensuremath{-}1.21}  % TM gamma, short days, four distinct HAR-set: min
\providecommand{\rhebTmGammaShortHarMax}{}\renewcommand{\rhebTmGammaShortHarMax}{\ensuremath{-}0.53}  % ... max
\providecommand{\rhebTmGammaShortTHarMin}{}\renewcommand{\rhebTmGammaShortTHarMin}{\ensuremath{-}1.01}  % ... its HAC t: min
\providecommand{\rhebTmGammaShortTHarMax}{}\renewcommand{\rhebTmGammaShortTHarMax}{\ensuremath{-}0.40}  % ... max
\providecommand{\rhebThirtyBetaHarMin}{}\renewcommand{\rhebThirtyBetaHarMin}{14.45}  % beta on r30, mid, four distinct HAR-set: min
\providecommand{\rhebThirtyBetaHarMax}{}\renewcommand{\rhebThirtyBetaHarMax}{55.15}  % ... max
\providecommand{\rhebThirtyBetaTHarMin}{}\renewcommand{\rhebThirtyBetaTHarMin}{0.52}  % ... its HAC t: min
\providecommand{\rhebThirtyBetaTHarMax}{}\renewcommand{\rhebThirtyBetaTHarMax}{1.93}  % ... max
\providecommand{\rhebThirtyTmGammaHarMin}{}\renewcommand{\rhebThirtyTmGammaHarMin}{\ensuremath{-}149.56}  % TM gamma on r30^2, four distinct HAR-set: min
\providecommand{\rhebThirtyTmGammaHarMax}{}\renewcommand{\rhebThirtyTmGammaHarMax}{\ensuremath{-}96.47}  % ... max
\providecommand{\rhebThirtyTmGammaTHarMin}{}\renewcommand{\rhebThirtyTmGammaTHarMin}{\ensuremath{-}3.09}  % ... its HAC t: min
\providecommand{\rhebThirtyTmGammaTHarMax}{}\renewcommand{\rhebThirtyTmGammaTHarMax}{\ensuremath{-}2.07}  % ... max
\providecommand{\rhebThirtyTmGammaLongHarMin}{}\renewcommand{\rhebThirtyTmGammaLongHarMin}{250.00}  % ... long days, four distinct HAR-set: min
\providecommand{\rhebThirtyTmGammaLongHarMax}{}\renewcommand{\rhebThirtyTmGammaLongHarMax}{396.11}  % ... max
\providecommand{\rhebThirtyTmGammaLongTHarMin}{}\renewcommand{\rhebThirtyTmGammaLongTHarMin}{2.28}  % ... its HAC t: min
\providecommand{\rhebThirtyTmGammaLongTHarMax}{}\renewcommand{\rhebThirtyTmGammaLongTHarMax}{2.92}  % ... max
\providecommand{\rhebThirtyTmGammaShortHarMin}{}\renewcommand{\rhebThirtyTmGammaShortHarMin}{\ensuremath{-}291.78}  % ... short days, four distinct HAR-set: min
\providecommand{\rhebThirtyTmGammaShortHarMax}{}\renewcommand{\rhebThirtyTmGammaShortHarMax}{\ensuremath{-}266.36}  % ... max
\providecommand{\rhebThirtyTmGammaShortTHarMin}{}\renewcommand{\rhebThirtyTmGammaShortTHarMin}{\ensuremath{-}4.42}  % ... its HAC t: min
\providecommand{\rhebThirtyTmGammaShortTHarMax}{}\renewcommand{\rhebThirtyTmGammaShortTHarMax}{\ensuremath{-}3.80}  % ... max
\providecommand{\rhebDownMeanHarMin}{}\renewcommand{\rhebDownMeanHarMin}{\ensuremath{-}0.5}  % mean on down days %, four distinct HAR-set: min
\providecommand{\rhebDownMeanHarMax}{}\renewcommand{\rhebDownMeanHarMax}{10.6}  % ... max
\providecommand{\rhebDownRestHarMin}{}\renewcommand{\rhebDownRestHarMin}{3.8}  % ... on the other days: min
\providecommand{\rhebDownRestHarMax}{}\renewcommand{\rhebDownRestHarMax}{7.3}  % ... max
\providecommand{\rhebDownDiffTHarMin}{}\renewcommand{\rhebDownDiffTHarMin}{\ensuremath{-}0.57}  % ... HAC t of the difference: min
\providecommand{\rhebDownDiffTHarMax}{}\renewcommand{\rhebDownDiffTHarMax}{0.67}  % ... max
\providecommand{\rhebWorstMeanHarMin}{}\renewcommand{\rhebWorstMeanHarMin}{\ensuremath{-}8.1}  % mean on worst-decile days, four distinct HAR-set: min
\providecommand{\rhebWorstMeanHarMax}{}\renewcommand{\rhebWorstMeanHarMax}{\ensuremath{-}2.3}  % ... max
\providecommand{\rhebWorstRestHarMin}{}\renewcommand{\rhebWorstRestHarMin}{3.3}  % ... on the other days: min
\providecommand{\rhebWorstRestHarMax}{}\renewcommand{\rhebWorstRestHarMax}{8.9}  % ... max
\providecommand{\rhebWorstDiffTHarMin}{}\renewcommand{\rhebWorstDiffTHarMin}{\ensuremath{-}1.21}  % ... HAC t of the difference: min
\providecommand{\rhebWorstDiffTHarMax}{}\renewcommand{\rhebWorstDiffTHarMax}{\ensuremath{-}0.53}  % ... max
\providecommand{\rhebTmSigHar}{}\renewcommand{\rhebTmSigHar}{0}  % four distinct HAR-set: TM gamma sig. 5%
\providecommand{\rhebTmSigHarWord}{}\renewcommand{\rhebTmSigHarWord}{zero}  % ... in words
\providecommand{\rhebThirtyBetaSigHar}{}\renewcommand{\rhebThirtyBetaSigHar}{0}  % four distinct HAR-set: r30 beta sig.
\providecommand{\rhebThirtyBetaSigHarWord}{}\renewcommand{\rhebThirtyBetaSigHarWord}{zero}  % ... in words
\providecommand{\rhebDownSigHar}{}\renewcommand{\rhebDownSigHar}{0}  % four distinct HAR-set: down-day diff sig.
\providecommand{\rhebDownSigHarWord}{}\renewcommand{\rhebDownSigHarWord}{zero}  % ... in words
\providecommand{\rhebWorstSigHar}{}\renewcommand{\rhebWorstSigHar}{0}  % four distinct HAR-set: worst-decile sig.
\providecommand{\rhebWorstSigHarWord}{}\renewcommand{\rhebWorstSigHarWord}{zero}  % ... in words
\providecommand{\rhebBetaFullXMin}{}\renewcommand{\rhebBetaFullXMin}{\ensuremath{-}0.45}  % CAPM beta, bid/ask, six full-set: min
\providecommand{\rhebBetaFullXMax}{}\renewcommand{\rhebBetaFullXMax}{2.99}  % ... max
\providecommand{\rhebBetaTFullXMin}{}\renewcommand{\rhebBetaTFullXMin}{\ensuremath{-}0.11}  % ... its HAC t, bid/ask, six full-set: min
\providecommand{\rhebBetaTFullXMax}{}\renewcommand{\rhebBetaTFullXMax}{0.76}  % ... max
\providecommand{\rhebBetaPFullXMin}{}\renewcommand{\rhebBetaPFullXMin}{0.467}  % ... its fixed-b p, bid/ask, six full-set: min
\providecommand{\rhebBetaPFullXMax}{}\renewcommand{\rhebBetaPFullXMax}{0.913}  % ... max
\providecommand{\rhebCorrFullXMin}{}\renewcommand{\rhebCorrFullXMin}{0.00}  % correlation, bid/ask, six full-set: min
\providecommand{\rhebCorrFullXMax}{}\renewcommand{\rhebCorrFullXMax}{0.03}  % ... max
\providecommand{\rhebRsqFullXMin}{}\renewcommand{\rhebRsqFullXMin}{0.00}  % CAPM R^2, bid/ask, six full-set: min
\providecommand{\rhebRsqFullXMax}{}\renewcommand{\rhebRsqFullXMax}{0.00}  % ... max
\providecommand{\rhebAlphaFullXMin}{}\renewcommand{\rhebAlphaFullXMin}{8.9}  % alpha % per day, bid/ask, six full-set: min
\providecommand{\rhebAlphaFullXMax}{}\renewcommand{\rhebAlphaFullXMax}{11.4}  % ... max
\providecommand{\rhebAlphaTFullXMin}{}\renewcommand{\rhebAlphaTFullXMin}{2.32}  % alpha HAC t, bid/ask, six full-set: min
\providecommand{\rhebAlphaTFullXMax}{}\renewcommand{\rhebAlphaTFullXMax}{3.09}  % ... max
\providecommand{\rhebAlphaPFullXMin}{}\renewcommand{\rhebAlphaPFullXMin}{0.005}  % alpha fixed-b p, bid/ask, six full-set: min
\providecommand{\rhebAlphaPFullXMax}{}\renewcommand{\rhebAlphaPFullXMax}{0.032}  % ... max
\providecommand{\rhebARFullXMin}{}\renewcommand{\rhebARFullXMin}{1.18}  % appraisal ratio (sqrt q, D19), bid/ask, six full-set: min
\providecommand{\rhebARFullXMax}{}\renewcommand{\rhebARFullXMax}{1.52}  % ... max
\providecommand{\rhebSRmaxFullXMin}{}\renewcommand{\rhebSRmaxFullXMin}{1.19}  % max SR STANDARD, bid/ask, six full-set: min
\providecommand{\rhebSRmaxFullXMax}{}\renewcommand{\rhebSRmaxFullXMax}{1.53}  % ... max
\providecommand{\rhebDSRmaxFullXMin}{}\renewcommand{\rhebDSRmaxFullXMin}{1.04}  % max SR gain STANDARD, bid/ask, six full-set: min
\providecommand{\rhebDSRmaxFullXMax}{}\renewcommand{\rhebDSRmaxFullXMax}{1.37}  % ... max
\providecommand{\rhebDSRmaxPFullXMin}{}\renewcommand{\rhebDSRmaxPFullXMin}{0.004}  % its bootstrap p, bid/ask, six full-set: min
\providecommand{\rhebDSRmaxPFullXMax}{}\renewcommand{\rhebDSRmaxPFullXMax}{0.014}  % ... max
\providecommand{\rhebDSRmaxLoFullXMin}{}\renewcommand{\rhebDSRmaxLoFullXMin}{0.03}  % its 95% interval, lower end, six full-set: min
\providecommand{\rhebDSRmaxLoFullXMax}{}\renewcommand{\rhebDSRmaxLoFullXMax}{0.18}  % ... max
\providecommand{\rhebDSRmaxHiFullXMin}{}\renewcommand{\rhebDSRmaxHiFullXMin}{1.94}  % its 95% interval, upper end, six full-set: min
\providecommand{\rhebDSRmaxHiFullXMax}{}\renewcommand{\rhebDSRmaxHiFullXMax}{2.34}  % ... max
\providecommand{\rhebGrsPFullXMin}{}\renewcommand{\rhebGrsPFullXMin}{0.013}  % GRS p, bid/ask, six full-set: min
\providecommand{\rhebGrsPFullXMax}{}\renewcommand{\rhebGrsPFullXMax}{0.053}  % ... max
\providecommand{\rhebTanShareFullXMin}{}\renewcommand{\rhebTanShareFullXMin}{99}  % tangency vol share %, bid/ask, six full-set: min
\providecommand{\rhebTanShareFullXMax}{}\renewcommand{\rhebTanShareFullXMax}{99}  % ... max
\providecommand{\rhebSRsFullXMin}{}\renewcommand{\rhebSRsFullXMin}{1.18}  % strategy SR STANDARD, bid/ask, six full-set: min
\providecommand{\rhebSRsFullXMax}{}\renewcommand{\rhebSRsFullXMax}{1.52}  % ... max
\providecommand{\rhebRtMvFullXMin}{}\renewcommand{\rhebRtMvFullXMin}{1.31}  % real-time Mv SR STANDARD, bid/ask, six full-set: min
\providecommand{\rhebRtMvFullXMax}{}\renewcommand{\rhebRtMvFullXMax}{1.59}  % ... max
\providecommand{\rhebRtMvDFullXMin}{}\renewcommand{\rhebRtMvDFullXMin}{0.81}  % ... minus market, STANDARD, six full-set: min
\providecommand{\rhebRtMvDFullXMax}{}\renewcommand{\rhebRtMvDFullXMax}{1.10}  % ... max
\providecommand{\rhebRtMvPFullXMin}{}\renewcommand{\rhebRtMvPFullXMin}{0.169}  % ... Ledoit-Wolf p (D19), six full-set: min
\providecommand{\rhebRtMvPFullXMax}{}\renewcommand{\rhebRtMvPFullXMax}{0.293}  % ... max
\providecommand{\rhebRtMvDSixFullXMin}{}\renewcommand{\rhebRtMvDSixFullXMin}{1.41}  % real-time Mv SR D6, bid/ask, six full-set: min
\providecommand{\rhebRtMvDSixFullXMax}{}\renewcommand{\rhebRtMvDSixFullXMax}{1.69}  % ... max
\providecommand{\rhebRtMvDSixDFullXMin}{}\renewcommand{\rhebRtMvDSixDFullXMin}{0.73}  % ... minus market, D6, six full-set: min
\providecommand{\rhebRtMvDSixDFullXMax}{}\renewcommand{\rhebRtMvDSixDFullXMax}{1.00}  % ... max
\providecommand{\rhebRtMvDSixPFullXMin}{}\renewcommand{\rhebRtMvDSixPFullXMin}{0.270}  % ... bootstrap p, D6, six full-set: min
\providecommand{\rhebRtMvDSixPFullXMax}{}\renewcommand{\rhebRtMvDSixPFullXMax}{0.420}  % ... max
\providecommand{\rhebExRtMvFullXMin}{}\renewcommand{\rhebExRtMvFullXMin}{0.64}  % w/o best days: RT Mv SR STANDARD: min
\providecommand{\rhebExRtMvFullXMax}{}\renewcommand{\rhebExRtMvFullXMax}{1.02}  % ... max
\providecommand{\rhebExRtMvPFullXMin}{}\renewcommand{\rhebExRtMvPFullXMin}{0.550}  % ... its Ledoit-Wolf p (D19): min
\providecommand{\rhebExRtMvPFullXMax}{}\renewcommand{\rhebExRtMvPFullXMax}{0.899}  % ... max
\providecommand{\rhebRtIvFullXMin}{}\renewcommand{\rhebRtIvFullXMin}{0.94}  % real-time Iv SR STANDARD, bid/ask, six full-set: min
\providecommand{\rhebRtIvFullXMax}{}\renewcommand{\rhebRtIvFullXMax}{1.33}  % ... max
\providecommand{\rhebRtIvDFullXMin}{}\renewcommand{\rhebRtIvDFullXMin}{0.44}  % ... minus market, STANDARD, six full-set: min
\providecommand{\rhebRtIvDFullXMax}{}\renewcommand{\rhebRtIvDFullXMax}{0.83}  % ... max
\providecommand{\rhebRtIvPFullXMin}{}\renewcommand{\rhebRtIvPFullXMin}{0.142}  % ... Ledoit-Wolf p (D19), six full-set: min
\providecommand{\rhebRtIvPFullXMax}{}\renewcommand{\rhebRtIvPFullXMax}{0.388}  % ... max
\providecommand{\rhebRtIvDSixFullXMin}{}\renewcommand{\rhebRtIvDSixFullXMin}{1.45}  % real-time Iv SR D6, bid/ask, six full-set: min
\providecommand{\rhebRtIvDSixFullXMax}{}\renewcommand{\rhebRtIvDSixFullXMax}{1.67}  % ... max
\providecommand{\rhebRtIvDSixDFullXMin}{}\renewcommand{\rhebRtIvDSixDFullXMin}{0.76}  % ... minus market, D6, six full-set: min
\providecommand{\rhebRtIvDSixDFullXMax}{}\renewcommand{\rhebRtIvDSixDFullXMax}{0.99}  % ... max
\providecommand{\rhebRtIvDSixPFullXMin}{}\renewcommand{\rhebRtIvDSixPFullXMin}{0.105}  % ... bootstrap p, D6, six full-set: min
\providecommand{\rhebRtIvDSixPFullXMax}{}\renewcommand{\rhebRtIvDSixPFullXMax}{0.218}  % ... max
\providecommand{\rhebExRtIvFullXMin}{}\renewcommand{\rhebExRtIvFullXMin}{0.68}  % w/o best days: RT Iv SR STANDARD: min
\providecommand{\rhebExRtIvFullXMax}{}\renewcommand{\rhebExRtIvFullXMax}{1.06}  % ... max
\providecommand{\rhebExRtIvPFullXMin}{}\renewcommand{\rhebExRtIvPFullXMin}{0.367}  % ... its Ledoit-Wolf p (D19): min
\providecommand{\rhebExRtIvPFullXMax}{}\renewcommand{\rhebExRtIvPFullXMax}{0.807}  % ... max
\providecommand{\rhebRtMvShareFullXMin}{}\renewcommand{\rhebRtMvShareFullXMin}{90}  % RT MV median vol share %, six full-set: min
\providecommand{\rhebRtMvShareFullXMax}{}\renewcommand{\rhebRtMvShareFullXMax}{90}  % ... max
\providecommand{\rhebRtMvCapDaysFullXMin}{}\renewcommand{\rhebRtMvCapDaysFullXMin}{100}  % RT MV: % of days at the cap, six full-set: min
\providecommand{\rhebRtMvCapDaysFullXMax}{}\renewcommand{\rhebRtMvCapDaysFullXMax}{100}  % ... max
\providecommand{\rhebRtMvWsFullXMin}{}\renewcommand{\rhebRtMvWsFullXMin}{3.3}  % RT MV median dollar weight %, six full-set: min
\providecommand{\rhebRtMvWsFullXMax}{}\renewcommand{\rhebRtMvWsFullXMax}{3.5}  % ... max
\providecommand{\rhebRtIvWsFullXMin}{}\renewcommand{\rhebRtIvWsFullXMin}{0.8}  % RT IV median dollar weight %, six full-set: min
\providecommand{\rhebRtIvWsFullXMax}{}\renewcommand{\rhebRtIvWsFullXMax}{0.9}  % ... max
\providecommand{\rhebExBetaFullXMin}{}\renewcommand{\rhebExBetaFullXMin}{\ensuremath{-}0.21}  % w/o best days: beta, bid/ask, six full-set: min
\providecommand{\rhebExBetaFullXMax}{}\renewcommand{\rhebExBetaFullXMax}{3.27}  % ... max
\providecommand{\rhebExBetaTFullXMin}{}\renewcommand{\rhebExBetaTFullXMin}{\ensuremath{-}0.06}  % ... its HAC t: min
\providecommand{\rhebExBetaTFullXMax}{}\renewcommand{\rhebExBetaTFullXMax}{0.94}  % ... max
\providecommand{\rhebExAlphaFullXMin}{}\renewcommand{\rhebExAlphaFullXMin}{3.4}  % w/o best days: alpha %, bid/ask, six full-set: min
\providecommand{\rhebExAlphaFullXMax}{}\renewcommand{\rhebExAlphaFullXMax}{6.2}  % ... max
\providecommand{\rhebExAlphaTFullXMin}{}\renewcommand{\rhebExAlphaTFullXMin}{1.25}  % ... its HAC t: min
\providecommand{\rhebExAlphaTFullXMax}{}\renewcommand{\rhebExAlphaTFullXMax}{2.28}  % ... max
\providecommand{\rhebExAlphaPFullXMin}{}\renewcommand{\rhebExAlphaPFullXMin}{0.034}  % ... its fixed-b p: min
\providecommand{\rhebExAlphaPFullXMax}{}\renewcommand{\rhebExAlphaPFullXMax}{0.235}  % ... max
\providecommand{\rhebExARFullXMin}{}\renewcommand{\rhebExARFullXMin}{0.52}  % ... appraisal ratio (sqrt q, D19): min
\providecommand{\rhebExARFullXMax}{}\renewcommand{\rhebExARFullXMax}{0.94}  % ... max
\providecommand{\rhebExSRmaxFullXMin}{}\renewcommand{\rhebExSRmaxFullXMin}{0.57}  % ... max SR STANDARD: min
\providecommand{\rhebExSRmaxFullXMax}{}\renewcommand{\rhebExSRmaxFullXMax}{0.96}  % ... max
\providecommand{\rhebExGrsPFullXMin}{}\renewcommand{\rhebExGrsPFullXMin}{0.124}  % ... GRS p: min
\providecommand{\rhebExGrsPFullXMax}{}\renewcommand{\rhebExGrsPFullXMax}{0.396}  % ... max
\providecommand{\rhebBetaSigFullX}{}\renewcommand{\rhebBetaSigFullX}{0}  % six full-set: beta sig. 5%, bid/ask
\providecommand{\rhebBetaSigFullXWord}{}\renewcommand{\rhebBetaSigFullXWord}{zero}  % ... in words
\providecommand{\rhebAlphaSigFullX}{}\renewcommand{\rhebAlphaSigFullX}{6}  % six full-set: alpha sig. 5% (HAC), bid/ask
\providecommand{\rhebAlphaSigFullXWord}{}\renewcommand{\rhebAlphaSigFullXWord}{six}  % ... in words
\providecommand{\rhebGrsSigFullX}{}\renewcommand{\rhebGrsSigFullX}{4}  % six full-set: GRS sig. 5%, bid/ask
\providecommand{\rhebGrsSigFullXWord}{}\renewcommand{\rhebGrsSigFullXWord}{four}  % ... in words
\providecommand{\rhebBootSigFullX}{}\renewcommand{\rhebBootSigFullX}{6}  % six full-set: max-SR bootstrap p < 0.05
\providecommand{\rhebBootSigFullXWord}{}\renewcommand{\rhebBootSigFullXWord}{six}  % ... in words
\providecommand{\rhebRtMvBeatFullX}{}\renewcommand{\rhebRtMvBeatFullX}{6}  % six full-set: RT Mv SR > mkt (STD)
\providecommand{\rhebRtMvBeatFullXWord}{}\renewcommand{\rhebRtMvBeatFullXWord}{six}  % ... in words
\providecommand{\rhebRtMvSigFullX}{}\renewcommand{\rhebRtMvSigFullX}{0}  % ... and Ledoit-Wolf p < 0.05 (D19; was the bootstrap p)
\providecommand{\rhebRtMvSigFullXWord}{}\renewcommand{\rhebRtMvSigFullXWord}{zero}  % ... in words
\providecommand{\rhebRtMvDSixSigFullX}{}\renewcommand{\rhebRtMvDSixSigFullX}{0}  % ... D6 difference > 0 with bootstrap p < 0.05
\providecommand{\rhebRtMvDSixSigFullXWord}{}\renewcommand{\rhebRtMvDSixSigFullXWord}{zero}  % ... in words
\providecommand{\rhebExRtMvSigFullX}{}\renewcommand{\rhebExRtMvSigFullX}{0}  % ... w/o best days: STANDARD difference > 0 with p < 0.05
\providecommand{\rhebExRtMvSigFullXWord}{}\renewcommand{\rhebExRtMvSigFullXWord}{zero}  % ... in words
\providecommand{\rhebRtIvBeatFullX}{}\renewcommand{\rhebRtIvBeatFullX}{6}  % six full-set: RT Iv SR > mkt (STD)
\providecommand{\rhebRtIvBeatFullXWord}{}\renewcommand{\rhebRtIvBeatFullXWord}{six}  % ... in words
\providecommand{\rhebRtIvSigFullX}{}\renewcommand{\rhebRtIvSigFullX}{0}  % ... and Ledoit-Wolf p < 0.05 (D19; was the bootstrap p)
\providecommand{\rhebRtIvSigFullXWord}{}\renewcommand{\rhebRtIvSigFullXWord}{zero}  % ... in words
\providecommand{\rhebRtIvDSixSigFullX}{}\renewcommand{\rhebRtIvDSixSigFullX}{0}  % ... D6 difference > 0 with bootstrap p < 0.05
\providecommand{\rhebRtIvDSixSigFullXWord}{}\renewcommand{\rhebRtIvDSixSigFullXWord}{zero}  % ... in words
\providecommand{\rhebExRtIvSigFullX}{}\renewcommand{\rhebExRtIvSigFullX}{0}  % ... w/o best days: STANDARD difference > 0 with p < 0.05
\providecommand{\rhebExRtIvSigFullXWord}{}\renewcommand{\rhebExRtIvSigFullXWord}{zero}  % ... in words
\providecommand{\rhebExAlphaSigFullX}{}\renewcommand{\rhebExAlphaSigFullX}{1}  % ... w/o best days: alpha sig.
\providecommand{\rhebExAlphaSigFullXWord}{}\renewcommand{\rhebExAlphaSigFullXWord}{one}  % ... in words
\providecommand{\rhebExGrsSigFullX}{}\renewcommand{\rhebExGrsSigFullX}{0}  % ... w/o best days: GRS sig.
\providecommand{\rhebExGrsSigFullXWord}{}\renewcommand{\rhebExGrsSigFullXWord}{zero}  % ... in words
\providecommand{\rhebBetaHarXMin}{}\renewcommand{\rhebBetaHarXMin}{0.22}  % CAPM beta, bid/ask, four distinct HAR-set: min
\providecommand{\rhebBetaHarXMax}{}\renewcommand{\rhebBetaHarXMax}{2.53}  % ... max
\providecommand{\rhebBetaTHarXMin}{}\renewcommand{\rhebBetaTHarXMin}{0.07}  % ... its HAC t, bid/ask, four distinct HAR-set: min
\providecommand{\rhebBetaTHarXMax}{}\renewcommand{\rhebBetaTHarXMax}{0.74}  % ... max
\providecommand{\rhebBetaPHarXMin}{}\renewcommand{\rhebBetaPHarXMin}{0.477}  % ... its fixed-b p, bid/ask, four distinct HAR-set: min
\providecommand{\rhebBetaPHarXMax}{}\renewcommand{\rhebBetaPHarXMax}{0.950}  % ... max
\providecommand{\rhebCorrHarXMin}{}\renewcommand{\rhebCorrHarXMin}{0.00}  % correlation, bid/ask, four distinct HAR-set: min
\providecommand{\rhebCorrHarXMax}{}\renewcommand{\rhebCorrHarXMax}{0.02}  % ... max
\providecommand{\rhebRsqHarXMin}{}\renewcommand{\rhebRsqHarXMin}{0.00}  % CAPM R^2, bid/ask, four distinct HAR-set: min
\providecommand{\rhebRsqHarXMax}{}\renewcommand{\rhebRsqHarXMax}{0.00}  % ... max
\providecommand{\rhebAlphaHarXMin}{}\renewcommand{\rhebAlphaHarXMin}{\ensuremath{-}0.3}  % alpha % per day, bid/ask, four distinct HAR-set: min
\providecommand{\rhebAlphaHarXMax}{}\renewcommand{\rhebAlphaHarXMax}{4.8}  % ... max
\providecommand{\rhebAlphaTHarXMin}{}\renewcommand{\rhebAlphaTHarXMin}{\ensuremath{-}0.07}  % alpha HAC t, bid/ask, four distinct HAR-set: min
\providecommand{\rhebAlphaTHarXMax}{}\renewcommand{\rhebAlphaTHarXMax}{1.31}  % ... max
\providecommand{\rhebAlphaPHarXMin}{}\renewcommand{\rhebAlphaPHarXMin}{0.213}  % alpha fixed-b p, bid/ask, four distinct HAR-set: min
\providecommand{\rhebAlphaPHarXMax}{}\renewcommand{\rhebAlphaPHarXMax}{0.949}  % ... max
\providecommand{\rhebARHarXMin}{}\renewcommand{\rhebARHarXMin}{\ensuremath{-}0.03}  % appraisal ratio (sqrt q, D19), bid/ask, four distinct HAR-set: min
\providecommand{\rhebARHarXMax}{}\renewcommand{\rhebARHarXMax}{0.63}  % ... max
\providecommand{\rhebSRmaxHarXMin}{}\renewcommand{\rhebSRmaxHarXMin}{0.15}  % max SR STANDARD, bid/ask, four distinct HAR-set: min
\providecommand{\rhebSRmaxHarXMax}{}\renewcommand{\rhebSRmaxHarXMax}{0.65}  % ... max
\providecommand{\rhebDSRmaxHarXMin}{}\renewcommand{\rhebDSRmaxHarXMin}{0.00}  % max SR gain STANDARD, bid/ask, four distinct HAR-set: min
\providecommand{\rhebDSRmaxHarXMax}{}\renewcommand{\rhebDSRmaxHarXMax}{0.50}  % ... max
\providecommand{\rhebDSRmaxPHarXMin}{}\renewcommand{\rhebDSRmaxPHarXMin}{0.106}  % its bootstrap p, bid/ask, four distinct HAR-set: min
\providecommand{\rhebDSRmaxPHarXMax}{}\renewcommand{\rhebDSRmaxPHarXMax}{0.892}  % ... max
\providecommand{\rhebDSRmaxLoHarXMin}{}\renewcommand{\rhebDSRmaxLoHarXMin}{0.00}  % its 95% interval, lower end, four distinct HAR-set: min
\providecommand{\rhebDSRmaxLoHarXMax}{}\renewcommand{\rhebDSRmaxLoHarXMax}{0.00}  % ... max
\providecommand{\rhebDSRmaxHiHarXMin}{}\renewcommand{\rhebDSRmaxHiHarXMin}{0.98}  % its 95% interval, upper end, four distinct HAR-set: min
\providecommand{\rhebDSRmaxHiHarXMax}{}\renewcommand{\rhebDSRmaxHiHarXMax}{1.43}  % ... max
\providecommand{\rhebGrsPHarXMin}{}\renewcommand{\rhebGrsPHarXMin}{0.297}  % GRS p, bid/ask, four distinct HAR-set: min
\providecommand{\rhebGrsPHarXMax}{}\renewcommand{\rhebGrsPHarXMax}{0.955}  % ... max
\providecommand{\rhebTanShareHarXMin}{}\renewcommand{\rhebTanShareHarXMin}{4}  % tangency vol share %, bid/ask, four distinct HAR-set: min
\providecommand{\rhebTanShareHarXMax}{}\renewcommand{\rhebTanShareHarXMax}{95}  % ... max
\providecommand{\rhebSRsHarXMin}{}\renewcommand{\rhebSRsHarXMin}{\ensuremath{-}0.03}  % strategy SR STANDARD, bid/ask, four distinct HAR-set: min
\providecommand{\rhebSRsHarXMax}{}\renewcommand{\rhebSRsHarXMax}{0.63}  % ... max
\providecommand{\rhebRtMvHarXMin}{}\renewcommand{\rhebRtMvHarXMin}{0.26}  % real-time Mv SR STANDARD, bid/ask, four distinct HAR-set: min
\providecommand{\rhebRtMvHarXMax}{}\renewcommand{\rhebRtMvHarXMax}{0.89}  % ... max
\providecommand{\rhebRtMvDHarXMin}{}\renewcommand{\rhebRtMvDHarXMin}{\ensuremath{-}0.24}  % ... minus market, STANDARD, four distinct HAR-set: min
\providecommand{\rhebRtMvDHarXMax}{}\renewcommand{\rhebRtMvDHarXMax}{0.40}  % ... max
\providecommand{\rhebRtMvPHarXMin}{}\renewcommand{\rhebRtMvPHarXMin}{0.605}  % ... Ledoit-Wolf p (D19), four distinct HAR-set: min
\providecommand{\rhebRtMvPHarXMax}{}\renewcommand{\rhebRtMvPHarXMax}{0.908}  % ... max
\providecommand{\rhebRtMvDSixHarXMin}{}\renewcommand{\rhebRtMvDSixHarXMin}{\ensuremath{-}0.02}  % real-time Mv SR D6, bid/ask, four distinct HAR-set: min
\providecommand{\rhebRtMvDSixHarXMax}{}\renewcommand{\rhebRtMvDSixHarXMax}{0.80}  % ... max
\providecommand{\rhebRtMvDSixDHarXMin}{}\renewcommand{\rhebRtMvDSixDHarXMin}{\ensuremath{-}0.71}  % ... minus market, D6, four distinct HAR-set: min
\providecommand{\rhebRtMvDSixDHarXMax}{}\renewcommand{\rhebRtMvDSixDHarXMax}{0.12}  % ... max
\providecommand{\rhebRtMvDSixPHarXMin}{}\renewcommand{\rhebRtMvDSixPHarXMin}{0.408}  % ... bootstrap p, D6, four distinct HAR-set: min
\providecommand{\rhebRtMvDSixPHarXMax}{}\renewcommand{\rhebRtMvDSixPHarXMax}{0.945}  % ... max
\providecommand{\rhebExRtMvHarXMin}{}\renewcommand{\rhebExRtMvHarXMin}{0.19}  % w/o best days: RT Mv SR STANDARD: min
\providecommand{\rhebExRtMvHarXMax}{}\renewcommand{\rhebExRtMvHarXMax}{0.64}  % ... max
\providecommand{\rhebExRtMvPHarXMin}{}\renewcommand{\rhebExRtMvPHarXMin}{0.665}  % ... its Ledoit-Wolf p (D19): min
\providecommand{\rhebExRtMvPHarXMax}{}\renewcommand{\rhebExRtMvPHarXMax}{0.994}  % ... max
\providecommand{\rhebRtIvHarXMin}{}\renewcommand{\rhebRtIvHarXMin}{0.61}  % real-time Iv SR STANDARD, bid/ask, four distinct HAR-set: min
\providecommand{\rhebRtIvHarXMax}{}\renewcommand{\rhebRtIvHarXMax}{0.84}  % ... max
\providecommand{\rhebRtIvDHarXMin}{}\renewcommand{\rhebRtIvDHarXMin}{0.11}  % ... minus market, STANDARD, four distinct HAR-set: min
\providecommand{\rhebRtIvDHarXMax}{}\renewcommand{\rhebRtIvDHarXMax}{0.34}  % ... max
\providecommand{\rhebRtIvPHarXMin}{}\renewcommand{\rhebRtIvPHarXMin}{0.528}  % ... Ledoit-Wolf p (D19), four distinct HAR-set: min
\providecommand{\rhebRtIvPHarXMax}{}\renewcommand{\rhebRtIvPHarXMax}{0.845}  % ... max
\providecommand{\rhebRtIvDSixHarXMin}{}\renewcommand{\rhebRtIvDSixHarXMin}{0.66}  % real-time Iv SR D6, bid/ask, four distinct HAR-set: min
\providecommand{\rhebRtIvDSixHarXMax}{}\renewcommand{\rhebRtIvDSixHarXMax}{0.99}  % ... max
\providecommand{\rhebRtIvDSixDHarXMin}{}\renewcommand{\rhebRtIvDSixDHarXMin}{\ensuremath{-}0.03}  % ... minus market, D6, four distinct HAR-set: min
\providecommand{\rhebRtIvDSixDHarXMax}{}\renewcommand{\rhebRtIvDSixDHarXMax}{0.31}  % ... max
\providecommand{\rhebRtIvDSixPHarXMin}{}\renewcommand{\rhebRtIvDSixPHarXMin}{0.626}  % ... bootstrap p, D6, four distinct HAR-set: min
\providecommand{\rhebRtIvDSixPHarXMax}{}\renewcommand{\rhebRtIvDSixPHarXMax}{0.978}  % ... max
\providecommand{\rhebExRtIvHarXMin}{}\renewcommand{\rhebExRtIvHarXMin}{0.66}  % w/o best days: RT Iv SR STANDARD: min
\providecommand{\rhebExRtIvHarXMax}{}\renewcommand{\rhebExRtIvHarXMax}{0.86}  % ... max
\providecommand{\rhebExRtIvPHarXMin}{}\renewcommand{\rhebExRtIvPHarXMin}{0.571}  % ... its Ledoit-Wolf p (D19): min
\providecommand{\rhebExRtIvPHarXMax}{}\renewcommand{\rhebExRtIvPHarXMax}{0.835}  % ... max
\providecommand{\rhebRtMvShareHarXMin}{}\renewcommand{\rhebRtMvShareHarXMin}{90}  % RT MV median vol share %, four distinct HAR-set: min
\providecommand{\rhebRtMvShareHarXMax}{}\renewcommand{\rhebRtMvShareHarXMax}{90}  % ... max
\providecommand{\rhebRtMvCapDaysHarXMin}{}\renewcommand{\rhebRtMvCapDaysHarXMin}{85}  % RT MV: % of days at the cap, four distinct HAR-set: min
\providecommand{\rhebRtMvCapDaysHarXMax}{}\renewcommand{\rhebRtMvCapDaysHarXMax}{97}  % ... max
\providecommand{\rhebRtMvWsHarXMin}{}\renewcommand{\rhebRtMvWsHarXMin}{3.3}  % RT MV median dollar weight %, four distinct HAR-set: min
\providecommand{\rhebRtMvWsHarXMax}{}\renewcommand{\rhebRtMvWsHarXMax}{3.6}  % ... max
\providecommand{\rhebRtIvWsHarXMin}{}\renewcommand{\rhebRtIvWsHarXMin}{0.8}  % RT IV median dollar weight %, four distinct HAR-set: min
\providecommand{\rhebRtIvWsHarXMax}{}\renewcommand{\rhebRtIvWsHarXMax}{0.9}  % ... max
\providecommand{\rhebExBetaHarXMin}{}\renewcommand{\rhebExBetaHarXMin}{\ensuremath{-}0.31}  % w/o best days: beta, bid/ask, four distinct HAR-set: min
\providecommand{\rhebExBetaHarXMax}{}\renewcommand{\rhebExBetaHarXMax}{1.59}  % ... max
\providecommand{\rhebExBetaTHarXMin}{}\renewcommand{\rhebExBetaTHarXMin}{\ensuremath{-}0.10}  % ... its HAC t: min
\providecommand{\rhebExBetaTHarXMax}{}\renewcommand{\rhebExBetaTHarXMax}{0.49}  % ... max
\providecommand{\rhebExAlphaHarXMin}{}\renewcommand{\rhebExAlphaHarXMin}{0.0}  % w/o best days: alpha %, bid/ask, four distinct HAR-set: min
\providecommand{\rhebExAlphaHarXMax}{}\renewcommand{\rhebExAlphaHarXMax}{2.2}  % ... max
\providecommand{\rhebExAlphaTHarXMin}{}\renewcommand{\rhebExAlphaTHarXMin}{\ensuremath{-}0.01}  % ... its HAC t: min
\providecommand{\rhebExAlphaTHarXMax}{}\renewcommand{\rhebExAlphaTHarXMax}{0.66}  % ... max
\providecommand{\rhebExAlphaPHarXMin}{}\renewcommand{\rhebExAlphaPHarXMin}{0.527}  % ... its fixed-b p: min
\providecommand{\rhebExAlphaPHarXMax}{}\renewcommand{\rhebExAlphaPHarXMax}{0.991}  % ... max
\providecommand{\rhebExARHarXMin}{}\renewcommand{\rhebExARHarXMin}{\ensuremath{-}0.01}  % ... appraisal ratio (sqrt q, D19): min
\providecommand{\rhebExARHarXMax}{}\renewcommand{\rhebExARHarXMax}{0.33}  % ... max
\providecommand{\rhebExSRmaxHarXMin}{}\renewcommand{\rhebExSRmaxHarXMin}{0.23}  % ... max SR STANDARD: min
\providecommand{\rhebExSRmaxHarXMax}{}\renewcommand{\rhebExSRmaxHarXMax}{0.41}  % ... max
\providecommand{\rhebExGrsPHarXMin}{}\renewcommand{\rhebExGrsPHarXMin}{0.583}  % ... GRS p: min
\providecommand{\rhebExGrsPHarXMax}{}\renewcommand{\rhebExGrsPHarXMax}{0.992}  % ... max
\providecommand{\rhebBetaSigHarX}{}\renewcommand{\rhebBetaSigHarX}{0}  % four distinct HAR-set: beta sig. 5%, bid/ask
\providecommand{\rhebBetaSigHarXWord}{}\renewcommand{\rhebBetaSigHarXWord}{zero}  % ... in words
\providecommand{\rhebAlphaSigHarX}{}\renewcommand{\rhebAlphaSigHarX}{0}  % four distinct HAR-set: alpha sig. 5% (HAC), bid/ask
\providecommand{\rhebAlphaSigHarXWord}{}\renewcommand{\rhebAlphaSigHarXWord}{zero}  % ... in words
\providecommand{\rhebGrsSigHarX}{}\renewcommand{\rhebGrsSigHarX}{0}  % four distinct HAR-set: GRS sig. 5%, bid/ask
\providecommand{\rhebGrsSigHarXWord}{}\renewcommand{\rhebGrsSigHarXWord}{zero}  % ... in words
\providecommand{\rhebBootSigHarX}{}\renewcommand{\rhebBootSigHarX}{0}  % four distinct HAR-set: max-SR bootstrap p < 0.05
\providecommand{\rhebBootSigHarXWord}{}\renewcommand{\rhebBootSigHarXWord}{zero}  % ... in words
\providecommand{\rhebRtMvBeatHarX}{}\renewcommand{\rhebRtMvBeatHarX}{2}  % four distinct HAR-set: RT Mv SR > mkt (STD)
\providecommand{\rhebRtMvBeatHarXWord}{}\renewcommand{\rhebRtMvBeatHarXWord}{two}  % ... in words
\providecommand{\rhebRtMvSigHarX}{}\renewcommand{\rhebRtMvSigHarX}{0}  % ... and Ledoit-Wolf p < 0.05 (D19; was the bootstrap p)
\providecommand{\rhebRtMvSigHarXWord}{}\renewcommand{\rhebRtMvSigHarXWord}{zero}  % ... in words
\providecommand{\rhebRtMvDSixSigHarX}{}\renewcommand{\rhebRtMvDSixSigHarX}{0}  % ... D6 difference > 0 with bootstrap p < 0.05
\providecommand{\rhebRtMvDSixSigHarXWord}{}\renewcommand{\rhebRtMvDSixSigHarXWord}{zero}  % ... in words
\providecommand{\rhebExRtMvSigHarX}{}\renewcommand{\rhebExRtMvSigHarX}{0}  % ... w/o best days: STANDARD difference > 0 with p < 0.05
\providecommand{\rhebExRtMvSigHarXWord}{}\renewcommand{\rhebExRtMvSigHarXWord}{zero}  % ... in words
\providecommand{\rhebRtIvBeatHarX}{}\renewcommand{\rhebRtIvBeatHarX}{4}  % four distinct HAR-set: RT Iv SR > mkt (STD)
\providecommand{\rhebRtIvBeatHarXWord}{}\renewcommand{\rhebRtIvBeatHarXWord}{four}  % ... in words
\providecommand{\rhebRtIvSigHarX}{}\renewcommand{\rhebRtIvSigHarX}{0}  % ... and Ledoit-Wolf p < 0.05 (D19; was the bootstrap p)
\providecommand{\rhebRtIvSigHarXWord}{}\renewcommand{\rhebRtIvSigHarXWord}{zero}  % ... in words
\providecommand{\rhebRtIvDSixSigHarX}{}\renewcommand{\rhebRtIvDSixSigHarX}{0}  % ... D6 difference > 0 with bootstrap p < 0.05
\providecommand{\rhebRtIvDSixSigHarXWord}{}\renewcommand{\rhebRtIvDSixSigHarXWord}{zero}  % ... in words
\providecommand{\rhebExRtIvSigHarX}{}\renewcommand{\rhebExRtIvSigHarX}{0}  % ... w/o best days: STANDARD difference > 0 with p < 0.05
\providecommand{\rhebExRtIvSigHarXWord}{}\renewcommand{\rhebExRtIvSigHarXWord}{zero}  % ... in words
\providecommand{\rhebExAlphaSigHarX}{}\renewcommand{\rhebExAlphaSigHarX}{0}  % ... w/o best days: alpha sig.
\providecommand{\rhebExAlphaSigHarXWord}{}\renewcommand{\rhebExAlphaSigHarXWord}{zero}  % ... in words
\providecommand{\rhebExGrsSigHarX}{}\renewcommand{\rhebExGrsSigHarX}{0}  % ... w/o best days: GRS sig.
\providecommand{\rhebExGrsSigHarXWord}{}\renewcommand{\rhebExGrsSigHarXWord}{zero}  % ... in words
 in the preamble of paper_rehaul.tex; put the subsection below in
%          writeup/rehaul/10_appendix.tex, section "Robustness of the trade" (app:rh_robust_trade); \input the two
%          floats in rehaul/10_appendix_floats.tex. It cites Table tab:rh_econ_bestdays
%          (generated/rehaul_econ_decomp_bestdays.tex, the 2026-10-25 decomposition draft) and \rheHACS
%          (rehaul_econ_numbers.tex, already loaded).
% BibTeX entries to add to writeup/references.bib (not edited here; the table note and this draft cite them):
%   @article{TreynorMazuy1966,
%     author  = {Treynor, Jack L. and Mazuy, Kay K.},
%     title   = {Can Mutual Funds Outguess the Market?},
%     journal = {Harvard Business Review}, volume = {44}, number = {4}, pages = {131--136}, year = {1966}}
%   @article{GibbonsRossShanken1989,
%     author  = {Gibbons, Michael R. and Ross, Stephen A. and Shanken, Jay},
%     title   = {A Test of the Efficiency of a Given Portfolio},
%     journal = {Econometrica}, volume = {57}, number = {5}, pages = {1121--1152}, year = {1989}}
% Every number below is a macro (\rheb..., or \rheSzDays / \rheHACS); nothing is typed. The words "none", "all",
% "every" and the signs stated in words are asserted in make_rehaul_econ_beta.py (DRAFT CLAIMS), so a rerun whose
% data no longer bear them out stops.

\subsection{Market exposure and allocation with the market}
\label{app:rh_econ_beta}
The sign strategy of Section~\ref{sec:rh_econ} holds one straddle over the last half hour of the session. We ask
whether its returns are a disguised position in the stock market, and what an investor who holds the S\&P~500 gains
by adding the strategy. The specification below was fixed before any of these statistics was computed.
\begin{itemize}
\item \textbf{Returns.} The strategy's daily return $r^{s}_t$ (in percent of the premium, at mid-quotes unless noted)
  is treated as an excess return, as in Section~\ref{sec:rh_econ_measures}. The market return $r^{m}_t$ is the
  S\&P~500 close-to-close return of option day $t$, from the previous trading day's close, minus the return of the
  three-month Treasury bill (FRED series DTB3, actual/360). It is a price return: the repository has no total-return
  series, so the market's mean excess return omits the dividend yield, which leaves the regression slopes and
  residuals unchanged. We also use $r^{30}_t$, the index's own return from 15:30 ET to the close, the window of the
  trade.
  \begin{itemize}
  \item Over the \rhebNDays{} post-warmup days the market's excess return averages \rhebMktMean\% a day with a
    standard deviation of \rhebMktSd\%, a Sharpe ratio of \rhebSRm{} (\rhebSRmAll{} over all \rhebNAll{} trading days
    of the same span). Under the paper's ratio $\SR$ of Section~\ref{sec:rh_econ_measures}, which gives each calendar
    year, including the partial first and last ones, the same weight, it is \rhebSRmDSix.
  \item The decomposition of the maximum Sharpe ratio below holds for the ratio of the pooled daily returns, so every
    Sharpe ratio in this appendix is $\sqrt{252}$ times the ratio of the mean to the standard deviation of the daily
    returns unless it is marked as $\SR$. All $t$-statistics use Newey--West standard errors with bandwidth
    $S = \rheHACS$ and fixed-$b$ critical values (Section~\ref{sec:rh_methods_tests}).
  \end{itemize}
\item \textbf{No market beta} (Table~\ref{tab:rh_econ_beta}, Panel A). The correlation of the six full-set strategies
  with the market is \rhebCorrFullMin{} to \rhebCorrFullMax, and their CAPM betas are \rhebBetaFullMin{} to
  \rhebBetaFullMax{} percent of the premium per percent of market return ($t$ from \rhebBetaTFullMin{} to
  \rhebBetaTFullMax); none of the six is significant at the 5\% level, at mid-quotes or at quoted bid/ask prices
  ($R^2$ at most \rhebRsqFullMax). The closing-interval HAR strategy ($\beta = \rhebBetaCiHar$,
  $t = \rhebBetaTCiHar$) and the short straddle ($\beta = \rhebBetaShort$, $t = \rhebBetaTShort$) are no different,
  nor are the \rhebNHarDistinctWord{} distinct HAR-set forecasts ($\beta$ from \rhebBetaHarMin{} to
  \rhebBetaHarMax).
  \begin{itemize}
  \item On the \rhebNDown{} days on which the market falls, the full-set strategies earn \rhebDownMeanFullMin{} to
    \rhebDownMeanFullMax\% of the premium, against \rhebDownRestFullMin{} to \rhebDownRestFullMax\% on the other days;
    none of the six differences is significant.
  \item On the \rhebNWorst{} days with the lowest market return (at most \rhebWorstCut\%) every full-set strategy
    has a negative mean return, \rhebWorstMeanFullMin{} to \rhebWorstMeanFullMax\% of the premium, against
    \rhebWorstRestFullMin{} to \rhebWorstRestFullMax\% on the other days ($t$ of the difference from
    \rhebWorstDiffTFullMin{} to \rhebWorstDiffTFullMax; significant at the 5\% level for \rhebWorstSigFullWord{} of
    the six). The short straddle returns \rhebWorstMeanShort\% on these days ($t = \rhebWorstDiffTShort$). The
    strategy has no linear exposure to the market, but like the short straddle it mostly is, it loses on the market's
    worst days.
  \end{itemize}
\item \textbf{The convexity is in the trading window, not in the daily return.} In the regression of $r^{s}_t$ on
  $r^{m}_t$ and $(r^{m}_t)^2$ \citep{TreynorMazuy1966}, the coefficient $\gamma$ on the squared market return is
  \rhebTmGammaFullMin{} to \rhebTmGammaFullMax{} for the full-set strategies ($t$ from \rhebTmGammaTFullMin{} to
  \rhebTmGammaTFullMax; none significant), and on the strategies' own long days
  \rhebTmGammaLongFullMin{} to \rhebTmGammaLongFullMax{} ($t$ at most \rhebTmGammaLongTFullMax).
  \begin{itemize}
  \item Against the index's return over the trade's own half hour the straddle's payoff shows. On their long days
    every full-set strategy loads positively on $(r^{30}_t)^2$, $\gamma$ from \rhebThirtyTmGammaLongFullMin{} to
    \rhebThirtyTmGammaLongFullMax{} ($t$ from \rhebThirtyTmGammaLongTFullMin{} to \rhebThirtyTmGammaLongTFullMax),
    and on their short days negatively, \rhebThirtyTmGammaShortFullMin{} to \rhebThirtyTmGammaShortFullMax{} ($t$
    from \rhebThirtyTmGammaShortTFullMin{} to \rhebThirtyTmGammaShortTFullMax; percent of the premium per squared
    percent of index return). These $t$-statistics are descriptive: the long and short days are not consecutive.
  \item Over all days the full-set strategies are short the squared half-hour move, $\gamma$ from
    \rhebThirtyTmGammaFullMin{} to \rhebThirtyTmGammaFullMax{} ($t$ from \rhebThirtyTmGammaTFullMin{} to
    \rhebThirtyTmGammaTFullMax), against \rhebThirtyTmGammaShort{} ($t = \rhebThirtyTmGammaTShort$) for the short
    straddle, and they carry no slope on its sign ($\beta$ on $r^{30}_t$ with $t$ at most
    \rhebThirtyBetaTFullMax).
  \end{itemize}
\item \textbf{Alpha} (Panel B). The CAPM intercept of the full-set strategies is \rhebAlphaFullMin{} to
  \rhebAlphaFullMax\% of the premium a day ($t$ from \rhebAlphaTFullMin{} to \rhebAlphaTFullMax, fixed-$b$ $p$ at most
  \rhebAlphaPFullMax), with appraisal ratios $\alpha/\sigma_\varepsilon$ of \rhebARFullMin{} to \rhebARFullMax{}
  (annualized by $\sqrt{252}$). With a beta near zero, $\alpha$ is close to the mean return and the appraisal ratio
  to the strategy's own Sharpe ratio (\rhebSRsFullMin{} to \rhebSRsFullMax).
  \begin{itemize}
  \item The closing-interval HAR strategy has $\alpha = \rhebAlphaCiHar\%$ ($t = \rhebAlphaTCiHar$) and an appraisal
    ratio of \rhebARCiHar; the short straddle $\alpha = \rhebAlphaShort\%$ ($t = \rhebAlphaTShort$). For the
    \rhebNHarDistinctWord{} distinct HAR-set forecasts $\alpha$ is \rhebAlphaHarMin{} to \rhebAlphaHarMax\% ($t$ at
    most \rhebAlphaTHarMax; none significant).
  \end{itemize}
\item \textbf{Maximum Sharpe ratio of the market and the strategy.} The tangency portfolio of the two reaches
  $(\SR_m^2 + (\alpha/\sigma_\varepsilon)^2)^{1/2}$, \rhebSRmaxFullMin{} to \rhebSRmaxFullMax{} for the full-set
  strategies against \rhebSRm{} for the market alone. Its gain over the market is significant exactly when $\alpha$
  is: the one-asset test of \citet{GibbonsRossShanken1989} rejects $\alpha = 0$ for \rhebGrsSigFullWord{} of the six at
  the 5\% level ($p$ from \rhebGrsPFullMin{} to \rhebGrsPFullMax), and for the closing-interval HAR strategy
  ($p = \rhebGrsPCiHar$), but not for the short straddle ($p = \rhebGrsPShort$).
  \begin{itemize}
  \item In the tangency portfolio the strategy carries \rhebTanShareFullMin{} to \rhebTanShareFullMax\% of the
    volatility. The market's Sharpe ratio over the sample is low, so the in-sample optimum is almost the strategy
    alone; the strategy's leverage is arbitrary, so this share, not a dollar weight, is the allocation.
  \item The paired block bootstrap of the gain, with the weights re-estimated in each resample, gives $p$ from
    \rhebDSRmaxPFullMin{} to \rhebDSRmaxPFullMax{} (95\% intervals with lower ends from \rhebDSRmaxLoFullMin{} to
    \rhebDSRmaxLoFullMax). The gain cannot be negative, so this test rejects in one tail only and its $p$ lies below
    the GRS $p$: for the \rhebNHarDistinctWord{} distinct HAR-set forecasts it is below 5\% for
    \rhebBootSigHarWord{} and the GRS $p$ for none. The GRS test and the $t$-statistic of $\alpha$ are the tests of
    record.
  \end{itemize}
\item \textbf{Real-time allocation} (Panel C, Figure~\ref{fig:rh_econ_beta_cum}). Each day the weights of the market
  and the strategy are set from the earlier post-warmup days only. After a burn-in of \rhebNBurn{} days the
  combinations are evaluated on the \rhebNRt{} days from \rhebRtFirst, over which the market alone has a Sharpe ratio
  of \rhebSRmRt{} ($\SR = \rhebSRmRtDSix$). The largest of the strategies' best days, \rhebTopDate, falls in the
  burn-in: it moves the estimated weights but not the evaluation.
  \begin{itemize}
  \item Mean--variance weights (expanding window, no short positions, neither leg above \rhebCap\% of the estimated
    portfolio volatility) reach Sharpe ratios of \rhebRtMvFullMin{} to \rhebRtMvFullMax{} with the full-set
    strategies. The estimated means put the strategy at the cap on every evaluation day, which amounts to a median of
    \rhebRtMvWsFullMin{} to \rhebRtMvWsFullMax\% of capital in straddle premium; the rule is in effect a fixed split of
    the risk at the cap. The gain over the market is
    significant at the 5\% level for none of the six (bootstrap $p$ from \rhebRtMvPFullMin{} to
    \rhebRtMvPFullMax; with $\SR$, \rhebRtMvDSixFullMin{} to \rhebRtMvDSixFullMax{} and $p$ from
    \rhebRtMvDSixPFullMin{} to \rhebRtMvDSixPFullMax).
  \item Inverse-volatility weights over the trailing \rhebNTrail{} days, which estimate no means and split the risk
    equally (a median of at most \rhebRtIvWsFullMax\% of capital in premium), reach \rhebRtIvFullMin{} to
    \rhebRtIvFullMax; the gain is significant for \rhebRtIvSigFullWord{} of the six ($p$ from \rhebRtIvPFullMin{} to
    \rhebRtIvPFullMax; with $\SR$, for \rhebRtIvDSixSigFullWord).
  \item With the closing-interval HAR strategy the two rules reach \rhebRtMvCiHar{} ($p = \rhebRtMvPCiHar$) and
    \rhebRtIvCiHar{} ($p = \rhebRtIvPCiHar$), with the short straddle \rhebRtMvShort{} and \rhebRtIvShort{}.
  \end{itemize}
\item \textbf{Quoted bid/ask prices} (Panel D). The betas are as at mid-quotes (\rhebBetaFullXMin{} to
  \rhebBetaFullXMax, none significant). The intercept falls to \rhebAlphaFullXMin{} to \rhebAlphaFullXMax\% a day ($t$ from
  \rhebAlphaTFullXMin{} to \rhebAlphaTFullXMax), the GRS test rejects for \rhebGrsSigFullXWord{} of the six ($p$ up to
  \rhebGrsPFullXMax), and the maximum Sharpe ratio is \rhebSRmaxFullXMin{} to \rhebSRmaxFullXMax. In real time the
  mean--variance combinations reach \rhebRtMvFullXMin{} to \rhebRtMvFullXMax{} and the inverse-volatility ones
  \rhebRtIvFullXMin{} to \rhebRtIvFullXMax, and none is significantly above the market ($p$ at least
  \rhebRtMvPFullXMin{} and \rhebRtIvPFullXMin).
\item \textbf{Without the days that carry the gain} (Panel E). We remove the \rhebNBest{} days in the union of the six
  full-set forecasts' best days (Table~\ref{tab:rh_econ_bestdays}) from every series, the market's included. The
  selection uses the strategies' realized returns, so this is a stress test against them.
  \begin{itemize}
  \item The betas stay near zero (\rhebExBetaFullMin{} to \rhebExBetaFullMax, $t$ at most \rhebExBetaTFullMax). The
    intercept falls to \rhebExAlphaFullMin{} to \rhebExAlphaFullMax\% a day ($t$ from \rhebExAlphaTFullMin{} to
    \rhebExAlphaTFullMax), the maximum Sharpe ratio to \rhebExSRmaxFullMin{} to \rhebExSRmaxFullMax{} against
    \rhebSRmEx{} for the market, and the GRS test rejects for \rhebExGrsSigFullWord{} of the six ($p$ from
    \rhebExGrsPFullMin{} to \rhebExGrsPFullMax; at quoted bid/ask prices the $t$-test of $\alpha$ rejects for
    \rhebExAlphaSigFullXWord).
  \item The removed days are the short straddle's largest losses, so its intercept rises, to \rhebExAlphaShort\% a day
    ($t = \rhebExAlphaTShort$), above that of the full-set strategies. What remains of their $\alpha$ is the premium
    the short straddle earns on ordinary days, in line with Table~\ref{tab:rh_econ_bestdays}.
  \item In real time no combination is significantly above the market (mean--variance \rhebExRtMvFullMin{} to
    \rhebExRtMvFullMax, inverse volatility \rhebExRtIvFullMin{} to \rhebExRtIvFullMax, against \rhebSRmExRt).
  \end{itemize}
\item \textbf{Caveats.} The sample is short, \rheSzDays{} option days with one bear market, and the market's Sharpe
  ratio over it is low, which makes the tangency portfolio almost all strategy. The strategies' gains over the short
  straddle rest on the \rhebNBest{} days with large closing moves; the $t$-statistics allow for serial correlation, not
  for a mean carried by a few days. Bid/ask prices lower every number. The market return omits dividends. The
  strategies are tested without a correction for multiple testing.
\end{itemize}

\paragraph{Market exposure and allocation.} The sign strategy is not a disguised position in the market: its betas
are near zero and insignificant, and the convexity of the straddle shows in the trade's half hour, not in the daily
market return. Its intercept is significant, so in this sample a holder of the S\&P~500 would have given almost all of
the risk budget to the strategy. In real time, the combinations reach Sharpe ratios of \rhebRtIvFullMin{} to
\rhebRtMvFullMax{} at mid-quotes against \rhebSRmRt{} for the market, but over \rhebNRt{} days that gain is not
statistically significant for the mean--variance weights, significant for \rhebRtIvSigFullWord{} of the six
inverse-volatility combinations, and for none at quoted
bid/ask prices or without the \rhebNBest{} days that carry the strategies' gain over the short straddle.
